% Generated mechanically from frozen input canon/ko-kr/integration-20260908-r10-p03-register/inputs/p03/ko/cohomology.tex.
% Do not edit this generated file; edit the bound locale source instead.

% OK, start here.
%

\setcounter{chapter}{19}
\chapter{층 코호몰로지}



\phantomsection
\label{cohomology-section-phantom}


\section{도입}
\label{cohomology-section-introduction}

\noindent
이 장에서는 위상공간 위 층의 코호몰로지에 관한 몇 가지 주제를 전개한다.
주로 환의 층 위 가군의 일반성에서 작업하며 환 달린 공간의 사상을 다룬다.
사이트 위의 층에서는 어떤 일이 일어나는지 보려면 Cohomology on Sites,
Section \ref{sites-cohomology-section-introduction}을 보라. 기본 참고문헌은
\cite{Godement}와 \cite{Iversen}이다.





\section{층의 코호몰로지}
\label{cohomology-section-cohomology-sheaves}

\noindent
$X$를 위상공간, $\mathcal{F}$를 아벨 층이라 하자. $X$ 위의 아벨 층의
범주에는 단사 대상이 충분히 많음을 알고 있다. Injectives, Lemma
\ref{injectives-lemma-abelian-sheaves-space}를 보라. 따라서 단사 분해
$\mathcal{F}[0] \to \mathcal{I}^\bullet$을 택할 수 있다. 관례에 따라
\begin{equation}
\label{cohomology-equation-cohomology}
H^i(X, \mathcal{F}) = H^i(\Gamma(X, \mathcal{I}^\bullet))
\end{equation}
로 정의하고 이를 아벨 층 $\mathcal{F}$의 {\it 제$i$ 코호몰로지군}이라
한다. 함자족 $H^i(X, -)$는 $\textit{Ab}(X) \to \textit{Ab}$인 보편
$\delta$-함자를 이룬다.

\medskip\noindent
$f : X \to Y$를 위상공간의 연속 사상이라 하자. 위와 같은
$\mathcal{F}[0] \to \mathcal{I}^\bullet$을 사용하여
\begin{equation}
\label{cohomology-equation-higher-direct-image}
R^if_*\mathcal{F} = H^i(f_*\mathcal{I}^\bullet)
\end{equation}
로 정의하고 이를 $\mathcal{F}$의 {\it 제$i$ 고차 직상}이라 한다.
함자족 $R^if_*$는 $\textit{Ab}(X) \to \textit{Ab}(Y)$인 보편
$\delta$-함자를 이룬다.

\medskip\noindent
$(X, \mathcal{O}_X)$를 환 달린 공간, $\mathcal{F}$를
$\mathcal{O}_X$-가군이라 하자. $X$ 위의 $\mathcal{O}_X$-가군의 범주에는
단사 대상이 충분히 많음을 알고 있다. Injectives, Lemma
\ref{injectives-lemma-sheaves-modules-space}를 보라. 따라서 단사 분해
$\mathcal{F}[0] \to \mathcal{I}^\bullet$을 택할 수 있다. 관례에 따라
\begin{equation}
\label{cohomology-equation-cohomology-modules}
H^i(X, \mathcal{F}) = H^i(\Gamma(X, \mathcal{I}^\bullet))
\end{equation}
로 정의하고 이를 $\mathcal{F}$의 {\it 제$i$ 코호몰로지군}이라 한다.
함자족 $H^i(X, -)$는
$\textit{Mod}(\mathcal{O}_X) \to \text{Mod}_{\mathcal{O}_X(X)}$인 보편
$\delta$-함자를 이룬다.

\medskip\noindent
$f : (X, \mathcal{O}_X) \to (Y, \mathcal{O}_Y)$를 환 달린 공간의
사상이라 하자. 위와 같은 $\mathcal{F}[0] \to \mathcal{I}^\bullet$을
사용하여
\begin{equation}
\label{cohomology-equation-higher-direct-image-modules}
R^if_*\mathcal{F} = H^i(f_*\mathcal{I}^\bullet)
\end{equation}
로 정의하고 이를 $\mathcal{F}$의 {\it 제$i$ 고차 직상}이라 한다.
함자족 $R^if_*$는
$\textit{Mod}(\mathcal{O}_X) \to \textit{Mod}(\mathcal{O}_Y)$인 보편
$\delta$-함자를 이룬다.





\section{유도함자}
\label{cohomology-section-derived-functors}

\noindent
분해 함자를 사용하여 오른쪽 유도함자를 얻는 방법을 간단히 설명한다.
비유계 유도함자에 대해서는 Section \ref{cohomology-section-unbounded}를 보라.

\medskip\noindent
$(X, \mathcal{O}_X)$를 환 달린 공간이라 하자. 범주
$\textit{Mod}(\mathcal{O}_X)$는 아벨 범주이다. Modules, Lemma
\ref{modules-lemma-abelian}를 보라. 이 장에서는
$$
K(\mathcal{O}_X) = K(\textit{Mod}(\mathcal{O}_X))
\quad
\text{and}
\quad
D(\mathcal{O}_X) = D(\textit{Mod}(\mathcal{O}_X)).
$$
로 쓰며, Derived Categories, Definition
\ref{derived-definition-complexes-notation}과 Definition
\ref{derived-definition-unbounded-derived-category}에서 도입한
삼각범주의 유계 형태에도 같은 표기를 사용한다. Derived Categories,
Remark \ref{derived-remark-big-abelian-category}에 의해 분해 함자
$$
j = j_X :
K^{+}(\textit{Mod}(\mathcal{O}_X))
\longrightarrow
K^{+}(\mathcal{I})
$$
가 존재한다. 여기서 $\mathcal{I}$는 단사 층들로 이루어진
$\textit{Mod}(\mathcal{O}_X)$의 엄밀 충만 가법 부분범주이다. 임의의
아벨 범주 $\mathcal{B}$로 가는 왼쪽 완전 함자
$F : \textit{Mod}(\mathcal{O}_X) \to \mathcal{B}$에 대해, Derived
Categories, Section \ref{derived-section-right-derived-functor}에
설명되어 있고 방금의 분해 함자 $j_X$를 사용하여 구성되는 오른쪽 유도함자를
$RF$로 나타낸다.
\begin{equation}
\label{cohomology-equation-RF}
RF = F \circ j_X' : D^{+}(X) \longrightarrow D^{+}(\mathcal{B})
\end{equation}
표기는 Derived Categories, Lemma
\ref{derived-lemma-right-derived-functor}를 보라. 상황에 따라 $RF$는
$\textit{Mod}(\mathcal{O}_X)$,
$\text{Comp}^{+}(\textit{Mod}(\mathcal{O}_X))$,
$K^{+}(X)$ 또는 $D^{+}(X)$ 위에 정의된 것으로 생각할 수 있다.
Derived Categories, Definition
\ref{derived-definition-higher-derived-functors}에 따라 제$i$ 오른쪽
유도함자
\begin{equation}
\label{cohomology-equation-RFi}
R^iF = H^i \circ RF : \textit{Mod}(\mathcal{O}_X) \longrightarrow \mathcal{B}
\end{equation}
를 얻는다. 따라서 $R^0F = F$이고
$\{R^iF, \delta\}_{i \geq 0}$는 보편 $\delta$-함자이다. Derived
Categories, Lemma \ref{derived-lemma-higher-derived-functors}를 보라.

\medskip\noindent
이 구성의 두 가지 특수한 경우를 들자. 환 $R$이 주어지면
$K(R) = K(\text{Mod}_R)$과 $D(R) = D(\text{Mod}_R)$로 쓰고 유계
형태에도 같은 표기를 사용한다. 임의의 열린집합 $U \subset X$에 대해
왼쪽 완전 함자
$
\Gamma(U, -) :
\textit{Mod}(\mathcal{O}_X)
\longrightarrow
\text{Mod}_{\mathcal{O}_X(U)}
$
가 있고, 위의 논의로부터
\begin{equation}
\label{cohomology-equation-total-derived-cohomology}
R\Gamma(U, -) :
D^{+}(X)
\longrightarrow
D^{+}(\mathcal{O}_X(U))
\end{equation}
을 얻는다. $H^i(U, -) = R^i\Gamma(U, -)$로 둔다. $U = X$이면
(\ref{cohomology-equation-cohomology-modules})를 되찾는다. $f : X \to Y$가 환 달린
공간의 사상이면 왼쪽 완전 함자
$
f_* :
\textit{Mod}(\mathcal{O}_X)
\longrightarrow
\textit{Mod}(\mathcal{O}_Y)
$
가 있고, 이로부터 {\it 유도 직상}
\begin{equation}
\label{cohomology-equation-total-derived-direct-image}
Rf_* :
D^{+}(X)
\longrightarrow
D^{+}(Y)
\end{equation}
$Rf_*\mathcal{F}^\bullet$의 제$i$ 코호몰로지 층은
$R^if_*\mathcal{F}^\bullet$로 나타내고
(\ref{cohomology-equation-higher-direct-image-modules})에 따라 제$i$
{\it 고차 직상}이라 한다. 위에 표시한 두 함자는 유도범주 사이의 완전
함자이다.

\medskip\noindent
{\bf 표기의 남용:} 함자 $Rf_*$나 다른 유도함자를 $X$ 위의 층
$\mathcal{F}$ 또는 층의 복합체에 적용할 때는 $\mathcal{F}$를
$\mathcal{F}$의 적절한 분해로 바꾼 것으로 이해한다. 이런 연산을 편리하게
하기 위해 대상
$\mathcal{F}^\bullet \in D(\mathcal{O}_X)$가 주어졌을 때,
$\textit{Mod}(\mathcal{O}_X)$의 단사 대상들로 이루어진 아래로 유계인
복합체 $\mathcal{I}^\bullet$에 대해 준동형사상
$\mathcal{F}^\bullet \to \mathcal{I}^\bullet$이 존재하면 그 복합체가
{\it 유도범주에서 $\mathcal{F}^\bullet$을
표현한다}고 한다. 같은 취지에서 ``$\alpha :
\mathcal{F}^\bullet \to \mathcal{G}^\bullet$를
$D(\mathcal{O}_X)$의 사상이라 하자''라는 말은 $\alpha$가 복합체의
사상으로 표현된다는 뜻이 아니다. 실제 복합체의 사상이 있을 때는 그렇게
명시한다.









\section{제1 코호몰로지와 토서}
\label{cohomology-section-h1-torsors}

\begin{definition}
\label{cohomology-definition-torsor}
$X$를 위상공간이라 하고,
$\mathcal{G}$를 $X$ 위의 (가환하지 않을 수도 있는) 군의 층이라 하자.
{\it 유사 토서}, 더 정확히는 {\it 유사 $\mathcal{G}$-토서}란 작용
$\mathcal{G} \times \mathcal{F} \to \mathcal{F}$가 주어진 $X$ 위의
집합의 층 $\mathcal{F}$로서 다음을 만족하는 것을 말한다.
\begin{enumerate}
\item $\mathcal{F}(U)$가 공집합이 아닐 때마다 작용
$\mathcal{G}(U) \times \mathcal{F}(U) \to \mathcal{F}(U)$
은 단순 추이적이다.
\end{enumerate}
{\it 유사 $\mathcal{G}$-토서의 사상} $\mathcal{F} \to \mathcal{F}'$은
$\mathcal{G}$-작용과 양립하는 집합의 층의 사상이다. {\it 토서}, 더 정확히는
{\it $\mathcal{G}$-토서}란 다음 조건도 만족하는 유사 $\mathcal{G}$-토서이다.
\begin{enumerate}
\item[(2)] 모든 $x \in X$에 대해 줄기 $\mathcal{F}_x$가 공집합이 아니다.
\end{enumerate}
{\it $\mathcal{G}$-토서의 사상}은 유사 $\mathcal{G}$-토서의 사상이다.
{\it 자명한 $\mathcal{G}$-토서}는 자명한 왼쪽 $\mathcal{G}$-작용이 주어진
층 $\mathcal{G}$이다.
\end{definition}

\noindent
토서의 사상은 자동으로 동형사상임이 분명하다.

\begin{lemma}
\label{cohomology-lemma-trivial-torsor}
$X$를 위상공간이라 하고,
$\mathcal{G}$를 $X$ 위의 (가환하지 않을 수도 있는) 군의 층이라 하자.
$\mathcal{G}$-토서 $\mathcal{F}$가 자명할 필요충분조건은
$\mathcal{F}(X) \not = \emptyset$인 것이다.
\end{lemma}

\begin{proof}
생략한다.
\end{proof}

\begin{lemma}
\label{cohomology-lemma-torsors-h1}
$X$를 위상공간이라 하고, $\mathcal{H}$를 $X$ 위의 아벨 층이라 하자.
$\mathcal{H}$-토서의 동형류 집합과 $H^1(X, \mathcal{H})$ 사이에는
표준 전단사가 있다.
\end{lemma}

\begin{proof}
$\mathcal{F}$를 $\mathcal{H}$-토서라 하자. $\mathcal{F}$ 위의 자유 아벨 층
$\mathbf{Z}[\mathcal{F}]$을 생각하자. 이것은 열린집합 $U \subset X$에
$n_i \in \mathbf{Z}$이고 $s_i \in \mathcal{F}(U)$인 유한 형식합
$\sum n_i[s_i]$들의 모임을 대응시키는 규칙의 층화이다. 자연스러운 사상
$$
\sigma : \mathbf{Z}[\mathcal{F}] \longrightarrow \underline{\mathbf{Z}}
$$
이 있으며, 국소 단면 $\sum n_i[s_i]$를 $\sum n_i$로 보낸다.
$\sigma$의 핵은 $[s] - [s']$ 꼴의 국소 단면들로 생성된다. 표준 사상
$a : \Ker(\sigma) \to \mathcal{H}$
이 있으며, $[s] - [s'] \mapsto h$로 보내는데, 여기서 $h$는
$h \cdot s' = s$를 만족하는 $\mathcal{H}$의 국소 단면이다. 밀어내기 도식
$$
\xymatrix{
0 \ar[r] &
\Ker(\sigma) \ar[r] \ar[d]^a &
\mathbf{Z}[\mathcal{F}] \ar[r] \ar[d] &
\underline{\mathbf{Z}} \ar[r] \ar[d] &
0 \\
0 \ar[r] &
\mathcal{H} \ar[r] &
\mathcal{E} \ar[r] &
\underline{\mathbf{Z}} \ar[r] &
0
}
$$
을 생각하자. 여기서 $\mathcal{E}$는 밀어내기로 얻은 확대이다. 아래의
짧은 완전열에 딸린 코호몰로지 긴 완전열에서
$\xi = \xi_\mathcal{F} \in H^1(X, \mathcal{H})$
를 얻는데, 이는 경계 사상을 $1 \in H^0(X, \underline{\mathbf{Z}})$에
적용하여 얻는다.

\medskip\noindent
거꾸로 $\xi \in H^1(X, \mathcal{H})$가 주어지면 다음과 같이 $\xi$에
토서를 대응시킬 수 있다. $\mathcal{H}$를 단사 아벨 층 $\mathcal{I}$에
넣는 매장 $\mathcal{H} \to \mathcal{I}$를 택하자.
$\mathcal{Q} = \mathcal{I}/\mathcal{H}$로 두면 다음 짧은 완전열을 얻는다.
$$
\xymatrix{
0 \ar[r] &
\mathcal{H} \ar[r] &
\mathcal{I} \ar[r] &
\mathcal{Q} \ar[r] &
0
}
$$
$H^1(X, \mathcal{I}) = 0$이므로(Derived Categories, Lemma
\ref{derived-lemma-higher-derived-functors} 참조), 원소 $\xi$는 전체 단면
$q \in H^0(X, \mathcal{Q})$의 상이다. 층 $\mathcal{Q}$에서 $q$로 가는
단면들로 이루어진 부분층(집합의 층으로서)을
$\mathcal{F} \subset \mathcal{I}$라 하자. $\mathcal{F}$가 토서임은
쉽게 확인할 수 있다.

\medskip\noindent
위에서 준 두 구성이 서로 역이라는 확인은 생략한다.
\end{proof}







\section{제1 코호몰로지와 확대}
\label{cohomology-section-h1-extensions}

\begin{lemma}
\label{cohomology-lemma-h1-extensions}
$(X, \mathcal{O}_X)$를 환 달린 공간이라 하고, $\mathcal{F}$를
$\mathcal{O}_X$-가군층이라 하자. 표준 전단사
$$
\Ext^1_{\textit{Mod}(\mathcal{O}_X)}(\mathcal{O}_X, \mathcal{F})
\longrightarrow
H^1(X, \mathcal{F})
$$
가 있으며, 이 전단사는 확대
$$
0 \to \mathcal{F} \to \mathcal{E} \to \mathcal{O}_X \to 0
$$
에 $1 \in \Gamma(X, \mathcal{O}_X)$의 $H^1(X, \mathcal{F})$에서의 상을
대응시킨다.
\end{lemma}

\begin{proof}
보조정리에서 준 사상의 역을 구성하자. $\xi \in H^1(X, \mathcal{F})$라 하자.
$\textit{Mod}(\mathcal{O}_X)$에서 $\mathcal{I}$가 단사 대상이 되도록
단사 사상 $\mathcal{F} \subset \mathcal{I}$를 택한다.
$\mathcal{Q} = \mathcal{I}/\mathcal{F}$로 두자. 코호몰로지 긴 완전열에 의해
$\xi$는 단면
$\tilde \xi \in \Gamma(X, \mathcal{Q}) =
\Hom_{\mathcal{O}_X}(\mathcal{O}_X, \mathcal{Q})$
의 상이다. 이제 다음 당김을 만들면 된다.
$$
\xymatrix{
0 \ar[r] &
\mathcal{F} \ar[r] \ar@{=}[d] &
\mathcal{E} \ar[r] \ar[d] &
\mathcal{O}_X \ar[r] \ar[d]^{\tilde \xi} &
0 \\
0 \ar[r] &
\mathcal{F} \ar[r] &
\mathcal{I} \ar[r] &
\mathcal{Q} \ar[r] &
0
}
$$
Homology, Section \ref{homology-section-extensions}을 보라.
\end{proof}








\section{제1 코호몰로지와 가역층}
\label{cohomology-section-invertible-sheaves}

\noindent
환 달린 공간의 피카르 군은 Modules, Section
\ref{modules-section-invertible}에서 정의하였다.

\begin{lemma}
\label{cohomology-lemma-h1-invertible}
$(X, \mathcal{O}_X)$를 환 달린 공간이라 하자. 모든 줄기
$\mathcal{O}_{X, x}$가 국소환이면 아벨군의 표준 동형
$$
H^1(X, \mathcal{O}_X^*) = \Pic(X).
$$
이 있다.
\end{lemma}

\begin{proof}
$\mathcal{L}$을 가역 $\mathcal{O}_X$-가군이라 하자. 다음 규칙으로 정의된
전층 $\mathcal{L}^*$를 생각하자.
$$
U \longmapsto \{s \in \mathcal{L}(U)
\text{ 중 } \mathcal{O}_U \xrightarrow{s \cdot -} \mathcal{L}_U
\text{ 가 동형사상인 것}\}
$$
이 전층은 층 조건을 만족한다. 또한 $f \in \mathcal{O}_X^*(U)$이고
$s \in \mathcal{L}^*(U)$이면 분명히 $fs \in \mathcal{L}^*(U)$이다.
마찬가지로 $s, s' \in \mathcal{L}^*(U)$이면 $fs = s'$를 만족하는
$f \in \mathcal{O}_X^*(U)$가 유일하게 존재한다. 더욱이 Modules, Lemma
\ref{modules-lemma-invertible-is-locally-free-rank-1}에 의해 층
$\mathcal{L}^*$는 국소적으로 단면을 가진다. 다시 말해
$\mathcal{L}^*$는 $\mathcal{O}_X^*$-토서이다. 따라서 다음 사상을 얻는다.
$$
\begin{matrix}
\text{환 달린 공간 }(X, \mathcal{O}_X)\text{ 위의 가역층} \\
\text{의 동형류}
\end{matrix}
\longrightarrow
\begin{matrix}
\mathcal{O}_X^*\text{-토서} \\
\text{의 동형류}
\end{matrix}
$$
이 사상이 아벨군의 준동형이라는 확인은 생략한다. 보조정리
\ref{cohomology-lemma-torsors-h1}에 의해 오른쪽은 $H^1(X, \mathcal{O}_X^*)$와
표준적으로 전단사이다. 따라서 이 사상이 단사이고 전사임을 보이면 된다.

\medskip\noindent
단사성. 토서 $\mathcal{L}^*$가 자명하면 보조정리
\ref{cohomology-lemma-trivial-torsor}에 의해 $\mathcal{L}^*$는 전체 단면을 가진다.
따라서 이는 정확히 $\mathcal{L} \cong \mathcal{O}_X$가
$\Pic(X)$의 항등원이라는 뜻이다.

\medskip\noindent
전사성. $\mathcal{F}$를 $\mathcal{O}_X^*$-토서라 하자. 집합의 전층
$$
\mathcal{L}_1 : U \longmapsto
(\mathcal{F}(U) \times \mathcal{O}_X(U))/\mathcal{O}_X^*(U)
$$
을 생각하자. 여기서 $f \in \mathcal{O}_X^*(U)$가 $(s, g)$에 작용한 결과는
$(fs, f^{-1}g)$이다. $s'/s$를 $fs = s'$를 만족하는
$\mathcal{O}_X^*$의 국소 단면 $f$라 하고,
$(s, g) + (s', g') = (s, g + (s'/s)g')$로 두며, $\mathcal{O}_X$의
국소 단면 $h$에 대해 $h(s, g) = (s, hg)$로 두면 $\mathcal{L}_1$은
$\mathcal{O}_X$-가군의 전층이 된다. 그 층화
$\mathcal{L} = \mathcal{L}_1^\#$가 가역 $\mathcal{O}_X$-가군이고,
이에 딸린 $\mathcal{O}_X^*$-토서 $\mathcal{L}^*$가 $\mathcal{F}$와
동형이라는 확인은 생략한다.
\end{proof}













\section{코호몰로지의 국소성}
\label{cohomology-section-locality}

\noindent
다음 보조정리는 열린집합 위에서 층 $\mathcal{F}$의 코호몰로지를 정의하는 데
모호함이 없음을 말한다.

\begin{lemma}
\label{cohomology-lemma-cohomology-of-open}
$X$를 환 달린 공간이라 하고, $U \subset X$를 열린 부분공간이라 하자.
\begin{enumerate}
\item $\mathcal{I}$가 단사 $\mathcal{O}_X$-가군이면
$\mathcal{I}|_U$는 단사 $\mathcal{O}_U$-가군이다.
\item 임의의 $\mathcal{O}_X$-가군층 $\mathcal{F}$에 대해
$H^p(U, \mathcal{F}) = H^p(U, \mathcal{F}|_U)$.
\end{enumerate}
\end{lemma}

\begin{proof}
$j : U \to X$를 열린 몰입이라 하자. $U$로 제한하는 함자 $j^{-1}$는
$0$에 의한 연장 함자 $j_!$의 오른쪽 수반임을 상기하라. Sheaves, Lemma
\ref{sheaves-lemma-j-shriek-modules}를 보라. 또한 $j_!$는 완전하다.
따라서 (1)은 Homology, Lemma
\ref{homology-lemma-adjoint-preserve-injectives}에서 따른다.

\medskip\noindent
정의에 따라 $H^p(U, \mathcal{F}) = H^p(\Gamma(U, \mathcal{I}^\bullet))$이다.
여기서 $\mathcal{F} \to \mathcal{I}^\bullet$은
$\textit{Mod}(\mathcal{O}_X)$에서의 단사 분해이다. 위에서 본 바에 따라
$\mathcal{F}|_U \to \mathcal{I}^\bullet|_U$는
$\textit{Mod}(\mathcal{O}_U)$에서의 단사 분해이다. 따라서
$H^p(U, \mathcal{F}|_U)$는
$H^p(\Gamma(U, \mathcal{I}^\bullet|_U))$와 같다. 물론 $X$ 위의 임의의 층
$\mathcal{F}$에 대해
$\Gamma(U, \mathcal{F}) = \Gamma(U, \mathcal{F}|_U)$이다.
따라서 (2)의 등식이 따른다.
\end{proof}

\noindent
$X$를 환 달린 공간, $\mathcal{F}$를 $\mathcal{O}_X$-가군층,
$U \subset V \subset X$를 열린집합들이라 하자. 그러면 표준적인
{\it 제한 사상}
\begin{equation}
\label{cohomology-equation-restriction-mapping}
H^n(V, \mathcal{F})
\longrightarrow
H^n(U, \mathcal{F}), \quad
\xi \longmapsto \xi|_U
\end{equation}
이 있으며 이는 $\mathcal{F}$에 관해 함자적이다. 실제로 임의의 단사 분해
$\mathcal{F} \to \mathcal{I}^\bullet$을 택하자. 층 $\mathcal{I}^p$의
제한 사상들은 복합체의 사상
$$
\Gamma(V, \mathcal{I}^\bullet)
\longrightarrow
\Gamma(U, \mathcal{I}^\bullet)
$$
을 준다. 왼쪽은 $R\Gamma(V, \mathcal{F})$를 나타내는 복합체이고,
오른쪽은 $R\Gamma(U, \mathcal{F})$를 나타내는 복합체이다. 함자 $H^n$을
적용하여 코호몰로지군 사이의 사상을 얻는다. 앞서 표시했듯이 이 사상을
$\xi \mapsto \xi|_U$로 나타낸다. 따라서 규칙
$U \mapsto H^n(U, \mathcal{F})$는 $\mathcal{O}_X$-가군의 준층이다.
이 준층은 관례상 $\underline{H}^n(\mathcal{F})$로 나타낸다. 보조정리
\ref{cohomology-lemma-include}에서 이 준층을 달리 해석할 것이다.

\begin{lemma}
\label{cohomology-lemma-kill-cohomology-class-on-covering}
$X$를 환 달린 공간, $\mathcal{F}$를 $\mathcal{O}_X$-가군층,
$U \subset X$를 열린 부분공간이라 하자. $n > 0$이고
$\xi \in H^n(U, \mathcal{F})$라 하자. 그러면 모든 $i \in I$에 대해
$\xi|_{U_i} = 0$이 되는 열린 덮개 $U = \bigcup_{i\in I} U_i$가 존재한다.
\end{lemma}

\begin{proof}
$\mathcal{F} \to \mathcal{I}^\bullet$을 단사 분해라 하자. 그러면
$$
H^n(U, \mathcal{F}) =
\frac{\Ker(\mathcal{I}^n(U) \to \mathcal{I}^{n + 1}(U))}
{\Im(\mathcal{I}^{n - 1}(U) \to \mathcal{I}^n(U))}.
$$
위의 표시에 나오는 코호몰로지류를 나타내는 원소
$\tilde \xi \in \mathcal{I}^n(U)$를 택하자. $\mathcal{I}^\bullet$은
$\mathcal{F}$의 단사 분해이고 $n > 0$이므로, 복합체
$\mathcal{I}^\bullet$은 차수 $n$에서 완전하다. 따라서 층으로서
$\Im(\mathcal{I}^{n - 1} \to \mathcal{I}^n) =
\Ker(\mathcal{I}^n \to \mathcal{I}^{n + 1})$이다. $\tilde \xi$는
$U$ 위에서 핵층의 단면이므로, $\tilde \xi|_{U_i}$가 어떤 단면
$\xi_i \in \mathcal{I}^{n - 1}(U_i)$의 $d$에 의한 상이 되는 열린 덮개
$U = \bigcup_{i \in I} U_i$가 존재한다. 제한 $\xi|_{U_i}$를
$\tilde \xi|_{U_i}$의 류에 대응하는 것으로 정의했으므로 결론이 따른다.
\end{proof}

\begin{lemma}
\label{cohomology-lemma-describe-higher-direct-images}
$f : X \to Y$를 환 달린 공간의 사상이라 하고,
$\mathcal{F}$를 $\mathcal{O}_X$-가군이라 하자. 층 $R^if_*\mathcal{F}$는 준층
$$
V \longmapsto H^i(f^{-1}(V), \mathcal{F})
$$
에 결부된 층이며, 제한 사상은 식 (\ref{cohomology-equation-restriction-mapping})의
것이다. 아래로 유계인 복합체 $\mathcal{F}^\bullet$에 $R^if_*$를 적용하는
경우에도 같은 명제가 성립한다.
\end{lemma}

\begin{proof}
$\mathcal{F} \to \mathcal{I}^\bullet$을 단사 분해라 하자. 정의에 따라
$R^if_*\mathcal{F}$는 복합체
$$
f_*\mathcal{I}^0 \to f_*\mathcal{I}^1 \to f_*\mathcal{I}^2 \to \ldots
$$
의 제$i$ 코호몰로지층이다. $\mathcal{O}_Y$-가군층의 아벨 범주 구조의
정의에 따라 이 코호몰로지층은 준층
$$
V
\longmapsto
\frac{\Ker(f_*\mathcal{I}^i(V) \to f_*\mathcal{I}^{i + 1}(V))}
{\Im(f_*\mathcal{I}^{i - 1}(V) \to f_*\mathcal{I}^i(V))}
$$
에 결부된 층이다. 이것은 분명히
$$
\frac{\Ker(\mathcal{I}^i(f^{-1}(V)) \to \mathcal{I}^{i + 1}(f^{-1}(V)))}
{\Im(\mathcal{I}^{i - 1}(f^{-1}(V)) \to \mathcal{I}^i(f^{-1}(V)))}
$$
와 같고, 이는 $H^i(f^{-1}(V), \mathcal{F})$와 같다. 이로써 끝난다.
\end{proof}

\begin{lemma}
\label{cohomology-lemma-localize-higher-direct-images}
$f : X \to Y$를 환 달린 공간의 사상, $\mathcal{F}$를
$\mathcal{O}_X$-가군, $V \subset Y$를 열린 부분공간이라 하자.
$g : f^{-1}(V) \to V$를 $f$의 제한이라 하자. 그러면
$$
R^pg_*(\mathcal{F}|_{f^{-1}(V)}) = (R^pf_*\mathcal{F})|_V
$$
이다. $\mathcal{F}^\bullet$이 아래로 유계인 $\mathcal{O}_X$-가군의
복합체일 때 유도상 $Rf_*\mathcal{F}^\bullet$에 대해서도 같은 명제가
성립한다.
\end{lemma}

\begin{proof}
첫째 증명. 보조정리 \ref{cohomology-lemma-describe-higher-direct-images}와
\ref{cohomology-lemma-cohomology-of-open}을 적용하면 표시된 등식을 얻는다.
둘째 증명. 단사 분해 $\mathcal{F} \to \mathcal{I}^\bullet$을 택하고
$\mathcal{F}|_{f^{-1}(V)} \to \mathcal{I}^\bullet|_{f^{-1}(V)}$도
단사 분해라는 사실을 사용한다.
\end{proof}

\begin{remark}
\label{cohomology-remark-daniel}
위의 보조정리 \ref{cohomology-lemma-kill-cohomology-class-on-covering}과
\ref{cohomology-lemma-describe-higher-direct-images}의 증명에 대한 다른 접근을 설명한다.
$(X, \mathcal{O}_X)$를 환 달린 공간이라 하자.
$i_X : \textit{Mod}(\mathcal{O}_X) \to \textit{PMod}(\mathcal{O}_X)$를
포함 함자, $\#$를 층화 함자라 하자. $i_X$는 왼쪽 완전이고 $\#$는
완전함을 상기하라.
\begin{enumerate}
\item 먼저 $i_X$의 오른쪽 유도함자가
$R^pi_X\mathcal{F} = \underline{H}^p(\mathcal{F})$로 주어진다는 아래의
보조정리 \ref{cohomology-lemma-include}를 증명한다. 다른 증명은 다음과 같다.
$p = 0$일 때 등식은 분명하다. $(R^pi_X)_{p \geq 0}$와
$(\underline{H}^p)_{p \geq 0}$는 모두 단사 대상에서 사라지는
델타 함자이므로 둘 다 보편적이고, 따라서 서로 동형이다. Homology,
Section \ref{homology-section-cohomological-delta-functor}을 보라.
\item 보조정리 \ref{cohomology-lemma-kill-cohomology-class-on-covering}은 임의의
$\mathcal{O}_X$-가군층 $\mathcal{F}$와 $p > 0$에 대해
$(\underline{H}^p(\mathcal{F}))^\# = 0$이라고 다시 말할 수 있다.
이를 보려면 ${}^\#$가 완전하므로
$$
(\underline{H}^p(\mathcal{F}))^\# =
(R^pi_X\mathcal{F})^\# =
R^p(\# \circ i_X)(\mathcal{F}) = 0
$$
을 얻으면 된다. 실제로 $\# \circ i_X$는 항등함자이다.
\item $f : X \to Y$를 환 달린 공간의 사상, $\mathcal{F}$를
$\mathcal{O}_X$-가군이라 하자. 준층
$V \mapsto H^p(f^{-1}V, \mathcal{F})$는
$R^p (i_Y \circ f_*)\mathcal{F}$와 같다. 위의 (1)의 논증에서와 같이
둘 다 보편 델타 함자를 준다는 것을 관찰하여 이를 증명할 수 있다.
따라서 보조정리 \ref{cohomology-lemma-describe-higher-direct-images}는
$R^p f_* \mathcal{F}= (R^p (i_Y \circ f_*)\mathcal{F})^\#$라고 말한다.
다시 $\#$가 완전하고 $\# \circ i_Y$가 항등함자라는 사실을 사용하면
$$
R^p f_* \mathcal{F} =
R^p(\# \circ i_Y \circ f_*)\mathcal{F} =
(R^p (i_Y \circ f_*)\mathcal{F})^\#
$$
을 얻는다. 이것이 원하는 결과이다.
\end{enumerate}
\end{remark}






\section{Mayer-Vietoris}
\label{cohomology-section-mayer-vietoris}

\noindent
아래에서 체흐-코호몰로지 스펙트럼열을 구성할 것이다. 보조정리
\ref{cohomology-lemma-cech-spectral-sequence}를 보라. 그 스펙트럼열의 특수한 경우가
Mayer--Vietoris 긴 완전열이다. 이는 스펙트럼열의 매우 기본적이고 유용하며
이해하기 쉬운 변형이므로 여기서 따로 다룬다.

\begin{lemma}
\label{cohomology-lemma-injective-restriction-surjective}
\begin{slogan}
단사 가군층은 플라스크이다.
\end{slogan}
$X$를 환 달린 공간이라 하고, $U' \subset U \subset X$를 열린
부분공간들이라 하자. 임의의 단사 $\mathcal{O}_X$-가군 $\mathcal{I}$에 대해
제한 사상 $\mathcal{I}(U) \to \mathcal{I}(U')$는 전사이다.
\end{lemma}

\begin{proof}
$j : U \to X$와 $j' : U' \to X$를 열린 몰입들이라 하자.
$j_!\mathcal{O}_U$는 $\mathcal{O}_U = \mathcal{O}_X|_U$의 영에 의한
연장임을 상기하라. Sheaves, Section
\ref{sheaves-section-open-immersions}을 보라. $j_!$는 제한의 왼쪽
수반이므로 임의의 $\mathcal{O}_X$-가군층 $\mathcal{F}$에 대해
$$
\Hom_{\mathcal{O}_X}(j_!\mathcal{O}_U, \mathcal{F})
=
\Hom_{\mathcal{O}_U}(\mathcal{O}_U, \mathcal{F}|_U)
=
\mathcal{F}(U)
$$
이다. Sheaves, Lemma \ref{sheaves-lemma-j-shriek-modules}를 보라.
마찬가지로 층 $j'_!\mathcal{O}_{U'}$는 함자
$\mathcal{F} \mapsto \mathcal{F}(U')$를 표현한다. 또한 자명한 표준
$\mathcal{O}_X$-가군 사상
$$
j'_!\mathcal{O}_{U'} \longrightarrow j_!\mathcal{O}_U
$$
이 있는데, 이는 Yoneda 보조정리(Categories, Lemma
\ref{categories-lemma-yoneda})를 통해 제한 사상
$\mathcal{F}(U) \to \mathcal{F}(U')$에 대응한다. 층
$j'_!\mathcal{O}_{U'}$, $j_!\mathcal{O}_U$의 줄기에 대한 설명으로부터
위에 표시한 사상이 단사임을 알 수 있다(위에서 인용한 보조정리 참조).
따라서 $\mathcal{I}$가 단사 $\mathcal{O}_X$-가군이면 사상
$$
\Hom_{\mathcal{O}_X}(j_!\mathcal{O}_U, \mathcal{I})
\longrightarrow
\Hom_{\mathcal{O}_X}(j'_!\mathcal{O}_{U'}, \mathcal{I})
$$
은 전사이다. Homology, Lemma
\ref{homology-lemma-characterize-injectives}를 보라. 이들을 모두 합치면
보조정리를 얻는다.
\end{proof}

\begin{lemma}[Mayer-Vietoris]
\label{cohomology-lemma-mayer-vietoris}
$X$를 환 달린 공간이라 하자. $X = U \cup V$가 두 열린 부분집합의
합집합이라고 가정하자. 모든 $\mathcal{O}_X$-가군 $\mathcal{F}$에 대해
코호몰로지 긴 완전열
$$
0 \to
H^0(X, \mathcal{F}) \to
H^0(U, \mathcal{F}) \oplus H^0(V, \mathcal{F}) \to
H^0(U \cap V, \mathcal{F}) \to
H^1(X, \mathcal{F}) \to \ldots
$$
이 존재한다. 이 긴 완전열은 $\mathcal{F}$에 대해 함자적이다.
\end{lemma}

\begin{proof}
층 조건에 따르면 임의의 아벨 층 $\mathcal{F}$에 대해
$(1, -1) : \mathcal{F}(U) \oplus \mathcal{F}(V) \to \mathcal{F}(U \cap V)$
의 핵은 첫 번째 사상에 의한 $\mathcal{F}(X)$의 상과 같다. 위 보조정리
\ref{cohomology-lemma-injective-restriction-surjective}에 따르면 사상
$(1, -1) : \mathcal{I}(U) \oplus \mathcal{I}(V) \to \mathcal{I}(U \cap V)$
은 $\mathcal{I}$가 단사 $\mathcal{O}_X$-가군일 때마다 전사이다.
따라서 $\mathcal{F} \to \mathcal{I}^\bullet$이 $\mathcal{F}$의 단사
분해이면 복합체의 짧은 완전열
$$
0 \to
\mathcal{I}^\bullet(X) \to
\mathcal{I}^\bullet(U) \oplus \mathcal{I}^\bullet(V) \to
\mathcal{I}^\bullet(U \cap V) \to
0.
$$
을 얻는다. 코호몰로지를 취하면 결과를 얻는다(Homology, Lemma
\ref{homology-lemma-long-exact-sequence-cochain}를 사용하라).
이 열의 함자성에 대한 증명은 생략한다.
\end{proof}

\begin{lemma}[상대 Mayer--Vietoris]
\label{cohomology-lemma-relative-mayer-vietoris}
$f : X \to Y$를 환 달린 공간의 사상이라 하자.
$X = U \cup V$가 두 열린 부분집합의 합집합이라고 가정하자.
$a = f|_U : U \to Y$, $b = f|_V : V \to Y$, 그리고
$c = f|_{U \cap V} : U \cap V \to Y$로 쓰자.
모든 $\mathcal{O}_X$-가군 $\mathcal{F}$에 대해 긴 완전열
$$
0 \to
f_*\mathcal{F} \to
a_*(\mathcal{F}|_U) \oplus b_*(\mathcal{F}|_V) \to
c_*(\mathcal{F}|_{U \cap V}) \to
R^1f_*\mathcal{F} \to \ldots
$$
이 존재한다. 이 긴 완전열은 $\mathcal{F}$에 대해 함자적이다.
\end{lemma}

\begin{proof}
$\mathcal{F} \to \mathcal{I}^\bullet$을 $\mathcal{F}$의 단사 분해라 하자.
복합체의 짧은 완전열
$$
0 \to
f_*\mathcal{I}^\bullet \to
a_*\mathcal{I}^\bullet|_U \oplus b_*\mathcal{I}^\bullet|_V \to
c_*\mathcal{I}^\bullet|_{U \cap V} \to
0.
$$
을 얻는다고 주장한다. 실제로 임의의 열린집합 $W \subset Y$와 임의의
$n \geq 0$에 대해 이에 대응하는 $W$ 위의 절단군의 열
$$
0 \to
\mathcal{I}^n(f^{-1}(W)) \to
\mathcal{I}^n(U \cap f^{-1}(W))
\oplus \mathcal{I}^n(V \cap f^{-1}(W)) \to
\mathcal{I}^n(U \cap V \cap f^{-1}(W)) \to
0
$$
이 짧은 완전열임은 보조정리 \ref{cohomology-lemma-mayer-vietoris}의 증명에서 보였다.
코호몰로지층을 취하고 $\mathcal{I}^\bullet|_U$가
$\mathcal{F}|_U$의 단사 분해이며 $\mathcal{I}^\bullet|_V$와
$\mathcal{I}^\bullet|_{U \cap V}$에 대해서도 마찬가지라는 사실을 사용하면
보조정리가 따른다. 보조정리 \ref{cohomology-lemma-cohomology-of-open}을 보라.
\end{proof}




















\section{체흐 복합체와 체흐 코호몰로지}
\label{cohomology-section-cech}

\noindent
$X$를 위상공간이라 하자.
$\mathcal{U} : U = \bigcup_{i \in I} U_i$를 열린 덮개라 하자. Topology,
Basic notion (\ref{topology-item-covering})을 보라. 관례에 따라
$\mathcal{U}$의 원소들의 $(p + 1)$중 교집합을
$U_{i_0\ldots i_p} = U_{i_0} \cap \ldots \cap U_{i_p}$로 나타낸다.
$\mathcal{F}$를 $X$ 위의 아벨 준층이라 하자. 다음과 같이 놓는다.
$$
\check{\mathcal{C}}^p(\mathcal{U}, \mathcal{F})
=
\prod\nolimits_{(i_0, \ldots, i_p) \in I^{p + 1}}
\mathcal{F}(U_{i_0\ldots i_p}).
$$
이는 아벨군이다. $s \in \check{\mathcal{C}}^p(\mathcal{U}, \mathcal{F})$에
대해 $\mathcal{F}(U_{i_0\ldots i_p})$에서의 그 값을
$s_{i_0\ldots i_p}$로 나타낸다. $s \in
\check{\mathcal{C}}^1(\mathcal{U}, \mathcal{F})$이고 $i, j \in I$이면
$s_{ij}$와 $s_{ji}$는 모두 $\mathcal{F}(U_i \cap U_j)$의 원소이지만,
$s_{ij}$와 $s_{ji}$ 사이에는 아무 관계도 부과하지 않음에 유의하자.
다시 말해 우리는 교대 코체인을 다루는 것이 {\it 아니다}(이는
Section \ref{cohomology-section-alternating-cech}에서 정의한다). 다음을 정의한다.
$$
d : \check{\mathcal{C}}^p(\mathcal{U}, \mathcal{F})
\longrightarrow
\check{\mathcal{C}}^{p + 1}(\mathcal{U}, \mathcal{F})
$$
그 공식은 다음과 같다.
\begin{equation}
\label{cohomology-equation-d-cech}
d(s)_{i_0\ldots i_{p + 1}}
=
\sum\nolimits_{j = 0}^{p + 1}
(-1)^j
s_{i_0\ldots \hat i_j \ldots i_{p + 1}}|_{U_{i_0\ldots i_{p + 1}}}
\end{equation}
$d \circ d = 0$임은 곧바로 알 수 있다. 다시 말해
$\check{\mathcal{C}}^\bullet(\mathcal{U}, \mathcal{F})$는 복합체이다.

\begin{definition}
\label{cohomology-definition-cech-complex}
$X$를 위상공간, $\mathcal{U} : U = \bigcup_{i \in I} U_i$를 열린 덮개,
$\mathcal{F}$를 $X$ 위의 아벨 준층이라 하자. 복합체
$\check{\mathcal{C}}^\bullet(\mathcal{U}, \mathcal{F})$를
$\mathcal{F}$와 열린 덮개 $\mathcal{U}$에 결부된 {\it 체흐 복합체}라고 한다.
그 코호몰로지군
$H^i(\check{\mathcal{C}}^\bullet(\mathcal{U}, \mathcal{F}))$를
$\mathcal{F}$와 덮개 $\mathcal{U}$에 결부된 {\it 체흐 코호몰로지군}이라고
하며 $\check H^i(\mathcal{U}, \mathcal{F})$로 나타낸다.
\end{definition}

\begin{lemma}
\label{cohomology-lemma-cech-h0}
$X$를 위상공간이라 하고, $\mathcal{F}$를 $X$ 위의 아벨 준층이라 하자.
다음은 서로 동치이다.
\begin{enumerate}
\item $\mathcal{F}$는 아벨 층이다.
\item 모든 열린 덮개 $\mathcal{U} : U = \bigcup_{i \in I} U_i$에 대해
자연스러운 사상
$$
\mathcal{F}(U) \to \check{H}^0(\mathcal{U}, \mathcal{F})
$$
은 전단사이다.
\end{enumerate}
\end{lemma}

\begin{proof}
층 조건은 정확히 모든 열린 덮개에 대해
$\mathcal{F}(U) \to \check{H}^0(\mathcal{U}, \mathcal{F})$가
전단사라는 조건이므로 성립한다.
\end{proof}

\begin{lemma}
\label{cohomology-lemma-cech-trivial}
$X$를 위상공간이라 하고, $\mathcal{F}$를 $X$ 위의 아벨 준층이라 하자.
$\mathcal{U} : U = \bigcup_{i \in I} U_i$를 열린 덮개라 하자. 어떤
$i \in I$에 대해 $U_i = U$이면, $\mathcal{F}(U)$를 차수 $-1$에 두고
$\mathcal{F}(U)$에서 $\check{\mathcal{C}}^0(\mathcal{U}, \mathcal{F})$로
가는 표준 사상을 미분으로 삼아 얻는 확대 체흐 복합체
$$
\mathcal{F}(U) \to \check{\mathcal{C}}^\bullet(\mathcal{U}, \mathcal{F})
$$
는 $0$과 호모토피 동치이다.
\end{lemma}

\begin{proof}
$U = U_i$인 원소 $i \in I$를 고정하자. $i_j = i$이면
$U_{i_0 \ldots i_p} = U_{i_0 \ldots \hat i_j \ldots i_p}$임에 유의하자.
호모토피
$$
h :
\prod\nolimits_{i_0 \ldots i_{p + 1}} \mathcal{F}(U_{i_0 \ldots i_{p + 1}})
\longrightarrow
\prod\nolimits_{i_0 \ldots i_p} \mathcal{F}(U_{i_0 \ldots i_p})
$$
를 다음 규칙으로 정의한다.
$$
h(s)_{i_0 \ldots i_p} = s_{i i_0 \ldots i_p}
$$
다시 말해 $h : \prod_{i_0} \mathcal{F}(U_{i_0}) \to \mathcal{F}(U)$는
인자 $\mathcal{F}(U_i) = \mathcal{F}(U)$로의 사영이고, 일반적으로
사상 $h$는 인자
$\mathcal{F}(U_{i i_1 \ldots i_{p + 1}}) =
\mathcal{F}(U_{i_1 \ldots i_{p + 1}})$들로의 사영과 같다. 다음을 계산한다.
\begin{align*}
(dh + hd)(s)_{i_0 \ldots i_p}
& =
\sum\nolimits_{j = 0}^p
(-1)^j
h(s)_{i_0 \ldots \hat i_j \ldots i_p}
+
d(s)_{i i_0 \ldots i_p}\\
& =
\sum\nolimits_{j = 0}^p
(-1)^j
s_{i i_0 \ldots \hat i_j \ldots i_p}
+
s_{i_0 \ldots i_p}
+
\sum\nolimits_{j = 0}^p
(-1)^{j + 1}
s_{i i_0 \ldots \hat i_j \ldots i_p} \\
& =
s_{i_0 \ldots i_p}
\end{align*}
이로써 항등 사상이 영 사상과 호모토픽임이 증명되고, 원하는 결론을 얻는다.
\end{proof}






\section{준층 위의 함자로서의 체흐 코호몰로지}
\label{cohomology-section-cech-functor}

\noindent
경고: 이 절에서는 거의 전적으로 준층과 준층의 범주를 다룬다. 여기의
결과들은 층과 층의 범주라는 설정에서는 완전히 틀린다!

\medskip\noindent
$X$를 환 달린 공간, $\mathcal{U} : U = \bigcup_{i \in I} U_i$를 열린
덮개라 하자. $\mathcal{F}$를 $\mathcal{O}_X$-가군의 준층이라 하자.
$\mathcal{F}$를 아벨군의 준층으로 보기만 하면 $\mathcal{F}$의 체흐 복합체
$\check{\mathcal{C}}^\bullet(\mathcal{U}, \mathcal{F})$
를 얻는다. 그러나 각 항 $\check{\mathcal{C}}^p(\mathcal{U}, \mathcal{F})$에는
자연스러운 $\mathcal{O}_X(U)$-가군 구조가 있고, 미분은
$\mathcal{O}_X(U)$-가군 사상들로 주어진다. 또한 구성
$$
\mathcal{F} \longmapsto \check{\mathcal{C}}^\bullet(\mathcal{U}, \mathcal{F})
$$
은 $\mathcal{F}$에 대해 함자적임이 분명하다. 실제로 이는 함자
\begin{equation}
\label{cohomology-equation-cech-functor}
\check{\mathcal{C}}^\bullet(\mathcal{U}, -) :
\textit{PMod}(\mathcal{O}_X)
\longrightarrow
\text{Comp}^{+}(\text{Mod}_{\mathcal{O}_X(U)})
\end{equation}
이다. 표기는 Derived Categories, Definition
\ref{derived-definition-complexes-notation}을 보라. 아벨 범주에서 아래로
유계인 복합체들의 범주는 아벨 범주임을 상기하라. Homology, Lemma
\ref{homology-lemma-cat-cochain-abelian}을 보라.

\begin{lemma}
\label{cohomology-lemma-cech-exact-presheaves}
Equation (\ref{cohomology-equation-cech-functor})으로 주어진 함자는 완전함자이다
(Homology, Lemma \ref{homology-lemma-exact-functor} 참조).
\end{lemma}

\begin{proof}
임의의 열린집합 $W \subset U$에 대해 함자
$\mathcal{F} \mapsto \mathcal{F}(W)$는 $\textit{PMod}(\mathcal{O}_X)$에서
$\text{Mod}_{\mathcal{O}_X(U)}$로 가는 가법 완전함자이다. 복합체의 항
$\check{\mathcal{C}}^p(\mathcal{U}, \mathcal{F})$는 이러한 완전함자들의
곱이므로 완전하다. 또한 복합체의 열이 완전일 필요충분조건은 각 차수에서
항들의 열이 완전인 것이다. 따라서 보조정리가 따른다.
\end{proof}

\begin{lemma}
\label{cohomology-lemma-cech-cohomology-delta-functor-presheaves}
$X$를 환 달린 공간, $\mathcal{U} : U = \bigcup_{i \in I} U_i$를 열린
덮개라 하자. 함자들
$\mathcal{F} \mapsto \check{H}^n(\mathcal{U}, \mathcal{F})$는
$\mathcal{O}_X$-가군의 준층들로 이루어진 아벨 범주에서
$\mathcal{O}_X(U)$-가군의 범주로 가는 $\delta$-함자를 이룬다(Homology,
Definition \ref{homology-definition-cohomological-delta-functor} 참조).
\end{lemma}

\begin{proof}
Lemma \ref{cohomology-lemma-cech-exact-presheaves}에 의해 $\mathcal{O}_X$-가군의
준층들의 짧은 완전열
$0 \to \mathcal{F}_1 \to \mathcal{F}_2 \to \mathcal{F}_3 \to 0$
은 $\mathcal{O}_X(U)$-가군 복합체들의 짧은 완전열로 보내진다. 따라서
Homology, Lemma \ref{homology-lemma-long-exact-sequence-cochain}을 사용하여
경계 사상
$\delta_{\mathcal{F}_1 \to \mathcal{F}_2 \to \mathcal{F}_3} :
\check{H}^n(\mathcal{U}, \mathcal{F}_3) \to
\check{H}^{n + 1}(\mathcal{U}, \mathcal{F}_1)$
과 이에 대응하는 긴 완전열을 얻는다. 이 사상들이 준층들의 짧은 완전열
사이의 사상들과 양립한다는 확인은 생략한다.
\end{proof}


\noindent
다음 보조정리의 서술에서는 열린 몰입 $j : U \to X$에 상대적인 가군
준층의 $0$에 의한 연장 함자 $j_{p!}$를 사용한다. Sheaves, Section
\ref{sheaves-section-open-immersions}을 보라. 임의의 열린집합
$W \subset X$와 임의의 $\mathcal{O}_X|_U$-가군 준층 $\mathcal{G}$에 대해
$$
(j_{p!}\mathcal{G})(W) =
\left\{
\begin{matrix}
\mathcal{G}(W) & \text{if } W \subset U \\
0 & \text{else.}
\end{matrix}
\right.
$$
이다. 또한 함자 $j_{p!}$는 제한 함자의 왼쪽 수반이다. Sheaves, Lemma
\ref{sheaves-lemma-j-shriek-modules}를 보라. 특히 다음 공식을 얻는다.
$$
\Hom_{\mathcal{O}_X}(j_{p!}\mathcal{O}_U, \mathcal{F})
=
\Hom_{\mathcal{O}_U}(\mathcal{O}_U, \mathcal{F}|_U)
=
\mathcal{F}(U).
$$
함자 $\mathcal{F} \mapsto \mathcal{F}(U)$는 준층의 범주에서 완전함자이므로,
준층 $j_{p!}\mathcal{O}_U$가 범주 $\textit{PMod}(\mathcal{O}_X)$의 사영
대상임을 알 수 있다. Homology, Lemma
\ref{homology-lemma-characterize-projectives}를 보라.

\medskip\noindent
결부된 열린 몰입이 $j_U, j_V$인 열린 부분집합 $U \subset V \subset X$가
주어지면 표준 사상 $(j_U)_{p!}\mathcal{O}_U \to
(j_V)_{p!}\mathcal{O}_V$가 있다. 이는 임의의 열린집합 $W \subset U$ 위의
절단에서는 항등이고 그 밖에서는 $0$이다. 식별
$\Hom_{\mathcal{O}_X}((j_U)_{p!}\mathcal{O}_U, (j_V)_{p!}\mathcal{O}_V) =
(j_V)_{p!}\mathcal{O}_V(U) = \mathcal{O}_V(U)$ 아래에서 이는 원소
$1 \in \mathcal{O}_V(U)$에 대응함에 유의하자.

\begin{lemma}
\label{cohomology-lemma-cech-map-into}
$X$를 환 달린 공간, $\mathcal{U} : U = \bigcup_{i \in I} U_i$를 열린
덮개라 하자. 열린 몰입을
$j_{i_0\ldots i_p} : U_{i_0 \ldots i_p} \to X$로 나타내자.
$\mathcal{O}_X$-가군의 준층들로 이루어진 사슬 복합체
$K(\mathcal{U})_\bullet$
$$
\ldots
\to
\bigoplus_{i_0i_1i_2} (j_{i_0i_1i_2})_{p!}\mathcal{O}_{U_{i_0i_1i_2}}
\to
\bigoplus_{i_0i_1} (j_{i_0i_1})_{p!}\mathcal{O}_{U_{i_0i_1}}
\to
\bigoplus_{i_0} (j_{i_0})_{p!}\mathcal{O}_{U_{i_0}}
\to 0 \to \ldots
$$
를 생각하자. 여기서 마지막 영이 아닌 항은 차수 $0$에 놓이고, 사상
$$
(j_{i_0\ldots i_{p + 1}})_{p!}\mathcal{O}_{U_{i_0\ldots i_{p + 1}}}
\longrightarrow
(j_{i_0\ldots \hat i_j \ldots i_{p + 1}})_{p!}
\mathcal{O}_{U_{i_0\ldots \hat i_j \ldots i_{p + 1}}}
$$
은 표준 사상의 $(-1)^j$배로 주어진다. 그러면 동형
$$
\Hom_{\mathcal{O}_X}(K(\mathcal{U})_\bullet, \mathcal{F})
=
\check{\mathcal{C}}^\bullet(\mathcal{U}, \mathcal{F})
$$
은 $\mathcal{F} \in \Ob(\textit{PMod}(\mathcal{O}_X))$에 대해 함자적이다.
\end{lemma}

\begin{proof}
보조정리 바로 위의 논의에서 다음을 보았다.
$$
\Hom_{\mathcal{O}_X}(
(j_{i_0\ldots i_p})_{p!}\mathcal{O}_{U_{i_0\ldots i_p}},
\mathcal{F})
=
\mathcal{F}(U_{i_0\ldots i_p}).
$$
따라서 실제로 직합
$$
\bigoplus\nolimits_{i_0 \ldots i_p}
(j_{i_0 \ldots i_p})_{p!}\mathcal{O}_{U_{i_0 \ldots i_p}}
$$
은 함자
$$
\mathcal{F}
\longmapsto
\prod\nolimits_{i_0\ldots i_p} \mathcal{F}(U_{i_0\ldots i_p}).
$$
를 표현함을 알 수 있다. 따라서 Categories, Yoneda Lemma
\ref{categories-lemma-yoneda}에 의해 주어진 항들을 갖는 복합체
$K(\mathcal{U})_\bullet$가 있음을 알 수 있다. 사상들이 보조정리에 주어진
것과 같다는 것은 쉽게 확인된다.
\end{proof}

\begin{lemma}
\label{cohomology-lemma-homology-complex}
$X$를 환 달린 공간, $\mathcal{U} : U = \bigcup_{i \in I} U_i$를 열린
덮개라 하자. $\mathcal{O}_\mathcal{U} \subset \mathcal{O}_X$를 사상
$\bigoplus j_{p!}\mathcal{O}_{U_i} \to \mathcal{O}_X$의 상 준층이라 하자.
위 보조정리 \ref{cohomology-lemma-cech-map-into}의 준층 사슬 복합체
$K(\mathcal{U})_\bullet$의 호몰로지 준층은 다음과 같다.
$$
H_i(K(\mathcal{U})_\bullet) =
\left\{
\begin{matrix}
0 & \text{if} & i \not = 0 \\
\mathcal{O}_\mathcal{U} & \text{if} & i = 0
\end{matrix}
\right.
$$
\end{lemma}

\begin{proof}
자명한 사상
$K(\mathcal{U})_0 =
\bigoplus_{i_0} (j_{i_0})_{p!}\mathcal{O}_{U_{i_0}} \to
\mathcal{O}_\mathcal{U}$와 함께 $\mathcal{O}_\mathcal{U}$를 차수 $-1$에
놓아 얻는 확대 복합체 $K^{ext}_\bullet$를 생각하자. 임의의 열린집합
$W \subset X$ 위에서 이 확대 복합체의 절단을 취한 것이 비순환 복합체임을
보이면 충분하다. 실제로 모든 $W \subset X$에 대해 복합체
$K^{ext}_\bullet(W)$가 영 복합체와 호모토피 동치라고 주장한다.
$W \subset U_i$일 필요충분조건이 $i \in I_1$이 되도록
$I = I_1 \amalg I_2$로 쓰자.

\medskip\noindent
$I_1 = \emptyset$이면 복합체 $K^{ext}_\bullet(W) = 0$이므로 증명할 것이 없다.

\medskip\noindent
$I_1 \not = \emptyset$이면
$\mathcal{O}_\mathcal{U}(W) = \mathcal{O}_X(W)$
이고
$$
K^{ext}_p(W) =
\bigoplus\nolimits_{i_0 \ldots i_p \in I_1} \mathcal{O}_X(W).
$$
이다. 이는 준층
$(j_{i_0 \ldots i_p})_{p!}\mathcal{O}_{U_{i_0 \ldots i_p}}$에 대한 간단한
기술에서 따른다. 또한 복합체 $K^{ext}_\bullet(W)$의 미분은 다음과 같다.
$$
d(s)_{i_0 \ldots i_p} =
\sum\nolimits_{j = 0, \ldots, p + 1} \sum\nolimits_{i \in I_1}
(-1)^j s_{i_0 \ldots i_{j - 1} i i_j \ldots i_p}.
$$
원소 $s$의 지지가 유한하므로 이 합은 유한하다. 원소
$i_{\text{fix}} \in I_1$을 고정하자. 사상
$$
h : K^{ext}_p(W) \longrightarrow K^{ext}_{p + 1}(W)
$$
를 다음 규칙으로 정의한다.
$$
h(s)_{i_0 \ldots i_{p + 1}} =
\left\{
\begin{matrix}
0 & \text{if} & i_0 \not = i_{\text{fix}} \\
s_{i_1 \ldots i_{p + 1}} & \text{if} & i_0 = i_{\text{fix}}
\end{matrix}
\right.
$$
이를 줄여서
$h(s)_{i_0 \ldots i_{p + 1}} = (i_0 = i_{\text{fix}}) s_{i_1 \ldots i_p}$
로 쓸 것이다. 그러면 다음을 계산한다.
\begin{eqnarray*}
& & (dh + hd)(s)_{i_0 \ldots i_p} \\
& = &
\sum_j \sum_{i \in I_1} (-1)^j h(s)_{i_0 \ldots i_{j - 1} i i_j \ldots i_p}
+
(i = i_0) d(s)_{i_1 \ldots i_p} \\
& = &
s_{i_0 \ldots i_p} +
\sum_{j \geq 1}\sum_{i \in I_1}
(-1)^j (i_0 = i_{\text{fix}}) s_{i_1 \ldots i_{j - 1} i i_j \ldots i_p}
+
(i_0 = i_{\text{fix}}) d(s)_{i_1 \ldots i_p}
\end{eqnarray*}
이는 원하는 대로 $s_{i_0 \ldots i_p}$와 같다.
\end{proof}

\begin{lemma}
\label{cohomology-lemma-cech-cohomology-derived-presheaves}
$X$를 환 달린 공간이라 하고, $\mathcal{U} : U = \bigcup_{i \in I} U_i$를
$U \subset X$의 열린 덮개라 하자. 체흐 코호몰로지 함자
$\check{H}^p(\mathcal{U}, -)$들은 하나의 $\delta$-함자로서 다음 함자의
오른쪽 유도함자들과 표준적으로 동형이다.
$$
\check{H}^0(\mathcal{U}, -) :
\textit{PMod}(\mathcal{O}_X)
\longrightarrow
\text{Mod}_{\mathcal{O}_X(U)}.
$$
또한 함자적 유사동형
$$
\check{\mathcal{C}}^\bullet(\mathcal{U}, \mathcal{F})
\longrightarrow
R\check{H}^0(\mathcal{U}, \mathcal{F})
$$
이 있다. 여기서 오른쪽은 왼쪽 완전함자
$\check{H}^0(\mathcal{U}, -)$의 오른쪽 유도함자
$$
R\check{H}^0(\mathcal{U}, -) :
D^{+}(\textit{PMod}(\mathcal{O}_X))
\longrightarrow
D^{+}(\mathcal{O}_X(U))
$$
를 나타낸다.
\end{lemma}

\begin{proof}
$\mathcal{O}_X$-가군 준층의 범주에는 단사가 충분히 있다. Injectives,
Proposition \ref{injectives-proposition-presheaves-modules}를 보라. 또한
$\check{H}^0(\mathcal{U}, -)$는 $\mathcal{O}_X$-가군 준층의 범주에서
$\mathcal{O}_X(U)$-가군의 범주로 가는 왼쪽 완전함자이다. 따라서 유도함자와
오른쪽 유도함자가 존재한다. Derived Categories, Section
\ref{derived-section-right-derived-functor}을 보라.

\medskip\noindent
$\mathcal{I}$를 단사 $\mathcal{O}_X$-가군 준층이라 하자. 이 경우 함자
$\Hom_{\mathcal{O}_X}(-, \mathcal{I})$는
$\textit{PMod}(\mathcal{O}_X)$에서 완전하다. Lemma
\ref{cohomology-lemma-cech-map-into}에 의해
$$
\Hom_{\mathcal{O}_X}(K(\mathcal{U})_\bullet, \mathcal{I})
=
\check{\mathcal{C}}^\bullet(\mathcal{U}, \mathcal{I}).
$$
이다. Lemma \ref{cohomology-lemma-homology-complex}에 의해 $K(\mathcal{U})_\bullet$는
$\mathcal{O}_\mathcal{U}[0]$과 유사동형이다. 따라서 위에서 언급한
$\mathcal{I}$로 들어가는 Hom의 완전성으로부터 모든 $i > 0$에 대해
$\check{H}^i(\mathcal{U}, \mathcal{I}) = 0$임을 알 수 있다. 그러므로
$\delta$-함자 $(\check{H}^n, \delta)$(Lemma
\ref{cohomology-lemma-cech-cohomology-delta-functor-presheaves} 참조)는 Homology,
Lemma \ref{homology-lemma-efface-implies-universal}의 가정을 만족하며,
따라서 보편 $\delta$-함자이다.

\medskip\noindent
Derived Categories, Lemma \ref{derived-lemma-higher-derived-functors}에 의해
열 $R^i\check{H}^0(\mathcal{U}, -)$도 보편 $\delta$-함자를 이룬다. 보편
$\delta$-함자의 유일성(Homology, Lemma
\ref{homology-lemma-uniqueness-universal-delta-functor} 참조)에 의해
$R^i\check{H}^0(\mathcal{U}, -) = \check{H}^i(\mathcal{U}, -)$를 얻는다.
이는 대부분의 응용에 충분하므로 독자는 증명의 나머지를 건너뛰어도 좋다.

\medskip\noindent
$\mathcal{F}$를 임의의 $\mathcal{O}_X$-가군 준층이라 하자. 범주
$\textit{PMod}(\mathcal{O}_X)$에서 단사 분해
$\mathcal{F} \to \mathcal{I}^\bullet$을 택하자. 항들이
$\check{\mathcal{C}}^p(\mathcal{U}, \mathcal{I}^q)$인 이중 복합체
$\check{\mathcal{C}}^\bullet(\mathcal{U}, \mathcal{I}^\bullet)$와 결부된
전체 복합체
$\text{Tot}(\check{\mathcal{C}}^\bullet(\mathcal{U}, \mathcal{I}^\bullet))$를
생각하자. Homology, Definition
\ref{homology-definition-associated-simple-complex}을 보라. 사상들
$\check{\mathcal{C}}^p(\mathcal{U}, \mathcal{F})
\to \check{\mathcal{C}}^p(\mathcal{U}, \mathcal{I}^0)$에서 오는 복합체의 사상
$$
\check{\mathcal{C}}^\bullet(\mathcal{U}, \mathcal{F})
\longrightarrow
\text{Tot}(\check{\mathcal{C}}^\bullet(\mathcal{U}, \mathcal{I}^\bullet))
$$
이 있고, 사상들
$\check{H}^0(\mathcal{U}, \mathcal{I}^q) \to
\check{\mathcal{C}}^0(\mathcal{U}, \mathcal{I}^q)$에서 오는 복합체의 사상
$$
\check{H}^0(\mathcal{U}, \mathcal{I}^\bullet)
\longrightarrow
\text{Tot}(\check{\mathcal{C}}^\bullet(\mathcal{U}, \mathcal{I}^\bullet))
$$
도 있다. Homology, Lemma
\ref{homology-lemma-double-complex-gives-resolution}을 적용하면 이 두 사상은
모두 유사동형이다. 실제로 함자로서의 체흐 복합체가 완전하므로(Lemma
\ref{cohomology-lemma-cech-exact-presheaves}) 이중 복합체의 열들은 양의 차수에서
완전하고, 방금 본 바와 같이 단사 준층 $\mathcal{I}^q$의 고차 체흐
코호몰로지군들이 영이므로 이중 복합체의 행들도 양의 차수에서 완전하다.
유사동형은 $D^{+}(\mathcal{O}_X(U))$에서 가역이 되므로 보조정리의 마지막에
표시한 사상을 얻는다. 이 사상이 함자적이라는 확인은 생략한다.
\end{proof}





\section{체흐 코호몰로지와 코호몰로지}
\label{cohomology-section-cech-cohomology-cohomology}

\begin{lemma}
\label{cohomology-lemma-injective-trivial-cech}
$X$를 환 달린 공간, $\mathcal{U} : U = \bigcup_{i \in I} U_i$를 열린
덮개라 하자. $\mathcal{I}$를 단사 $\mathcal{O}_X$-가군이라 하자. 그러면
$$
\check{H}^p(\mathcal{U}, \mathcal{I}) =
\left\{
\begin{matrix}
\mathcal{I}(U) & \text{if} & p = 0 \\
0 & \text{if} & p > 0
\end{matrix}
\right.
$$
\end{lemma}

\begin{proof}
단사 $\mathcal{O}_X$-가군은 범주 $\textit{PMod}(\mathcal{O}_X)$의
대상으로서도 단사이다(예를 들어 층화가 포함 함자의 완전한 왼쪽 수반이기
때문이다. Homology, Lemma
\ref{homology-lemma-adjoint-preserve-injectives}를 사용하라). 따라서 Lemma
\ref{cohomology-lemma-cech-cohomology-derived-presheaves}(또는 그 증명)를 적용하면
결과를 알 수 있다.
\end{proof}

\begin{lemma}
\label{cohomology-lemma-cech-cohomology}
$X$를 환 달린 공간, $\mathcal{U} : U = \bigcup_{i \in I} U_i$를 열린
덮개라 하자. 함자들 $\textit{Mod}(\mathcal{O}_X) \to
D^{+}(\mathcal{O}_X(U))$ 사이의 변환
$$
\check{\mathcal{C}}^\bullet(\mathcal{U}, -)
\longrightarrow
R\Gamma(U, -)
$$
이 있다. 특히 $\mathcal{F}$가 $\textit{Mod}(\mathcal{O}_X)$를 달릴 때
이는 표준 사상
$\check{H}^p(\mathcal{U}, \mathcal{F}) \to H^p(U, \mathcal{F})$를 준다.
\end{lemma}

\begin{proof}
$\mathcal{F}$를 $\mathcal{O}_X$-가군이라 하자. 단사 분해
$\mathcal{F} \to \mathcal{I}^\bullet$을 택하자. 항들이
$\check{\mathcal{C}}^p(\mathcal{U}, \mathcal{I}^q)$인 이중 복합체
$\check{\mathcal{C}}^\bullet(\mathcal{U}, \mathcal{I}^\bullet)$를 생각하자.
사상들 $\mathcal{I}^q(U) \to
\check{H}^0(\mathcal{U}, \mathcal{I}^q)$에서 오는 복합체의 사상
$$
\alpha :
\Gamma(U, \mathcal{I}^\bullet)
\longrightarrow
\text{Tot}(\check{\mathcal{C}}^\bullet(\mathcal{U}, \mathcal{I}^\bullet))
$$
과 사상 $\mathcal{F} \to \mathcal{I}^0$에서 오는 복합체의 사상
$$
\beta :
\check{\mathcal{C}}^\bullet(\mathcal{U}, \mathcal{F})
\longrightarrow
\text{Tot}(\check{\mathcal{C}}^\bullet(\mathcal{U}, \mathcal{I}^\bullet))
$$
이 있다. Homology, Lemma
\ref{homology-lemma-double-complex-gives-resolution}을 적용하면 $\alpha$가
유사동형임을 알 수 있다. 실제로 Lemma
\ref{cohomology-lemma-injective-trivial-cech}에 의해 이중 복합체
$\check{\mathcal{C}}^\bullet(\mathcal{U}, \mathcal{I}^\bullet)$의
$q$번째 행은 $\Gamma(U, \mathcal{I}^q)$의 분해이다. 따라서 $\alpha$는
$D^{+}(\mathcal{O}_X(U))$에서 가역이 되고, 보조정리의 변환은 $\beta$에
$\alpha$의 역을 뒤이어 합성한 것이다. 이것이 함자적이라는 확인은 생략한다.
\end{proof}

\begin{lemma}
\label{cohomology-lemma-cech-h1}
$X$를 위상공간, $\mathcal{H}$를 $X$ 위의 아벨 층,
$\mathcal{U} : X = \bigcup_{i \in I} U_i$를 열린 덮개라 하자. 사상
$$
\check{H}^1(\mathcal{U}, \mathcal{H}) \longrightarrow H^1(X, \mathcal{H})
$$
은 단사이며, Lemma \ref{cohomology-lemma-torsors-h1}의 전단사를 통해
$\check{H}^1(\mathcal{U}, \mathcal{H})$를 각 $U_i$ 위로 제한했을 때
자명한 토서가 되는 $\mathcal{H}$-토서들의 동형류 집합과 식별한다.
\end{lemma}

\begin{proof}
이를 보이기 위해 역 사상을 구성한다. 즉, $\mathcal{F}$를 각 $U_i$로의
제한이 자명한 $\mathcal{H}$-토서라 하자. Lemma
\ref{cohomology-lemma-trivial-torsor}에 의해 이는 절단 $s_i \in \mathcal{F}(U_i)$가
존재한다는 뜻이다. $U_{i_0} \cap U_{i_1}$ 위에는 다음을 만족하는
$\mathcal{H}$의 유일한 절단 $s_{i_0i_1}$이 있다.
$s_{i_0i_1} \cdot s_{i_0}|_{U_{i_0} \cap U_{i_1}} =
s_{i_1}|_{U_{i_0} \cap U_{i_1}}$. 계산하면 $s_{i_0i_1}$이 체흐
코사이클이고 그 류가 잘 정의됨을 알 수 있다(즉, 절단 $s_i$의 선택에
의존하지 않는다). 역 사상은 $\mathcal{F}$의 동형류를 코사이클
$(s_{i_0i_1})$의 코호몰로지류로 보낸다. 이 사상이 실제로 역이라는 확인은
생략한다.
\end{proof}

\begin{lemma}
\label{cohomology-lemma-include}
$X$를 환 달린 공간이라 하자. 함자
$i : \textit{Mod}(\mathcal{O}_X) \to \textit{PMod}(\mathcal{O}_X)$를
생각하자. 이는 왼쪽 완전함자이고 그 오른쪽 유도함자는 다음과 같다.
$$
R^pi(\mathcal{F}) = \underline{H}^p(\mathcal{F}) :
U \longmapsto H^p(U, \mathcal{F})
$$
Section \ref{cohomology-section-locality}의 논의를 보라.
\end{lemma}

\begin{proof}
$i$가 왼쪽 완전임은 분명하다. 단사 분해
$\mathcal{F} \to \mathcal{I}^\bullet$을 택하자. 정의에 따라 $R^pi$는
복합체 $\mathcal{I}^\bullet$의 $p$번째 코호몰로지 {\it 준층}이다.
다시 말해 열린집합 $U$ 위에서 $R^pi(\mathcal{F})$의 절단은 다음과 같다.
$$
\frac{\Ker(\mathcal{I}^p(U) \to \mathcal{I}^{p + 1}(U))}
{\Im(\mathcal{I}^{p - 1}(U) \to \mathcal{I}^p(U))}.
$$
이는 $H^p(U, \mathcal{F})$의 정의이다.
\end{proof}

\begin{lemma}
\label{cohomology-lemma-cech-spectral-sequence}
$X$를 환 달린 공간, $\mathcal{U} : U = \bigcup_{i \in I} U_i$를 열린
덮개라 하자. 임의의 $\mathcal{O}_X$-가군층 $\mathcal{F}$에 대해
$$
E_2^{p, q} = \check{H}^p(\mathcal{U}, \underline{H}^q(\mathcal{F}))
$$
를 가지며 $H^{p + q}(U, \mathcal{F})$로 수렴하는 스펙트럼열
$(E_r, d_r)_{r \geq 0}$이 있다. 이 스펙트럼열은 $\mathcal{F}$에 대해
함자적이다.
\end{lemma}

\begin{proof}
이는 함자들
$$
i :  \textit{Mod}(\mathcal{O}_X) \to \textit{PMod}(\mathcal{O}_X)
\quad\text{and}\quad
\check{H}^0(\mathcal{U}, - ) : \textit{PMod}(\mathcal{O}_X)
\to \text{Mod}_{\mathcal{O}_X(U)}.
$$
에 대한 Grothendieck 스펙트럼열이다(Derived Categories, Lemma
\ref{derived-lemma-grothendieck-spectral-sequence} 참조). 실제로 Lemma
\ref{cohomology-lemma-cech-h0}에 의해
$\check{H}^0(\mathcal{U}, i(\mathcal{F})) = \mathcal{F}(U)$이다. Lemma
\ref{cohomology-lemma-injective-trivial-cech}에 의해 $i(\mathcal{I})$는 체흐
비순환이다. 그리고 Lemma
\ref{cohomology-lemma-cech-cohomology-derived-presheaves}에 의해
$\textit{PMod}(\mathcal{O}_X)$ 위의 함자로서
$\check{H}^p(\mathcal{U}, -) = R^p\check{H}^0(\mathcal{U}, -)$이다.
이들을 모두 합치면 보조정리를 얻는다.
\end{proof}

\begin{lemma}
\label{cohomology-lemma-cech-spectral-sequence-application}
$X$를 환 달린 공간, $\mathcal{U} : U = \bigcup_{i \in I} U_i$를 열린
덮개, $\mathcal{F}$를 $\mathcal{O}_X$-가군이라 하자. 모든 $i > 0$,
모든 $p \geq 0$, 모든 $i_0, \ldots, i_p \in I$에 대해
$H^i(U_{i_0 \ldots i_p}, \mathcal{F}) = 0$이라고 가정하자. 그러면
$\mathcal{O}_X(U)$-가군으로서
$\check{H}^p(\mathcal{U}, \mathcal{F}) = H^p(U, \mathcal{F})$이다.
\end{lemma}

\begin{proof}
Lemma \ref{cohomology-lemma-cech-spectral-sequence}의 스펙트럼열을 사용한다. 가정은
$q \not = 0$인 모든 $(p, q)$에 대해 $E_2^{p, q} = 0$이라는 뜻이다.
따라서 스펙트럼열은 $E_2$에서 퇴화하고 결과가 따른다.
\end{proof}

\begin{lemma}
\label{cohomology-lemma-ses-cech-h1}
$X$를 환 달린 공간이라 하자.
$$
0 \to \mathcal{F} \to \mathcal{G} \to \mathcal{H} \to 0
$$
을 $\mathcal{O}_X$-가군들의 짧은 완전열이라 하자. $U \subset X$를 열린
부분집합이라 하자. $\check{H}^1(\mathcal{U}, \mathcal{F}) = 0$을 만족하는
$U$의 열린 덮개 $\mathcal{U}$들의 공종계가 존재하면 사상
$\mathcal{G}(U) \to \mathcal{H}(U)$는 전사이다.
\end{lemma}

\begin{proof}
원소 $s \in \mathcal{H}(U)$를 택하자. (a)
$\check{H}^1(\mathcal{U}, \mathcal{F}) = 0$이고 (b) $s|_{U_i}$가 어떤 절단
$s_i \in \mathcal{G}(U_i)$의 상이 되도록 열린 덮개
$\mathcal{U} : U = \bigcup_{i \in I} U_i$를 택하자. (b)를 만족하는
$\mathcal{U}$는 확실히 찾을 수 있으므로, 보조정리의 가정에 의해 (a)와
(b)를 모두 만족하는 $\mathcal{U}$를 찾을 수 있다. 절단들
$$
s_{i_0i_1} = s_{i_1}|_{U_{i_0i_1}} - s_{i_0}|_{U_{i_0i_1}}.
$$
을 생각하자. $s_i$가 $s$의 올림이므로
$s_{i_0i_1} \in \mathcal{F}(U_{i_0i_1})$임을 알 수 있다.
$\check{H}^1(\mathcal{U}, \mathcal{F})$가 영이므로 다음을 만족하는 절단
$t_i \in \mathcal{F}(U_i)$들을 찾을 수 있다.
$$
s_{i_0i_1} = t_{i_1}|_{U_{i_0i_1}} - t_{i_0}|_{U_{i_0i_1}}.
$$
그러면 절단들 $s_i - t_i$는 분명 층 조건을 만족하며 서로 붙어서 $s$로
가는 $U$ 위의 $\mathcal{G}$의 절단이 된다. 따라서 결론을 얻는다.
\end{proof}

\begin{lemma}
\label{cohomology-lemma-cech-vanish}
\begin{slogan}
아벨 층의 고차 체흐 코호몰로지가 모든 열린 덮개에 대해 영이면,
고차 코호몰로지도 영이다.
\end{slogan}
$X$를 환 달린 공간이라 하자. $\mathcal{F}$를 다음을 만족하는
$\mathcal{O}_X$-가군이라 하자.
$$
\check{H}^p(\mathcal{U}, \mathcal{F}) = 0
$$
이는 $X$의 열린 부분집합에 대한 임의의 열린 덮개
$\mathcal{U} : U = \bigcup_{i \in I} U_i$와 모든 $p > 0$에 대해 성립한다고
하자. 그러면 임의의 열린집합 $U \subset X$와 모든 $p > 0$에 대해
$H^p(U, \mathcal{F}) = 0$이다.
\end{lemma}

\begin{proof}
$\mathcal{F}$를 보조정리의 가정을 만족하는 층이라 하자. 이를
``$\mathcal{F}$의 고차 체흐 코호몰로지는 임의의 열린 덮개에 대해
영이다''라고 말하겠다. $\mathcal{F} \to \mathcal{I}$를 단사
$\mathcal{O}_X$-가군으로 가는 매입으로 택하자. Lemma
\ref{cohomology-lemma-injective-trivial-cech}에 의해 $\mathcal{I}$의 고차 체흐
코호몰로지는 임의의 열린 덮개에 대해 영이다.
$\mathcal{Q} = \mathcal{I}/\mathcal{F}$로 놓으면 짧은 완전열
$$
0 \to \mathcal{F} \to \mathcal{I} \to \mathcal{Q} \to 0.
$$
을 얻는다. Lemma \ref{cohomology-lemma-ses-cech-h1}과 우리의 가정에 의해 이 열은
실제로 준층의 열로서 완전하다! 특히 임의의 열린 덮개 $\mathcal{U}$에 대해
체흐 코호몰로지군의 긴 완전열을 얻는다. 예를 들어 Lemma
\ref{cohomology-lemma-cech-cohomology-delta-functor-presheaves}를 보라. 이로부터
$\mathcal{Q}$도 모든 열린 덮개에 대해 고차 체흐 코호몰로지가 영인
$\mathcal{O}_X$-가군임을 알 수 있다.

\medskip\noindent
이제 임의의 열린집합 $U \subset X$에 대한 코호몰로지 긴 완전열
$$
\xymatrix{
0 \ar[r] &
H^0(U, \mathcal{F}) \ar[r] &
H^0(U, \mathcal{I}) \ar[r] &
H^0(U, \mathcal{Q}) \ar[lld] \\
&
H^1(U, \mathcal{F}) \ar[r] &
H^1(U, \mathcal{I}) \ar[r] &
H^1(U, \mathcal{Q}) \ar[lld] \\
&
\ldots & \ldots & \ldots \\
}
$$
을 살펴보자. $\mathcal{I}$가 단사이므로 $n > 0$에 대해
$H^n(U, \mathcal{I}) = 0$이다(Derived Categories, Lemma
\ref{derived-lemma-higher-derived-functors} 참조). 위에서
$H^0(U, \mathcal{I}) \to H^0(U, \mathcal{Q})$가 전사임을 보았으므로
$H^1(U, \mathcal{F}) = 0$이다. $\mathcal{F}$는 고차 체흐 코호몰로지가
영인 임의의 $\mathcal{O}_X$-가군이었고 $\mathcal{Q}$도 그러한 층이므로
(위 참조) $H^1(U, \mathcal{Q}) = 0$도 얻는다. 그러면 긴 완전열에 의해
$H^2(U, \mathcal{F}) = 0$이 따른다. 같은 과정을 계속한다.
\end{proof}

\begin{lemma}
\label{cohomology-lemma-cech-vanish-basis}
(Lemma \ref{cohomology-lemma-cech-vanish}의 변형.)
$X$를 환 달린 공간, $\mathcal{B}$를 $X$의 위상에 대한 기저,
$\mathcal{F}$를 $\mathcal{O}_X$-가군이라 하자. 다음 성질을 갖는 열린
덮개들의 집합 $\text{Cov}$가 존재한다고 가정하자.
\begin{enumerate}
\item $\mathcal{U} : U = \bigcup_{i \in I} U_i$인 모든
$\mathcal{U} \in \text{Cov}$에 대해 $U, U_i \in \mathcal{B}$이고 모든
$U_{i_0 \ldots i_p} \in \mathcal{B}$이다.
\item 모든 $U \in \mathcal{B}$에 대해 $\text{Cov}$에 나타나는 $U$의 열린
덮개들은 $U$의 열린 덮개들의 공종계를 이룬다.
\item 모든 $\mathcal{U} \in \text{Cov}$와 모든 $p > 0$에 대해
$\check{H}^p(\mathcal{U}, \mathcal{F}) = 0$이다.
\end{enumerate}
그러면 모든 $p > 0$과 모든 $U \in \mathcal{B}$에 대해
$H^p(U, \mathcal{F}) = 0$이다.
\end{lemma}

\begin{proof}
$\mathcal{F}$와 $\text{Cov}$가 보조정리에서와 같다고 하자. 이를
``$\mathcal{F}$의 고차 체흐 코호몰로지는 임의의
$\mathcal{U} \in \text{Cov}$에 대해 영이다''라고 말하겠다.
$\mathcal{F} \to \mathcal{I}$를 단사 $\mathcal{O}_X$-가군으로 가는
매입으로 택하자. Lemma \ref{cohomology-lemma-injective-trivial-cech}에 의해
$\mathcal{I}$의 고차 체흐 코호몰로지는 임의의
$\mathcal{U} \in \text{Cov}$에 대해 영이다.
$\mathcal{Q} = \mathcal{I}/\mathcal{F}$로 놓으면 짧은 완전열
$$
0 \to \mathcal{F} \to \mathcal{I} \to \mathcal{Q} \to 0.
$$
을 얻는다. Lemma \ref{cohomology-lemma-ses-cech-h1}과 가정 (2)에 의해 이 열은 모든
$U \in \mathcal{B}$에 대해 완전열
$$
0 \to \mathcal{F}(U) \to \mathcal{I}(U) \to \mathcal{Q}(U) \to 0.
$$
을 낳는다. 따라서 임의의 $\mathcal{U} \in \text{Cov}$에 대해 체흐
복합체들의 짧은 완전열
$$
0 \to
\check{\mathcal{C}}^\bullet(\mathcal{U}, \mathcal{F}) \to
\check{\mathcal{C}}^\bullet(\mathcal{U}, \mathcal{I}) \to
\check{\mathcal{C}}^\bullet(\mathcal{U}, \mathcal{Q}) \to 0
$$
을 얻는다. 가정 (1)에 의해 체흐 복합체의 각 항은 $\mathcal{B}$의 원소들
위에서 취한 값들의 곱으로 이루어지기 때문이다. 특히 임의의 열린 덮개
$\mathcal{U} \in \text{Cov}$에 대해 체흐 코호몰로지군들의 긴 완전열을
얻는다. 이로부터 $\mathcal{Q}$도 모든 $\mathcal{U} \in \text{Cov}$에 대해
고차 체흐 코호몰로지가 영인 $\mathcal{O}_X$-가군임을 알 수 있다.

\medskip\noindent
이제 임의의 $U \in \mathcal{B}$에 대한 코호몰로지 긴 완전열
$$
\xymatrix{
0 \ar[r] &
H^0(U, \mathcal{F}) \ar[r] &
H^0(U, \mathcal{I}) \ar[r] &
H^0(U, \mathcal{Q}) \ar[lld] \\
&
H^1(U, \mathcal{F}) \ar[r] &
H^1(U, \mathcal{I}) \ar[r] &
H^1(U, \mathcal{Q}) \ar[lld] \\
&
\ldots & \ldots & \ldots \\
}
$$
을 살펴보자. $\mathcal{I}$가 단사이므로 $n > 0$에 대해
$H^n(U, \mathcal{I}) = 0$이다(Derived Categories, Lemma
\ref{derived-lemma-higher-derived-functors} 참조). 위에서
$H^0(U, \mathcal{I}) \to H^0(U, \mathcal{Q})$가 전사임을 보았으므로
$H^1(U, \mathcal{F}) = 0$이다. $\mathcal{F}$는 모든
$\mathcal{U} \in \text{Cov}$에 대해 고차 체흐 코호몰로지가 영인 임의의
$\mathcal{O}_X$-가군이었고 $\mathcal{Q}$도 그러한 층이므로(위 참조)
$H^1(U, \mathcal{Q}) = 0$도 얻는다. 그러면 긴 완전열에 의해
$H^2(U, \mathcal{F}) = 0$이 따른다. 같은 과정을 계속한다.
\end{proof}

\begin{lemma}
\label{cohomology-lemma-pushforward-injective}
$f : X \to Y$를 환 달린 공간의 사상이라 하자. $\mathcal{I}$를 단사
$\mathcal{O}_X$-가군이라 하자. 그러면 다음이 성립한다.
\begin{enumerate}
\item $\check{H}^p(\mathcal{V}, f_*\mathcal{I}) = 0$
은 모든 $p > 0$과 $Y$의 임의의 열린 덮개
$\mathcal{V} : V = \bigcup_{j \in J} V_j$에 대해 성립한다.
\item $H^p(V, f_*\mathcal{I}) = 0$은 모든 $p > 0$과 모든 열린집합
$V \subset Y$에 대해 성립한다.
\end{enumerate}
다시 말해 임의의 열린집합 $V \subset Y$에 대해 $f_*\mathcal{I}$는
$\Gamma(V, -)$에 대해 오른쪽 비순환이다(Derived Categories, Definition
\ref{derived-definition-derived-functor} 참조).
\end{lemma}

\begin{proof}
$\mathcal{U} : f^{-1}(V) = \bigcup_{j \in J} f^{-1}(V_j)$로 놓자. 이는
$X$의 열린 덮개이고
$$
\check{\mathcal{C}}^\bullet(\mathcal{V}, f_*\mathcal{I}) =
\check{\mathcal{C}}^\bullet(\mathcal{U}, \mathcal{I}).
$$
이다. 실제로
$$
f_*\mathcal{I}(V_{j_0 \ldots j_p})
= \mathcal{I}(f^{-1}(V_{j_0 \ldots j_p})) =
\mathcal{I}(f^{-1}(V_{j_0}) \cap \ldots \cap f^{-1}(V_{j_p}))
= \mathcal{I}(U_{j_0 \ldots j_p}).
$$
이기 때문이다. 따라서 보조정리의 첫째 명제는 Lemma
\ref{cohomology-lemma-injective-trivial-cech}에서 따른다. 둘째 명제는 첫째 명제와
Lemma \ref{cohomology-lemma-cech-vanish}에서 따른다.
\end{proof}

\noindent
다음 보조정리는 특히 $f_* : \textit{Ab}(X) \to \textit{Ab}(Y)$가 단사
아벨 층을 단사 아벨 층으로 보냄을 함의한다.

\begin{lemma}
\label{cohomology-lemma-pushforward-injective-flat}
$f : X \to Y$를 환 달린 공간의 사상이라 하고 $f$가 평탄하다고 가정하자.
그러면 임의의 단사 $\mathcal{O}_X$-가군 $\mathcal{I}$에 대해
$f_*\mathcal{I}$는 단사 $\mathcal{O}_Y$-가군이다.
\end{lemma}

\begin{proof}
이 경우 함자 $f^*$는 단사 사상을 단사 사상으로 보낸다(Modules, Lemma
\ref{modules-lemma-pullback-flat}). 따라서 결과는 Homology, Lemma
\ref{homology-lemma-adjoint-preserve-injectives}에서 따른다.
\end{proof}

\begin{lemma}
\label{cohomology-lemma-cohomology-products}
$(X, \mathcal{O}_X)$를 환 달린 공간, $I$를 집합이라 하자. 각
$i \in I$에 대해 $\mathcal{F}_i$를 $\mathcal{O}_X$-가군이라 하자.
$U \subset X$를 열린집합이라 하자. 표준 사상
$$
H^p(U, \prod\nolimits_{i \in I} \mathcal{F}_i)
\longrightarrow
\prod\nolimits_{i \in I} H^p(U, \mathcal{F}_i)
$$
은 $p = 0$일 때 동형이고 $p = 1$일 때 단사이다.
\end{lemma}

\begin{proof}
$p = 0$에 대한 명제는 층들의 곱이 밑에 놓인 준층들의 곱과 같기 때문에
성립한다. Sheaves, Section \ref{sheaves-section-limits-sheaves}를 보라.
$p = 1$인 경우를 증명하자. $\mathcal{F} = \prod \mathcal{F}_i$로 놓자.
$\xi \in H^1(U, \mathcal{F})$가 $\prod H^1(U, \mathcal{F}_i)$에서 영으로
간다고 하자. 코호몰로지의 국소성(Lemma
\ref{cohomology-lemma-kill-cohomology-class-on-covering} 참조)에 의해 모든 $j$에 대해
$\xi|_{U_j} = 0$이 되도록 하는 열린 덮개
$\mathcal{U} : U = \bigcup U_j$가 존재한다. Lemma \ref{cohomology-lemma-cech-h1}에
의해 이는 $\xi$가 어떤 원소
$\check \xi \in \check H^1(\mathcal{U}, \mathcal{F})$에서 온다는 뜻이다.
모든 $i$에 대해 사상
$\check H^1(\mathcal{U}, \mathcal{F}_i) \to H^1(U, \mathcal{F}_i)$
은 단사이고(Lemma \ref{cohomology-lemma-cech-h1}에 의해), $\xi$의 상은
$\prod H^1(U, \mathcal{F}_i)$에서 영이므로 그 상
$\check \xi_i = 0$은 $\check H^1(\mathcal{U}, \mathcal{F}_i)$에서 영이다.
한편 $\mathcal{F} = \prod \mathcal{F}_i$이므로
$\check{\mathcal{C}}^\bullet(\mathcal{U}, \mathcal{F})$는 복합체들
$\check{\mathcal{C}}^\bullet(\mathcal{U}, \mathcal{F}_i)$의 곱이다.
따라서 Homology, Lemma
\ref{homology-lemma-product-abelian-groups-exact}에 의해 원하는 대로
$\check \xi = 0$을 얻는다.
\end{proof}







\section{플라스크 층}
\label{cohomology-section-flasque}

\noindent
정의는 다음과 같다.

\begin{definition}
\label{cohomology-definition-flasque}
$X$를 위상공간이라 하자. 집합의 준층 $\mathcal{F}$가 {\it 플라스크} 또는
{\it 플래비}라는 것은 $X$의 모든 열린집합 $U \subset V$에 대해 제한 사상
$\mathcal{F}(V) \to \mathcal{F}(U)$가 전사라는 뜻이다.
\end{definition}

\noindent
$X$가 환 달린 공간일 때 아벨 층과 가군층에도 이 용어를 사용할 것이다.
분명 모든 열린집합 $U \subset X$에 대해 제한 사상
$\mathcal{F}(X) \to \mathcal{F}(U)$가 전사라고 가정하는 것으로 충분하다.

\begin{lemma}
\label{cohomology-lemma-injective-flasque}
$(X, \mathcal{O}_X)$를 환 달린 공간이라 하자. 그러면 모든 단사
$\mathcal{O}_X$-가군은 플라스크이다.
\end{lemma}

\begin{proof}
이는 Lemma \ref{cohomology-lemma-injective-restriction-surjective}를 다시 서술한 것이다.
\end{proof}

\begin{lemma}
\label{cohomology-lemma-flasque-acyclic}
$(X, \mathcal{O}_X)$를 환 달린 공간이라 하자. 모든 플라스크
$\mathcal{O}_X$-가군은 $R\Gamma(X, -)$에 대해 비순환이며, $X$의 임의의
열린집합 $U$에 대해 $R\Gamma(U, -)$에 대해서도 비순환이다.
\end{lemma}

\begin{proof}
Derived Categories, Lemma
\ref{derived-lemma-subcategory-right-acyclics}를 사용하여 증명한다. 모든
단사 가군은 플라스크이므로 모든 $\mathcal{O}_X$-가군을 플라스크 가군에
매입할 수 있다. Injectives, Lemma
\ref{injectives-lemma-abelian-sheaves-space}를 보라. 따라서 짧은 완전열
$$
0 \to \mathcal{F} \to \mathcal{G} \to \mathcal{H} \to 0
$$
이 주어지고 $\mathcal{F}$와 $\mathcal{G}$가 플라스크이면
$\mathcal{H}$도 플라스크이고, $X$의 임의의 열린집합 위에서 절단을 취한
뒤에도 이 열이 짧은 완전열임을 보이면 충분하다. 실제로 둘째 명제가 첫째를
함의한다. $U \subset X$를 열린 부분공간, $s \in \mathcal{H}(U)$라 하자.
$s$를 $U$ 위의 $\mathcal{G}$의 절단으로 올릴 수 있음을 보이겠다. 이를 위해
쌍 $(V, t)$들의 집합 $T$를 생각하자. 여기서 $V \subset U$는 열려 있고,
$t \in \mathcal{G}(V)$는 $\mathcal{H}$에서 $s|_V$로 가는 절단이다.
$V \subset V'$이고 $t'|_V = t$일 필요충분조건으로
$(V, t) \leq (V', t')$를 정하여 $T$에 부분순서를 준다.
$(V_\alpha, t_\alpha)$, $\alpha \in A$가 $T$의 전순서 부분집합이면
$V = \bigcup V_\alpha$는 열려 있고, $\mathcal{G}$의 층 조건에 의해 각
$V_\alpha$ 위에서 $t_\alpha$로 제한되는 유일한 절단
$t \in \mathcal{G}(V)$가 있다. 따라서 Zorn 보조정리에 의해 $T$에는 극대
원소 $(V, t)$가 존재한다. $V = U$임을 보여 증명을 끝내자. 임의의
$x \in U$를 택하자. $x$의 작은 열린 근방 $W \subset U$와
$\mathcal{H}$에서 $s|_W$로 가는 $t' \in \mathcal{G}(W)$를 찾을 수 있다.
그러면 $t'|_{W \cap V} - t|_{W \cap V}$는 $\mathcal{H}$에서 영으로 가므로
어떤 절단 $r' \in \mathcal{F}(W \cap V)$에서 온다. $\mathcal{F}$가
플라스크이므로 $W \cap V$ 위에서 $r'$로 제한되는 절단
$r \in \mathcal{F}(W)$를 찾을 수 있다. $t'$를 $r$의 상만큼 바꾸면
$t$와 $t'$가 $W \cap V$ 위에서 같은 절단으로 제한된다고 가정할 수 있다.
$\mathcal{G}$의 층 조건에 의해 $t$와 $t'$로 제한되는 $W \cup V$ 위의
$\mathcal{G}$의 절단 $\tilde t$를 찾을 수 있다. $(V, t)$의 극대성에 의해
$V \cup W = V$이다. 따라서 $x \in V$이고 증명이 끝난다.
\end{proof}

\noindent
다음 보조정리는 플라스크 준층에 대해서는 성립하지 않는다.

\begin{lemma}
\label{cohomology-lemma-flasque-acyclic-cech}
$(X, \mathcal{O}_X)$를 환 달린 공간, $\mathcal{F}$를
$\mathcal{O}_X$-가군층, $\mathcal{U} : U = \bigcup U_i$를 열린 덮개라 하자.
$\mathcal{F}$가 플라스크이면 $p > 0$에 대해
$\check{H}^p(\mathcal{U}, \mathcal{F}) = 0$이다.
\end{lemma}

\begin{proof}
Lemma \ref{cohomology-lemma-cech-spectral-sequence}의 서술에 사용된 준층
$\underline{H}^q(\mathcal{F})$는 Lemma \ref{cohomology-lemma-flasque-acyclic}에 의해
영이다. 따라서 다시 Lemma \ref{cohomology-lemma-flasque-acyclic}에 의해
$\check{H}^p(U, \mathcal{F}) = H^p(U, \mathcal{F}) = 0$이다.
\end{proof}

\begin{lemma}
\label{cohomology-lemma-flasque-acyclic-pushforward}
$f : (X, \mathcal{O}_X) \to (Y, \mathcal{O}_Y)$를 환 달린 공간의 사상,
$\mathcal{F}$를 $\mathcal{O}_X$-가군층이라 하자. $\mathcal{F}$가
플라스크이면 $p > 0$에 대해 $R^pf_*\mathcal{F} = 0$이다.
\end{lemma}

\begin{proof}
Lemma \ref{cohomology-lemma-describe-higher-direct-images}와 Lemma
\ref{cohomology-lemma-flasque-acyclic}에서 즉시 따른다.
\end{proof}

\noindent
다음 보조정리는 유한 덮개에 대해서는 초보적인 귀납 논증으로 증명할 수 있다.
\cite{FOAG}의 체흐 코호몰로지에 대한 논의와 비교하라.

\begin{lemma}
\label{cohomology-lemma-vanishing-ravi}
$X$를 위상공간, $\mathcal{F}$를 $X$ 위의 아벨 층,
$\mathcal{U} : U = \bigcup_{i \in I} U_i$를 열린 덮개라 하자.
$U'$가 $U_{i_0 \ldots i_p}$ 꼴의 열린집합들의 임의의 합집합일 때 제한 사상
$\mathcal{F}(U) \to \mathcal{F}(U')$가 전사라고 가정하자. 그러면
$p > 0$에 대해 $\check{H}^p(\mathcal{U}, \mathcal{F})$는 영이다.
\end{lemma}

\begin{proof}
$Y$를 $I$의 공집합이 아닌 부분집합들의 집합이라 하자. 문자
$A, B, C, \ldots$로 $Y$의 원소, 즉 $I$의 공집합이 아닌 부분집합을
나타낼 것이다. 유한하고 공집합이 아닌 부분집합 $J \subset I$에 대해
$$
V_J = \{A \in Y \mid J \subset A\}
$$
로 놓자. 이는 $V_{\{i\}} = \{A \in Y \mid i \in A\}$이고
$V_J = \bigcap_{j \in J} V_{\{j\}}$라는 뜻이다. 이때
$V_J \subset V_K$일 필요충분조건은 $J \supset K$이다. 부분집합들
$V_J$의 모임이 $Y$의 위상에 대한 기저가 되게 하는 $Y$ 위의 위상은
유일하다. 모든 열린집합은 어떤 유한 부분집합들의 족 $J_t$에 대해
$$
V = \bigcup\nolimits_{t \in T} V_{J_t}
$$
꼴이다. $J_t \subset J_{t'}$이면 $V$를 바꾸지 않고 족에서 $J_{t'}$를
제거할 수 있다. 따라서 $J_t$들 사이에 포함관계가 없다고 가정할 수 있다.
이 경우 $V$의 극소 원소는 집합들 $A = J_t$이다. 그러므로 열린집합
$V$에서 족 $(J_t)_{t \in T}$를 읽어낼 수 있다.

\medskip\noindent
$Y$의 열린 덮개들을 완전히 이해할 수 있다. 먼저 $A \in Y$의 원소들이
$I$의 공집합이 아닌 부분집합이므로
$$
Y = \bigcup\nolimits_{i \in I} V_{\{i\}}
$$
이다. 다른 덮개들을 이해하기 위해 $V$가 위와 같고 $V_s \subset Y$가
족 $(J_{s, t})_{t \in T_s}$에 대응하는 열린집합이라고 하자. 그러면
$$
V = \bigcup\nolimits_{s \in S} V_s
$$
일 필요충분조건은 각 $t \in T$에 대해 $J_t = J_{s, t_s}$가 되도록 하는
$s \in S$와 $t_s \in T_s$가 존재하는 것이다. 실제로 족
$(J_t)_{t \in T}$가 극소이므로 극소 원소 $A = J_t$는 어떤 $s$에 대해
$V_s$에 속해야 하고, 따라서 어떤 $t_s \in T_s$에 대해
$A \in V_{J_{t_s}}$이다. 그런데 $A$는 $V_s$에서도 극소이므로
$J_{t_s} = J_t$를 얻는다.

\medskip\noindent
이제 $Y$의 열린집합들의 집합을 $X$의 열린집합들로 보낸다. 구체적으로
$Y$를 $U$로 보내고, 열린집합 $V_J$에는 규칙
$$
V_J \mapsto U_J = \bigcap\nolimits_{i \in J} U_i
$$
을 사용하며, 임의의 열린집합 $V$에는 규칙
$$
V = \bigcup\nolimits_{t \in T} V_{J_t}
\mapsto
\bigcup\nolimits_{t \in T} U_{J_t}
$$
으로 확장한다. 위에서 얻은 $Y$의 열린 덮개 분류에 의해 이 규칙은 열린
덮개를 열린 덮개로 보낸다. 따라서 $\mathcal{G}(Y) = \mathcal{F}(U)$로
놓고, $V = \bigcup\nolimits_{t \in T} V_{J_t}$에 대해
$$
\mathcal{G}(V) = \mathcal{F}\left(\bigcup\nolimits_{t \in T} U_{J_t}\right)
$$
로 놓으며 $\mathcal{F}$의 제한 사상들을 사용하면 $Y$ 위의 아벨 층
$\mathcal{G}$를 얻는다.

\medskip\noindent
이 준비를 마쳤으므로 다음과 같이 보조정리를 증명할 수 있다. $Y$의 열린
덮개 $\mathcal{V} : Y = \bigcup_{i \in I} V_{\{i\}}$가 있다. 구성에 의해
체흐 복합체들의 등식
$$
\check{C}^\bullet(\mathcal{V}, \mathcal{G}) =
\check{C}^\bullet(\mathcal{U}, \mathcal{F})
$$
을 얻는다. 층 $\mathcal{G}$는 $Y$ 위에서 플라스크이므로(보조정리의
$\mathcal{F}$에 대한 가정에 의해), 영이 된다는 결론은 Lemma
\ref{cohomology-lemma-flasque-acyclic-cech}에서 따른다.
\end{proof}




\section{Leray 스펙트럼열}
\label{cohomology-section-Leray}

\begin{lemma}
\label{cohomology-lemma-before-Leray}
$f : X \to Y$를 환 달린 공간의 사상이라 하자. 가환 도식
$$
\xymatrix{
D^{+}(X) \ar[rr]_-{R\Gamma(X, -)} \ar[d]_{Rf_*} & &
D^{+}(\mathcal{O}_X(X)) \ar[d]^{\text{restriction}} \\
D^{+}(Y) \ar[rr]^-{R\Gamma(Y, -)} & &
D^{+}(\mathcal{O}_Y(Y))
}
$$
이 있다. 더 일반적으로 임의의 열린집합 $V \subset Y$와
$U = f^{-1}(V)$에 대해 가환 도식
$$
\xymatrix{
D^{+}(X) \ar[rr]_-{R\Gamma(U, -)} \ar[d]_{Rf_*} & &
D^{+}(\mathcal{O}_X(U)) \ar[d]^{\text{restriction}} \\
D^{+}(Y) \ar[rr]^-{R\Gamma(V, -)} & &
D^{+}(\mathcal{O}_Y(V))
}
$$
이 있다. 더 자세한 설명은 Remark \ref{cohomology-remark-elucidate-lemma}도 보라.
\end{lemma}

\begin{proof}
$\Gamma_{res} : \textit{Mod}(\mathcal{O}_X) \to
\text{Mod}_{\mathcal{O}_Y(Y)}$를 $\mathcal{O}_X$-가군 $\mathcal{F}$에
$f^\sharp : \mathcal{O}_Y(Y) \to \mathcal{O}_X(X)$를 통해
$\mathcal{O}_Y(Y)$-가군으로 본 $\mathcal{F}$의 대역 절단을 대응시키는
함자라 하자. 또한 $f^\sharp : \mathcal{O}_Y(Y) \to
\mathcal{O}_X(X)$가 유도하는 스칼라 제한 함자를
$restriction : \text{Mod}_{\mathcal{O}_X(X)} \to
\text{Mod}_{\mathcal{O}_Y(Y)}$라 하자. $restriction$은 완전하므로 그 오른쪽
유도함자는 제한 함자를 그대로 적용하여 계산된다. Derived Categories,
Lemma \ref{derived-lemma-right-derived-exact-functor}를 보라. 분명
$$
\Gamma_{res}
=
restriction \circ \Gamma(X, -)
=
\Gamma(Y, -) \circ f_*
$$
이다. Derived Categories, Lemma
\ref{derived-lemma-compose-derived-functors}가 두 합성 모두에 적용된다고
주장한다. 첫째 합성에 대해서는 위의 설명으로 분명하다. 둘째 합성에
대해서는 Lemma \ref{cohomology-lemma-pushforward-injective}에서 따르는데, 그 보조정리는
단사 $\mathcal{O}_X$-가군들이 $Y$ 위의 $\Gamma(Y, -)$-비순환 층들로
보내짐을 함의한다.
\end{proof}

\begin{remark}
\label{cohomology-remark-elucidate-lemma}
Lemma \ref{cohomology-lemma-before-Leray}의 의미를 구체적으로 설명하자. 그 보조정리는
$f : X \to Y$, $\mathcal{F} \in \textit{Mod}(\mathcal{O}_X)$, 그리고 단사
분해 $\mathcal{F} \to \mathcal{I}^\bullet$이 주어졌을 때 다음이 성립한다는
말이다.
$$
\begin{matrix}
R\Gamma(X, \mathcal{F}) & \text{is represented by} &
\Gamma(X, \mathcal{I}^\bullet) \\
Rf_*\mathcal{F} & \text{is represented by} & f_*\mathcal{I}^\bullet \\
R\Gamma(Y, Rf_*\mathcal{F}) & \text{is represented by} &
\Gamma(Y, f_*\mathcal{I}^\bullet)
\end{matrix}
$$
마지막 사실은 Leray의 비순환성 보조정리(Derived Categories, Lemma
\ref{derived-lemma-leray-acyclicity})와 Lemma
\ref{cohomology-lemma-pushforward-injective}에서 따른다. 마지막으로 이를 자명한 관찰
$$
\Gamma(X, \mathcal{I}^\bullet)
=
\Gamma(Y, f_*\mathcal{I}^\bullet).
$$
과 결합하여 보조정리의 도식이 가환임을 얻는다.
\end{remark}

\begin{lemma}
\label{cohomology-lemma-modules-abelian}
$X$를 환 달린 공간, $\mathcal{F}$를 $\mathcal{O}_X$-가군이라 하자.
\begin{enumerate}
\item 열린집합 $U \subset X$에 대해 $\mathcal{F}$를
$\mathcal{O}_X$-가군으로 보아 계산한 코호몰로지군
$H^i(U, \mathcal{F})$와 아벨 층으로 보아 계산한 것은 같다.
\item $f : X \to Y$를 환 달린 공간의 사상이라 하자. $\mathcal{F}$를
$\mathcal{O}_X$-가군으로 보아 계산한 고차 직접상
$R^if_*\mathcal{F}$와 아벨 층으로 보아 계산한 것은 같다.
\end{enumerate}
$\mathcal{O}_X$-가군의 아래로 유계인 복합체에 대해서도 비슷한 명제가
성립한다.
\end{lemma}

\begin{proof}
밑에 놓인 위상공간에서는 항등이고 환의 층의 유일한 사상
$\underline{\mathbf{Z}}_X \to \mathcal{O}_X$로 주어지는 환 달린 공간의 사상
$(X, \mathcal{O}_X) \to (X, \underline{\mathbf{Z}}_X)$
을 생각하자. $\mathcal{F}$를 $\mathcal{O}_X$-가군이라 하자. 같은 층을
$\underline{\mathbf{Z}}_X$-가군, 즉 아벨군의 층으로 본 것을
$\mathcal{F}_{ab}$로 나타내자. $\mathcal{F} \to \mathcal{I}^\bullet$을 단사
분해라 하자. Remark \ref{cohomology-remark-elucidate-lemma}에 의해
$\Gamma(X, \mathcal{I}^\bullet)$은 $R\Gamma(X, \mathcal{F})$와
$R\Gamma(X, \mathcal{F}_{ab})$를 모두 계산한다. 이로써 (1)이 증명된다.

\medskip\noindent
(2)를 증명하기 위해 (1)과 Lemma
\ref{cohomology-lemma-describe-higher-direct-images}를 사용한다. 결과는 즉시 따른다.
\end{proof}

\begin{lemma}[Leray spectral sequence]
\label{cohomology-lemma-Leray}
$f : X \to Y$를 환 달린 공간의 사상, $\mathcal{F}^\bullet$을
$\mathcal{O}_X$-가군의 아래로 유계인 복합체라 하자. 스펙트럼열
$$
E_2^{p, q} = H^p(Y, R^qf_*(\mathcal{F}^\bullet))
$$
은 $H^{p + q}(X, \mathcal{F}^\bullet)$로 수렴한다.
\end{lemma}

\begin{proof}
이는 함자 합성 $\Gamma_{res} = \Gamma(Y, -) \circ f_*$에서 오는
Grothendieck 스펙트럼열(Derived Categories, Lemma
\ref{derived-lemma-grothendieck-spectral-sequence})일 뿐이다. 여기서
$\Gamma_{res}$는 Lemma \ref{cohomology-lemma-before-Leray}의 증명에서와 같다. Derived
Categories, Lemma \ref{derived-lemma-grothendieck-spectral-sequence}의 가정이
만족됨을 보려면 Lemma \ref{cohomology-lemma-before-Leray}의 증명이나 Remark
\ref{cohomology-remark-elucidate-lemma}를 보라.
\end{proof}

\begin{remark}
\label{cohomology-remark-Leray-ss-more-structure}
Lemma \ref{cohomology-lemma-Leray}에서 증명한 방식에 따른 Leray 스펙트럼열은
$\Gamma(Y, \mathcal{O}_Y)$-가군의 스펙트럼열이다. 그러나 실제로는
$\Gamma(X, \mathcal{O}_X)$-가군의 스펙트럼열임을 쉽게 알 수 있다. 예를
들어 $f$는 환 달린 공간의 사상
$f' :  (X, \mathcal{O}_X) \to (Y, f_*\mathcal{O}_X)$을 낳는다. Lemma
\ref{cohomology-lemma-modules-abelian}에 의해 $\mathcal{O}_X$-가군 $\mathcal{F}$와
$f$에 대한 Leray 스펙트럼열의 항 $E_r^{p, q}$는 적어도 $r \geq 2$에 대해
$\mathcal{F}$와 $f'$에 대한 항과 같다. 실제로 둘 다 $\mathcal{F}$를 아벨
층으로 보았을 때의 Leray 스펙트럼열의 항과 같다. 그리고
$(f_*\mathcal{O}_X)(Y) = \mathcal{O}_X(X)$이므로 결과를 얻는다. Leray
스펙트럼열이 추가 구조를 지니는 경우가 흔히 있다.
\end{remark}

\begin{lemma}
\label{cohomology-lemma-apply-Leray}
$f : X \to Y$를 환 달린 공간의 사상, $\mathcal{F}$를
$\mathcal{O}_X$-가군이라 하자.
\begin{enumerate}
\item $q > 0$에 대해 $R^qf_*\mathcal{F} = 0$이면 모든 $p$에 대해
$H^p(X, \mathcal{F}) = H^p(Y, f_*\mathcal{F})$이다.
\item 모든 $q$와 $p > 0$에 대해 $H^p(Y, R^qf_*\mathcal{F}) = 0$이면
모든 $q$에 대해 $H^q(X, \mathcal{F}) = H^0(Y, R^qf_*\mathcal{F})$이다.
\end{enumerate}
\end{lemma}

\begin{proof}
이는 Leray 스펙트럼열을 $E_2$에서 퇴화시키는 두 가지 간단한 조건이다.
이 사실들을 스펙트럼열을 사용하지 않고 직접 증명할 수도 있으며, 이는 층의
코호몰로지에 관한 좋은 연습문제이다.
\end{proof}

\begin{lemma}
\label{cohomology-lemma-higher-direct-images-compose}
\begin{slogan}
합성의 전체 유도함자는 전체 유도함자들의 합성이다.
\end{slogan}
$f : X \to Y$와 $g : Y \to Z$를 환 달린 공간의 사상들이라 하자.
이 경우 $D^{+}(X) \to D^{+}(Z)$인 함자로서
$Rg_* \circ Rf_* = R(g \circ f)_*$이다.
\end{lemma}

\begin{proof}
Derived Categories, Lemma
\ref{derived-lemma-compose-derived-functors}를 적용한다. 분명
$g_* \circ f_* = (g \circ f)_*$이다. Sheaves, Lemma
\ref{sheaves-lemma-pushforward-composition}을 보라. 이제
$f_*\mathcal{I}$가 $g_*$-비순환임을 보이면 된다. 이는 Lemma
\ref{cohomology-lemma-pushforward-injective}와 Lemma
\ref{cohomology-lemma-describe-higher-direct-images}의 고차 직접상 $R^ig_*$에 대한
기술에서 따른다.
\end{proof}

\begin{lemma}[상대 Leray 스펙트럼열]
\label{cohomology-lemma-relative-Leray}
$f : X \to Y$와 $g : Y \to Z$를 환 달린 공간의 사상들이라 하고,
$\mathcal{F}$를 $\mathcal{O}_X$-가군이라 하자. 다음을 갖는 스펙트럼열이 있다.
$$
E_2^{p, q} = R^pg_*(R^qf_*\mathcal{F})
$$
이는 $R^{p + q}(g \circ f)_*\mathcal{F}$로 수렴한다. 이 스펙트럼열은
$\mathcal{F}$에 대해 함자적이며, $\mathcal{O}_X$-가군의 아래로 유계인
복합체에 대한 판도 있다.
\end{lemma}

\begin{proof}
이는 함자 합성에 대한 Grothendieck 스펙트럼열이며, Lemma
\ref{cohomology-lemma-higher-direct-images-compose}와 Derived Categories, Lemma
\ref{derived-lemma-grothendieck-spectral-sequence}에서 따른다.
\end{proof}














\section{코호몰로지의 함자성}
\label{cohomology-section-functoriality}

\begin{lemma}
\label{cohomology-lemma-functoriality}
$f : X \to Y$를 환 달린 공간의 사상이라 하자.
$\mathcal{G}^\bullet$과 $\mathcal{F}^\bullet$을 각각
$\mathcal{O}_Y$-가군과 $\mathcal{O}_X$-가군의 아래로 유계인 복합체라 하자.
복합체의 사상
$\varphi : \mathcal{G}^\bullet \to f_*\mathcal{F}^\bullet$
가 주어졌다고 하자. 그러면 $D^{+}(Y)$에 표준 사상
$$
\mathcal{G}^\bullet
\longrightarrow
Rf_*(\mathcal{F}^\bullet)
$$
이 존재한다. 더구나 이 구성은 삼중항
$(\mathcal{G}^\bullet, \mathcal{F}^\bullet, \varphi)$에 대해 함자적이다.
\end{lemma}

\begin{proof}
단사 분해 $\mathcal{F}^\bullet \to \mathcal{I}^\bullet$을 택하자.
정의에 따라 $Rf_*(\mathcal{F}^\bullet)$은
$K^{+}(\mathcal{O}_Y)$에서 $f_*\mathcal{I}^\bullet$로 표현된다. 합성
$$
\mathcal{G}^\bullet \to f_*\mathcal{F}^\bullet \to f_*\mathcal{I}^\bullet
$$
은 $K^{+}(Y)$의 사상이며, 국소화 함자
$j_Y : K^{+}(Y) \to D^{+}(Y)$를 적용하면 보조정리의 사상이 된다.
\end{proof}

\noindent
$f : X \to Y$를 환 달린 공간의 사상이라 하자. $\mathcal{G}$를
$\mathcal{O}_Y$-가군, $\mathcal{F}$를 $\mathcal{O}_X$-가군이라 하자.
$\mathcal{G}$에서 $\mathcal{F}$로 가는 $f$-사상 $\varphi$란 사상
$\varphi : \mathcal{G} \to f_*\mathcal{F}$, 또는 이와 같은 말로 사상
$\varphi : f^*\mathcal{G} \to \mathcal{F}$임을 상기하자.
Sheaves, Definition \ref{sheaves-definition-f-map}을 보라.
이러한 $f$-사상은 $D^{+}(\mathcal{O}_Y(Y))$에서 복합체의 사상
\begin{equation}
\label{cohomology-equation-functorial-derived}
\varphi :
R\Gamma(Y, \mathcal{G})
\longrightarrow
R\Gamma(X, \mathcal{F})
\end{equation}
을 낳는다. 즉 Lemma \ref{cohomology-lemma-functoriality}의
$D^{+}(Y)$ 안의 사상 $\mathcal{G} \to Rf_*\mathcal{F}$를 사용하고
$R\Gamma(Y, -)$를 적용한다. Lemma \ref{cohomology-lemma-before-Leray}에 의해
$R\Gamma(X, \mathcal{F}) = R\Gamma(Y, Rf_*\mathcal{F})$
임을 알 수 있으므로 표시한 화살표를 얻는다. 아래의 Remark
\ref{cohomology-remark-explain-arrow}에서 이를 완전히 풀어 쓴다. 특히 코호몰로지 위의 사상
\begin{equation}
\label{cohomology-equation-functorial}
\varphi : H^i(Y, \mathcal{G}) \longrightarrow H^i(X, \mathcal{F}).
\end{equation}
들을 얻는다.

\begin{remark}
\label{cohomology-remark-explain-arrow}
$f : X \to Y$를 환 달린 공간의 사상이라 하자. $\mathcal{G}$를
$\mathcal{O}_Y$-가군, $\mathcal{F}$를 $\mathcal{O}_X$-가군이라 하고,
$\varphi$를 $\mathcal{G}$에서 $\mathcal{F}$로 가는 $f$-사상이라 하자.
단사 $\mathcal{O}_X$-가군들의 복합체에 의한 분해
$\mathcal{F} \to \mathcal{I}^\bullet$을 택하자. 또한 단사
$\mathcal{O}_Y$-가군들의 복합체에 의한 분해
$\mathcal{G} \to \mathcal{J}^\bullet$과
$f_*\mathcal{I}^\bullet \to (\mathcal{J}')^\bullet$을 택하자.
Derived Categories, Lemma \ref{derived-lemma-morphisms-lift}에 의해 도식
\begin{equation}
\label{cohomology-equation-choice}
\xymatrix{
\mathcal{G} \ar[d] \ar[r] &
f_*\mathcal{F} \ar[r] &
f_*\mathcal{I}^\bullet \ar[d] \\
\mathcal{J}^\bullet \ar[rr]^\beta & &
(\mathcal{J}')^\bullet
}
\end{equation}
을 가환으로 만드는 복합체의 사상 $\beta$가 존재한다. 대역 절단 함자를 적용하면 도식
$$
\xymatrix{
 & & \Gamma(Y, f_*\mathcal{I}^\bullet) \ar[d]_{qis} \ar@{=}[r] &
\Gamma(X, \mathcal{I}^\bullet) \\
\Gamma(Y, \mathcal{J}^\bullet) \ar[rr]^\beta & &
\Gamma(Y, (\mathcal{J}')^\bullet) &
}
$$
을 얻는다. 왼쪽 아래의 복합체는 $R\Gamma(Y, \mathcal{G})$를 표현하고,
오른쪽 위의 복합체는 $R\Gamma(X, \mathcal{F})$를 표현한다.
수직 화살표는 Lemma \ref{cohomology-lemma-before-Leray}에 의해 유사동형이며, 국소화 함자
$K^{+}(\mathcal{O}_Y(Y)) \to D^{+}(\mathcal{O}_Y(Y))$.
를 적용하면 가역이 된다. 화살표 (\ref{cohomology-equation-functorial-derived})는
수평 사상에 이어 수직 사상의 역을 합성한 것으로 주어진다.
\end{remark}





\section{세분과 체흐 코호몰로지}
\label{cohomology-section-refinements-cech}

\noindent
$(X, \mathcal{O}_X)$를 환 달린 공간이라 하자.
$\mathcal{U} : X = \bigcup_{i \in I} U_i$와
$\mathcal{V} : X = \bigcup_{j \in J} V_j$를 열린 덮개라 하자.
$\mathcal{U}$가 $\mathcal{V}$의 세분이라고 가정하자. 모든 $i \in I$에 대해
$U_i \subset V_{c(i)}$가 되도록 사상 $c : I \to J$를 택하자.
이는 체흐 복합체의 사상
$$
\gamma :
\check{\mathcal{C}}^\bullet(\mathcal{V}, \mathcal{F})
\longrightarrow
\check{\mathcal{C}}^\bullet(\mathcal{U}, \mathcal{F}),
\quad
(\xi_{j_0 \ldots j_p})
\longmapsto
(\xi_{c(i_0) \ldots c(i_p)}|_{U_{i_0 \ldots i_p}})
$$
을 유도하며, 이는 $\mathcal{O}_X$-가군층 $\mathcal{F}$에 대해 함자적이다.
$c' : I \to J$도 모든 $i \in I$에 대해
$U_i \subset V_{c'(i)}$를 만족하는 사상이라고 하자. 그러면 대응하는 사상
$\gamma$와 $\gamma'$은 서로 호모토픽이다. 구체적으로,
$\gamma - \gamma' = \text{d} \circ h + h \circ \text{d}$
이며, 여기서
$h : \check{\mathcal{C}}^{p + 1}(\mathcal{V}, \mathcal{F}) \to
\check{\mathcal{C}}^p(\mathcal{U}, \mathcal{F})$
는 공식
$$
h(\alpha)_{i_0 \ldots i_p} =
\sum\nolimits_{a = 0}^{p}
(-1)^a
\alpha_{c(i_0)\ldots c(i_a) c'(i_a) \ldots c'(i_p)}
$$
로 주어진다. 이것이 성립함을 보이는 계산은 생략한다. 더 일반적인 경우의
증명은 (\ref{cohomology-equation-transformation}) 뒤의 논의를 보라. 특히 체흐
코호몰로지군 위의 사상은 $c$의 선택과 무관하다. 더구나
$\mathcal{W} : X = \bigcup_{k \in K} W_k$가 셋째 열린 덮개이고
$\mathcal{V}$가 $\mathcal{W}$의 세분이면, 사상들의 합성
$$
\check{\mathcal{C}}^\bullet(\mathcal{W}, \mathcal{F})
\longrightarrow
\check{\mathcal{C}}^\bullet(\mathcal{V}, \mathcal{F})
\longrightarrow
\check{\mathcal{C}}^\bullet(\mathcal{U}, \mathcal{F})
$$
은 사상 $I \to J$와 $J \to K$에 결부된 사상이며, 이는 합성
$I \to K$에 결부된 사상임이 분명하다. 따라서 체흐 코호몰로지군을
$$
\check{H}^p(X, \mathcal{F}) =
\colim_\mathcal{U} \check{H}^p(\mathcal{U}, \mathcal{F})
$$
로 정의할 수 있다. 여기서 여극한은 세분에 의해 준순서가 주어진 $X$의 모든
열린 덮개 위에서 취한다.

\medskip\noindent
위에서 정의한 사상 $\gamma$는 코호몰로지로 가는 사상과 양립한다. 즉 합성
$$
\check{H}^p(\mathcal{V}, \mathcal{F}) \to
\check{H}^p(\mathcal{U}, \mathcal{F})
\xrightarrow{\text{보조정리 \ref{cohomology-lemma-cech-cohomology}}}
H^p(X, \mathcal{F})
$$
은 Lemma \ref{cohomology-lemma-cech-cohomology}의 첫째 군에서 코호몰로지로 가는 표준
사상이다. 아래 보조정리에서는 이를 조금 더 일반적인 상황에서 증명한다.
그 결과 체흐 코호몰로지에서 코호몰로지로 가는 잘 정의된 사상
\begin{equation}
\label{cohomology-equation-cech-to-cohomology}
\check{H}^p(X, \mathcal{F}) =
\colim_\mathcal{U} \check{H}^p(\mathcal{U}, \mathcal{F})
\longrightarrow
H^p(X, \mathcal{F})
\end{equation}
을 얻는다.

\begin{lemma}
\label{cohomology-lemma-functoriality-cech}
$f : X \to Y$를 환 달린 공간의 사상이라 하자.
$\varphi : f^*\mathcal{G} \to \mathcal{F}$를
$\mathcal{O}_Y$-가군 $\mathcal{G}$에서 $\mathcal{O}_X$-가군
$\mathcal{F}$로 가는 $f$-사상이라 하자.
$\mathcal{U} : X = \bigcup_{i \in I} U_i$와
$\mathcal{V} : Y = \bigcup_{j \in J} V_j$를 열린 덮개라 하자.
$\mathcal{U}$가
$f^{-1}\mathcal{V} : X = \bigcup_{j \in J} f^{-1}(V_j)$의 세분이라고
가정하자. 이 경우 $D^{+}(\mathcal{O}_X(X))$에 가환 도식
$$
\xymatrix{
\check{\mathcal{C}}^\bullet(\mathcal{U}, \mathcal{F}) \ar[r] &
R\Gamma(X, \mathcal{F}) \\
\check{\mathcal{C}}^\bullet(\mathcal{V}, \mathcal{G}) \ar[r]
\ar[u]^\gamma &
R\Gamma(Y, \mathcal{G}) \ar[u]
}
$$
이 존재한다. 여기서 수평 화살표들은 Lemma
\ref{cohomology-lemma-cech-cohomology}에 의해, 오른쪽 수직 화살표는
(\ref{cohomology-equation-functorial-derived})에 의해 주어진다. 특히 코호몰로지군의 가환 도식
$$
\xymatrix{
\check{H}^p(\mathcal{U}, \mathcal{F}) \ar[r] &
H^p(X, \mathcal{F}) \\
\check{H}^p(\mathcal{V}, \mathcal{G}) \ar[r]
\ar[u]^\gamma &
H^p(Y, \mathcal{G}) \ar[u]
}
$$
을 얻으며, 여기서 오른쪽 수직 화살표는 (\ref{cohomology-equation-functorial})이다.
\end{lemma}

\begin{proof}
먼저 왼쪽 수직 화살표를 정의하자. 즉 모든 $i \in I$에 대해
$U_i \subset f^{-1}(V_{c(i)})$가 되도록 사상 $c : I \to J$를 택한다.
차수 $p$에서 사상을 공식
$$
\gamma(s)_{i_0 \ldots i_p} = \varphi(s)_{c(i_0) \ldots c(i_p)}
$$
로 정의한다. 가정에 따라 $\varphi$가 실제로 사상
$\mathcal{G}(V_{c(i_0) \ldots c(i_p)}) \to \mathcal{F}(U_{i_0 \ldots i_p})$
을 유도하므로 이는 의미가 있다. 또한 이것이 복합체의 사상을 정의함도
분명하다. $X$ 위의 단사 분해
$\mathcal{F} \to \mathcal{I}^\bullet$과 $Y$ 위의 단사 분해
$\mathcal{G} \to J^\bullet$을 택하자. Lemma
\ref{cohomology-lemma-cech-cohomology}의 증명에 따라 다음 항들을 갖는 이중 복합체
$A^{\bullet, \bullet}$과 $B^{\bullet, \bullet}$을 도입한다.
$$
B^{p, q} = \check{\mathcal{C}}^p(\mathcal{V}, \mathcal{J}^q)
\quad
\text{및}
\quad
A^{p, q} = \check{\mathcal{C}}^p(\mathcal{U}, \mathcal{I}^q).
$$
위의 Remark \ref{cohomology-remark-explain-arrow}에서처럼 $Y$ 위의 단사 분해
$f_*\mathcal{I} \to (\mathcal{J}')^\bullet$과
(\ref{cohomology-equation-choice})를 가환으로 만드는 복합체의 사상
$\beta : \mathcal{J} \to (\mathcal{J}')^\bullet$을 택한다.
또 다음 항들을 갖는 이중 복합체 $(B')^{\bullet, \bullet}$과
$(B'')^{\bullet, \bullet}$을 도입한다.
$$
(B')^{p, q} = \check{\mathcal{C}}^p(\mathcal{V}, (\mathcal{J}')^q)
\quad
\text{및}
\quad
(B'')^{p, q} = \check{\mathcal{C}}^p(\mathcal{V}, f_*\mathcal{I}^q).
$$
$f_*\mathcal{I}^\bullet$에서 $\mathcal{I}^\bullet$로 가는 복합체의
$f$-사상이 있음에 유의하라. 따라서 위와 같은 공식이 이중 복합체의 사상
$$
\gamma : (B'')^{\bullet, \bullet} \longrightarrow A^{\bullet, \bullet}.
$$
을 정의함이 분명하다. 복합체의 도식
$$
\xymatrix{
\check{\mathcal{C}}^\bullet(\mathcal{U}, \mathcal{F})
\ar[r] &
\text{Tot}(A^{\bullet, \bullet}) & & &
\Gamma(X, \mathcal{I}^\bullet) \ar[lll]^{qis}
\ar@{=}[ddl]\\
\check{\mathcal{C}}^\bullet(\mathcal{V}, \mathcal{G})
\ar[r] \ar[u]^\gamma &
\text{Tot}(B^{\bullet, \bullet}) \ar[r]^\beta &
\text{Tot}((B')^{\bullet, \bullet}) &
\text{Tot}((B'')^{\bullet, \bullet}) \ar[l] \ar[llu]_{s\gamma} \\
& \Gamma(Y, \mathcal{J}^\bullet) \ar[u]^{qis} \ar[r]^\beta &
\Gamma(Y, (\mathcal{J}')^\bullet) \ar[u] &
\Gamma(Y, f_*\mathcal{I}^\bullet) \ar[u] \ar[l]_{qis}
}
$$
$\text{Tot}(A^{\bullet, \bullet})$과
$\text{Tot}(B^{\bullet, \bullet})$을 종점으로 하는 두 수평 화살표는 Lemma
\ref{cohomology-lemma-cech-cohomology}에서 설명한 것들이다. 왼쪽 위 도형(오각형)은
(\ref{cohomology-equation-choice})가 가환이므로 가환이다. 아래의 두 정사각형은 자명하게
가환한다. 오른쪽 위 도형(정사각형)도 정의에서 바로 가환임을 알 수 있다.
이제 정의와 다음 사실에서 보조정리의 결과가 따른다. 도식의 바깥쪽을 따라
$\check{\mathcal{C}}^\bullet(\mathcal{V}, \mathcal{G})$에서
$\Gamma(X, \mathcal{I}^\bullet)$까지 위쪽으로 가든 아래쪽으로 가든 결과가 같다
(도중에 만나는 모든 유사동형의 역을 취해야 한다).
\end{proof}





\section{하우스도르프 준콤팩트 공간 위의 코호몰로지}
\label{cohomology-section-cohomology-LC}

\noindent
그런 공간에서는 체흐 코호몰로지와 코호몰로지가 일치한다.

\begin{lemma}
\label{cohomology-lemma-cech-always}
$X$를 위상공간, $\mathcal{F}$를 아벨 층이라 하자. 그러면
(\ref{cohomology-equation-cech-to-cohomology})에서 정의한 사상
$\check{H}^1(X, \mathcal{F}) \to H^1(X, \mathcal{F})$은 동형이다.
\end{lemma}

\begin{proof}
$\mathcal{U}$를 $X$의 열린 덮개라 하자. Lemma
\ref{cohomology-lemma-cech-spectral-sequence}에 의해 완전열
$$
0 \to \check{H}^1(\mathcal{U}, \mathcal{F}) \to H^1(X, \mathcal{F})
\to \check{H}^0(\mathcal{U}, \underline{H}^1(\mathcal{F}))
$$
이 있다. 따라서 이 사상은 단사이다. 전사임을 보이려면
$\check{H}^0(\mathcal{U}, \underline{H}^1(\mathcal{F}))$의 임의의 원소가
$\mathcal{U}$를 어떤 세분으로 바꾼 뒤 영으로 보내짐을 보이면 충분하다.
이는 정의와 다음 사실에서 바로 따른다. $\underline{H}^1(\mathcal{F})$는
아벨군의 준층이고 코호몰로지의 국소성에 의해 그 층화가 영이다. Lemma
\ref{cohomology-lemma-kill-cohomology-class-on-covering}을 보라.
\end{proof}

\begin{lemma}
\label{cohomology-lemma-cech-Hausdorff-quasi-compact}
$X$를 하우스도르프 준콤팩트 위상공간이라 하고,
$\mathcal{F}$를 $X$ 위의 아벨 층이라 하자. 그러면 모든 $n$에 대해
(\ref{cohomology-equation-cech-to-cohomology})에서 정의한 사상
$\check{H}^n(X, \mathcal{F}) \to H^n(X, \mathcal{F})$은 동형이다.
\end{lemma}

\begin{proof}
$n = 0, 1$일 때 $\check{H}^n(X, -) \to H^n(X, -)$가 함자들의 동형임은
이미 안다. Lemma \ref{cohomology-lemma-cech-always}를 보라. 함자 $H^n(X, -)$들은
보편 $\delta$-함자를 이룬다. Derived Categories, Lemma
\ref{derived-lemma-higher-derived-functors}를 보라.
$\check{H}^n(X, -)$가 보편 $\delta$-함자를 이루고
$\check{H}^n(X, -) \to H^n(X, -)$가 경계 사상들과 양립함을 보이면,
보편 $\delta$-함자의 유일성에 의해 이 사상은 자동으로 동형이 된다.
Homology, Lemma \ref{homology-lemma-uniqueness-universal-delta-functor}를 보라.

\medskip\noindent
$0 \to \mathcal{F} \to \mathcal{G} \to \mathcal{H} \to 0$을
$X$ 위의 아벨 층들의 짧은 완전열이라 하자.
$\mathcal{U} : X = \bigcup_{i \in I} U_i$를 열린 덮개라 하자.
이로부터 복합체들의 복합체
$$
0 \to \check{\mathcal{C}}^\bullet(\mathcal{U}, \mathcal{F}) \to
\check{\mathcal{C}}^\bullet(\mathcal{U}, \mathcal{G}) \to
\check{\mathcal{C}}^\bullet(\mathcal{U}, \mathcal{H}) \to 0
$$
를 얻는데, 이는 일반적으로 오른쪽에서 완전하지 않다. 이 열은 사상들
$$
\check{H}^n(\mathcal{U}, \mathcal{F}) \to
\check{H}^n(\mathcal{U}, \mathcal{G}) \to
\check{H}^n(\mathcal{U}, \mathcal{H})
$$
을 정의하지만 경계 연산자
$\delta : \check{H}^n(\mathcal{U}, \mathcal{H}) \to
\check{H}^{n + 1}(\mathcal{U}, \mathcal{F})$를 정의하기에는 충분하지 않다.
실제로 그런 연산자는 일반적으로 존재하지 않는다. 그러나 코사이클
$h = (h_{i_0 \ldots i_n})$의 코호몰로지류인 원소
$\overline{h} \in \check{H}^n(\mathcal{U}, \mathcal{H})$가 주어지면,
열린 덮개들
$$
U_{i_0 \ldots i_n} = \bigcup W_{i_0 \ldots i_n, k}
$$
를 택하여 $h_{i_0 \ldots i_n}|_{W_{i_0 \ldots i_n, k}}$가
$W_{i_0 \ldots i_n, k}$ 위의 $\mathcal{G}$의 절단으로 올라가게 할 수 있다.
Topology, Lemma \ref{topology-lemma-refine-covering}에 의해(여기서 $X$가
하우스도르프이고 준콤팩트라는 가정을 사용한다) 열린 덮개
$\mathcal{V} : X = \bigcup_{j \in J} V_j$와 사상 $\alpha : J \to I$를
택하여 $V_j \subset U_{\alpha(j)}$가 되게 하고(이는 세분이다), 모든
$j_0, \ldots, j_n \in J$에 대해 어떤 $k$가 존재하여
$V_{j_0 \ldots j_n} \subset W_{\alpha(j_0) \ldots \alpha(j_n), k}$.
가 되게 할 수 있다. 복합체들의 사상
$$
\xymatrix{
0 \ar[r] &
\check{\mathcal{C}}^\bullet(\mathcal{U}, \mathcal{F}) \ar[d] \ar[r] &
\check{\mathcal{C}}^\bullet(\mathcal{U}, \mathcal{G}) \ar[d] \ar[r] &
\check{\mathcal{C}}^\bullet(\mathcal{U}, \mathcal{H}) \ar[d] \ar[r] &
0 \\
0 \ar[r] &
\check{\mathcal{C}}^\bullet(\mathcal{V}, \mathcal{F}) \ar[r] &
\check{\mathcal{C}}^\bullet(\mathcal{V}, \mathcal{G}) \ar[r] &
\check{\mathcal{C}}^\bullet(\mathcal{V}, \mathcal{H}) \ar[r] &
0
}
$$
을 얻는다. 실제로 수직 화살표들은 체흐 코호몰로지군 사이의 전이 사상들을
정의하는 데 사용한 복합체의 사상들이다. 우리가 택한 세분에 의해
$$
g_{j_0 \ldots j_n} \in
\mathcal{G}(V_{j_0 \ldots j_n}),\quad
g_{j_0 \ldots j_n} \longmapsto
h_{\alpha(j_0) \ldots \alpha(j_n)}|_{V_{j_0 \ldots j_n}}
$$
를 택할 수 있다. 코체인 $g = (g_{j_0 \ldots j_n})$는 일반적으로 코사이클이
아니지만, 그 체흐 경계 $\text{d}(g)$는
$\check{\mathcal{C}}^{n + 1}(\mathcal{V}, \mathcal{H})$에서 영으로 보내진다
(위의 가환 도식과 $h$가 코사이클이라는 사실에 의한다). 따라서
$\text{d}(g)$는 $\check{\mathcal{C}}^\bullet(\mathcal{V}, \mathcal{F})$의
코사이클이다. 이로써 다음을 정의할 수 있다.
$$
\delta(\overline{h}) = \text{d}(g)\text{의 류, 소속 군은 }
\check{H}^{n + 1}(\mathcal{V}, \mathcal{F})
$$
이제 원소 $\xi \in \check{H}^n(X, \mathcal{G})$가 주어지면 열린 덮개
$\mathcal{U}$와 원소
$\overline{h} \in \check{H}^n(\mathcal{U}, \mathcal{G})$를 택하여 체흐
코호몰로지를 정의하는 여극한에서 $\xi$로 보내지게 한다. 그런 다음 위와 같이
$\mathcal{V}$와 $g$를 택하고, $\delta(\xi)$를
$\check{H}^n(X, \mathcal{F})$ 안의 $\delta(\overline{h})$의 상으로 정한다.
이제 여러 성질을 확인해야 하는데 모두 직접적이다. 예를 들어 이 구성이
$\mathcal{U}, \overline{h}, \mathcal{V}, \alpha : J \to I, g$의 선택과
무관함을 확인해야 한다. $\text{d}(g)$의 류는 올림
$g_{i_0 \ldots i_n}$의 선택과 무관하다. 두 선택의 차이는 코바운더리이기
때문이다. $\alpha$에 대한 독립성도 성립한다.\footnote{이는 중요한 확인이다.
$\alpha$의 비유일성만이 $X$의 모든 열린 덮개 위에서 체흐 복합체의 여극한을
취하여 체흐 코호몰로지를 계산하는 복합체의 짧은 완전열을 얻는 것을 막기
때문이다.} 실제로 $\alpha$를 달리 택하면 위 도식의 수직 복합체 사상들은 서로
호모토픽이다. Section \ref{cohomology-section-refinements-cech}를 보라. 다른 선택들에
대해서는 $X$의 열린 덮개가 유한 개 주어졌을 때 이들 모두의 세분인 열린 덮개를
항상 찾을 수 있다는 사실을 사용한다. 가법성도 확인해야 하며 같은 방식으로
보인다. 마지막으로 사상 $\check{H}^n(X, -) \to H^n(X, -)$가 경계 사상들과
양립함을 확인해야 한다. 이를 위해 단사 분해들
$$
\xymatrix{
0 \ar[r] &
\mathcal{F} \ar[r] \ar[d] &
\mathcal{G} \ar[r] \ar[d] &
\mathcal{H} \ar[r] \ar[d] &
0 \\
0 \ar[r] &
\mathcal{I}_1^\bullet \ar[r] &
\mathcal{I}_2^\bullet \ar[r] &
\mathcal{I}_3^\bullet \ar[r] &
0
}
$$
을 Derived Categories, Lemma \ref{derived-lemma-injective-resolution-ses}에서처럼
택한다. 이로부터 가환 도식
$$
\xymatrix{
0 \ar[r] &
\check{\mathcal{C}}^\bullet(\mathcal{U}, \mathcal{F}) \ar[r] \ar[d] &
\check{\mathcal{C}}^\bullet(\mathcal{U}, \mathcal{F}) \ar[r] \ar[d] &
\check{\mathcal{C}}^\bullet(\mathcal{U}, \mathcal{F}) \ar[r] \ar[d] &
0 \\
0 \ar[r] &
\text{Tot}(\check{\mathcal{C}}^\bullet(\mathcal{U}, \mathcal{I}_1^\bullet))
\ar[r] &
\text{Tot}(\check{\mathcal{C}}^\bullet(\mathcal{U}, \mathcal{I}_2^\bullet))
\ar[r] &
\text{Tot}(\check{\mathcal{C}}^\bullet(\mathcal{U}, \mathcal{I}_3^\bullet))
\ar[r] &
0
}
$$
을 얻는다. 여기서 $\mathcal{U}$는 위와 같은 열린 덮개이고, 수직 사상들은
$\check{H}^n(\mathcal{U}, -) \to H^n(X, -)$를 정의하는 데 사용한 것들이다.
Lemma \ref{cohomology-lemma-cech-cohomology}를 보라. 단사 대상들의 복합체로 이루어진 열은
항별 분할 완전이므로 아래쪽 복합체는 완전하다. 따라서 코호몰로지의 경계
사상은 이 아래쪽 완전열에 대한 통상적인 절차로 계산된다. Homology, Lemma
\ref{homology-lemma-long-exact-sequence-cochain}을 보라. 체흐 코호몰로지의 경계
사상을 정의한 세분 $\mathcal{V}$로 넘어간 뒤에도 마찬가지이다. 따라서 두 경계
사상은 같은 구성을 사용하므로 일치한다(첫째 사상이 주어진 열린 덮개 위의 체흐
코호몰로지 원소에 정의되는 경우). 이로써 체흐 코호몰로지 위의
$\delta$-함자 구조의 구성과 이 구조가 통상적인 코호몰로지 위의 주어진
$\delta$-함자 구조와 양립하는 이유에 대한 논의를 마친다.

\medskip\noindent
끝으로 Lemma \ref{cohomology-lemma-injective-trivial-cech}를 적용하면 단사 층 위에서 고차
체흐 코호몰로지가 자명함을 알 수 있다. 따라서 Homology, Lemma
\ref{homology-lemma-efface-implies-universal}에 의해 체흐 코호몰로지는 보편
$\delta$-함자이다.
\end{proof}

\begin{lemma}
\label{cohomology-lemma-cohomology-of-closed}
\begin{reference}
\cite[Expose V bis, 4.1.3]{SGA4}
\end{reference}
$X$를 위상공간이라 하자. $Z \subset X$를 준콤팩트 부분집합이라 하고,
$Z$의 서로 다른 두 점은 $X$에서 서로소인 열린 근방들을 가진다고 하자.
$X$ 위의 모든 아벨 층 $\mathcal{F}$에 대해 표준 사상
$$
\colim H^p(U, \mathcal{F})
\longrightarrow
H^p(Z, \mathcal{F}|_Z)
$$
은 동형이다. 여기서 여극한은 $X$ 안에서 $Z$의 열린 근방 $U$들 위에서 취한다.
\end{lemma}

\begin{proof}
먼저 $p = 0$인 경우를 증명한다. 단사성은 $\mathcal{F}|_Z$의 정의에서 따르며
일반적으로 성립한다(임의의 위상공간 $X$의 임의의 부분집합에 대해).
이제 $s \in H^0(Z, \mathcal{F}|_Z)$라 하자. 열린집합 $U_i \subset X$를
택하여 $Z \subset \bigcup U_i$가 되게 하고, $s|_{Z \cap U_i}$가
$s_i \in \mathcal{F}(U_i)$에서 오게 할 수 있다. 따라서 열린집합
$W_{ij} \subset U_i \cap U_j$가 존재하여
$W_{ij} \cap Z = U_i \cap U_j \cap Z$이고
$s_i|_{W_{ij}} = s_j|_{W_{ij}}$이다. Topology, Lemma
\ref{topology-lemma-lift-covering-of-quasi-compact-hausdorff-subset}를 적용하면
$X$의 열린집합 $V_i$를 택하여 $V_i \subset U_i$이고
$V_i \cap V_j \subset W_{ij}$가 되게 할 수 있다. 그러므로 $s_i|_{V_i}$들은
$Z$의 열린 근방 $\bigcup V_i$ 위에서 $\mathcal{F}$의 한 절단으로 이어붙는다.

\medskip\noindent
증명을 끝내려면 $\mathcal{I}$가 $X$ 위의 단사 아벨 층일 때 $p > 0$에 대해
$H^p(Z, \mathcal{I}|_Z) = 0$임을 보이면 충분하다. 이는 짧은 완전열과 차원
이동을 사용하여 따르며 세부사항은 생략한다. 따라서 어떤 $p > 0$에 대해
$\overline{\xi}$가 $H^p(Z, \mathcal{I}|_Z)$의 원소라고 하자. Lemma
\ref{cohomology-lemma-cech-Hausdorff-quasi-compact}에 의해 $\overline{\xi}$는 $Z$의 어떤
열린 덮개 $\mathcal{V} : Z = \bigcup V_i$에 대한
$\check{H}^p(\mathcal{V}, \mathcal{I}|_Z)$에서 온다. $\overline{\xi}$가
$\check{\mathcal{C}}^p(\mathcal{V}, \mathcal{I}|_Z)$ 안의 코사이클
$\xi = (\xi_{i_0 \ldots i_p})$의 류의 상이라고 하자.

\medskip\noindent
다음 규칙으로 정의되는 부분준층 $\mathcal{I}' \subset \mathcal{I}|_Z$를 생각하자.
$$
\mathcal{I}'(V) =
\{s \in \mathcal{I}|_Z(V) \mid
\exists (U, t),\ U \subset X\text{ 열린집합},
\ t \in \mathcal{I}(U),\ V = Z \cap U,\ s = t|_{Z \cap U} \}
$$
그러면 $\mathcal{I}|_Z$는 $\mathcal{I}'$의 층화이다. 따라서 모든
$(p + 1)$-튜플 $i_0 \ldots i_p$에 대해 열린 덮개
$V_{i_0 \ldots i_p} = \bigcup W_{i_0 \ldots i_p, k}$를 택하여
$\xi_{i_0 \ldots i_p}|_{W_{i_0 \ldots i_p, k}}$가 $\mathcal{I}'$의 절단이
되게 할 수 있다. Topology, Lemma \ref{topology-lemma-refine-covering}를 적용하고
$\mathcal{V}$를 세분한 뒤에는 각 $\xi_{i_0 \ldots i_p}$가 준층
$\mathcal{I}'$의 절단이라고 가정할 수 있다.

\medskip\noindent
어떤 열린집합 $U_i \subset X$에 대해 $V_i = Z \cap U_i$로 쓰자.
$\mathcal{I}$가 플라스크이고(Lemma \ref{cohomology-lemma-injective-flasque}), 모든
$(p + 1)$-튜플 $i_0 \ldots i_p$에 대해 $\xi_{i_0 \ldots i_p}$가
$\mathcal{I}'$의 절단이므로, $V_{i_0 \ldots i_p} =
Z \cap U_{i_0 \ldots i_p}$ 위에서 $\xi_{i_0 \ldots i_p}$로 제한되는 절단
$s_{i_0 \ldots i_p} \in \mathcal{I}(U_{i_0 \ldots i_p})$를 택할 수 있다.
(단사 대상이 플라스크라는 사실을 사용하지 않고도 Topology, Lemma
\ref{topology-lemma-lift-covering-of-quasi-compact-hausdorff-subset}를 한 번 더
적용하면 된다.) $s = (s_{i_0 \ldots i_p})$를 열린 덮개
$U = \bigcup U_i$에 대응하는 코체인이라 하자. $\text{d}(\xi) = 0$이므로 절단
$\text{d}(s)_{i_0 \ldots i_{p + 1}}$은
$Z \cap U_{i_0 \ldots i_{p + 1}}$ 위에서 영으로 제한된다. 따라서 증명 처음의
설명에 의해 열린 부분집합
$W_{i_0 \ldots i_{p + 1}} \subset U_{i_0 \ldots i_{p + 1}}$을 택하여
$Z \cap W_{i_0 \ldots i_{p + 1}} = Z \cap U_{i_0 \ldots i_{p + 1}}$이고
$\text{d}(s)_{i_0 \ldots i_{p + 1}}|_{W_{i_0 \ldots i_{p + 1}}} = 0$이 되게
할 수 있다. Topology, Lemma
\ref{topology-lemma-lift-covering-of-quasi-compact-hausdorff-subset}에 의해
$U'_i \subset U_i$를 택하여 $Z \subset \bigcup U'_i$이고
$U'_{i_0 \ldots i_{p + 1}} \subset W_{i_0 \ldots i_{p + 1}}$가 되게 할 수 있다.
그러면 $s'_{i_0 \ldots i_p} = s_{i_0 \ldots i_p}|_{U'_{i_0 \ldots i_p}}$인
$s' = (s'_{i_0 \ldots i_p})$는 $Z$의 열린 근방의 열린 덮개
$U' = \bigcup U'_i$에 대한 $\mathcal{I}$의 코사이클이다. $\mathcal{I}$의 고차
체흐 코호몰로지군은 자명하므로(Lemma \ref{cohomology-lemma-injective-trivial-cech}),
$s'$은 코바운더리이다. 따라서 열린 덮개 $Z = \bigcup Z \cap U'_i$에 대한 체흐
복합체에서 $\xi$의 상도 코바운더리이며 증명이 끝난다.
\end{proof}






\section{기저변환 사상}
\label{cohomology-section-base-change-map}

\noindent
몇몇 경우에는 기저변환 사상을 구성하는 방법이 필요하다. 아직 유도 당김을
논의하지 않았으므로, 여기서는 환 달린 공간의 평탄 사상에 의한 기저변환만
다룬다. 결과를 서술하기 전에 유도 범주 위의 평탄 당김을 논의하자. 즉
$g : X \to Y$가 환 달린 공간의 평탄 사상이라고 하자. Modules, Lemma
\ref{modules-lemma-pullback-flat}에 의해 함자
$g^* : \textit{Mod}(\mathcal{O}_Y) \to \textit{Mod}(\mathcal{O}_X)$는
완전하다. 따라서 유도함자
$$
g^* : D^{+}(Y) \to D^{+}(X)
$$
를 가지며, 이는 $D^{+}(Y)$의 주어진 대상의 대표 복합체를 단순히 당겨서
계산한다. Derived Categories, Lemma
\ref{derived-lemma-right-derived-exact-functor}를 보라. 따라서 앞서 나타낸 대로
이 함자를 $Lg^*$가 아니라 $g^*$로 나타낸다.

\begin{lemma}
\label{cohomology-lemma-base-change-map-flat-case}
다음 가환 도식의 환 달린 공간들이 주어졌다고 하자.
$$
\xymatrix{
X' \ar[r]_{g'} \ar[d]_{f'} &
X \ar[d]^f \\
S' \ar[r]^g &
S
}
$$
$\mathcal{F}^\bullet$을 $\mathcal{O}_X$-가군의 아래로 유계인 복합체라 하자.
$g$와 $g'$이 모두 평탄하다고 가정하자. 그러면 $D^{+}(S')$에 표준 기저변환 사상
$$
g^*Rf_*\mathcal{F}^\bullet
\longrightarrow
R(f')_*(g')^*\mathcal{F}^\bullet
$$
이 존재한다.
\end{lemma}

\begin{proof}
단사 분해 $\mathcal{F}^\bullet \to \mathcal{I}^\bullet$과
$(g')^*\mathcal{F}^\bullet \to \mathcal{J}^\bullet$을 택하자. Lemma
\ref{cohomology-lemma-pushforward-injective-flat}에 의해 $(g')_*\mathcal{J}^\bullet$은
$R(g')_*(g')^*\mathcal{F}^\bullet$을 표현하는 단사 대상들의 복합체이다. 따라서
Derived Categories, Lemmas \ref{derived-lemma-morphisms-lift} 및
\ref{derived-lemma-morphisms-equal-up-to-homotopy}에 의해 도식의 화살표 $\beta$
$$
\xymatrix{
(g')_*(g')^*\mathcal{F}^\bullet \ar[r] &
(g')_*\mathcal{J}^\bullet \\
\mathcal{F}^\bullet \ar[u]^{adjunction} \ar[r] &
\mathcal{I}^\bullet \ar[u]_\beta
}
$$
가 존재하며 호모토피까지 유일하다. $S$로 직접상을 취하면
$$
f_*\beta :
f_*\mathcal{I}^\bullet
\longrightarrow
f_*(g')_*\mathcal{J}^\bullet
=
g_*(f')_*\mathcal{J}^\bullet
$$
을 얻는다. $g^*$와 $g_*$의 수반성에 의해 복합체의 사상
$g^*f_*\mathcal{I}^\bullet \to (f')_*\mathcal{J}^\bullet$을 얻는다. 전체 과정에서
유일한 선택은 사상 $\beta$의 선택이었고 모든 일을 복합체 수준에서 했으므로,
이 사상은 호모토피까지 유일함에 유의하라.
\end{proof}

\begin{remark}
\label{cohomology-remark-correct-version-base-change-map}
기저변환 사상의 ``올바른'' 판은 사상
$$
Lg^* Rf_* \mathcal{F}^\bullet
\longrightarrow
R(f')_* L(g')^*\mathcal{F}^\bullet.
$$
이다. 이 사상의 구성에는 비유계 복합체가 관여한다. Remark
\ref{cohomology-remark-base-change}를 보라.
\end{remark}





\section{위상수학에서의 고유 기저변환}
\label{cohomology-section-proper-base-change}

\noindent
이 절에서는 위상수학의 고유 기저변환 정리의 매우 일반적인 판을 증명한다.
이 정리에 따르면 고차 직접상 $R^pf_*$의 줄기는 섬유 위에서 계산할 수 있다.

\begin{lemma}
\label{cohomology-lemma-proper-base-change}
$f : (X, \mathcal{O}_X) \to (Y, \mathcal{O}_Y)$를 환 달린 공간의 사상이라 하고,
$y \in Y$라 하자. 다음을 가정하자.
\begin{enumerate}
\item $f$는 닫힌 사상이다.
\item $f$는 분리 사상이다.
\item $f^{-1}(y)$는 준콤팩트이다.
\end{enumerate}
그러면 $D^+(\mathcal{O}_X)$의 $E$에 대해 $D^+(\mathcal{O}_{Y, y})$에서
$(Rf_*E)_y = R\Gamma(f^{-1}(y), E|_{f^{-1}(y)})$이다.
\end{lemma}

\begin{proof}
Lemma \ref{cohomology-lemma-base-change-map-flat-case}의 기저변환 사상은 표준 사상
$(Rf_*E)_y \to R\Gamma(f^{-1}(y), E|_{f^{-1}(y)})$를 준다. 이 사상이
동형임을 증명하기 위해 $E$를 단사 대상들의 아래로 유계인 복합체
$\mathcal{I}^\bullet$로 표현하자. $Z = f^{-1}(\{y\})$로 두자. Lemma
\ref{cohomology-lemma-cohomology-of-closed}의 가정들이 만족된다. Topology, Lemma
\ref{topology-lemma-separated}를 보라. 따라서 제한 $\mathcal{I}^n|_Z$는
$\Gamma(Z, -)$에 대해 비순환이다. 그러므로 $R\Gamma(Z, E|_Z)$는 복합체
$\Gamma(Z, \mathcal{I}^\bullet|_Z)$로 표현된다. Derived Categories, Lemma
\ref{derived-lemma-leray-acyclicity}를 보라. 다시 말해 사상
$$
\colim_V \mathcal{I}^\bullet(f^{-1}(V))
\longrightarrow
\Gamma(Z, \mathcal{I}^\bullet|_Z)
$$
이 동형임을 보여야 한다. Lemma \ref{cohomology-lemma-cohomology-of-closed}를 사용하면
$Z = f^{-1}(\{y\})$의 열린 근방 $f^{-1}(V)$들의 모임이 모든 열린 근방의 계에서
공종임을 보이면 충분하다. $f^{-1}(\{y\}) \subset U$가 열린 근방이면 $f$가
닫힌 사상이므로 $V = Y \setminus f(X \setminus U)$는 $y$의 열린 근방이고
$f^{-1}(V) \subset U$이다. 이로써 보조정리가 증명된다.
\end{proof}

\begin{theorem}[고유 기저변환]
\label{cohomology-theorem-proper-base-change}
\begin{reference}
\cite[Expose V bis, 4.1.1]{SGA4}
\end{reference}
위상공간들의 데카르트 정사각형
$$
\xymatrix{
X' = Y' \times_Y X \ar[d]_{f'} \ar[r]_-{g'} & X \ar[d]^f \\
Y' \ar[r]^g & Y
}
$$
을 생각하자. $f$가 고유 사상이라고 가정하자. $E$를 $D^+(X)$의 대상이라 하자.
그러면 Lemma \ref{cohomology-lemma-base-change-map-flat-case}의 기저변환 사상
$$
g^{-1}Rf_*E \longrightarrow Rf'_*(g')^{-1}E
$$
은 $D^+(Y')$에서 동형이다.
\end{theorem}

\begin{proof}
$y' \in Y'$을 상이 $y \in Y$인 점이라 하자. 기저변환 사상이 $y'$에서의 줄기
위에 동형을 유도함을 보이면 충분하다. $f$가 고유 사상이므로 $f'$도 고유 사상이고,
$f$와 $f'$의 섬유들은 준콤팩트이며 $f$와 $f'$은 닫힌 사상이다. Topology,
Theorem \ref{topology-theorem-characterize-proper} 및 Lemma
\ref{topology-lemma-base-change-separated}를 보라. 따라서 Lemma
\ref{cohomology-lemma-proper-base-change}를 두 번 적용하여
$$
(Rf'_*(g')^{-1}E)_{y'} = R\Gamma((f')^{-1}(y'), (g')^{-1}E|_{(f')^{-1}(y')})
$$
그리고
$$
(Rf_*E)_y = R\Gamma(f^{-1}(y), E|_{f^{-1}(y)})
$$
임을 알 수 있다. 유도된 섬유 사상 $(f')^{-1}(y') \to f^{-1}(y)$는 위상공간의
위상동형이고, $E|_{f^{-1}(y)}$의 당김은
$(g')^{-1}E|_{(f')^{-1}(y')}$이다. 원하는 결과가 따른다.
\end{proof}

\begin{lemma}[집합의 층에 대한 고유 기저변환]
\label{cohomology-lemma-proper-base-change-sheaves-of-sets}
위상공간들의 데카르트 정사각형
$$
\xymatrix{
X' \ar[d]_{f'} \ar[r]_-{g'} & X \ar[d]^f \\
Y' \ar[r]^g & Y
}
$$
을 생각하자. $f$가 고유 사상이라고 가정하자. 그러면 $X$ 위의 임의의 집합층
$\mathcal{F}$에 대해
$g^{-1}f_*\mathcal{F} = f'_*(g')^{-1}\mathcal{F}$
이다.
\end{lemma}

\begin{proof}
Theorem \ref{cohomology-theorem-proper-base-change}의 증명과 똑같이 논증하면
$(f_*\mathcal{F})_y = \Gamma(X_y, \mathcal{F}|_{X_y})$임을 보이는 것으로
충분하다. 이어 Lemma \ref{cohomology-lemma-proper-base-change}에서처럼 논증하면 이는
집합의 층에 대한 Lemma \ref{cohomology-lemma-cohomology-of-closed}의 $p = 0$인 경우로
환원된다. Lemma \ref{cohomology-lemma-cohomology-of-closed}의 증명의 첫 부분은 집합의
층에 대해서도 성립하므로 증명이 끝난다. 몇 가지 세부사항은 생략한다.
\end{proof}



\section{코호몰로지와 여극한}
\label{cohomology-section-limits}

\noindent
$X$를 환 달린 공간이라 하자. $(\mathcal{F}_i, \varphi_{ii'})$를 유향집합
$I$ 위의 $\mathcal{O}_X$-가군층의 계라 하자. Categories, Section
\ref{categories-section-posets-limits}를 보라. 각 $i$에 대해 표준 사상
$\mathcal{F}_i \to \colim_i \mathcal{F}_i$가 있으므로, 모든 $p \geq 0$에 대해
표준 사상
$$
\colim_i H^p(X, \mathcal{F}_i)
\longrightarrow
H^p(X, \colim_i \mathcal{F}_i)
$$
을 얻는다. 물론 모든 열린집합 $U \subset X$에 대해서도 비슷한 사상이 있다.
이 사상들은 일반적으로 동형이 아니며 $p = 0$일 때조차 그렇다. 이 절에서는
Sheaves, Lemma \ref{sheaves-lemma-directed-colimits-sections}의 결과를
일반화한다. 특수한 경우 $\mathcal{G} = \mathcal{O}_X$에 대해서는 Modules,
Lemma \ref{modules-lemma-finite-presentation-quasi-compact-colimit}도 보라.

\begin{lemma}
\label{cohomology-lemma-quasi-separated-cohomology-colimit}
$X$를 환 달린 공간이라 하자. $X$의 바탕 위상공간이 다음 성질들을 만족한다고
가정하자.
\begin{enumerate}
\item 준콤팩트 열린 부분집합들의 기저가 존재한다.
\item 임의의 두 준콤팩트 열린집합의 교집합은 준콤팩트이다.
\end{enumerate}
그러면 $\mathcal{O}_X$-가군층의 임의의 유향계
$(\mathcal{F}_i, \varphi_{ii'})$와 임의의 준콤팩트 열린집합 $U \subset X$에 대해
표준 사상
$$
\colim_i H^q(U, \mathcal{F}_i)
\longrightarrow
H^q(U, \colim_i \mathcal{F}_i)
$$
은 모든 $q \geq 0$에 대해 동형이다.
\end{lemma}

\begin{proof}
이 증명에서는 모든 준콤팩트 열린집합 $U \subset X$에 대해 동시에 논증하는
것이 중요하다. Sheaves, Lemma
\ref{sheaves-lemma-directed-colimits-sections}와 Topology, Lemma
\ref{topology-lemma-topology-quasi-separated-scheme}를 함께 사용하면 $q = 0$과
임의의 준콤팩트 열린집합 $U \subset X$에 대해 결과가 성립한다. 모든
$q \leq q_0$에 대해 결과를 증명했다고 가정하고 $q = q_0 + 1$인 경우를 증명하자.

\medskip\noindent
유향계에 대한 우리의 관례에 따라 지표집합 $I$는 유향집합이고, $I$ 위의
임의의 $\mathcal{O}_X$-가군의 계 $(\mathcal{F}_i, \varphi_{ii'})$는 유향계이다.
Injectives, Lemma \ref{injectives-lemma-sheaves-modules-space}에 의해
$\mathcal{O}_X$-가군의 범주에는 함자적 단사 매입들이 있다. 따라서 임의의 계
$(\mathcal{F}_i, \varphi_{ii'})$에 대해 계 $(\mathcal{I}_i, \varphi_{ii'})$와
단사 $\mathcal{O}_X$-가군 사상 $\mathcal{F}_i \to \mathcal{I}_i$들로 주어지는
계의 사상이 존재하며, 각 $\mathcal{I}_i$는 단사 $\mathcal{O}_X$-가군이다.
여핵을 $\mathcal{Q}_i$로 나타내면 짧은 완전열들
$$
0 \to
\mathcal{F}_i \to
\mathcal{I}_i \to
\mathcal{Q}_i \to 0.
$$
을 얻는다. 또한 열
$$
0 \to
\colim_i \mathcal{F}_i \to
\colim_i \mathcal{I}_i \to
\colim_i \mathcal{Q}_i \to 0.
$$
도 $\mathcal{O}_X$-가군의 짧은 완전열이라고 주장한다. 이는 줄기에서 확인할 수
있다. Sheaves, Sections \ref{sheaves-section-limits-presheaves} 및
\ref{sheaves-section-limits-sheaves}에 의해 줄기를 취하는 일은 여극한과
가환한다. 아벨군의 짧은 완전열들의 유향 여극한은 짧은 완전열이므로(Algebra,
Lemma \ref{algebra-lemma-directed-colimit-exact} 참조) 결과가 따른다. 모든
준콤팩트 열린집합 $U \subset X$와 모든 $q \geq 1$에 대해
$H^q(U, \colim_i \mathcal{I}_i) = 0$이라고 주장한다. 우선 이 주장을 받아들이고
도식
$$
\xymatrix{
\colim_i H^{q_0}(U, \mathcal{I}_i) \ar[d] \ar[r] &
\colim_i H^{q_0}(U, \mathcal{Q}_i) \ar[d] \ar[r] &
\colim_i H^{q_0 + 1}(U, \mathcal{F}_i) \ar[d] \ar[r] &
0 \ar[d] \\
H^{q_0}(U, \colim_i \mathcal{I}_i) \ar[r] &
H^{q_0}(U, \colim_i \mathcal{Q}_i) \ar[r] &
H^{q_0 + 1}(U, \colim_i \mathcal{F}_i) \ar[r] &
0
}
$$
을 생각하자. 오른쪽 아래의 영은 위 주장으로부터, 오른쪽 위의 영은 층
$\mathcal{I}_i$들이 단사라는 사실로부터 나온다. Algebra, Lemma
\ref{algebra-lemma-directed-colimit-exact}를 적용하면 위 행은 완전하다. 따라서
뱀 보조정리에 의해 $q = q_0 + 1$인 경우의 결과를 얻는다.

\medskip\noindent
이제 주장을 증명하면 된다. Lemma \ref{cohomology-lemma-cech-vanish-basis}를 사용할 것이다.
$\mathcal{B}$를 $X$의 모든 준콤팩트 열린 부분집합의 모임이라 하자. 가정에 따라
이는 $X$의 위상의 기저이다. $\text{Cov}$를 유한 열린 덮개
$\mathcal{U} : U = \bigcup_{j = 1, \ldots, m} U_j$ 가운데 $U$와 각 $U_j$가
$X$의 준콤팩트 열린집합인 것들의 모임이라 하자. $q = 0$인 경우의 결과에 의해
$\mathcal{U} \in \text{Cov}$에 대해
$$
\check{\mathcal{C}}^\bullet(\mathcal{U}, \colim_i \mathcal{I}_i)
=
\colim_i \check{\mathcal{C}}^\bullet(\mathcal{U}, \mathcal{I}_i)
$$
이다. 모든 다중 교집합 $U_{j_0 \ldots j_p}$가 준콤팩트이기 때문이다. Lemma
\ref{cohomology-lemma-injective-trivial-cech}에 의해 체흐 복합체의 여극한 안의 각 복합체는
차수 $\geq 1$에서 비순환이다. 따라서 Algebra, Lemma
\ref{algebra-lemma-directed-colimit-exact}에 의해 체흐 복합체
$\check{\mathcal{C}}^\bullet(\mathcal{U}, \colim_i \mathcal{I}_i)$도 차수
$\geq 1$에서 비순환이다. 다시 말해 모든 $p \geq 1$에 대해
$\check{H}^p(\mathcal{U},  \colim_i \mathcal{I}_i) = 0$이다. 그러므로 Lemma
\ref{cohomology-lemma-cech-vanish-basis}의 가정들이 만족되어 주장이 따른다.
\end{proof}

\begin{lemma}
\label{cohomology-lemma-higher-direct-image-colimit}
$f : X \to Y$를 위상공간의 연속 사상이라 하자.
$(\mathcal{F}_i, \varphi_{ii'})$를 $X$ 위의 아벨 층의 계라 하고
$\mathcal{F} = \colim \mathcal{F}_i$로 두자. $p \geq 0$을 정수라 하자.
$H^p(f^{-1}(V), \mathcal{F}) = \colim H^p(f^{-1}(V), \mathcal{F}_i)$를 만족하는
열린집합 $V \subset Y$들의 모임이 $Y$의 위상의 기저라고 가정하자. 그러면
$R^pf_*\mathcal{F} = \colim R^pf_*\mathcal{F}_i$이다.
\end{lemma}

\begin{proof}
$R^pf_*\mathcal{F}$는 $V$에 $H^p(f^{-1}(V), \mathcal{F})$를 대응시키는 준층
$\mathcal{G}$의 층화임을 상기하자. Lemma
\ref{cohomology-lemma-describe-higher-direct-images}를 보라. 마찬가지로
$R^pf_*\mathcal{F}_i$는 $V$에 $H^p(f^{-1}(V), \mathcal{F}_i)$를 대응시키는
준층 $\mathcal{G}_i$의 층화이다. 층화는 층에서 준층으로 가는 포함의 왼쪽
수반임을 상기하자. Sheaves, Section \ref{sheaves-section-sheafification}을 보라.
따라서 층화는 여극한과 가환한다. Categories, Lemma
\ref{categories-lemma-adjoint-exact}를 보라. 그러므로 준층의 사상(여극한은
준층의 범주에서 취한다)
$$
\colim \mathcal{G}_i \longrightarrow \mathcal{G}
$$
이 층화 위에 동형을 유도함을 보이면 충분하다. 이를 위해서는 준층
$\mathcal{G}$와 $\colim \mathcal{G}_i$가 $Y$의 위상의 한 기저 위에서
일치함을 보이면 충분하다. 실제로 이 경우 그 층화들의 줄기는 기저 원소들 위의
준층 값에서 직접 계산할 수 있으므로 일치한다. 필요한 일치는 정확히 보조정리의
가정이다.
\end{proof}

\noindent
다음으로 Sheaves, Lemma \ref{sheaves-lemma-descend-opens}의 코호몰로지에 대한
대응물을 서술한다. $X$를 스펙트럴 공간이라 하고, Topology, Lemma
\ref{topology-lemma-directed-inverse-limit-spectral-spaces}에서처럼 스펙트럴 전이
사상들을 갖는 도식에 대해 스펙트럴 공간 $X_i$들의 공여과 극한으로 주어졌다고
하자. 다음이 주어졌다고 가정하자.
\begin{enumerate}
\item 모든 $i \in \Ob(\mathcal{I})$에 대해 $X_i$ 위의 아벨 층
$\mathcal{F}_i$가 주어진다.
\item $a : j \to i$에 대해 아벨 층의 $f_a$-사상
$\varphi_a : \mathcal{F}_i \to \mathcal{F}_j$가 주어진다(Sheaves, Definition
\ref{sheaves-definition-f-map} 참조).
\end{enumerate}
$c = a \circ b$일 때마다 $\varphi_c = \varphi_b \circ \varphi_a$라고 가정하고,
$X$ 위에서 $\mathcal{F} = \colim p_i^{-1}\mathcal{F}_i$로 두자.

\begin{lemma}
\label{cohomology-lemma-colimit}
위에서 논의한 상황에서 $i \in \Ob(\mathcal{I})$라 하고,
$U_i \subset X_i$를 준콤팩트 열린집합이라 하자. 그러면
$$
\colim_{a : j \to i} H^p(f_a^{-1}(U_i), \mathcal{F}_j) =
H^p(p_i^{-1}(U_i), \mathcal{F})
$$
이 모든 $p \geq 0$에 대해 성립한다. 특히
$H^p(X, \mathcal{F}) = \colim H^p(X_i, \mathcal{F}_i)$이다.
\end{lemma}

\begin{proof}
$p = 0$인 경우는 Sheaves, Lemma \ref{sheaves-lemma-descend-opens}이다.

\medskip\noindent
이 문단에서는 $\mathcal{G}_i$가 단사 아벨 층이고 $\gamma_i$가 단사인 계의 사상
$(\gamma_i) : (\mathcal{F}_i, \varphi_a) \to (\mathcal{G}_i, \psi_a)$를 찾을 수
있음을 보인다. 각 $i$에 대해 단사 사상
$\mathcal{F}_i \to \mathcal{I}_i$를 택하자. 여기서 $\mathcal{I}_i$는 $X_i$ 위의
단사 아벨 층이다. 그러면 사상들의 족
$$
\gamma_i :
\mathcal{F}_i
\longrightarrow
\prod\nolimits_{b : k \to i} f_{b, *}\mathcal{I}_k = \mathcal{G}_i
$$
을 생각할 수 있다. 그 성분 사상들은 사상
$f_b^{-1}\mathcal{F}_i \to \mathcal{F}_k \to \mathcal{I}_k$에 수반인 사상들이다.
$\mathcal{I}$ 안의 $a : j \to i$에 대해 표준 사상
$$
\psi_a : f_a^{-1}\mathcal{G}_i \to \mathcal{G}_j
$$
이 존재하며, 그 성분들은 표준 사상
$f_b^{-1}f_{a \circ b, *}\mathcal{I}_k \to f_{b, *}\mathcal{I}_k$
들이다. 여기서 $b : k \to j$이다. 따라서 계들의 단사 사상
$\{\gamma_i\} : \{\mathcal{F}_i, \varphi_a) \to (\mathcal{G}_i, \psi_a)$
을 얻는다. $\mathcal{G}_i$는 $X_i$ 위의 단사 아벨군층이다. Lemma
\ref{cohomology-lemma-pushforward-injective-flat}과 Homology, Lemma
\ref{homology-lemma-product-injectives}를 보라. 이로써 구성이 끝난다.

\medskip\noindent
Lemma \ref{cohomology-lemma-quasi-separated-cohomology-colimit}의 증명과 똑같이 논증하면
$p > 0$에 대해 $H^p(X, \colim f_i^{-1}\mathcal{G}_i) = 0$임을 증명하는 것으로
충분함을 알 수 있다.

\medskip\noindent
$\mathcal{G} = \colim f_i^{-1}\mathcal{G}_i$로 두자. $X$의 모든 준콤팩트
열린집합 위에서 $\mathcal{G}$의 코호몰로지가 소멸함을 보이려면, 준콤팩트
열린집합을 유한 개의 준콤팩트 열린집합으로 덮는 $X$의 임의의 열린 덮개
$\mathcal{U}$에 대해 $\mathcal{G}$의 체흐 코호몰로지가 영임을 보이면 충분하다.
Lemma \ref{cohomology-lemma-cech-vanish-basis}를 보라. Topology, Lemma
\ref{topology-lemma-descend-opens}에 의해 이러한 덮개는 어떤 $i$에 대해 공간
$X_i$ 위의 같은 종류의 덮개 $\mathcal{U}_i$의 $p_i$에 의한 역상이다. $p = 0$인
경우에 의해
$$
\check{\mathcal{C}}^\bullet(\mathcal{U}, \mathcal{G}) =
\colim_{a : j \to i}
\check{\mathcal{C}}^\bullet(f_a^{-1}(\mathcal{U}_i), \mathcal{G}_j)
$$
이다. 오른쪽은 양의 차수에서 비순환인 복합체들의 여과 여극한이다. Lemma
\ref{cohomology-lemma-injective-trivial-cech}를 보라. 따라서 Algebra, Lemma
\ref{algebra-lemma-directed-colimit-exact}에 의해 결론을 얻는다.
\end{proof}









\section{뇌터 위상공간에서의 소멸}
\label{cohomology-section-vanishing-Noetherian}

\noindent
목표는 Grothendieck의 한 정리, 즉
Proposition \ref{cohomology-proposition-vanishing-Noetherian}을 증명하는 것이다.
\cite{Tohoku}를 보라.

\begin{lemma}
\label{cohomology-lemma-cohomology-and-closed-immersions}
$i : Z \to X$를 위상공간의 닫힌 몰입이라 하자.
$Z$ 위의 임의의 아벨 층 $\mathcal{F}$에 대해
$H^p(Z, \mathcal{F}) = H^p(X, i_*\mathcal{F})$이다.
\end{lemma}

\begin{proof}
$i_*$가 완전하기 때문에 이는 참이다(Modules, Lemma
\ref{modules-lemma-i-star-exact}를 보라). 따라서 함자로서
$R^pi_* = 0$이다(Derived Categories, Lemma
\ref{derived-lemma-right-derived-exact-functor}).
그러므로 Lemma \ref{cohomology-lemma-apply-Leray}를 적용할 수 있다.
\end{proof}

\begin{lemma}
\label{cohomology-lemma-irreducible-constant-cohomology-zero}
$X$를 기약 위상공간이라 하자. 그러면 임의의 아벨 군 $A$와
모든 $p > 0$에 대해 $H^p(X, \underline{A}) = 0$이다.
\end{lemma}

\begin{proof}
$\underline{A}$는 Sheaves, Definition
\ref{sheaves-definition-constant-sheaf}에서 정의한 상수층임을 상기하자.
$X$가 기약이므로 공집합이 아닌 임의의 열린집합 $U$는 기약이고, 특히
연결이다. 따라서 공집합이 아닌 열린집합 $U \subset X$에 대해
$\underline{A}(U) = A$이다. 또한 $\underline{A}(\emptyset) = 0$이다.
그러므로 $\underline{A}$는 $X$ 위의 플라스크 아벨 층이다.
소멸은 Lemma \ref{cohomology-lemma-flasque-acyclic}에서 따른다.
\end{proof}

\begin{lemma}
\label{cohomology-lemma-subsheaf-of-constant-sheaf}
\begin{reference}
\cite[Page 168]{Tohoku}.
\end{reference}
$X$를 임의의 두 준콤팩트 열린집합의 교집합이 준콤팩트인 위상공간이라
하자. $\mathcal{F} \subset \underline{\mathbf{Z}}$를 준콤팩트
열린집합들 위의 유한 개 절단으로 생성되는 부분층이라 하자. 그러면
아벨 부분층들에 의한 유한 여과
$$
(0) = \mathcal{F}_0 \subset \mathcal{F}_1 \subset \ldots \subset
\mathcal{F}_n = \mathcal{F}
$$
가 존재하며, 각 $0 < i \leq n$에 대해 짧은 완전열
$$
0 \to j'_!\underline{\mathbf{Z}}_V \to j_!\underline{\mathbf{Z}}_U \to
\mathcal{F}_i/\mathcal{F}_{i - 1} \to 0
$$
이 존재한다. 여기서 $j : U \to X$와 $j' : V \to X$는 준콤팩트
열린집합을 $X$에 넣는 포함이다.
\end{lemma}

\begin{proof}
$\mathcal{F}$가 준콤팩트 열린집합 $U_1, \ldots, U_t$ 위의 절단
$s_1, \ldots, s_t$로 생성된다고 하자. $U_i$는 준콤팩트이고 $s_i$는
$\mathbf{Z}$로 가는 국소 상수 함수이므로, 필요하면 $U_i$를 열린닫힌
부분집합들로 이루어진 유한 분해의 조각들로 바꾸어 $s_i$가 상수 절단이라고
가정할 수 있다. $n_i \in \mathbf{Z}$에 대해 $s_i = n_i$라 하자.
물론 $n_i = 0$이면 목록에서 $(U_i, n_i)$를 제거할 수 있다. 필요하면
부호를 뒤집어 $n_i > 0$이라고 가정할 수도 있다. 다음으로 임의의 부분집합
$I \subset \{1, \ldots, t\}$에 대해
$\bigcap_{i \in I} U_i$와 $\gcd(n_i, i \in I)$를 목록에 더할 수 있다.
이 작업을 마치면 목록 $(U_1, n_1), \ldots, (U_t, n_t)$가 다음 성질을
만족함을 알 수 있다. $x \in X$에 대해
$I_x = \{i \in \{1, \ldots, t\} \mid x \in U_i\}$로 두자. 그러면
$\gcd(n_i, i \in I_x)$는 어떤 $i \in I_x$에 대한 $n_i$로 실현된다.

\medskip\noindent
여과로 $\mathcal{F}_0 = (0)$을 취하고, $n_i \leq n$인 $i$들에 대한
$U_i$ 위의 절단 $n_i$들로 생성되는 층을 $\mathcal{F}_n$으로 취한다.
$n \gg 0$이면 $\mathcal{F}_n = \mathcal{F}$임은 명백하다. 또한 몫
$\mathcal{F}_n/\mathcal{F}_{n - 1}$은
$U = \bigcup_{n_i \leq n} U_i$ 위의 절단 $n$으로 생성되고,
사상 $j_!\underline{\mathbf{Z}}_U \to
\mathcal{F}_n/\mathcal{F}_{n - 1}$의 핵은
$V = \bigcup_{n_i \leq n - 1} U_i$ 위의 절단 $n$으로 생성된다.
따라서 보조정리의 명제에 나오는 짧은 완전열을 얻는다.
\end{proof}

\begin{lemma}
\label{cohomology-lemma-vanishing-generated-one-section}
\begin{reference}
이는 \cite[Proposition 3.6.1]{Tohoku}의 특수한 경우이다.
\end{reference}
$X$를 위상공간이라 하고 $d \geq 0$을 정수라 하자. 다음을 가정하자.
\begin{enumerate}
\item $X$는 준콤팩트이다.
\item 준콤팩트 열린집합들은 $X$의 기저를 이룬다.
\item 두 준콤팩트 열린집합의 교집합은 준콤팩트이다.
\item 임의의 준콤팩트 열린집합 $j : U \to X$와 모든 $p > d$에 대해
$H^p(X, j_!\underline{\mathbf{Z}}_U) = 0$이다.
\end{enumerate}
그러면 $X$ 위의 임의의 아벨 층 $\mathcal{F}$와 모든 $p > d$에 대해
$H^p(X, \mathcal{F}) = 0$이다.
\end{lemma}

\begin{proof}
$U$가 $X$의 준콤팩트 열린집합들을 달릴 때
$S = \coprod_{U \subset X} \mathcal{F}(U)$로 두자.
각 유한 부분집합 $A = \{s_1, \ldots, s_n\} \subset S$에 대해
모든 $s_i$로 생성되는 $\mathcal{F}$의 부분층을 $\mathcal{F}_A$라 하자
(Modules, Definition
\ref{modules-definition-generated-by-local-sections}를 보라).
$A \subset A'$이면 $\mathcal{F}_A \subset \mathcal{F}_{A'}$임에
유의하자. 따라서 $\{\mathcal{F}_A\}$는 $S$의 유한 부분집합들로 이루어진
유향 부분순서집합 위의 계를 이룬다. Modules, Lemma
\ref{modules-lemma-generated-by-local-sections-stalk}에 의해 줄기에서
보면
$$
\colim_A \mathcal{F}_A = \mathcal{F}
$$
임이 명백하다. Lemma \ref{cohomology-lemma-quasi-separated-cohomology-colimit}에
의해
$$
H^p(X, \mathcal{F}) =
\colim_A H^p(X, \mathcal{F}_A)
$$
이다. 따라서 아벨 층 $\mathcal{F}_A$에 대해 소멸을 증명하면 충분하다.
다시 말해, $\mathcal{F}$가 $X$의 준콤팩트 열린집합 위의 유한 개 국소
절단으로 생성되는 경우에 결과를 증명하면 충분하다.

\medskip\noindent
$\mathcal{F}$가 국소 절단 $s_1, \ldots, s_n$으로 생성된다고 하자.
$s_1, \ldots, s_{n - 1}$로 생성되는 부분층을
$\mathcal{F}' \subset \mathcal{F}$라 하자. 그러면 짧은 완전열
$$
0 \to \mathcal{F}' \to \mathcal{F} \to \mathcal{F}/\mathcal{F}' \to 0
$$
을 얻는다. 코호몰로지의 긴 완전열에서, $n$개보다 적은 국소 절단으로
생성되는 아벨 층 $\mathcal{F}'$와 $\mathcal{F}/\mathcal{F}'$에 대해
소멸을 증명하면 충분함을 알 수 있다. 따라서 많아야 하나의 국소 절단으로
생성되는 층에 대해 소멸을 증명하면 충분하다. 이러한 층들은 정확히
$U$가 $X$의 준콤팩트 열린집합일 때 층
$j_!\underline{\mathbf{Z}}_U$의 몫들이다.

\medskip\noindent
이제 $U$가 $X$의 준콤팩트 열린집합이고 짧은 완전열
$$
0 \to \mathcal{K} \to j_!\underline{\mathbf{Z}}_U \to \mathcal{F} \to 0
$$
이 주어졌다고 하자. $q \geq d + 1$에 대해
$H^q(X, \mathcal{K})$가 영임을 보이면 충분하다. 위와 같이
$\mathcal{K}$를 준콤팩트 열린집합 위의 유한 개 절단으로 생성되는
부분층 $\mathcal{K}'$들의 여과 여극한으로 쓸 수 있다. 그러면
$\mathcal{F}$는 층
$j_!\underline{\mathbf{Z}}_U/\mathcal{K}'$들의 여과 여극한이다.
이렇게 하여 $\mathcal{K}$가 준콤팩트 열린집합 위의 유한 개 절단으로
생성되는 경우로 환원된다. $\mathcal{K}$가
$\underline{\mathbf{Z}}_X$의 부분층임에 유의하자. 따라서 Lemma
\ref{cohomology-lemma-subsheaf-of-constant-sheaf}에 의해 $\mathcal{K}$에는
연속 몫 $\mathcal{Q}$들이 짧은 완전열
$$
0 \to j''_!\underline{\mathbf{Z}}_W \to
j'_!\underline{\mathbf{Z}}_V \to \mathcal{Q} \to 0
$$
에 들어가는 유한 여과가 존재한다. 여기서 $j'' : W \to X$와
$j' : V \to X$는 준콤팩트 열린집합의 포함이다. 그러므로
$p > d$에 대한 $H^p(X, \mathcal{Q})$의 소멸은 보조정리에서 가정한
$j''_!\underline{\mathbf{Z}}_W$와
$j'_!\underline{\mathbf{Z}}_V$의 코호몰로지 군의 소멸에서 따른다.
$\mathcal{K}$로 돌아가면, 긴 코호몰로지 완전열을 사용하는 귀납 논증을
통해 이것은 원하는 소멸도 함의한다.
\end{proof}

\begin{example}
\label{cohomology-example-datta}
$X = \mathbf{N}$에 다음 위상을 주자. 열린집합들은
$\emptyset$, $X$, 그리고 $n \geq 1$에 대해
$U_n = \{i \mid i \leq n\}$이다. $X$ 위의 아벨 층 $\mathcal{F}$는
아벨 군의 역계 $A_n = \mathcal{F}(U_n)$와 같은 것이며,
$\Gamma(X, \mathcal{F}) = \lim A_n$이다. 역극한 함자는 역계의
범주에서 완전함자가 아니므로, 영이 아닌 $H^1$을 갖는 아벨 층이
존재함을 알 수 있다. 끝으로 독자는 $j : U = U_n \to X$가 포함일 때
$p \geq 1$에 대해 $H^p(X, j_!\mathbf{Z}_U) = 0$임을 확인할 수 있다.
따라서 $X$는 $d = 0$에 대해 Lemma
\ref{cohomology-lemma-vanishing-generated-one-section}의 조건 (2), (3), (4)를
만족하지만 결론은 만족하지 않는 공간의 예이다.
\end{example}

\begin{lemma}
\label{cohomology-lemma-subsheaf-irreducible}
$X$를 기약 위상공간이라 하자. $\mathcal{H} \subset
\underline{\mathbf{Z}}$를 상수층의 아벨 부분층이라 하자.
그러면 어떤 $d \in \mathbf{Z}$에 대해
$\mathcal{H}|_U = \underline{d\mathbf{Z}}_U$를 만족하는 공집합이 아닌
열린집합 $U \subset X$가 존재한다.
\end{lemma}

\begin{proof}
$X$의 공집합이 아닌 임의의 열린집합 $V$에 대해
$\underline{\mathbf{Z}}(V) = \mathbf{Z}$임을 상기하자(Lemma
\ref{cohomology-lemma-irreducible-constant-cohomology-zero}의 증명을 보라).
$\mathcal{H} = 0$이면 보조정리는 $d = 0$으로 성립한다.
$\mathcal{H} \not = 0$이면 $\mathcal{H}(U) \not = 0$인 공집합이 아닌
열린집합 $U \subset X$가 존재한다. 어떤 $n \geq 1$에 대해
$\mathcal{H}(U) = n\mathbf{Z}$라 하자. 그러면
$\underline{n\mathbf{Z}}_U
\subset \mathcal{H}|_U \subset
\underline{\mathbf{Z}}_U$임을 알 수 있다. 첫 번째 포함이 진포함이면,
어떤 정수 $1 \leq n' < n$과 공집합이 아닌 $U' \subset U$를 찾아
$\underline{n'\mathbf{Z}}_{U'}
\subset \mathcal{H}|_{U'} \subset
\underline{\mathbf{Z}}_{U'}$로 둘 수 있다. 이 과정은 유한 단계 뒤에는
멈추어야 하므로 보조정리를 얻는다.
\end{proof}

\begin{proposition}[Grothendieck]
\label{cohomology-proposition-vanishing-Noetherian}
\begin{reference}
\cite[Theorem 3.6.5]{Tohoku}.
\end{reference}
$X$를 뇌터 위상공간이라 하자. $\dim(X) \leq d$이면,
$X$ 위의 임의의 아벨 층 $\mathcal{F}$와 모든 $p > d$에 대해
$H^p(X, \mathcal{F}) = 0$이다.
\end{proposition}

\begin{proof}
이 보조정리를 $d$에 대한 귀납법으로 증명한다. $d$를 고정하고,
차원이 $< d$인 모든 뇌터 위상공간에 대해 보조정리가 성립한다고 가정하자.

\medskip\noindent
$\mathcal{F}$를 $X$ 위의 아벨 층이라 하자. $U \subset X$가
열린집합이라고 하고, 닫힌 여집합을 $Z \subset X$로 나타내자.
포함 사상을 각각 $j : U \to X$와 $i : Z \to X$로 나타내자.
그러면 짧은 완전열
$$
0 \to j_{!}j^*\mathcal{F} \to \mathcal{F} \to i_*i^*\mathcal{F} \to 0
$$
이 존재한다. Modules, Lemma
\ref{modules-lemma-canonical-exact-sequence}를 보라.
$j_!j^*\mathcal{F}$는 $U$의 위상적 폐포 $Z'$ 위에 지지된다는 점에
유의하자. 즉, $i' : Z' \to X$가 포함일 때, $Z'$ 위의 어떤 아벨 층
$\mathcal{F}'$에 대해 $i'_*\mathcal{F}'$의 꼴이다.

\medskip\noindent
이를 사용하여 $X$가 기약인 경우로 환원할 수 있다. 실제로 Topology,
Lemma \ref{topology-lemma-Noetherian}에 따르면 $X$에는 유한 개의
기약 성분이 있다. $X$에 기약 성분이 둘 이상이면 $Z \subset X$를
$X$의 한 기약 성분이라 하고 $U = X \setminus Z$로 두자. 위의 내용과
코호몰로지의 긴 완전열에 의해, $p > d$에 대해
$H^p(X, i_*i^*\mathcal{F})$와 $H^p(X, i'_*\mathcal{F}')$의 소멸을
증명하면 충분하다. Lemma
\ref{cohomology-lemma-cohomology-and-closed-immersions}에 의해 $p > d$에 대해
$H^p(Z, i^*\mathcal{F})$와 $H^p(Z', \mathcal{F}')$가 소멸함을
증명하면 충분하다. $Z'$와 $Z$는 기약 성분의 수가 더 적으므로,
실제로 $X$가 기약인 경우로 환원된다.

\medskip\noindent
$d = 0$이고 $X$가 기약이면 $X$는 $X$의 유일한 공집합 아닌
열린 부분집합이다. 따라서 모든 층은 상수이고 고차 코호몰로지 군은
소멸한다(예를 들어 Lemma
\ref{cohomology-lemma-irreducible-constant-cohomology-zero}에 의한다).

\medskip\noindent
$X$가 차원 $d > 0$인 기약공간이라고 하자. Lemma
\ref{cohomology-lemma-vanishing-generated-one-section}에 의해
어떤 열린집합 $U \subset X$에 대해
$\mathcal{F} = j_!\underline{\mathbf{Z}}_U$인 경우로 환원된다.
이 경우 짧은 완전열
$$
0 \to j_!(\underline{\mathbf{Z}}_U) \to
\underline{\mathbf{Z}}_X \to i_*\underline{\mathbf{Z}}_Z \to 0
$$
을 살펴본다. 여기서 $Z = X \setminus U$이다. Lemma
\ref{cohomology-lemma-irreducible-constant-cohomology-zero}에 의해 모든
$p \geq 1$에 대해 $H^p(X, \underline{\mathbf{Z}}_X)$가 소멸한다.
귀납가정에 의해 $p \geq d$에 대해
$H^p(X, i_*\underline{\mathbf{Z}}_Z) =
H^p(Z, \underline{\mathbf{Z}}_Z) = 0$이다.
따라서 코호몰로지의 긴 완전열로부터 결론을 얻는다.
\end{proof}






\section{닫힌 부분집합에 지지를 갖는 코호몰로지}
\label{cohomology-section-cohomology-support}

\noindent
이 절에서는 꼭 필요한 최소한만 논의한다. 논의는 Section
\ref{cohomology-section-cohomology-support-bis}에서 계속된다.

\medskip\noindent
$X$를 위상공간이라 하고 $Z \subset X$를 닫힌 부분집합이라 하자.
$\mathcal{F}$를 $X$ 위의 아벨 층이라 하자. 지지가 $Z$에 포함되는
절단들의 부분집합을
$$
\Gamma_Z(X, \mathcal{F}) =
\{s \in \mathcal{F}(X) \mid \text{Supp}(s) \subset Z\}
$$
로 둔다. 절단의 지지는 Modules, Definition
\ref{modules-definition-support}에서 정의한다. Modules, Lemma
\ref{modules-lemma-support-section-closed}에 의해
$\Gamma_Z(X, \mathcal{F})$는 $\Gamma(X, \mathcal{F})$의 부분군이다.
같은 보조정리에 의해 대응 $\mathcal{F} \mapsto
\Gamma_Z(X, \mathcal{F})$는 $\mathcal{F}$에 대한 함자이다.
이 함자는 왼쪽 완전이지만 일반적으로 완전하지는 않다.

\medskip\noindent
아벨 층의 범주에는 충분히 많은 단사 대상이 있으므로(Injectives, Lemma
\ref{injectives-lemma-abelian-sheaves-space}), Derived Categories,
Lemma \ref{derived-lemma-enough-injectives-right-derived}에 의해
오른쪽 유도 함자
$$
R\Gamma_Z(X, -) : D^+(X) \longrightarrow D^+(\textit{Ab})
$$
를 얻는다. 대상 $K$에서 $R\Gamma_Z(X, -)$의 값은 $K$를 단사 아벨
층들로 이루어진 아래로 유계인 복합체 $\mathcal{I}^\bullet$로 나타낸 뒤
$\Gamma_Z(X, \mathcal{I}^\bullet)$를 취하여 계산한다. Derived
Categories, Lemma \ref{derived-lemma-injective-acyclic}를 보라.
$Z$에 지지를 갖는 아벨 층 $\mathcal{F}$의 코호몰로지 군은
$H^q_Z(X, \mathcal{F}) = R^q\Gamma_Z(X, \mathcal{F})$로 정의한다.

\medskip\noindent
$\mathcal{I}$를 $X$ 위의 단사 아벨 층이라 하고
$U = X \setminus Z$로 두자. 그러면 제한 사상
$\mathcal{I}(X) \to \mathcal{I}(U)$는 전사이며(Lemma
\ref{cohomology-lemma-injective-restriction-surjective}), 그 핵은
$\Gamma_Z(X, \mathcal{I})$이다. 곧바로 각 $K \in D^+(X)$에 대해
$D^+(\textit{Ab})$ 안의 특수 삼각형
$$
R\Gamma_Z(X, K) \to R\Gamma(X, K) \to R\Gamma(U, K) \to R\Gamma_Z(X, K)[1]
$$
을 얻는다. 그 결과 임의의 $D^+(X)$의 $K$에 대해 긴 코호몰로지 완전열
$$
\ldots \to H^i_Z(X, K) \to H^i(X, K) \to H^i(U, K) \to
H^{i + 1}_Z(X, K) \to \ldots
$$
을 얻는다.

\medskip\noindent
$X$ 위의 아벨 층 $\mathcal{F}$에 대해 {\it $Z$에 지지를 갖는
절단들의 부분층}을 생각할 수 있다. 이를 $\mathcal{H}_Z(\mathcal{F})$로
나타내고 다음 규칙으로 정의한다.
$$
\mathcal{H}_Z(\mathcal{F})(U) =
\{s \in \mathcal{F}(U) \mid \text{Supp}(s) \subset U \cap Z\} =
\Gamma_{Z \cap U}(U, \mathcal{F}|_U)
$$
Modules, Lemma \ref{modules-lemma-i-star-exact}의 동치를 사용하여
$\mathcal{H}_Z(\mathcal{F})$를 $Z$ 위의 아벨 층으로 볼 수 있다.
Modules, Remark
\ref{modules-remark-sections-support-in-closed}를 보라.
따라서 함자
$$
\textit{Ab}(X) \longrightarrow \textit{Ab}(Z),\quad
\mathcal{F} \longmapsto
\mathcal{H}_Z(\mathcal{F})\text{ viewed as a sheaf on }Z
$$
를 얻는다. 이 함자는 왼쪽 완전이지만 일반적으로 완전하지는 않다.
위와 똑같이 오른쪽 유도 함자
$$
R\mathcal{H}_Z : D^+(X) \longrightarrow D^+(Z)
$$
를 얻는다. $\mathcal{H}^0_Z(\mathcal{F}) =
\mathcal{H}_Z(\mathcal{F})$가 되도록
$\mathcal{H}^q_Z(\mathcal{F}) = R^q\mathcal{H}_Z(\mathcal{F})$로 둔다.

\medskip\noindent
임의의 아벨 층 $\mathcal{F}$에 대해
$\Gamma_Z(X, \mathcal{F}) = \Gamma(Z,
\mathcal{H}_Z(\mathcal{F}))$임에 유의하자. 아래의 Lemma
\ref{cohomology-lemma-sections-with-support-acyclic}에 의해 함자
$\mathcal{H}_Z$는 단사 아벨 층을 $\Gamma(Z, -)$에 대해 오른쪽
비순환인 층으로 보낸다. 따라서 Derived Categories, Lemma
\ref{derived-lemma-grothendieck-spectral-sequence}에 의해
$D^+(X)$의 $K$에 대해 함자적인 수렴하는 Grothendieck 스펙트럼열
$$
E_2^{p, q} = H^p(Z, \mathcal{H}^q_Z(K)) \Rightarrow H^{p + q}_Z(X, K)
$$
을 얻는다.

\begin{lemma}
\label{cohomology-lemma-sections-with-support-acyclic}
$i : Z \to X$를 닫힌 부분집합의 포함이라 하자.
$\mathcal{I}$를 $X$ 위의 단사 아벨 층이라 하자. 그러면
$\mathcal{H}_Z(\mathcal{I})$는 $Z$ 위의 단사 아벨 층이다.
\end{lemma}

\begin{proof}
$\mathcal{H}_Z(-)$가 완전함자 $i_*$의 오른쪽 수반이므로
Homology, Lemma \ref{homology-lemma-adjoint-preserve-injectives}에서
따른다. Modules, Lemmas \ref{modules-lemma-i-star-exact}와
\ref{modules-lemma-i-star-right-adjoint}를 보라.
\end{proof}





\section{스펙트럴 공간 위의 코호몰로지}
\label{cohomology-section-spectral}

\noindent
스펙트럴 공간의 코호몰로지에 관한 핵심 결과는 Lemma
\ref{cohomology-lemma-colimit}이다. 대략 말해, 이는 Topology, Definition
\ref{topology-definition-spectral-space}에서 정의한 스펙트럴 공간의
범주에서 코호몰로지가 공여과 극한과 교환한다는 뜻이다. 이를 적용하면
Lemmas \ref{cohomology-lemma-cohomology-of-closed}와
\ref{cohomology-lemma-proper-base-change}의 유사한 결과를 다음과 같이 얻을 수 있다.

\begin{lemma}
\label{cohomology-lemma-cohomology-of-neighbourhoods-of-closed}
$X$를 스펙트럴 공간이라 하고 $\mathcal{F}$를 $X$ 위의 아벨 층이라 하자.
$E \subset X$를 준콤팩트 부분집합이라 하자. $W \subset X$를 $E$의 한
점으로 특수화되는 $X$의 점들의 집합이라 하자.
\begin{enumerate}
\item $H^p(W, \mathcal{F}|_W) = \colim H^p(U, \mathcal{F})$이며, 여기서
여극한은 $E$의 준콤팩트 열린 근방들에 걸쳐 취한다.
\item $E$가 구성가능 부분집합이면
$H^p(W \setminus E, \mathcal{F}|_{W \setminus E}) =
\colim H^p(U \setminus E, \mathcal{F}|_{U \setminus E})$이다.
\end{enumerate}
\end{lemma}

\begin{proof}
Topology, Lemma \ref{topology-lemma-make-spectral-space}에 의해
$W = \lim U$임을 알 수 있다. 여기서 극한은 $E$를 포함하는 준콤팩트
열린집합들에 걸쳐 취한다. Topology, Lemma
\ref{topology-lemma-spectral-sub}에 의해 각 $U$는 스펙트럴 공간이다.
따라서 Lemma \ref{cohomology-lemma-colimit}를 적용하여 (1)이 성립함을 얻는다.
(2)에도 같은 증명이 적용되지만, 이때는 Topology, Lemma
\ref{topology-lemma-make-spectral-space-minus}를 사용한다.
\end{proof}

\begin{lemma}
\label{cohomology-lemma-proper-base-change-spectral}
$f : X \to Y$를 스펙트럴 공간 사이의 스펙트럴 사상이라 하자.
$y \in Y$라 하고, $E \subset Y$를 $y$로 특수화되는 점들의 집합이라
하자. $\mathcal{F}$를 $X$ 위의 아벨 층이라 하자. 그러면
$(R^pf_*\mathcal{F})_y =
H^p(f^{-1}(E), \mathcal{F}|_{f^{-1}(E)})$이다.
\end{lemma}

\begin{proof}
$V$가 $Y$에서 $y$의 준콤팩트 열린 근방들을 달릴 때
$E = \bigcap V$임에 유의하자. 따라서
$f^{-1}(E) = \bigcap f^{-1}(V)$이다. 이는 위상공간으로서
$f^{-1}(E) = \lim f^{-1}(V)$임을 뜻한다. $f$가 스펙트럴이므로 각
$f^{-1}(V)$도 스펙트럴 공간이다(Topology, Lemma
\ref{topology-lemma-spectral-sub}). 그러므로 $f^{-1}(E)$가 스펙트럴
공간임을 얻으며, Lemma \ref{cohomology-lemma-colimit}에 의해
$$
H^p(f^{-1}(E), \mathcal{F}|_{f^{-1}(E)}) =
\colim H^p(f^{-1}(V), \mathcal{F})
$$
이다. 한편 $y$에서 $R^pf_*\mathcal{F}$의 줄기는 오른쪽의 여극한으로
주어진다.
\end{proof}

\begin{lemma}
\label{cohomology-lemma-vanishing-for-profinite}
$X$를 프로유한 위상공간이라 하자. 그러면 모든 $q > 0$과 모든 아벨 층
$\mathcal{F}$에 대해 $H^q(X, \mathcal{F}) = 0$이다.
\end{lemma}

\begin{proof}
$X$의 임의의 열린 덮개는 열린 조각들로 이루어진 유한 서로소 합 분해로
세분될 수 있다. Topology, Lemma
\ref{topology-lemma-profinite-refine-open-covering}를 보라. 따라서
$\mathcal{F} \to \mathcal{G}$가 $X$ 위의 아벨 층의 전사이면
$\mathcal{F}(X) \to \mathcal{G}(X)$도 전사이다. 다시 말해 대역 절단
함자는 완전함자이다. 그러므로 그 고차 유도 함자들은 영이다. Derived
Categories, Lemma
\ref{derived-lemma-right-derived-exact-functor}를 보라.
\end{proof}

\noindent
다음 코호몰로지 소멸 결과는 Grothendieck의 결과(Proposition
\ref{cohomology-proposition-vanishing-Noetherian})를 개선하며,
\cite{Scheiderer}에서 찾을 수 있다.

\begin{proposition}
\label{cohomology-proposition-cohomological-dimension-spectral}
\begin{reference}
(1)은 \cite{Scheiderer}의 주정리이다.
\end{reference}
$X$를 크룰 차원 $d$인 스펙트럴 공간이라 하고,
$\mathcal{F}$를 $X$ 위의 아벨 층이라 하자.
\begin{enumerate}
\item $q > d$이면 $H^q(X, \mathcal{F}) = 0$이다.
\item 모든 준콤팩트 열린집합 $U \subset X$에 대해
$H^d(X, \mathcal{F}) \to H^d(U, \mathcal{F})$는 전사이다.
\item 임의의 구성가능 닫힌 부분집합 $Z \subset X$와 $q > d$에 대해
$H^q_Z(X, \mathcal{F}) = 0$이다.
\end{enumerate}
\end{proposition}

\begin{proof}
이 결과를 $d$에 대한 귀납법으로 증명한다.

\medskip\noindent
$d = 0$이면 $X$는 프로유한 공간이다. Topology, Lemma
\ref{topology-lemma-characterize-profinite-spectral}를 보라.
따라서 (1)은 Lemma \ref{cohomology-lemma-vanishing-for-profinite}에서 따른다.
$U \subset X$가 준콤팩트 열린집합이면, $U$는 하우스도르프 공간의
준콤팩트 부분집합이므로 닫힌집합이기도 하다. 따라서 위상공간으로서
$X = U \amalg (X \setminus U)$이고 (2)가 성립함을 알 수 있다.
(3)의 $Z$가 주어지면 긴 완전열
$$
H^{q - 1}(X, \mathcal{F}) \to
H^{q - 1}(X \setminus Z, \mathcal{F}) \to
H^q_Z(X, \mathcal{F}) \to H^q(X, \mathcal{F})
$$
을 생각한다. $X$와 $U = X \setminus Z$가 프로유한이고(실제로
$Z$가 구성가능이므로 $U$는 준콤팩트이다), (2)와 (1)이 성립하므로
$Z$에 지지를 갖는 코호몰로지 군의 원하는 소멸을 얻는다.

\medskip\noindent
귀납 단계를 보자. $d \geq 1$이고 차원이 $< d$인 모든 스펙트럴 공간에
대해 명제가 성립한다고 가정하자. 먼저 $X$에 대해 (2)를 증명한다.
$U$를 준콤팩트 열린집합이라 하고
$\xi \in H^d(U, \mathcal{F})$라 하자. $Z = X \setminus U$로 두고,
$W \subset X$를 $Z$로 특수화되는 점들의 집합이라 하자. Lemma
\ref{cohomology-lemma-cohomology-of-neighbourhoods-of-closed}에 의해
$$
H^d(W \setminus Z, \mathcal{F}|_{W \setminus Z}) =
\colim_{Z \subset V} H^d(V \setminus Z, \mathcal{F})
$$
이다. 여기서 여극한은 $X$에서 $Z$의 준콤팩트 열린 근방 $V$들에
걸쳐 취한다. Topology, Lemma
\ref{topology-lemma-make-spectral-space}에 의해 $W \setminus Z$가
스펙트럴 공간임을 알 수 있다. $W$의 모든 점이 $Z$의 한 점으로
특수화되므로 $W \setminus Z$는 크룰 차원이 $< d$인 스펙트럴
공간이다. 귀납가정에 의해 $\xi$의
$H^d(W \setminus Z, \mathcal{F}|_{W \setminus Z})$에서의 상은 영이다.
위의 공식에 의해 $\xi|_{V \setminus Z} = 0$을 만족하는 준콤팩트
열린집합 $Z \subset V \subset X$가 존재한다.
$V \setminus Z = V \cap U$이므로, 덮개 $X = U \cup V$에 대해
Mayer--Vietoris(Lemma \ref{cohomology-lemma-mayer-vietoris})를 적용하면
$U$에서는 $\xi$로, $V$에서는 영으로 제한되는
$\tilde \xi \in H^d(X, \mathcal{F})$가 존재한다. 다시 말해 (2)가
성립한다.

\medskip\noindent
(2)를 가정하고 (1)을 증명하자. 단사 분해
$\mathcal{F} \to \mathcal{I}^\bullet$을 택하고
$$
\mathcal{G} = \Im(\mathcal{I}^{d - 1} \to \mathcal{I}^d) =
\Ker(\mathcal{I}^d \to \mathcal{I}^{d + 1})
$$
로 두자. 준콤팩트 열린집합 $U \subset X$에 대해 다음 완전열들의
사상을 얻는다.
$$
\xymatrix{
\mathcal{I}^{d - 1}(X) \ar[r] \ar[d] &
\mathcal{G}(X) \ar[r] \ar[d] &
H^d(X, \mathcal{F}) \ar[d] \ar[r] & 0 \\
\mathcal{I}^{d - 1}(U) \ar[r] &
\mathcal{G}(U) \ar[r] &
H^d(U, \mathcal{F}) \ar[r] & 0
}
$$
층 $\mathcal{I}^{d - 1}$은 Lemma \ref{cohomology-lemma-injective-flasque}와
$d \geq 1$에 의해 플라스크이다. (2)에 의해 오른쪽 수직 화살표는
전사이다. 도표 추적을 하면 사상
$\mathcal{G}(X) \to \mathcal{G}(U)$가 전사임을 얻는다.
Lemma \ref{cohomology-lemma-vanishing-ravi}에 의해 준콤팩트 열린집합을
준콤팩트 열린집합들로 덮는 임의의 유한 덮개
$\mathcal{U} : U = U_1 \cup \ldots \cup U_n$와 모든 $q > 0$에 대해
$\check{H}^q(\mathcal{U}, \mathcal{G}) = 0$임을 얻는다.
Lemma \ref{cohomology-lemma-cech-vanish-basis}를 적용하면 $X$의 모든 준콤팩트
열린집합 $U$와 모든 $q > 0$에 대해 $H^q(U, \mathcal{G}) = 0$이다.
Leray의 비순환성 보조정리(Derived Categories, Lemma
\ref{derived-lemma-leray-acyclicity})에 의해
$$
H^q(X, \mathcal{F}) =
H^q\left(
\Gamma(X, \mathcal{I}^0) \to \ldots \to
\Gamma(X, \mathcal{I}^{d - 1}) \to \Gamma(X, \mathcal{G})
\right)
$$
임을 얻는다. 특히 $q > d$이면 코호몰로지 군이 소멸한다.

\medskip\noindent
(3)을 증명하자. (3)의 $Z$가 주어지면 긴 완전열
$$
H^{q - 1}(X, \mathcal{F}) \to
H^{q - 1}(X \setminus Z, \mathcal{F}) \to
H^q_Z(X, \mathcal{F}) \to H^q(X, \mathcal{F})
$$
을 생각한다. $X$와 $U = X \setminus Z$는 차원이 $\leq d$인
스펙트럴 공간이고(Topology, Lemma
\ref{topology-lemma-spectral-sub}), (2)와 (1)이 성립하므로 원하는
소멸을 얻는다.
\end{proof}















\section{교대 체흐 복합체}
\label{cohomology-section-alternating-cech}

\noindent
이 절에서는 체흐 복합체를 교대 체흐 복합체 및 몇몇 관련 복합체와
비교한다.

\medskip\noindent
$X$를 위상공간이라 하고 $\mathcal{U} : U = \bigcup_{i \in I} U_i$를
열린 덮개라 하자. $p \geq 0$에 대해
$$
\check{\mathcal{C}}_{alt}^p(\mathcal{U}, \mathcal{F})
=
\left\{
\begin{matrix}
s \in  \check{\mathcal{C}}^p(\mathcal{U}, \mathcal{F})
\text{ such that }
s_{i_0 \ldots i_p} = 0 \text{ if } i_n = i_m \text{ for some } n \not = m\\
\text{ and }
s_{i_0\ldots i_n \ldots i_m \ldots i_p}
=
-s_{i_0\ldots i_m \ldots i_n \ldots i_p}
\text{ in any case.}
\end{matrix}
\right\}
$$
로 두자. Equation (\ref{cohomology-equation-d-cech})의 미분 $d$가
$\check{\mathcal{C}}^p_{alt}(\mathcal{U}, \mathcal{F})$를
$\check{\mathcal{C}}^{p + 1}_{alt}(\mathcal{U}, \mathcal{F})$ 안으로
보낸다는 확인은 생략한다.

\begin{definition}
\label{cohomology-definition-alternating-cech-complex}
$X$를 위상공간이라 하고 $\mathcal{U} : U = \bigcup_{i \in I} U_i$를
열린 덮개라 하자. $\mathcal{F}$를 $X$ 위의 아벨 준층이라 하자.
복합체 $\check{\mathcal{C}}_{alt}^\bullet(\mathcal{U}, \mathcal{F})$를
$\mathcal{F}$와 열린 덮개 $\mathcal{U}$에 딸린
{\it 교대 체흐 복합체}라 한다.
\end{definition}

\noindent
따라서 복합체의 표준 사상
$$
\check{\mathcal{C}}_{alt}^\bullet(\mathcal{U}, \mathcal{F})
\longrightarrow
\check{\mathcal{C}}^\bullet(\mathcal{U}, \mathcal{F})
$$
이 존재한다. 이는 교대 체흐 복합체를 보통의 체흐 복합체에 넣는 포함이다.

\medskip\noindent
덮개 $\mathcal{U} : U = \bigcup_{i \in I} U_i$에 $I$ 위의 전순서
$<$가 주어졌다고 하자. 이 경우
$$
\check{\mathcal{C}}_{ord}^p(\mathcal{U}, \mathcal{F})
=
\prod\nolimits_{(i_0, \ldots, i_p) \in I^{p + 1}, i_0 < \ldots < i_p}
\mathcal{F}(U_{i_0\ldots i_p}).
$$
로 두자. 이는 아벨 군이다.
$s \in \check{\mathcal{C}}_{ord}^p(\mathcal{U}, \mathcal{F})$에 대해
$\mathcal{F}(U_{i_0\ldots i_p})$에서의 값을
$s_{i_0\ldots i_p}$로 나타낸다. 사상
$$
d : \check{\mathcal{C}}_{ord}^p(\mathcal{U}, \mathcal{F})
\longrightarrow
\check{\mathcal{C}}_{ord}^{p + 1}(\mathcal{U}, \mathcal{F})
$$
을 임의의 $i_0 < \ldots < i_{p + 1}$에 대해 공식
$$
d(s)_{i_0\ldots i_{p + 1}}
=
\sum\nolimits_{j = 0}^{p + 1}
(-1)^j
s_{i_0\ldots \hat i_j \ldots i_{p + 1}}|_{U_{i_0\ldots i_{p + 1}}}
$$
으로 정의한다. 이 공식은 Equation (\ref{cohomology-equation-d-cech})과 같다.
$d \circ d = 0$임은 바로 알 수 있다. 다시 말해
$\check{\mathcal{C}}_{ord}^\bullet(\mathcal{U}, \mathcal{F})$는
복합체이다.

\begin{definition}
\label{cohomology-definition-ordered-cech-complex}
$X$를 위상공간이라 하고
$\mathcal{U} : U = \bigcup_{i \in I} U_i$를 열린 덮개라 하자.
$I$에 전순서가 주어졌다고 가정하자. $\mathcal{F}$를 $X$ 위의
아벨 준층이라 하자. 복합체
$\check{\mathcal{C}}_{ord}^\bullet(\mathcal{U}, \mathcal{F})$를
$\mathcal{F}$, 열린 덮개 $\mathcal{U}$, 그리고 $I$에 주어진
전순서에 딸린 {\it 순서화 체흐 복합체}라 한다.
\end{definition}

\noindent
이 복합체를 교대 체흐 복합체라고 부르기도 한다. 그 까닭은 순서화 체흐
복합체와 교대 체흐 복합체 사이에 명백한 비교 사상이 있기 때문이다.
즉, 다음 규칙으로 주어지는 사상
$$
c :
\check{\mathcal{C}}_{ord}^\bullet(\mathcal{U}, \mathcal{F})
\longrightarrow
\check{\mathcal{C}}^\bullet(\mathcal{U}, \mathcal{F})
$$
을 생각하자.
$$
c(s)_{i_0\ldots i_p} =
\left\{
\begin{matrix}
0 &
\text{if} &
i_n = i_m \text{ for some } n \not = m\\
\text{sgn}(\sigma) s_{i_{\sigma(0)}\ldots i_{\sigma(p)}} &
\text{if} &
i_{\sigma(0)} < i_{\sigma(1)} < \ldots < i_{\sigma(p)}
\end{matrix}
\right.
$$
여기서 $\sigma$는 $\{0, \ldots, p\}$의 순열이고
$\text{sgn}(\sigma)$는 그 부호이다. 문헌에서는 흔히 사상 $c$를
통해 교대 체흐 복합체와 순서화 체흐 복합체를 동일시한다.
즉, 다음의 쉬운 보조정리가 성립한다.

\begin{lemma}
\label{cohomology-lemma-ordered-alternating}
$X$를 위상공간이라 하고
$\mathcal{U} : U = \bigcup_{i \in I} U_i$를 열린 덮개라 하자.
$I$에 전순서가 주어졌다고 가정하자. 사상 $c$는 복합체 사상이다.
실제로 이는 복합체의 동형사상
$$
c : \check{\mathcal{C}}_{ord}^\bullet(\mathcal{U}, \mathcal{F})
\to \check{\mathcal{C}}_{alt}^\bullet(\mathcal{U}, \mathcal{F})
$$
을 유도한다.
\end{lemma}

\begin{proof}
생략한다.
\end{proof}

\noindent
다음 규칙으로 기술되는 사상
$$
\pi :
\check{\mathcal{C}}^\bullet(\mathcal{U}, \mathcal{F})
\longrightarrow
\check{\mathcal{C}}_{ord}^\bullet(\mathcal{U}, \mathcal{F})
$$
도 있다.
$$
\pi(s)_{i_0\ldots i_p} = s_{i_0\ldots i_p}
$$
여기서 $i_0 < i_1 < \ldots < i_p$이다.

\begin{lemma}
\label{cohomology-lemma-project-to-ordered}
$X$를 위상공간이라 하고
$\mathcal{U} : U = \bigcup_{i \in I} U_i$를 열린 덮개라 하자.
$I$에 전순서가 주어졌다고 가정하자. 사상
$\pi : \check{\mathcal{C}}^\bullet(\mathcal{U}, \mathcal{F})
\to \check{\mathcal{C}}_{ord}^\bullet(\mathcal{U}, \mathcal{F})$는
복합체 사상이다. 이는 $c$의 왼쪽 역인 복합체의 동형사상
$$
\pi : \check{\mathcal{C}}_{alt}^\bullet(\mathcal{U}, \mathcal{F})
\to \check{\mathcal{C}}_{ord}^\bullet(\mathcal{U}, \mathcal{F})
$$
을 유도한다.
\end{lemma}

\begin{proof}
생략한다.
\end{proof}

\begin{remark}
\label{cohomology-remark-compared-ordered-complexes}
이는 지표집합 $I$ 위에 두 전순서 $<_1$과 $<_2$가 있으면 복합체의
동형사상 $\tau = \pi_2 \circ c_1 :
\check{\mathcal{C}}_{ord\text{-}1}(\mathcal{U}, \mathcal{F}) \to
\check{\mathcal{C}}_{ord\text{-}2}(\mathcal{U}, \mathcal{F})$를 얻는다는
뜻이다. 다음은 명백하다.
$$
\tau(s)_{i_0 \ldots i_p} =
\text{sign}(\sigma) s_{i_{\sigma(0)} \ldots i_{\sigma(p)}}
$$
여기서 $i_0 <_1 i_1 <_1 \ldots <_1 i_p$이고
$i_{\sigma(0)} <_2 i_{\sigma(1)} <_2 \ldots <_2 i_{\sigma(p)}$이다.
이런 의미에서 순서화 체흐 복합체는 선택한 전순서와 무관하다.
\end{remark}

\begin{lemma}
\label{cohomology-lemma-alternating-usual}
$X$를 위상공간이라 하고
$\mathcal{U} : U = \bigcup_{i \in I} U_i$를 열린 덮개라 하자.
$I$에 전순서가 주어졌다고 가정하자. 사상 $c \circ \pi$는
$\check{\mathcal{C}}^\bullet(\mathcal{U}, \mathcal{F})$ 위의
항등사상과 호모토픽이다. 특히 포함 사상
$\check{\mathcal{C}}_{alt}^\bullet(\mathcal{U}, \mathcal{F}) \to
\check{\mathcal{C}}^\bullet(\mathcal{U}, \mathcal{F})$는 호모토피
동치이다.
\end{lemma}

\begin{proof}
임의의 다중지표 $(i_0, \ldots, i_p) \in I^{p + 1}$에 대해 다음을
만족하는 유일한 순열
$\sigma : \{0, \ldots, p\} \to \{0, \ldots, p\}$가 존재한다.
$$
i_{\sigma(0)} \leq i_{\sigma(1)} \leq \ldots \leq i_{\sigma(p)}
\quad
\text{and}
\quad
\sigma(j) < \sigma(j + 1)
\quad
\text{if}
\quad
i_{\sigma(j)} = i_{\sigma(j + 1)}.
$$
이 순열을 $\sigma = \sigma^{i_0 \ldots i_p}$로 나타낸다.

\medskip\noindent
임의의 순열 $\sigma : \{0, \ldots, p\} \to \{0, \ldots, p\}$와
임의의 $a$, $0 \leq a \leq p$에 대해, $0 \leq j < a$이면
$\sigma_a(j) = \sigma(j)$이고
$\sigma_a(a) < \sigma_a(a + 1) < \ldots < \sigma_a(p)$인
$\{0, \ldots, p\}$의 유일한 순열을 $\sigma_a$로 나타낸다.
예를 들어 $p = 3$이고 $\sigma$, $\tau$가
$$
\begin{matrix}
\text{id} & 0 & 1 & 2 & 3 \\
\sigma & 3 & 2 & 1 & 0
\end{matrix}
\quad \text{and} \quad
\begin{matrix}
\text{id} & 0 & 1 & 2 & 3 \\
\tau & 3 & 0 & 2 & 1
\end{matrix}
$$
로 주어지면
$$
\begin{matrix}
\text{id} & 0 & 1 & 2 & 3 \\
\sigma_0 & 0 & 1 & 2 & 3 \\
\sigma_1 & 3 & 0 & 1 & 2 \\
\sigma_2 & 3 & 2 & 0 & 1 \\
\sigma_3 & 3 & 2 & 1 & 0 \\
\end{matrix}
\quad \text{and} \quad
\begin{matrix}
\text{id} & 0 & 1 & 2 & 3 \\
\tau_0 & 0 & 1 & 2 & 3 \\
\tau_1 & 3 & 0 & 1 & 2 \\
\tau_2 & 3 & 0 & 1 & 2 \\
\tau_3 & 3 & 0 & 2 & 1 \\
\end{matrix}
$$
을 얻는다. 항상 $\sigma_0 = \text{id}$이고
$\sigma_p = \sigma$임은 명백하다.

\medskip\noindent
이 표기를 도입했으므로
$s \in \check{\mathcal{C}}^{p + 1}(\mathcal{U}, \mathcal{F})$에 대해
$h(s) \in \check{\mathcal{C}}^p(\mathcal{U}, \mathcal{F})$를 다음
성분들을 갖는 원소로 정의한다.
\begin{equation}
\label{cohomology-equation-first-homotopy}
h(s)_{i_0\ldots i_p} =
\sum\nolimits_{0 \leq a \leq p}
(-1)^a \text{sign}(\sigma_a)
s_{i_{\sigma(0)} \ldots i_{\sigma(a)} i_{\sigma_a(a)} \ldots i_{\sigma_a(p)}}
\end{equation}
여기서 $\sigma = \sigma^{i_0 \ldots i_p}$이다. 지표
$i_{\sigma(a)}$는
$i_{\sigma(0)} \ldots i_{\sigma(a)} i_{\sigma_a(a)} \ldots
i_{\sigma_a(p)}$에 두 번 나타난다. 한 번은 처음 $a + 1$개 지표의
모임에, 다른 한 번은 두 번째 $p - a + 1$개 지표의 모임에 나타난다.
실제로 $\sigma_a$의 정의에 의해 어떤 $j \geq a$에 대해
$\sigma_a(j) = \sigma(a)$이다. 따라서 각 원소
$s_{i_{\sigma(0)} \ldots i_{\sigma(a)} i_{\sigma_a(a)} \ldots
i_{\sigma_a(p)}}$가 열린집합 $U_{i_0 \ldots i_p}$ 위에서 정의되므로
합은 의미가 있다. 또한 $a = 0$이면 $s_{i_0 \ldots i_p}$를 얻고,
$a = p$이면
$(-1)^p \text{sign}(\sigma) s_{i_{\sigma(0)} \ldots i_{\sigma(p)}}$를
얻는다.

\medskip\noindent
다음을 주장한다.
$$
(dh + hd)(s)_{i_0 \ldots i_p} =
s_{i_0 \ldots i_p} -
\text{sign}(\sigma) s_{i_{\sigma(0)} \ldots i_{\sigma(p)}}
$$
여기서 $\sigma = \sigma^{i_0 \ldots i_p}$이다. 이 주장의 확인은
생략한다. (stacks-project 하위 디렉터리 scripts에는 이 주장의
유한 개 사례를 확인하는 데 쓸 수 있는 first-homotopy.gp라는
PARI/gp 스크립트가 있다. 부호가 올바른지 확인하기 위해 이 스크립트를
작성했다.) 다음 규칙으로 주어지는 연산자를
$$
\kappa :
\check{\mathcal{C}}^\bullet(\mathcal{U}, \mathcal{F})
\longrightarrow
\check{\mathcal{C}}^\bullet(\mathcal{U}, \mathcal{F})
$$
라 하자.
$$
\kappa(s)_{i_0 \ldots i_p} =
\text{sign}(\sigma^{i_0 \ldots i_p}) s_{i_{\sigma(0)} \ldots i_{\sigma(p)}}.
$$
위 주장은 $\kappa$가 복합체 사상이며 $\kappa$가 체흐 복합체의
항등사상과 호모토픽임을 뜻한다. 그러나 연산자 $\kappa$의 상은 교대 부분복합체가
아니므로 이것만으로 보조정리가 곧바로 따라오지는 않는다. 실제로
$\kappa$의 상은 “준교대” 복합체
$\check{\mathcal{C}}_{semi\text{-}alt}^p(\mathcal{U}, \mathcal{F})$이다.
이 복합체의 $p$-코체인이 $s$라는 것은, 임의의
$(i_0, \ldots, i_p) \in I^{p + 1}$와
$\sigma = \sigma^{i_0 \ldots i_p}$에 대해
$$
s_{i_0 \ldots i_p} = \text{sign}(\sigma) s_{i_{\sigma(0)} \ldots i_{\sigma(p)}}
$$
인 것과 동치이다. 체흐 복합체의 또 다른 변형인 준순서화 체흐 복합체를
다음과 같이 정의한다.
$$
\check{\mathcal{C}}_{semi\text{-}ord}^p(\mathcal{U}, \mathcal{F})
=
\prod\nolimits_{i_0 \leq i_1 \leq \ldots \leq i_p}
\mathcal{F}(U_{i_0 \ldots i_p})
$$
Equation (\ref{cohomology-equation-d-cech})이 여기에서도 미분을 정의하여 복합체를
준다는 것은 쉽게 알 수 있다. 또한 (Lemma
\ref{cohomology-lemma-project-to-ordered}와 마찬가지로) 사영 사상
$$
\check{\mathcal{C}}_{semi\text{-}alt}^\bullet(\mathcal{U}, \mathcal{F})
\longrightarrow
\check{\mathcal{C}}_{semi\text{-}ord}^\bullet(\mathcal{U}, \mathcal{F})
$$
이 복합체의 동형사상임도 명백하다.

\medskip\noindent
그러므로 명백한 포함 사상
$$
\check{\mathcal{C}}_{ord}^p(\mathcal{U}, \mathcal{F})
\longrightarrow
\check{\mathcal{C}}_{semi\text{-}ord}^p(\mathcal{U}, \mathcal{F})
$$
이 호모토피 동치임을 보이면 보조정리가 따른다. 이를 위해 다음 호모토피를
사용한다.
\begin{equation}
\label{cohomology-equation-second-homotopy}
h(s)_{i_0 \ldots i_p} =
\left\{
\begin{matrix}
0 & \text{if} & i_0 < i_1 < \ldots < i_p \\
(-1)^a s_{i_0 \ldots i_{a - 1} i_a i_a i_{a + 1} \ldots i_p}
& \text{if} & i_0 < i_1 < \ldots < i_{a - 1} < i_a = i_{a + 1}
\end{matrix}
\right.
\end{equation}
다음을 주장한다.
$$
(dh + hd)(s)_{i_0 \ldots i_p} =
\left\{
\begin{matrix}
0 & \text{if} & i_0 < i_1 < \ldots < i_p \\
s_{i_0 \ldots i_p}
& \text{else} &
\end{matrix}
\right.
$$
확인은 생략한다. (stacks-project 하위 디렉터리 scripts에는 이 주장의
유한 개 사례를 확인하는 데 쓸 수 있는 second-homotopy.gp라는 PARI/gp
스크립트가 있다. 부호가 올바른지 확인하기 위해 이 스크립트를 작성했다.)
이 주장은 사영과 자연스러운 포함의 합성
$$
\check{\mathcal{C}}_{semi\text{-}ord}^\bullet(\mathcal{U}, \mathcal{F})
\longrightarrow
\check{\mathcal{C}}_{ord}^\bullet(\mathcal{U}, \mathcal{F})
\longrightarrow
\check{\mathcal{C}}_{semi\text{-}ord}^\bullet(\mathcal{U}, \mathcal{F})
$$
이 원하는 대로 항등사상과 호모토픽임을 분명히 보여 준다.
\end{proof}

\begin{lemma}
\label{cohomology-lemma-alternating-cech-trivial}
$X$를 위상공간이라 하고 $\mathcal{F}$를 $X$ 위의 아벨 준층이라 하자.
$\mathcal{U} : U = \bigcup_{i \in I} U_i$를 열린 덮개라 하자.
어떤 $i \in I$에 대해 $U_i = U$이면, $\mathcal{F}(U)$를 차수 $-1$에
놓고 $\mathcal{F}(U)$에서
$\check{\mathcal{C}}^0(\mathcal{U}, \mathcal{F})$로 가는 표준 사상을
미분으로 삼아 얻는 확대 교대 체흐 복합체
$$
\mathcal{F}(U) \to \check{\mathcal{C}}_{alt}^\bullet(\mathcal{U}, \mathcal{F})
$$
는 $0$과 호모토피 동치이다. 마찬가지로 $I$ 위의 임의의 전순서에 대해
확대 순서화 체흐 복합체
$$
\mathcal{F}(U) \to
\check{\mathcal{C}}_{ord}^\bullet(\mathcal{U}, \mathcal{F})
$$
는 $0$과 호모토피 동치이다.
\end{lemma}

\begin{proof}[첫째 증명]
Lemmas \ref{cohomology-lemma-cech-trivial}와 \ref{cohomology-lemma-alternating-usual}을
결합하면 된다.
\end{proof}

\begin{proof}[둘째 증명]
교대 체흐 복합체와 순서화 체흐 복합체는 동형이므로 순서화된 것에 대해
증명하면 충분하다. 표준 표기를 쓰자. 확대 순서화 체흐 복합체에서
차수 $p$인 코체인 $s$는 $s = (s_{i_0 \ldots i_p})$의 꼴이며,
여기서 $s_{i_0 \ldots i_p}$는
$\mathcal{F}(U_{i_0 \ldots i_p})$에 속하고
$i_0 < \ldots < i_p$이다. 이 표기로
$$
d(x)_{i_0 \ldots i_{p + 1}} =
\sum\nolimits_j (-1)^j x_{i_0 \ldots \hat i_j \ldots i_p}
$$
이다. $U = U_i$인 지표 $i \in I$를 고정하자. 호모토피로 다음 규칙에
의해 주어지는 사상들을 사용한다.
$$
h : \text{cochains of degree }p + 1 \to \text{cochains of degree }p
$$
$$
h(s)_{i_0 \ldots i_p} = 0 \text{ if } i \in \{i_0, \ldots, i_p\}
\text{ and }
h(s)_{i_0 \ldots i_p} = 
(-1)^j s_{i_0 \ldots i_j i i_{j + 1} \ldots i_p} \text{ if not}
$$
두 번째 경우 $j$는 $i_j < i < i_{j + 1}$을 만족하는 유일한 지표이다.
또한 $U = U_i$이므로 등식
$$
\mathcal{F}(U_{i_0 \ldots i_p}) =
\mathcal{F}(U_{i_0 \ldots i_j i i_{j + 1} \ldots i_p})
$$
이 성립하며, 이를 사용하면
$(-1)^j s_{i_0 \ldots i_j i i_{j + 1} \ldots i_p}$를
$\mathcal{F}(U_{i_0 \ldots i_p})$의 원소로 생각할 수 있다.
계산을 통해 $d h + h d$가 항등사상의 음수와 같음을 보일 것이다.
그러면 증명이 끝난다. 이를 위해 차수 $p$인 코체인 $s$를 고정하고
$I$의 원소 $i_0 < \ldots < i_p$를 택하자.

\medskip\noindent
경우 I: $i \in \{i_0, \ldots, i_p\}$라고 하자. $i = i_t$라 하자.
그러면 $h(d(s))_{i_0 \ldots i_p} = 0$이다. 한편
$$
d(h(s))_{i_0 \ldots i_p} =
\sum (-1)^j h(s)_{i_0 \ldots \hat i_j \ldots i_p} =
(-1)^t h(s)_{i_0 \ldots \hat i \ldots i_p} =
(-1)^t (-1)^{t - 1} s_{i_0 \ldots i_p}
$$
이다. 따라서 원하는 대로
$(dh + hd)(s)_{i_0 \ldots i_p} = -s_{i_0 \ldots i_p}$이다.

\medskip\noindent
경우 II: $i \not \in \{i_0, \ldots, i_p\}$라고 하자.
$i_j < i < i_{j + 1}$인 $j$를 택하면
\begin{align*}
h(d(s))_{i_0 \ldots i_p}
& =
(-1)^j d(s)_{i_0 \ldots i_j i i_{j + 1} \ldots i_p} \\
& =
\sum\nolimits_{j' \leq j} (-1)^{j + j'}
s_{i_0 \ldots \hat i_{j'} \ldots i_j i i_{j + 1} \ldots i_p} -
s_{i_0 \ldots i_p} \\
&
+ \sum\nolimits_{j' > j} (-1)^{j + j' + 1}
s_{i_0 \ldots i_j i i_{j + 1} \ldots \hat i_{j'} \ldots i_p}
\end{align*}
임을 알 수 있다. 한편
\begin{align*}
d(h(s))_{i_0 \ldots i_p}
& =
\sum\nolimits_{j'} (-1)^{j'} h(s)_{i_0 \ldots \hat i_{j'} \ldots i_p} \\
& =
\sum\nolimits_{j' \leq j} (-1)^{j' + j - 1}
s_{i_0 \ldots \hat i_{j'} \ldots i_j i i_{j + 1} \ldots i_p} \\
& +
\sum\nolimits_{j' > j} (-1)^{j' + j}
s_{i_0 \ldots i_j i i_{j + 1} \ldots \hat i_{j'} \ldots i_p}
\end{align*}
이다. 이들을 더하면 원하는 대로
$(dh + hd)(s)_{i_0 \ldots i_p} = - s_{i_0 \ldots i_p}$를 얻는다.
\end{proof}






\section{체흐 복합체에 대한 대안적 관점}
\label{cohomology-section-locally-finite-cech}

\noindent
이 절에서는 체흐 복합체와 코호몰로지의 관계를 확립하는 다른 방법을
논의한다.

\begin{lemma}
\label{cohomology-lemma-covering-resolution}
$X$를 환 달린 공간이라 하고
$\mathcal{U} : X = \bigcup_{i \in I} U_i$를 $X$의 열린 덮개라 하자.
$\mathcal{F}$를 $\mathcal{O}_X$-가군이라 하자. $\mathcal{F}$를
$U_{i_0 \ldots i_p}$에 제한한 것을
$\mathcal{F}_{i_0 \ldots i_p}$로 나타내자. 다음을 만족하는
$\mathcal{O}_X$-가군의 복합체
${\mathfrak C}^\bullet(\mathcal{U}, \mathcal{F})$가 존재한다.
$$
{\mathfrak C}^p(\mathcal{U}, \mathcal{F}) =
\prod\nolimits_{i_0 \ldots i_p}
(j_{i_0 \ldots i_p})_* \mathcal{F}_{i_0 \ldots i_p}
$$
그리고 미분
$d : {\mathfrak C}^p(\mathcal{U}, \mathcal{F})
\to {\mathfrak C}^{p + 1}(\mathcal{U}, \mathcal{F})$는
Equation (\ref{cohomology-equation-d-cech})과 같다. 또한 유사동형인 표준 사상
$$
\mathcal{F} \to {\mathfrak C}^\bullet(\mathcal{U}, \mathcal{F})
$$
이 존재한다. 즉, ${\mathfrak C}^\bullet(\mathcal{U}, \mathcal{F})$는
$\mathcal{F}$의 분해이다.
\end{lemma}

\begin{proof}
$$
0 \to \mathcal{F} \to \mathfrak{C}^0(\mathcal{U}, \mathcal{F}) \to
\mathfrak{C}^1(\mathcal{U}, \mathcal{F}) \to  \ldots
$$
이 줄기에서 완전함을 확인하자. $x \in X$라 하고
$x \in U_{i_{\text{fix}}}$인 $i_{\text{fix}} \in I$를 택하자.
그러면 다음과 같이
$$
h : \mathfrak{C}^p(\mathcal{U}, \mathcal{F})_x
\to \mathfrak{C}^{p - 1}(\mathcal{U}, \mathcal{F})_x
$$
를 정의한다. $s \in \mathfrak{C}^p(\mathcal{U}, \mathcal{F})_x$이면
$x$의 어떤 근방 $V$에서 정의된 대표
$$
\widetilde{s} \in
\mathfrak{C}^p(\mathcal{U}, \mathcal{F})(V) =
\prod\nolimits_{i_0 \ldots i_p}
\mathcal{F}(V \cap U_{i_0} \cap \ldots \cap U_{i_p})
$$
를 택하고
$$
h(s)_{i_0 \ldots i_{p - 1}} =
\widetilde{s}_{i_{\text{fix}} i_0 \ldots i_{p - 1}, x}.
$$
로 둔다. 같은 공식으로($p = 0$인 경우)
$\mathfrak{C}^{0}(\mathcal{U},\mathcal{F})_x \to \mathcal{F}_x$인 사상을
얻는다. 형식적으로 다음과 같이 계산한다.
\begin{align*}
(dh + hd)(s)_{i_0 \ldots i_p}
& =
\sum\nolimits_{j = 0}^p
(-1)^j
h(s)_{i_0 \ldots \hat i_j \ldots i_p}
+
d(s)_{i_{\text{fix}} i_0 \ldots i_p}\\
& =
\sum\nolimits_{j = 0}^p
(-1)^j
s_{i_{\text{fix}} i_0 \ldots \hat i_j \ldots i_p}
+
s_{i_0 \ldots i_p}
+
\sum\nolimits_{j = 0}^p
(-1)^{j + 1}
s_{i_{\text{fix}} i_0 \ldots \hat i_j \ldots i_p} \\
& =
s_{i_0 \ldots i_p}
\end{align*}
이는 $h$가 확대 복합체
$$
0 \to \mathcal{F}_x \to \mathfrak{C}^0(\mathcal{U}, \mathcal{F})_x
\to \mathfrak{C}^1(\mathcal{U}, \mathcal{F})_x \to \ldots
$$
의 항등사상에서 영으로 가는 호모토피임을 보여 주므로 결론을 얻는다.
\end{proof}

\noindent
이 보조정리를 사용하면 Lemma \ref{cohomology-lemma-cech-spectral-sequence}의
체흐-코호몰로지 스펙트럼열을 쉽게 다시 증명할 수 있다. 즉, $X$,
$\mathcal{U}$, $\mathcal{F}$를 Lemma
\ref{cohomology-lemma-covering-resolution}에서와 같이 두고 단사 분해
$\mathcal{F} \to \mathcal{I}^\bullet$을 택하자. 그러면 이중 복합체
$$
A^{\bullet, \bullet} =
\Gamma(X, {\mathfrak C}^\bullet(\mathcal{U}, \mathcal{I}^\bullet)).
$$
를 생각할 수 있다. 구성에 의해
$$
A^{p, q} = \prod\nolimits_{i_0 \ldots i_p} \mathcal{I}^q(U_{i_0 \ldots i_p})
$$
이다. 이 이중 복합체에 딸린 Homology, Section
\ref{homology-section-double-complex}의 두 스펙트럼열을 생각하자.
특히 Homology, Lemma \ref{homology-lemma-ss-double-complex}를 보라.
곱을 취하는 것이 완전하므로(Homology, Lemma
\ref{homology-lemma-product-abelian-groups-exact}), 스펙트럼열
$({}'E_r, {}'d_r)_{r \geq 0}$에 대해서는
${}'E_2^{p, q} =
\check{H}^p(\mathcal{U}, \underline{H}^q(\mathcal{F}))$를 얻는다.
스펙트럼열 $({}''E_r, {}''d_r)_{r \geq 0}$에 대해서는 $p > 0$이면
${}''E_2^{p, q} = 0$이고
${}''E_2^{0, q} = H^q(X, \mathcal{F})$임을 얻는다. 실제로 $q$를
고정하면 층의 복합체
${\mathfrak C}^\bullet(\mathcal{U}, \mathcal{I}^q)$는 단사층
$\mathcal{I}^q$의 분해(Lemma \ref{cohomology-lemma-covering-resolution})이며,
그 항들도 단사층이다(Lemmas \ref{cohomology-lemma-cohomology-of-open},
\ref{cohomology-lemma-pushforward-injective-flat}와 Homology, Lemma
\ref{homology-lemma-product-injectives}에 의한다). 따라서
$\Gamma(X, {\mathfrak C}^\bullet(\mathcal{U}, \mathcal{I}^q))$의
코호몰로지는 양의 차수에서는 영이고 차수 $0$에서는
$\Gamma(X, \mathcal{I}^q)$와 같다. 다음 미분의 코호몰로지를 취하면
스펙트럼열 $({}''E_r, {}''d_r)_{r \geq 0}$에 대한 주장을 얻는다.
두 스펙트럼열 모두 $A^{\bullet, \bullet}$에 딸린 전체 복합체의
코호몰로지로 수렴하므로 결과가 따른다.

\begin{definition}
\label{cohomology-definition-covering-locally-finite}
$X$를 위상공간이라 하자. 열린 덮개
$X = \bigcup_{i \in I} U_i$가 {\it 국소 유한}이라는 것은, 모든
$x \in X$에 대해 $x$의 열린 근방 $W$가 존재하여
$\{i \in I \mid W \cap U_i \not = \emptyset\}$가 유한이라는 뜻이다.
\end{definition}

\begin{remark}
\label{cohomology-remark-locally-finite-sections}
$X = \bigcup_{i \in I} U_i$를 국소 유한 열린 덮개라 하자. 포함 사상을
$j_i : U_i \to X$로 나타내자. 각 $i$에 대해 $U_i$ 위의 아벨 층
$\mathcal{F}_i$가 주어졌다고 하자. 아벨 층
$\mathcal{G} = \bigoplus_{i \in I} (j_i)_*\mathcal{F}_i$를 생각하자.
그러면 실제로 열린집합 $V \subset X$에 대해
$$
\Gamma(V, \mathcal{G}) = \prod\nolimits_{i \in I} \mathcal{F}_i(V \cap U_i).
$$
이다. 다시 말해
$$
\bigoplus\nolimits_{i \in I} (j_i)_*\mathcal{F}_i =
\prod\nolimits_{i \in I} (j_i)_*\mathcal{F}_i
$$
이다. 층들의 모임의 직합이 우리가 예상하는 것의 층화라는 점을 깨닫기
전에는 이것이 이상해 보인다. Modules, Section
\ref{modules-section-kernels}의 논의를 보라. 따라서 이 경우 Lemma
\ref{cohomology-lemma-covering-resolution}의 복합체의 항들은
$$
{\mathfrak C}^p(\mathcal{U}, \mathcal{F}) =
\bigoplus\nolimits_{i_0 \ldots i_p}
(j_{i_0 \ldots i_p})_* \mathcal{F}_{i_0 \ldots i_p}
$$
이고, 이는 때때로 유용하다.
\end{remark}












\section{복합체의 체흐 코호몰로지}
\label{cohomology-section-cech-cohomology-of-complexes}

\noindent
일반적으로 $X$ 위의 아벨 군의 층 ${\mathcal F}$와
${\mathcal G}$에 대해 컵곱 사상
$$
H^i(X, {\mathcal F}) \times H^j(X, {\mathcal G})
\longrightarrow
H^{i + j}(X, {\mathcal F} \otimes_{\mathbf Z} {\mathcal G}).
$$
이 존재한다. 이 절에서는 컵곱의 명시적 공식을 써서 체흐 코사이클로
이를 정의한다. 코호몰로지가 체흐 코호몰로지와 같지 않을 수 있다는
점이 염려된다면, 초덮개를 쓰면서도 코사이클 표기를 사용할 수 있다.
이는 임의의 토포스 위의 초코호몰로지에 대한 컵곱을 정의하는 데에도
쓸 수 있다는 장점이 있다(향후 참고문헌 삽입).

\medskip\noindent
${\mathcal F}^\bullet$을 $X$ 위의 아벨 군의 준층들로 이루어진 아래로
유계인 복합체라 하자. 흔히 체흐 코사이클을 사용하여
$H^n(X, {\mathcal F}^\bullet)$을 계산할 수 있다. 실제로
${\mathcal U} : X = \bigcup_{i \in I} U_i$를 $X$의 열린 덮개라 하자.
체흐 복합체 $\check{\mathcal{C}}^\bullet(\mathcal{U}, \mathcal{F})$
(Definition \ref{cohomology-definition-cech-complex})는 준층 $\mathcal{F}$에 대해
함자적이므로 이중 복합체
$\check{\mathcal{C}}^\bullet(\mathcal{U}, \mathcal{F}^\bullet)$을 얻는다.
$\check{\mathcal{C}}^\bullet({\mathcal U}, {\mathcal F}^\bullet)$에 딸린
전체 복합체는 차수 $n$의 항이
$$
\text{Tot}^n(\check{\mathcal{C}}^\bullet({\mathcal U}, {\mathcal F}^\bullet))
=
\bigoplus\nolimits_{p + q = n}
\prod\nolimits_{i_0\ldots i_p} {\mathcal F}^q(U_{i_0\ldots i_p})
$$
인 복합체이다. Homology, Definition
\ref{homology-definition-associated-simple-complex}을 보라.
$\text{Tot}^n$의 전형적인 원소를
$\alpha = \{\alpha_{i_0\ldots i_p}\}$로 나타내며, 여기서
$\alpha_{i_0 \ldots i_p} \in
\mathcal{F}^q(U_{i_0\ldots i_p})$이다. 다시 말해
$\alpha_{i_0\ldots i_p}$의 $\mathcal{F}$-차수는 $q = n - p$이다.
이 표기를 쓰려면 $\alpha$가 어느 차수에 속하는지를 항상 유념해야 한다.
이 상황을 공식
$\deg_{\mathcal F}(\alpha_{i_0\ldots i_p}) = q$로 나타낸다.
Homology, Definition
\ref{homology-definition-associated-simple-complex}의 규약에 따르면,
차수 $n$인 원소 $\alpha$의 미분은
$$
d(\alpha)_{i_0\ldots i_{p + 1}}
=
\sum\nolimits_{j = 0}^{p + 1}
(-1)^j \alpha_{i_0 \ldots \hat i_j \ldots i_{p + 1}} + 
(-1)^{p + 1}d_{{\mathcal F}}(\alpha_{i_0 \ldots i_{p + 1}})
$$
로 주어진다. 여기서 $d_\mathcal{F}$는 복합체
$\mathcal{F}^\bullet$의 미분이다. 표현
$\alpha_{i_0 \ldots \hat i_j \ldots i_{p + 1}}$는
$\alpha_{i_0 \ldots \hat i_j \ldots i_{p + 1}}
\in {\mathcal F}(U_{i_0\ldots\hat i_j\ldots i_{p + 1}})$를
$U_{i_0 \ldots i_{p + 1}}$에 제한한 것을 뜻한다.

\medskip\noindent
$\text{Tot}(\check{\mathcal{C}}^\bullet({\mathcal U},
{\mathcal F}^\bullet))$의 구성은 ${\mathcal F}^\bullet$에 대해 함자적이다.
또한 다음 복합체의 자연변환이 있다.
\begin{equation}
\label{cohomology-equation-global-sections-to-cech}
\Gamma(X, {\mathcal F}^\bullet)
\longrightarrow
\text{Tot}(\check{\mathcal{C}}^\bullet({\mathcal U}, {\mathcal F}^\bullet))
\end{equation}
이는 다음 규칙으로 정의된다. 절단
$s\in \Gamma(X, {\mathcal F}^n)$는
$\alpha_{i_0} = s|_{U_{i_0}}$이고 $p > 0$에 대해
$\alpha_{i_0\ldots i_p} = 0$인 원소
$\alpha = \{\alpha_{i_0\ldots i_p}\}$로 보내진다.

\medskip\noindent
세분을 보자. ${\mathcal V} = \{ V_j \}_{j\in J}$를
${\mathcal U}$의 세분이라 하자. 이는 모든 $j\in J$에 대해
$V_j \subset U_{t(j)}$인 사상 $t: J \to I$가 존재한다는 뜻이다.
이로부터 자연변환
\begin{equation}
\label{cohomology-equation-transformation}
T_t :
\text{Tot}(\check{\mathcal{C}}^\bullet({\mathcal U}, {\mathcal F}^\bullet))
\longrightarrow
\text{Tot}(\check{\mathcal{C}}^\bullet({\mathcal V}, {\mathcal F}^\bullet)).
\end{equation}
을 얻으며, 이는 규칙
$$
T_t(\alpha)_{j_0\ldots j_p}
=
\alpha_{t(j_0)\ldots t(j_p)}|_{V_{j_0\ldots j_p}}.
$$
으로 정의된다. 위와 같은 두 사상 $t, t' : J \to I$에 대해 위에서
구성한 사상 $T_t$와 $T_{t'}$는 호모토픽이다. 호모토피는 차수 $n$인
원소 $\alpha$에 대해
$$
h(\alpha)_{j_0\ldots j_p}
=
\sum\nolimits_{a = 0}^{p}
(-1)^a
\alpha_{t(j_0)\ldots t(j_a) t'(j_a) \ldots t'(j_p)}
$$
로 주어진다. 이는 다음 계산으로 확인된다. 여기서 다시 $\alpha$는
차수 $n$인 원소이다(따라서 $d(\alpha)$의 차수는 $n + 1$이고
$h(\alpha)$의 차수는 $n - 1$이다).
\begin{align*}
(
d(h(\alpha)) + h(d(\alpha))
)_{j_0\ldots j_p}
= &
\sum\nolimits_{k = 0}^p
(-1)^k
h(\alpha)_{j_0 \ldots \hat j_k \ldots j_p}
+ \\
&
(-1)^p
d_{\mathcal F}(h(\alpha)_{j_0 \ldots j_p})
+ \\
&
\sum\nolimits_{a = 0}^p
(-1)^a
d(\alpha)_{t(j_0) \ldots t(j_a) t'(j_a) \ldots t'(j_p)}
\\
= &
\sum\nolimits_{k = 0}^p
\sum\nolimits_{a = 0}^{k - 1}
(-1)^{k + a}
\alpha_{t(j_0)\ldots t(j_a)t'(j_a)\ldots \hat{t'(j_k)}\ldots t'(j_p)}
+ \\
&
\sum\nolimits_{k = 0}^p
\sum\nolimits_{a = k + 1}^p
(-1)^{k + a - 1}
\alpha_{t(j_0)\ldots \hat{t(j_k)}\ldots t(j_a)t'(j_a)\ldots t'(j_p)}
+ \\
&
\sum\nolimits_{a = 0}^p
(-1)^{p + a}
d_{\mathcal F}(\alpha_{t(j_0)\ldots t(j_a) t'(j_a) \ldots t'(j_p)})
+ \\
&
\sum\nolimits_{a = 0}^p
\sum\nolimits_{k = 0}^a
(-1)^{a + k}
\alpha_{t(j_0)\ldots\hat{t(j_k)}\ldots t(j_a)t'(j_a)\ldots t'(j_p)}
+ \\
&
\sum\nolimits_{a = 0}^p
\sum\nolimits_{k = a}^p
(-1)^{a + k + 1}
\alpha_{t(j_0) \ldots t(j_a) t'(j_a) \ldots \hat{t'(j_k)} \ldots t'(j_p)}
+ \\
&
\sum\nolimits_{a = 0}^p
(-1)^{a + p + 1}
d_{\mathcal F}(\alpha_{t(j_0)\ldots t(j_a) t'(j_a) \ldots t'(j_p)})
\\
= &
\alpha_{t'(j_0)\ldots t'(j_p)} +
(-1)^{2p + 1}\alpha_{t(j_0)\ldots t(j_p)}
\\
= &
T_{t'}(\alpha)_{j_0\ldots j_p} - T_t(\alpha)_{j_0\ldots j_p}
\end{align*}
상쇄의 확인은 독자에게 맡긴다. (첫째, 둘째, 넷째, 다섯째 합의 항들 가운데
$k$와 $a$를 모두 갖는 항들은 상쇄된다. 다만 넷째와 다섯째에만 나타나는
$a = k$인 항들은 서로 상쇄되며, 원하는 두 항은 예외이다.)
따라서 유도 사상
$$
H^n(T_t) :
H^n(
\text{Tot}(\check{\mathcal{C}}^\bullet({\mathcal U}, {\mathcal F}^\bullet))
)
\to
H^n(
\text{Tot}(\check{\mathcal{C}}^\bullet({\mathcal V}, {\mathcal F}^\bullet))
)
$$
은 $t$의 선택과 무관하다. 모든 세분에 걸친 체흐 코호몰로지 군의
여극한을 사상 $H^\bullet(T_t)$를 통해 취한 것을
{\it 체흐 초코호몰로지}라 정의한다.

\medskip\noindent
($X$의 모든 열린 덮개에 걸친) 여극한에서 다음 보조정리는 체흐
초코호몰로지에서 코호몰로지로 가는 사상을 준다. 이는 흔히
동형사상이며, 초덮개를 사용하면 항상 동형사상이다.

\begin{lemma}
\label{cohomology-lemma-cech-complex-complex}
$(X, \mathcal{O}_X)$를 환 달린 공간이라 하자.
$\mathcal{U} : X = \bigcup_{i \in I} U_i$를 열린 덮개라 하자.
$\mathcal{O}_X$-가군들로 이루어진 아래로 유계인 복합체
$\mathcal{F}^\bullet$에 대해 표준 사상
$$
\text{Tot}(\check{\mathcal{C}}^\bullet(\mathcal{U}, \mathcal{F}^\bullet))
\longrightarrow
R\Gamma(X, \mathcal{F}^\bullet)
$$
이 존재한다. 이는 $\mathcal{F}^\bullet$에 대해 함자적이며
(\ref{cohomology-equation-global-sections-to-cech}) 및
(\ref{cohomology-equation-transformation})과 양립한다. 또한
$H^{p + q}(X, \mathcal{F}^\bullet)$로 수렴하는 스펙트럼열
$(E_r, d_r)_{r \geq 0}$이 존재하며
$$
E_2^{p, q} =
H^p(\text{Tot}(\check{\mathcal{C}}^\bullet(\mathcal{U},
\underline{H}^q(\mathcal{F}^\bullet)))
$$
이다.
\end{lemma}

\begin{proof}
${\mathcal I}^\bullet$을 아래로 유계인 단사 대상들의 복합체라 하자.
$\mathcal{I}^\bullet$에 대한
(\ref{cohomology-equation-global-sections-to-cech})의 사상은
$\Gamma(X, {\mathcal I}^\bullet) \to
\text{Tot}(\check{\mathcal{C}}^\bullet({\mathcal U},
{\mathcal I}^\bullet))$이다. 이는 Homology, Lemma
\ref{homology-lemma-double-complex-gives-resolution}을 이중 복합체
$\check{\mathcal{C}}^\bullet({\mathcal U},
{\mathcal I}^\bullet)$에 적용하고 Lemma
\ref{cohomology-lemma-injective-trivial-cech}를 사용하면 아벨 군의 복합체의
유사동형임을 알 수 있다.
${\mathcal F}^\bullet \to {\mathcal I}^\bullet$을
${\mathcal F}^\bullet$에서 아래로 유계인 단사 대상들의 복합체로 가는
유사동형이라 하자. $R\Gamma(X, {\mathcal F}^\bullet)$은 복합체
$\Gamma(X, {\mathcal I}^\bullet)$로 표현되므로 다음 사상을 사용하여
보조정리의 사상을 얻는다.
$$
\text{Tot}(\check{\mathcal{C}}^\bullet({\mathcal U}, {\mathcal F}^\bullet))
\longrightarrow
\text{Tot}(\check{\mathcal{C}}^\bullet({\mathcal U}, {\mathcal I}^\bullet)).
$$
함자성과 양립성의 확인은 생략한다. 보조정리의 스펙트럼열을 구성하기
위해 Cartan--Eilenberg 분해
$\mathcal{F}^\bullet \to \mathcal{I}^{\bullet, \bullet}$을 택하자.
Derived Categories, Lemma
\ref{derived-lemma-cartan-eilenberg}를 보라. 이 경우
$\mathcal{F}^\bullet \to \text{Tot}(\mathcal{I}^{\bullet, \bullet})$은
단사 분해이고, 따라서 위에서 본 바와 같이
$$
\text{Tot}(\check{\mathcal{C}}^\bullet({\mathcal U},
\text{Tot}({\mathcal I}^{\bullet, \bullet})))
$$
가 $R\Gamma(X, \mathcal{F}^\bullet)$을 계산한다. Homology, Remark
\ref{homology-remark-triple-complex}에 의해 이를 삼중 복합체
$\check{\mathcal{C}}^\bullet({\mathcal U},
{\mathcal I}^{\bullet, \bullet})$에 딸린 전체 복합체로 볼 수 있다.
따라서 같은 비고를 다시 사용하면 이를 항들이
$$
A^{n, m} =
\bigoplus\nolimits_{p + q = n}
\check{\mathcal{C}}^p({\mathcal U}, \mathcal{I}^{q, m})
$$
인 이중 복합체 $A^{\bullet, \bullet}$에 딸린 전체 복합체로 볼 수 있다.
$\mathcal{I}^{q, \bullet}$은 $\mathcal{F}^q$의 단사 분해이므로
$A^{\bullet, \bullet}$에 딸린 첫 번째 스펙트럼열(Homology, Lemma
\ref{homology-lemma-ss-double-complex})을 적용하여
$$
E_1^{n, m} =
\bigoplus\nolimits_{p + q = n}
\check{\mathcal{C}}^p(\mathcal{U}, \underline{H}^m(\mathcal{F}^q))
$$
인 스펙트럼열을 얻는다. 이는 복합체
$\text{Tot}(\check{\mathcal{C}}^\bullet(\mathcal{U},
\underline{H}^m(\mathcal{F}^\bullet))$의 $n$번째 항이다.
따라서 보조정리에 기술된 $E_2$ 항들을 얻는다. 수렴성은 Homology,
Lemma \ref{homology-lemma-first-quadrant-ss}에서 따른다.
\end{proof}

\begin{lemma}
\label{cohomology-lemma-cech-complex-complex-computes}
$(X, \mathcal{O}_X)$를 환 달린 공간이라 하자.
$\mathcal{U} : X = \bigcup_{i \in I} U_i$를 열린 덮개라 하자.
$\mathcal{F}^\bullet$을 $\mathcal{O}_X$-가군들로 이루어진 아래로 유계인
복합체라 하자. 모든 $i > 0$과 모든 $p, i_0, \ldots, i_p, q$에 대해
$H^i(U_{i_0 \ldots i_p}, \mathcal{F}^q) = 0$이면, Lemma
\ref{cohomology-lemma-cech-complex-complex}의 사상
$
\text{Tot}(\check{\mathcal{C}}^\bullet(\mathcal{U}, \mathcal{F}^\bullet))
\to
R\Gamma(X, \mathcal{F}^\bullet)
$
은 동형사상이다.
\end{lemma}

\begin{proof}
Lemma \ref{cohomology-lemma-cech-complex-complex}의 스펙트럼열에서 곧바로 따른다.
\end{proof}

\begin{remark}
\label{cohomology-remark-shift-complex-cech-complex}
$(X, \mathcal{O}_X)$를 환 달린 공간이라 하고
$\mathcal{U} : X = \bigcup_{i \in I} U_i$를 열린 덮개라 하자.
$\mathcal{F}^\bullet$을 $\mathcal{O}_X$-가군들로 이루어진 아래로 유계인
복합체라 하고 $b$를 정수라 하자. 유도 범주에서 다음 가환 도표가
존재한다고 주장한다.
$$
\xymatrix{
\text{Tot}(\check{\mathcal{C}}^\bullet(\mathcal{U}, \mathcal{F}^\bullet))[b]
\ar[r] \ar[d]_\gamma &
R\Gamma(X, \mathcal{F}^\bullet)[b] \ar[d] \\
\text{Tot}(\check{\mathcal{C}}^\bullet(\mathcal{U}, \mathcal{F}^\bullet[b]))
\ar[r] &
R\Gamma(X, \mathcal{F}^\bullet[b])
}
$$
여기서 사상 $\gamma$는 Homology, Remark
\ref{homology-remark-shift-double-complex}에서 구성한 복합체의
사상이다. 이 말이 의미가 있는 까닭은 이중 복합체
$\check{\mathcal{C}}^\bullet(\mathcal{U}, \mathcal{F}^\bullet[b])$가
Homology, Remark \ref{homology-remark-shift-double-complex}에서 도입한
이중 복합체
$\check{\mathcal{C}}^\bullet(\mathcal{U}, \mathcal{F}^\bullet)[0, b]$와
분명히 같기 때문이다. 도표의 가환성을 확인하기 위해 Lemma
\ref{cohomology-lemma-cech-complex-complex}의 증명에서처럼 단사 분해
$\mathcal{F}^\bullet \to \mathcal{I}^\bullet$을 택할 수 있다.
도표를 추적하면 $\mathcal{F}^\bullet$을 $\mathcal{I}^\bullet$로
바꾼 경우에 도표가 가환함을 확인하는 것으로 충분함을 알 수 있다.
이제 확대된 도표
$$
\xymatrix{
\Gamma(X, \mathcal{I}^\bullet)[b] \ar[r] \ar[d] &
\text{Tot}(\check{\mathcal{C}}^\bullet(\mathcal{U}, \mathcal{I}^\bullet))[b]
\ar[r] \ar[d]_\gamma &
R\Gamma(X, \mathcal{I}^\bullet)[b] \ar[d] \\
\Gamma(X, \mathcal{I}^\bullet[b]) \ar[r] &
\text{Tot}(\check{\mathcal{C}}^\bullet(\mathcal{U}, \mathcal{I}^\bullet[b]))
\ar[r] &
R\Gamma(X, \mathcal{I}^\bullet[b])
}
$$
를 생각하자. 여기서 왼쪽 수평 화살표들은
(\ref{cohomology-equation-global-sections-to-cech})이다. 이 경우 수평 화살표들은
유도 범주에서 동형사상이므로(Lemma
\ref{cohomology-lemma-cech-complex-complex}의 증명을 보라), 왼쪽 정사각형이
가환함을 보이면 충분하다. 이는 사상 $\gamma$가 직합항
$\check{\mathcal{C}}^0(\mathcal{U}, \mathcal{I}^{q + b})$에서 부호
$1$을 사용하기 때문이다. Homology, Remark
\ref{homology-remark-shift-double-complex}의 공식을 보라.
\end{remark}

\noindent
$X$를 위상공간이라 하고
$\mathcal{U} : X = \bigcup_{i \in I} U_i$를 열린 덮개라 하며,
$\mathcal{F}^\bullet$을 아벨 군의 준층들로 이루어진 아래로 유계인
복합체라 하자. 다음과 같이 정의되는 사상
$\tau :
\text{Tot}(\check{\mathcal{C}}^\bullet({\mathcal U}, {\mathcal F}^\bullet))
\to
\text{Tot}(\check{\mathcal{C}}^\bullet({\mathcal U}, {\mathcal F}^\bullet))$
를 생각하자.
$$
\tau(\alpha)_{i_0 \ldots i_p} = (-1)^{p(p + 1)/2} \alpha_{i_p \ldots i_0}.
$$
그러면 차수 $n$인 원소 $\alpha$에 대해
\begin{align*}
& d(\tau(\alpha))_{i_0 \ldots i_{p + 1}} \\
& =
\sum\nolimits_{j = 0}^{p + 1}
(-1)^j
\tau(\alpha)_{i_0 \ldots \hat i_j \ldots i_{p + 1}}
+
(-1)^{p + 1}
d_{\mathcal F}(\tau(\alpha)_{i_0 \ldots i_{p + 1}})
\\
& =
\sum\nolimits_{j = 0}^{p + 1}
(-1)^{j + \frac{p(p + 1)}{2}}
\alpha_{i_{p + 1} \ldots \hat i_j \ldots i_0}
+
(-1)^{p + 1 + \frac{(p + 1)(p + 2)}{2}}
d_{\mathcal F}(\alpha_{i_{p + 1} \ldots i_0})
\end{align*}
이다. 한편
\begin{align*}
& \tau(d(\alpha))_{i_0\ldots i_{p + 1}} \\
& =
(-1)^{\frac{(p + 1)(p + 2)}{2}} d(\alpha)_{i_{p + 1} \ldots i_0}
\\
& =
(-1)^{\frac{(p + 1)(p + 2)}{2}}
\left(
\sum\nolimits_{j = 0}^{p + 1}
(-1)^j
\alpha_{i_{p + 1}\ldots \hat i_{p + 1 - j} \ldots i_0}
+
(-1)^{p + 1}
d_{\mathcal F}(\alpha_{i_{p + 1}\ldots i_0})
\right)
\end{align*}
이다. $p(p + 1)/2 \equiv (p + 1)(p + 2)/2 + p + 1 \bmod 2$이므로
$d(\tau(\alpha)) = \tau(d(\alpha))$임을 얻는다. 다시 말해 $\tau$는
복합체
$\text{Tot}(\check{\mathcal{C}}^\bullet({\mathcal U},
{\mathcal F}^\bullet))$의 자기 사상이다. 도표
$$
\begin{matrix}
\Gamma(X, {\mathcal F}^\bullet) &
\longrightarrow &
\text{Tot}(\check{\mathcal{C}}^\bullet({\mathcal U}, {\mathcal F}^\bullet)) \\
\downarrow \text{id} & & \downarrow \tau \\
\Gamma(X, {\mathcal F}^\bullet) &
\longrightarrow &
\text{Tot}(\check{\mathcal{C}}^\bullet({\mathcal U}, {\mathcal F}^\bullet))
\end{matrix}
$$
가 가환함에 유의하자. 또한 $\tau$는 분명히 세분과 양립한다.
이는 $\tau$가 체흐 코호몰로지 위에서 항등사상으로 작용함을 시사한다.
즉, 여극한에서 그러하다. 다만 체흐 초코호몰로지가 초코호몰로지와
일치한다고 가정하며, 초덮개를 쓰면 이는 항상 성립한다. 실제로 $\tau$가
전체 체흐 복합체
$\text{Tot}(\check{\mathcal{C}}^\bullet({\mathcal U},
{\mathcal F}^\bullet))$ 위에서 항등사상과 호모토픽이라고 주장한다.
이를 증명하기 위해 차수 $n$인 $\alpha$에 대해 다음을 호모토피로
사용한다.
$$
h(\alpha)_{i_0 \ldots i_p}
=
\sum\nolimits_{a = 0}^p
\epsilon_p(a)
\alpha_{i_0 \ldots i_a i_p \ldots i_a}
\quad\text{with}\quad
\epsilon_p(a) = (-1)^{\frac{(p - a)(p - a - 1)}{2} + p}
$$
보통처럼 $|_{U_{i_0 \ldots i_p}}$는 쓰지 않는다. 이는 다음 계산으로
확인된다. 여기서 다시 $\alpha$는 차수 $n$인 원소이다.
\begin{align*}
(d(h(\alpha)) + h(d(\alpha)))_{i_0 \ldots i_p}
= &
\sum\nolimits_{k = 0}^p
(-1)^k
h(\alpha)_{i_0 \ldots \hat i_k \ldots i_p}
+ \\
&
(-1)^p
d_{\mathcal F}(h(\alpha)_{i_0 \ldots i_p})
+ \\
&
\sum\nolimits_{a = 0}^p
\epsilon_p(a)
d(\alpha)_{i_0 \ldots i_a i_p \ldots i_a}
\\
= &
\sum\nolimits_{k = 0}^p
\sum\nolimits_{a = 0}^{k - 1}
(-1)^k \epsilon_{p - 1}(a)
\alpha_{i_0 \ldots i_a i_p \ldots \hat{i_k} \ldots i_a}
+ \\
&
\sum\nolimits_{k = 0}^p
\sum\nolimits_{a = k + 1}^p
(-1)^k \epsilon_{p - 1}(a - 1)
\alpha_{i_0 \ldots \hat{i_k} \ldots i_a i_p \ldots i_a}
+ \\
&
\sum\nolimits_{a = 0}^p
(-1)^p \epsilon_p(a)
d_{\mathcal F}(\alpha_{i_0 \ldots i_a i_p \ldots i_a})
+ \\
&
\sum\nolimits_{a = 0}^p
\sum\nolimits_{k = 0}^a
\epsilon_p(a) (-1)^k
\alpha_{i_0 \ldots \hat{i_k} \ldots i_a i_p \ldots i_a}
+ \\
&
\sum\nolimits_{a = 0}^p
\sum\nolimits_{k = a}^p
\epsilon_p(a) (-1)^{p + a + 1 - k}
\alpha_{i_0 \ldots i_a i_p \ldots \hat{i_k} \ldots i_a}
+ \\
&
\sum\nolimits_{a = 0}^p
\epsilon_p(a) (-1)^{p + 1}
d_{\mathcal F}(\alpha_{i_0 \ldots i_a i_p \ldots i_a})
\\
= &
\epsilon_p(0) \alpha_{i_p \ldots i_0} +
\epsilon_p(p) (-1)^{p + 1} \alpha_{i_0 \ldots i_p} \\
= &
(-1)^{\frac{p(p + 1)}{2}}\alpha_{i_p \ldots i_0}
- \alpha_{i_0 \ldots i_p}
\end{align*}
다음 등식들 때문에 상쇄가 일어난다.
$$
(-1)^k \epsilon_{p - 1}(a) + \epsilon_p(a)(-1)^{p + a + 1 - k} = 0
\quad\text{and}\quad
(-1)^k\epsilon_{p - 1}(a - 1) + \epsilon_p(a) (-1)^k = 0
$$
상쇄의 확인은 독자에게 맡긴다.

\medskip\noindent

아벨 층들로 이루어진 아래로 유계인 두 복합체
${\mathcal F}^\bullet$과 ${\mathcal G}^\bullet$이 있다고 하자.
복합체 $\text{Tot}({\mathcal F}^\bullet\otimes_{\mathbf Z}
{\mathcal G}^\bullet)$을 항들이
$\bigoplus_{p + q = n} {\mathcal F}^p \otimes {\mathcal G}^q$이고,
동차 원소 $\alpha$와 $\beta$에 대한 미분이 규칙
\begin{equation}
\label{cohomology-equation-differential-tensor-product-complexes}
d(\alpha \otimes \beta) =
d(\alpha)\otimes \beta + (-1)^{\deg(\alpha)} \alpha \otimes d(\beta)
\end{equation}
으로 주어지는 복합체라 정의한다. Homology, Definition
\ref{homology-definition-associated-simple-complex}을 보라.

\medskip\noindent
$M^\bullet$과 $N^\bullet$을 아벨 군으로 이루어진 아래로 유계인 두
복합체라 하자. $m$, 각각 $n$이 $M^\bullet$, 각각 $N^\bullet$의
코사이클이면, $m \otimes n$은
$\text{Tot}(M^\bullet\otimes N^\bullet)$의 코사이클임이 바로
따른다. 따라서 컵곱
$$
H^i(M^\bullet) \times H^j(N^\bullet)
\longrightarrow
H^{i + j}(Tot(M^\bullet\otimes N^\bullet)).
$$
을 얻는다. 이는 More on Algebra, Section
\ref{more-algebra-section-products-tor}에서도 논의한다.

\medskip\noindent
그러므로 복합체의 초코호몰로지에서 컵곱을 구성하는 일은 복합체 사상
\begin{equation}
\label{cohomology-equation-needs-signs}
\text{Tot}\left(
\text{Tot}(\check{\mathcal{C}}^\bullet({\mathcal U}, {\mathcal F}^\bullet))
\otimes_{\mathbf Z}
\text{Tot}(\check{\mathcal{C}}^\bullet({\mathcal U}, {\mathcal G}^\bullet))
\right)
\longrightarrow
\text{Tot}(
\check{\mathcal{C}}^\bullet({\mathcal U},
\text{Tot}({\mathcal F}^\bullet\otimes {\mathcal G}^\bullet)
))
\end{equation}
의 구성에 달려 있다. 이 사상을 $\cup$로 나타내며 다음 규칙으로
주어진다.
$$
(\alpha \cup \beta)_{i_0 \ldots i_p}
=
\sum\nolimits_{r = 0}^p
\epsilon(n, m, p, r)
\alpha_{i_0 \ldots i_r} \otimes \beta_{i_r \ldots i_p}.
$$
여기서 $\alpha$의 차수는 $n$이고 $\beta$의 차수는 $m$이며
$$
\epsilon(n, m, p, r) = (-1)^{(p + r)n + rp + r}.
$$
로 둔다. $\epsilon(n, m, p, n) = 1$임에 유의하자. 따라서
$\mathcal{F}^\bullet = \mathcal{F}[0]$가 차수 $0$에 놓인 하나의
아벨 층 $\mathcal{F}$로 이루어진 복합체이면 $\cup$의 공식에는
부호가 없다(이 경우 $r = n$이 아니면
$\alpha_{i_0 \ldots i_r} = 0$이기 때문이다). 부호가 있어야 하는
이유와 그 계산법은 Deligne의 \cite[Exposee XVII]{SGA4}를 보라.
(\ref{cohomology-equation-needs-signs})가 복합체 사상임을 확인하려면
$\text{Tot}(
\text{Tot}(\check{\mathcal{C}}^\bullet({\mathcal U}, {\mathcal F}^\bullet))
\otimes_{\mathbf Z}
\text{Tot}(\check{\mathcal{C}}^\bullet({\mathcal U}, {\mathcal G}^\bullet))
)$
위의 미분 정의(Homology, Definition
\ref{homology-definition-associated-simple-complex})에 따라
$$
d(\alpha \cup \beta) =
d(\alpha) \cup \beta +
(-1)^{\deg(\alpha)} \alpha \cup d(\beta)
$$
임을 보여야 한다. 먼저
\begin{align*}
d(\alpha \cup \beta)_{i_0 \ldots i_{p + 1}}
= &
\sum\nolimits_{j = 0}^{p + 1}
(-1)^j
(\alpha \cup \beta)_{i_0 \ldots \hat i_j \ldots i_{p + 1}}
+
(-1)^{p + 1}
d_{{\mathcal F} \otimes {\mathcal G}}
((\alpha \cup \beta)_{i_0 \ldots i_{p + 1}})
\\
= &
\sum\nolimits_{j = 0}^{p + 1}
\sum\nolimits_{r = 0}^{j - 1}
(-1)^j \epsilon(n, m, p, r)
\alpha_{i_0 \ldots i_r} \otimes \beta_{i_r \ldots \hat i_j \ldots i_{p + 1}}
+ \\
&
\sum\nolimits_{j = 0}^{p + 1}
\sum\nolimits_{r = j + 1}^{p + 1}
(-1)^j \epsilon(n, m, p, r - 1)
\alpha_{i_0 \ldots \hat i_j \ldots i_r} \otimes \beta_{i_r \ldots i_{p + 1}}
+ \\
&
\sum\nolimits_{r = 0}^{p + 1}
(-1)^{p + 1} \epsilon(n, m, p + 1, r)
d_{{\mathcal F} \otimes {\mathcal G}}
(\alpha_{i_0 \ldots i_r} \otimes \beta_{i_r \ldots i_{p + 1}})
\end{align*}
를 계산한다. 마지막 항의 직합항들은
$$
(-1)^{p + 1} \epsilon(n, m, p + 1, r)
\left(
d_{\mathcal F}(\alpha_{i_0 \ldots i_r}) \otimes 
\beta_{i_r \ldots i_{p + 1}} +
(-1)^{n - r}
\alpha_{i_0 \ldots i_r} \otimes d_{\mathcal G}(\beta_{i_r \ldots i_{p + 1}})
\right).
$$
와 같다. 실제로
$\deg_\mathcal{F}(\alpha_{i_0 \ldots i_r}) = n - r$이다. 한편
\begin{align*}
(d(\alpha) \cup \beta)_{i_0\ldots i_{p + 1}}
= &
\sum\nolimits_{r = 0}^{p + 1}
\epsilon(n + 1, m, p + 1, r)
d(\alpha)_{i_0\ldots i_r} \otimes \beta_{i_r\ldots i_{p + 1}}
\\
= &
\sum\nolimits_{r = 0}^{p + 1}
\sum\nolimits_{j = 0}^{r}
\epsilon(n + 1, m, p + 1, r) (-1)^j
\alpha_{i_0\ldots\hat{i_j}\ldots i_r} \otimes \beta_{i_r\ldots i_{p + 1}}
+ \\
&
\sum\nolimits_{r = 0}^{p + 1}
\epsilon(n + 1, m, p + 1, r) (-1)^r
d_{\mathcal F}(\alpha_{i_0 \ldots i_r}) \otimes \beta_{i_r\ldots i_{p + 1}}
\end{align*}
이고,
\begin{align*}
(\alpha \cup d(\beta))_{i_0\ldots i_{p + 1}}
= &
\sum\nolimits_{r = 0}^{p + 1}
\epsilon(n, m + 1, p + 1, r)
\alpha_{i_0 \ldots i_r} \otimes d(\beta)_{i_r \ldots i_{p + 1}}
\\
= &
\sum\nolimits_{r = 0}^{p + 1}
\sum\nolimits_{j = r}^{p + 1}
\epsilon(n, m + 1, p + 1, r) (-1)^{j - r}
\alpha_{i_0 \ldots i_r} \otimes \beta_{i_r \ldots \hat{i_j}\ldots i_{p + 1}}
+ \\
&
\sum\nolimits_{r = 0}^{p + 1}
\epsilon(n, m + 1, p + 1, r) (-1)^{p + 1 - r}
\alpha_{i_0 \ldots i_r} \otimes d_{\mathcal G}(\beta_{i_r \ldots i_{p + 1}})
\end{align*}
이다. 다음 등식들이 성립하면 원하는 등식이 성립한다.
\begin{align*}
(-1)^{p + 1} \epsilon(n, m, p + 1, r)
& =
\epsilon(n + 1, m, p + 1, r) (-1)^r \\
(-1)^{p + 1} \epsilon(n, m, p + 1, r) (-1)^{n - r}
& =
(-1)^n \epsilon(n, m + 1, p + 1, r) (-1)^{p + 1 - r} \\
\epsilon(n + 1, m, p + 1, r) (-1)^r
& =
(-1)^{1 + n} \epsilon(n, m + 1, p + 1, r - 1) \\
(-1)^j \epsilon(n, m, p, r)
& =
(-1)^n \epsilon(n, m + 1, p + 1, r) (-1)^{j - r} \\
(-1)^j \epsilon(n, m, p, r - 1)
& =
\epsilon(n + 1, m, p + 1, r) (-1)^j
\end{align*}
(셋째 등식은 $d(\alpha) \cup \beta$와
$(-1)^n \alpha \cup d(\beta)$에서 $r = j$인 항들이 서로 상쇄되게
하는 데 필요하다.) 확인은 독자에게 맡긴다. (또는 Stacks 프로젝트의
scripts 하위 디렉터리에 있는 signs.gp 스크립트를 확인하라.)

\medskip\noindent
컵곱의 결합법칙을 보자. ${\mathcal F}^\bullet$,
${\mathcal G}^\bullet$, ${\mathcal H}^\bullet$을 $X$ 위의 아벨 군으로
이루어진 아래로 유계인 복합체라 하자. 부호가 개입하지 않는 명백한
사상은 복합체의 동형사상
$$
\text{Tot}(
\text{Tot}({\mathcal F}^\bullet \otimes_{\mathbf Z} {\mathcal G}^\bullet)
\otimes_{\mathbf Z} {\mathcal H}^\bullet
)
\longrightarrow
\text{Tot}(
{\mathcal F}^\bullet \otimes_{\mathbf Z}
\text{Tot}({\mathcal G}^\bullet \otimes_{\mathbf Z} {\mathcal H}^\bullet)
).
$$
이다. 달리 말하면, 삼중 복합체
${\mathcal F}^\bullet \otimes_{\mathbf Z} {\mathcal G}^\bullet
\otimes_{\mathbf Z} {\mathcal H}^\bullet$은 잘 정의된 전체 복합체를
주며, 동차 원소에 대해 그 미분은
$$
d(\alpha \otimes \beta \otimes \gamma) =
d(\alpha) \otimes \beta \otimes \gamma +
(-1)^{\deg(\alpha)} \alpha \otimes d(\beta) \otimes \gamma +
(-1)^{\deg(\alpha) + \deg(\beta)} \alpha \otimes \beta \otimes d(\gamma)
$$
를 만족한다. 이 사상을 사용하면
$$
(\alpha \cup \beta) \cup \gamma = \alpha \cup ( \beta \cup \gamma)
$$
임을 쉽게 확인할 수 있다. 실제로 $\alpha$의 차수가 $a$이고
$\beta$의 차수가 $b$이며 $\gamma$의 차수가 $c$이면
\begin{align*}
((\alpha \cup \beta) \cup \gamma)_{i_0 \ldots i_p}
= &
\sum\nolimits_{r = 0}^p
\epsilon(a + b, c, p, r)
(\alpha \cup \beta)_{i_0 \ldots i_r} \otimes \gamma_{i_r \ldots i_p}
\\
= &
\sum\nolimits_{r = 0}^p
\sum\nolimits_{s = 0}^r
\epsilon(a + b, c, p, r) \epsilon(a, b, r, s)
\alpha_{i_0 \ldots i_s} \otimes
\beta_{i_s \ldots i_r} \otimes
\gamma_{i_r \ldots i_p}
\end{align*}
이고,
\begin{align*}
(\alpha \cup (\beta \cup \gamma)_{i_0\ldots i_p}
= &
\sum\nolimits_{s = 0}^p
\epsilon(a, b + c, p, s)
\alpha_{i_0 \ldots i_s} \otimes (\beta \cup \gamma)_{i_s \ldots i_p}
\\
= &
\sum\nolimits_{s = 0}^p
\sum\nolimits_{r = s}^p
\epsilon(a, b + c, p, s) \epsilon(b, c, p - s, r - s)
\alpha_{i_0 \ldots i_s} \otimes \beta_{i_s \ldots i_r} \otimes
\gamma_{i_r \ldots i_p}
\end{align*}
이다. 자명한 법 $2$ 계산으로 부호들이 일치함을 알 수 있다.
(또는 Stacks 프로젝트의 scripts 하위 디렉터리에 있는 signs.gp
스크립트를 확인하라.)

\medskip\noindent
끝으로 컵곱이 적어도 코호몰로지 수준에서 등급 가환 구조를 보존하는
이유를 설명한다. 이를 위해 위에서 도입한 연산자 $\tau$를 사용한다.
${\mathcal F}^\bullet$을 아벨 군으로 이루어진 아래로 유계인 복합체라
하고, 등급 가환 곱셈
$$
\wedge^\bullet :
\text{Tot}({\mathcal F}^\bullet\otimes {\mathcal F}^\bullet)
\longrightarrow
{\mathcal F}^\bullet.
$$
이 주어졌다고 하자. 이는 다음을 뜻한다. $s$가 ${\mathcal F}^a$의
국소 절단이고 $t$가 ${\mathcal F}^b$의 국소 절단이면
$s \wedge t$는 ${\mathcal F}^{a + b}$의 국소 절단이다.
등급 가환이라는 것은 $s \wedge t = (-1)^{ab} t \wedge s$라는 뜻이다.
$\wedge$가 복합체 사상이므로
$d(s\wedge t) = d(s) \wedge t + (-1)^a s \wedge d(t)$이다. 합성
$$
\text{Tot}(
\text{Tot}(\check{\mathcal{C}}^\bullet({\mathcal U}, {\mathcal F}^\bullet))
\otimes
\text{Tot}(\check{\mathcal{C}}^\bullet({\mathcal U}, {\mathcal F}^\bullet))
)
\to
\text{Tot}(
\check{\mathcal{C}}^\bullet({\mathcal U},
\text{Tot}({\mathcal F}^\bullet\otimes_{\mathbf Z}{\mathcal F}^\bullet))
)
\to
\text{Tot}(\check{\mathcal{C}}^\bullet({\mathcal U}, {\mathcal F}^\bullet))
$$
은 코호몰로지 위의 컵곱
$$
H^n(
\text{Tot}(\check{\mathcal{C}}^\bullet({\mathcal U}, {\mathcal F}^\bullet))
)
\times
H^m(
\text{Tot}(\check{\mathcal{C}}^\bullet({\mathcal U}, {\mathcal F}^\bullet))
)
\longrightarrow
H^{n + m}(
\text{Tot}(\check{\mathcal{C}}^\bullet({\mathcal U}, {\mathcal F}^\bullet))
)
$$
을 유도한다. 따라서 여극한에서 체흐 코호몰로지 위의 곱을 유도하고,
그 결과 (필요하면 초덮개를 사용하여) ${\mathcal F}^\bullet$의
코호몰로지 위의 곱을 유도한다. 이 곱도 코호몰로지 위에서 등급
가환이라고 주장한다. 이를 증명하기 위해 먼저
$\text{Tot}(\check{\mathcal{C}}^\bullet({\mathcal U},
{\mathcal F}^\bullet))$에서 차수 $n$인 원소 $\alpha$와
$\text{Tot}(\check{\mathcal{C}}^\bullet({\mathcal U},
{\mathcal F}^\bullet))$에서 차수 $m$인 원소 $\beta$를 택하여
다음을 계산한다.
\begin{align*}
\wedge^\bullet(\alpha \cup \beta)_{i_0 \ldots i_p}
= &
\sum\nolimits_{r = 0}^p
\epsilon(n, m, p, r)
\alpha_{i_0 \ldots i_r} \wedge \beta_{i_r \ldots i_p} \\
= &
\sum\nolimits_{r = 0}^p
\epsilon(n, m, p, r)
(-1)^{\deg(\alpha_{i_0 \ldots i_r})\deg(\beta_{i_r \ldots i_p})}
\beta_{i_r \ldots i_p} \wedge \alpha_{i_0 \ldots i_r}
\end{align*}
이는 $\wedge$가 등급 가환이기 때문이다. 한편
\begin{align*}
\tau(\wedge^\bullet(\tau(\beta) \cup \tau(\alpha)))_{i_0 \ldots i_p}
= &
\chi(p)
\sum\nolimits_{r = 0}^p
\epsilon(m, n, p, r)
\tau(\beta)_{i_p \ldots i_{p - r}} \wedge \tau(\alpha)_{i_{p - r} \ldots i_0}
\\
= &
\chi(p)
\sum\nolimits_{r = 0}^p
\epsilon(m, n, p, r) \chi(r) \chi(p - r)
\beta_{i_{p - r} \ldots i_p} \wedge \alpha_{i_0 \ldots i_{p - r}}
\\
= &
\chi(p)
\sum\nolimits_{r = 0}^p
\epsilon(m, n, p, p - r) \chi(r) \chi(p - r)
\beta_{i_r \ldots i_p} \wedge \alpha_{i_0 \ldots i_r}
\end{align*}
이다. 여기서 $\chi(t) = (-1)^{\frac{t(t + 1)}{2}}$이다. 앞에서
$\tau$가 코호몰로지 위에서 항등사상으로 작용함을 증명했으므로
$$
\epsilon(n, m, p, r)
(-1)^{(n - r)(m - (p - r))}
=
(-1)^{nm} \chi(p)\epsilon(m, n, p, p - r) \chi(r) \chi(p - r)
$$
임을 확인해야 한다. 자명한 법 $2$ 계산으로 부호들이 일치함을 알 수
있다. (또는 Stacks 프로젝트의 scripts 하위 디렉터리에 있는 signs.gp
스크립트를 확인하라.)

\medskip\noindent
끝으로 컵곱과 경계 사상의 양립성을 연구한다. 다음이
$$
0
\to
{\mathcal F}_1^\bullet
\to
{\mathcal F}_2^\bullet
\to
{\mathcal F}_3^\bullet
\to
0
\quad\text{and}\quad
0
\leftarrow
{\mathcal G}_1^\bullet
\leftarrow
{\mathcal G}_2^\bullet
\leftarrow
{\mathcal G}_3^\bullet
\leftarrow
0
$$
$X$ 위의 아벨 층으로 이루어진 아래로 유계인 복합체의 짧은 완전열이라고
하자. ${\mathcal H}^\bullet$을 아벨 층으로 이루어진 또 하나의 아래로
유계인 복합체라 하고, 복합체 사이의 사상과 양립하는 복합체 사상
$$
\gamma_i :
\text{Tot}({\mathcal F}_i^\bullet \otimes_{\mathbf Z} {\mathcal G}_i^\bullet)
\longrightarrow
{\mathcal H}^\bullet
$$
이 주어졌다고 하자. 즉, 도표
$$
\xymatrix{
\text{Tot}({\mathcal F}_1^\bullet \otimes_{\mathbf Z} {\mathcal G}_1^\bullet)
\ar[d]_{\gamma_1}
&
\text{Tot}({\mathcal F}_1^\bullet \otimes_{\mathbf Z} {\mathcal G}_2^\bullet)
\ar[l] \ar[d]
\\
\mathcal{H}^\bullet &
\text{Tot}({\mathcal F}_2^\bullet \otimes_{\mathbf Z} {\mathcal G}_2^\bullet)
\ar[l]_-{\gamma_2}
}
$$
와
$$
\xymatrix{
\text{Tot}({\mathcal F}_2^\bullet \otimes_{\mathbf Z} {\mathcal G}_2^\bullet)
\ar[d]_{\gamma_2}
&
\text{Tot}({\mathcal F}_2^\bullet \otimes_{\mathbf Z} {\mathcal G}_3^\bullet)
\ar[l] \ar[d]
\\
\mathcal{H}^\bullet &
\text{Tot}({\mathcal F}_3^\bullet \otimes_{\mathbf Z} {\mathcal G}_3^\bullet)
\ar[l]_-{\gamma_3}
}
$$
가 가환한다고 하자.

\begin{lemma}
\label{cohomology-lemma-compute-sign-cup-product-boundaries}
위의 상황에서 층 $\mathcal{F}_i^p$와 $\mathcal{G}_j^q$에 대해
체흐 코호몰로지가 코호몰로지와 일치한다고 가정하자.
$a_3 \in H^n(X, \mathcal{F}_3^\bullet)$ 및
$b_1 \in H^m(X, \mathcal{G}_1^\bullet)$라 하자. 그러면
$$
\gamma_1( \partial a_3 \cup b_1) =
(-1)^{n + 1} \gamma_3( a_3 \cup \partial b_1)
$$
이다. 이 등식은 $H^{n + m + 1}(X, \mathcal{H}^\bullet)$에서 성립하며,
여기서 $\partial$는 위의 복합체의 짧은 완전열에 딸린 코호몰로지의
경계 사상을 뜻한다.
\end{lemma}

\begin{proof}
다음 규약과 표기를 사용한다. ${\mathcal F}_1^p$를
${\mathcal F}_2^p$의 부분층으로 생각하고, ${\mathcal G}_3^q$를
${\mathcal G}_2^q$의 부분층으로 생각한다. 따라서 $s$가
${\mathcal F}_1^p$의 국소 절단이면 ${\mathcal F}_2^p$의 대응하는
절단도 $s$로 나타낸다. ${\mathcal G}_3^q$의 국소 절단에 대해서도
마찬가지이다. 또한 $s$가 ${\mathcal F}_2^p$의 국소 절단이면
${\mathcal F}_3^p$에서의 상을 $\bar s$로 나타낸다. 사상
${\mathcal G}_2^q \to {\mathcal G}^q_1$에 대해서도 마찬가지이다.
특히 $s$가 ${\mathcal F}_2^p$의 국소 절단이고 $\bar s = 0$이면
$s$는 ${\mathcal F}_1^p$의 국소 절단이다. 위 도표들의 가환성에서,
${\mathcal F}_2^p$의 국소 절단 $s$와 ${\mathcal G}_3^q$의 국소
절단 $t$에 대해
$\gamma_2(s \otimes t) = \gamma_3(\bar s \otimes t)$가
${\mathcal H}^{p + q}$의 절단으로서 성립한다.

\medskip\noindent
${\mathcal U} : X =  \bigcup_{i \in I} U_i$를 $X$의 열린 덮개라 하자.
$\alpha_3$, 각각 $\beta_1$이
$\text{Tot}(
\check{\mathcal{C}}^\bullet({\mathcal U}, {\mathcal F}_3^\bullet))$,
각각 $\text{Tot}(
\check{\mathcal{C}}^\bullet({\mathcal U}, {\mathcal G}_1^\bullet))$의
차수 $n$, 각각 $m$인 코사이클이고, $a_3$, 각각 $b_1$을 나타낸다고
하자. 필요하면 $\mathcal{U}$를 세분하여,
$\text{Tot}(
\check{\mathcal{C}}^\bullet({\mathcal U}, {\mathcal F}_2^\bullet))$,
각각 $\text{Tot}(
\check{\mathcal{C}}^\bullet({\mathcal U}, {\mathcal G}_2^\bullet))$ 안에서
차수 $n$, 각각 $m$이고 $\alpha_3$, 각각 $\beta_1$로 보내지는
코체인 $\alpha_2$, 각각 $\beta_2$를 찾을 수 있다. 그러면
$$
\overline{d(\alpha_2)} = d(\bar \alpha_2) = 0
\quad\text{and}\quad
\overline{d(\beta_2)} = d(\bar \beta_2) = 0.
$$
임을 알 수 있다. 이는 $\alpha_1 = d(\alpha_2)$가
$\text{Tot}(\check{\mathcal{C}}^\bullet({\mathcal U},
{\mathcal F}_1^\bullet))$ 안의 차수 $n + 1$인 코사이클이며
$\partial a_3$을 나타낸다는 뜻이다. 마찬가지로
$\beta_3 = d(\beta_2)$는
$\text{Tot}(\check{\mathcal{C}}^\bullet({\mathcal U},
{\mathcal G}_3^\bullet))$ 안의 차수 $m + 1$인 코사이클이며
$\partial b_1$을 나타낸다. 따라서
\begin{align*}
d(\gamma_2(\alpha_2 \cup \beta_2))
& =
\gamma_2(d(\alpha_2 \cup \beta_2))
\\
& =
\gamma_2(d(\alpha_2) \cup \beta_2 + (-1)^n \alpha_2 \cup d(\beta_2) )
\\
& =
\gamma_2( \alpha_1 \cup \beta_2)  + (-1)^n \gamma_2( \alpha_2 \cup \beta_3)
\\
& =
\gamma_1(\alpha_1 \cup \beta_1) + (-1)^n \gamma_3(\alpha_3 \cup \beta_3)
\end{align*}
로 계산할 수 있다. 이는 보조정리에 적힌 것처럼 부호가
$(-1)^{n + 1}$임까지 알려 준다.\footnote{부호는 $D(X)$의 삼각형에
딸린 코호몰로지 긴 완전열의 부호 규약에 의존한다. Stacks 프로젝트의
규약은 (a) 특수 삼각형이 항별 분할 완전열에 대응하고, (b) 긴 완전열의
경계 사상이 부호의 개입 없이 뱀 보조정리의 사상으로 주어진다는 것이다.
Derived Categories, Section
\ref{derived-section-homotopy-triangulated}를 보라.}
\end{proof}

\begin{lemma}
\label{cohomology-lemma-boundary-derivation-over-cup-product}
$X$를 위상공간이라 하고, $\mathcal{O}' \to \mathcal{O}$를 환의 층의
전사라 하자. 그 핵 $\mathcal{I} \subset \mathcal{O}'$의 제곱이
영이라고 하자. 그러면 $M = H^1(X, \mathcal{I})$은
$R = H^0(X, \mathcal{O})$-가군이고, 짧은 완전열
$$
0 \to \mathcal{I} \to \mathcal{O}' \to \mathcal{O} \to 0
$$
에 딸린 경계 사상 $\partial : R \to M$은 미분이다(Algebra,
Definition \ref{algebra-definition-derivation}).
\end{lemma}

\begin{proof}
가정에 의해 $\mathcal{I} \cdot \mathcal{I} = 0$이므로 사상
$\mathcal{O}' \to \SheafHom(\mathcal{I}, \mathcal{I})$은
$\mathcal{O}$를 통해 인수분해된다. 따라서 $\mathcal{I}$는
$\mathcal{O}$-가군의 층이고, 이것이 $M$ 위의 $R$-가군 구조를
정의한다. 경계 사상은 가법적이므로 라이프니츠 법칙을 증명하면 충분하다.
$f \in R$라 하자. $f|_{U_i} \in \mathcal{O}(U_i)$를 올리는
$f_i \in \mathcal{O}'(U_i)$가 존재하도록 열린 덮개
$\mathcal{U} : X = \bigcup U_i$를 택하자.
$f_i - f_j$는 $\mathcal{I}(U_i \cap U_j)$의 원소임에 유의하자.
그러면 $\partial(f)$는
$\alpha_{i_0i_1} = f_{i_0} - f_{i_1}$인 $1$-코사이클 $\alpha$의
체흐 코호몰로지류에 대응한다. (Lemma \ref{cohomology-lemma-cech-h1}에 의해
$\mathcal{U}$에 관한 첫째 체흐 코호몰로지 군은 $M$의 부분가군임에
유의하자.) 다음으로 $g \in R$를 두 번째 원소라 하고, 필요하면 열린
덮개를 세분하여 $g_i \in \mathcal{O}'(U_i)$가
$g|_{U_i} \in \mathcal{O}(U_i)$를 올린다고 하자. 그러면
$\partial(g)$는 $\beta_{i_0i_1} = g_{i_0} - g_{i_1}$인 코사이클
$\beta$로 주어진다. $f_ig_i \in \mathcal{O}'(U_i)$가
$fg|_{U_i}$를 올리므로 $\partial(fg)$는 다음 코사이클 $\gamma$로
주어진다.
$$
\gamma_{i_0i_1} = f_{i_0}g_{i_0} - f_{i_1}g_{i_1} =
(f_{i_0} - f_{i_1})g_{i_0} + f_{i_1}(g_{i_0} - g_{i_1}) =
\alpha_{i_0i_1}g + f\beta_{i_0i_1}
$$
이는 $\mathcal{I}$ 위의 $\mathcal{O}$-가군 구조의 정의에 따른다.
이로써 라이프니츠 법칙이 증명되고 증명이 끝난다.
\end{proof}










\section{평탄 분해}
\label{cohomology-section-flat}

\noindent
이 절의 내용은 \cite{Spaltenstein}을 참고할 수 있다.
$(X, \mathcal{O}_X)$를 환 달린 공간이라 하자.
Modules, Lemma \ref{modules-lemma-module-quotient-flat}에 의해
모든 $\mathcal{O}_X$-가군은 평탄 $\mathcal{O}_X$-가군의 몫이다.
Derived Categories, Lemma \ref{derived-lemma-subcategory-left-resolution}에
의해 모든 위로 유계인 $\mathcal{O}_X$-가군 복합체는 평탄
$\mathcal{O}_X$-가군들로 이루어진 위로 유계인 복합체에 의한 왼쪽
분해를 갖는다. 그러나 비유계 복합체의 경우에는 평탄 분해만으로는
충분하지 않음이 드러난다.

\begin{lemma}
\label{cohomology-lemma-derived-tor-exact}
$(X, \mathcal{O}_X)$를 환 달린 공간이라 하자.
$\mathcal{G}^\bullet$를 $\mathcal{O}_X$-가군 복합체라 하자. 함자
$$
K(\textit{Mod}(\mathcal{O}_X))
\longrightarrow
K(\textit{Mod}(\mathcal{O}_X)),
\quad
\mathcal{F}^\bullet \longmapsto
\text{Tot}(\mathcal{G}^\bullet \otimes_{\mathcal{O}_X} \mathcal{F}^\bullet)
$$
와
$$
K(\textit{Mod}(\mathcal{O}_X))
\longrightarrow
K(\textit{Mod}(\mathcal{O}_X)),
\quad
\mathcal{F}^\bullet \longmapsto
\text{Tot}(\mathcal{F}^\bullet \otimes_{\mathcal{O}_X} \mathcal{G}^\bullet)
$$
는 삼각범주의 완전 함자이다.
\end{lemma}

\begin{proof}
이는 Derived Categories, Remark
\ref{derived-remark-double-complex-as-tensor-product-of}에서 따른다.
\end{proof}

\begin{definition}
\label{cohomology-definition-K-flat}
$(X, \mathcal{O}_X)$를 환 달린 공간이라 하자.
$\mathcal{O}_X$-가군 복합체 $\mathcal{K}^\bullet$가
{\it K-평탄}이라는 것은 모든 비순환 $\mathcal{O}_X$-가군 복합체
$\mathcal{F}^\bullet$에 대해 복합체
$$
\text{Tot}(\mathcal{F}^\bullet \otimes_{\mathcal{O}_X} \mathcal{K}^\bullet)
$$
가 비순환이라는 뜻이다.
\end{definition}

\begin{lemma}
\label{cohomology-lemma-K-flat-quasi-isomorphism}
$(X, \mathcal{O}_X)$를 환 달린 공간이라 하자.
$\mathcal{K}^\bullet$를 K-평탄 복합체라 하자. 그러면 함자
$$
K(\textit{Mod}(\mathcal{O}_X))
\longrightarrow
K(\textit{Mod}(\mathcal{O}_X)), \quad
\mathcal{F}^\bullet
\longmapsto
\text{Tot}(\mathcal{F}^\bullet \otimes_{\mathcal{O}_X} \mathcal{K}^\bullet)
$$
는 유사동형을 유사동형으로 보낸다.
\end{lemma}

\begin{proof}
Lemma \ref{cohomology-lemma-derived-tor-exact}과 유사동형은 원뿔이 비순환인
사상으로 특징지어진다는 사실에서 따른다.
\end{proof}

\begin{lemma}
\label{cohomology-lemma-check-K-flat-stalks}
$(X, \mathcal{O}_X)$를 환 달린 공간이라 하자. $\mathcal{K}^\bullet$를
$\mathcal{O}_X$-가군 복합체라 하자. 그러면 $\mathcal{K}^\bullet$가
K-평탄일 필요충분조건은 모든 $x \in X$에 대해
$\mathcal{O}_{X, x}$-가군 복합체 $\mathcal{K}_x^\bullet$가
K-평탄인 것이다(More on Algebra, Definition
\ref{more-algebra-definition-K-flat}).
\end{lemma}

\begin{proof}
모든 $x \in X$에 대해 $\mathcal{K}_x^\bullet$가 K-평탄이라고 하자.
$\otimes$와 직합은 줄기를 취하는 것과 가환하고 완전성은 줄기에서
확인할 수 있으므로 $\mathcal{K}^\bullet$는 K-평탄이다. Modules,
Lemma \ref{modules-lemma-abelian}을 보라. 역으로
$\mathcal{K}^\bullet$가 K-평탄이라고 가정하자. $x \in X$를 고르고,
$M^\bullet$를 비순환 $\mathcal{O}_{X, x}$-가군 복합체라 하자. 그러면
$i_{x, *}M^\bullet$는 비순환 $\mathcal{O}_X$-가군 복합체이다. 따라서
$\text{Tot}(i_{x, *}M^\bullet \otimes_{\mathcal{O}_X} \mathcal{K}^\bullet)$
는 비순환이다. $x$에서 줄기를 취하면
$\text{Tot}(M^\bullet \otimes_{\mathcal{O}_{X, x}} \mathcal{K}_x^\bullet)$
가 비순환임을 알 수 있다.
\end{proof}

\begin{lemma}
\label{cohomology-lemma-tensor-product-K-flat}
$(X, \mathcal{O}_X)$를 환 달린 공간이라 하자.
$\mathcal{K}^\bullet$, $\mathcal{L}^\bullet$가 K-평탄
$\mathcal{O}_X$-가군 복합체이면
$\text{Tot}(\mathcal{K}^\bullet \otimes_{\mathcal{O}_X} \mathcal{L}^\bullet)$
도 K-평탄 $\mathcal{O}_X$-가군 복합체이다.
\end{lemma}

\begin{proof}
다음 동형과 정의에서 따른다.
$$
\text{Tot}(\mathcal{M}^\bullet \otimes_{\mathcal{O}_X}
\text{Tot}(\mathcal{K}^\bullet \otimes_{\mathcal{O}_X} \mathcal{L}^\bullet))
=
\text{Tot}(\text{Tot}(\mathcal{M}^\bullet \otimes_{\mathcal{O}_X}
\mathcal{K}^\bullet) \otimes_{\mathcal{O}_X} \mathcal{L}^\bullet)
$$
\end{proof}

\begin{lemma}
\label{cohomology-lemma-K-flat-two-out-of-three}
$(X, \mathcal{O}_X)$를 환 달린 공간이라 하자.
$(\mathcal{K}_1^\bullet, \mathcal{K}_2^\bullet, \mathcal{K}_3^\bullet)$를
$K(\textit{Mod}(\mathcal{O}_X))$ 안의 특수 삼각형이라 하자.
세 $\mathcal{K}_i^\bullet$ 가운데 둘이 K-평탄이면 나머지 하나도
K-평탄이다.
\end{lemma}

\begin{proof}
Lemma \ref{cohomology-lemma-derived-tor-exact}과
$K(\textit{Mod}(\mathcal{O}_X))$ 안의 특수 삼각형에서 세 대상 가운데
둘이 비순환이면 나머지 하나도 비순환이라는 사실에서 따른다.
\end{proof}

\begin{lemma}
\label{cohomology-lemma-K-flat-two-out-of-three-ses}
$(X, \mathcal{O}_X)$를 환 달린 공간이라 하자.
$0 \to \mathcal{K}_1^\bullet \to \mathcal{K}_2^\bullet \to
\mathcal{K}_3^\bullet \to 0$를 복합체들의 짧은 완전열이라 하고,
$\mathcal{K}_3^\bullet$의 각 항은 평탄 $\mathcal{O}_X$-가군이라고 하자.
세 $\mathcal{K}_i^\bullet$ 가운데 둘이 K-평탄이면 나머지 하나도
K-평탄이다.
\end{lemma}

\begin{proof}
Modules, Lemma \ref{modules-lemma-flat-tor-zero}에 의해 모든 복합체
$\mathcal{L}^\bullet$에 대해 복합체들의 짧은 완전열
$$
0 \to
\text{Tot}(\mathcal{L}^\bullet \otimes_{\mathcal{O}_X} \mathcal{K}_1^\bullet)
\to
\text{Tot}(\mathcal{L}^\bullet \otimes_{\mathcal{O}_X} \mathcal{K}_1^\bullet)
\to
\text{Tot}(\mathcal{L}^\bullet \otimes_{\mathcal{O}_X} \mathcal{K}_1^\bullet)
\to 0
$$
을 얻는다. 따라서 보조정리는 코호몰로지 층들의 긴 완전열과
K-평탄 복합체의 정의에서 따른다.
\end{proof}

\begin{lemma}
\label{cohomology-lemma-pullback-K-flat}
$f : (X, \mathcal{O}_X) \to (Y, \mathcal{O}_Y)$를 환 달린 공간의
사상이라 하자. K-평탄 $\mathcal{O}_Y$-가군 복합체의 당김은
K-평탄 $\mathcal{O}_X$-가군 복합체이다.
\end{lemma}

\begin{proof}
이는 줄기에서 확인할 수 있다. Lemma
\ref{cohomology-lemma-check-K-flat-stalks}를 보라. 따라서 Sheaves, Lemma
\ref{sheaves-lemma-stalk-pullback-modules}와 More on Algebra, Lemma
\ref{more-algebra-lemma-base-change-K-flat}에서 따른다.
\end{proof}

\begin{lemma}
\label{cohomology-lemma-bounded-flat-K-flat}
$(X, \mathcal{O}_X)$를 환 달린 공간이라 하자. 평탄
$\mathcal{O}_X$-가군들로 이루어진 위로 유계인 복합체는 K-평탄이다.
\end{lemma}

\begin{proof}
이는 줄기에서 확인할 수 있다. Lemma
\ref{cohomology-lemma-check-K-flat-stalks}를 보라. 그러므로 이 보조정리는 Modules,
Lemma \ref{modules-lemma-flat-stalks-flat}와 More on Algebra, Lemma
\ref{more-algebra-lemma-derived-tor-quasi-isomorphism}에서 따른다.
\end{proof}

\noindent
다음 보조정리에서 복합체들의 계의 여극한은 항별 여극한을 뜻한다.

\begin{lemma}
\label{cohomology-lemma-colimit-K-flat}
$(X, \mathcal{O}_X)$를 환 달린 공간이라 하자.
$\mathcal{K}_1^\bullet \to \mathcal{K}_2^\bullet \to \ldots$를
K-평탄 복합체들의 계라 하자. 그러면
$\colim_i \mathcal{K}_i^\bullet$는 K-평탄이다.
\end{lemma}

\begin{proof}
항별 여극한을 취하고 있으므로
$$
\colim_i \text{Tot}(
\mathcal{F}^\bullet \otimes_{\mathcal{O}_X} \mathcal{K}_i^\bullet)
=
\text{Tot}(\mathcal{F}^\bullet \otimes_{\mathcal{O}_X}
\colim_i \mathcal{K}_i^\bullet)
$$
임이 명백하다. 따라서 보조정리는 여과 여극한이 완전하다는 사실에서
따른다.
\end{proof}

\begin{lemma}
\label{cohomology-lemma-resolution-by-direct-sums-extensions-by-zero}
$(X, \mathcal{O}_X)$를 환 달린 공간이라 하자. 임의의
$\mathcal{O}_X$-가군 복합체 $\mathcal{G}^\bullet$에 대해
$\mathcal{O}_X$-가군 복합체들의 가환도식
$$
\xymatrix{
\mathcal{K}_1^\bullet \ar[d] \ar[r] &
\mathcal{K}_2^\bullet \ar[d] \ar[r] & \ldots \\
\tau_{\leq 1}\mathcal{G}^\bullet \ar[r] &
\tau_{\leq 2}\mathcal{G}^\bullet \ar[r] & \ldots
}
$$
이 존재하며 다음 성질을 만족한다. (1) 수직 화살표들은 유사동형이고
항별 전사이다. (2) 각 $\mathcal{K}_n^\bullet$는 위로 유계인
복합체이며, 그 항들은 $j_{U!}\mathcal{O}_U$ 꼴의
$\mathcal{O}_X$-가군들의 직합이다. (3) 사상
$\mathcal{K}_n^\bullet \to \mathcal{K}_{n + 1}^\bullet$는 항별 분할
단사이고, 그 여핵들은 $j_{U!}\mathcal{O}_U$ 꼴의
$\mathcal{O}_X$-가군들의 직합이다. 더욱이 사상
$\colim \mathcal{K}_n^\bullet \to \mathcal{G}^\bullet$는 유사동형이다.
\end{lemma}

\begin{proof}
도식의 존재와 성질 (1), (2), (3)은 Modules, Lemma
\ref{modules-lemma-module-quotient-flat}과 Derived Categories, Lemma
\ref{derived-lemma-special-direct-system}에서 곧바로 따른다. 유도된 사상
$\colim \mathcal{K}_n^\bullet \to \mathcal{G}^\bullet$는 여과 여극한이
완전하므로 유사동형이다.
\end{proof}

\begin{lemma}
\label{cohomology-lemma-K-flat-resolution}
$(X, \mathcal{O}_X)$를 환 달린 공간이라 하자. 임의의 복합체
$\mathcal{G}^\bullet$에 대해, 각 항이 평탄 $\mathcal{O}_X$-가군인
$K$-평탄 복합체 $\mathcal{K}^\bullet$와 항별 전사인 유사동형
$\mathcal{K}^\bullet \to \mathcal{G}^\bullet$가 존재한다.
\end{lemma}

\begin{proof}
Lemma \ref{cohomology-lemma-resolution-by-direct-sums-extensions-by-zero}의 도식을
고르자. 각 복합체 $\mathcal{K}_n^\bullet$는 평탄 가군들로 이루어진
위로 유계인 복합체이다. Modules, Lemma
\ref{modules-lemma-j-shriek-flat}을 보라. 따라서 Lemma
\ref{cohomology-lemma-bounded-flat-K-flat}에 의해 $\mathcal{K}_n^\bullet$는
K-평탄이다. 그러므로 Lemma \ref{cohomology-lemma-colimit-K-flat}에 의해
$\colim \mathcal{K}_n^\bullet$는 K-평탄이다. 구성상 유도된 사상
$\colim \mathcal{K}_n^\bullet \to \mathcal{G}^\bullet$는
유사동형이며 항별 전사이다. Lemma
\ref{cohomology-lemma-resolution-by-direct-sums-extensions-by-zero}의 성질 (3)에
의하면 $\colim \mathcal{K}_n^m$은 평탄 가군들의 직합이므로 평탄이다.
이로써 마지막 주장도 증명된다.
\end{proof}

\begin{lemma}
\label{cohomology-lemma-derived-tor-quasi-isomorphism-other-side}
$(X, \mathcal{O}_X)$를 환 달린 공간이라 하자.
$\alpha : \mathcal{P}^\bullet \to \mathcal{Q}^\bullet$를 K-평탄
$\mathcal{O}_X$-가군 복합체들의 유사동형이라 하자. 모든
$\mathcal{O}_X$-가군 복합체 $\mathcal{F}^\bullet$에 대해 유도된 사상
$$
\text{Tot}(\text{id}_{\mathcal{F}^\bullet} \otimes \alpha) :
\text{Tot}(\mathcal{F}^\bullet \otimes_{\mathcal{O}_X} \mathcal{P}^\bullet)
\longrightarrow
\text{Tot}(\mathcal{F}^\bullet \otimes_{\mathcal{O}_X} \mathcal{Q}^\bullet)
$$
은 유사동형이다.
\end{lemma}

\begin{proof}
Lemma \ref{cohomology-lemma-K-flat-resolution}에 따라 $\mathcal{K}^\bullet$가
K-평탄인 유사동형
$\mathcal{K}^\bullet \to \mathcal{F}^\bullet$를 고르자. 다음
가환도식을 생각하자.
$$
\xymatrix{
\text{Tot}(\mathcal{K}^\bullet
\otimes_{\mathcal{O}_X} \mathcal{P}^\bullet) \ar[r] \ar[d] &
\text{Tot}(\mathcal{K}^\bullet
\otimes_{\mathcal{O}_X} \mathcal{Q}^\bullet) \ar[d] \\
\text{Tot}(\mathcal{F}^\bullet
\otimes_{\mathcal{O}_X} \mathcal{P}^\bullet) \ar[r] &
\text{Tot}(\mathcal{F}^\bullet
\otimes_{\mathcal{O}_X} \mathcal{Q}^\bullet)
}
$$
Lemma \ref{cohomology-lemma-K-flat-quasi-isomorphism}에 의해 수직 화살표들과 위쪽
수평 화살표가 유사동형이므로 결과가 따른다.
\end{proof}

\noindent
$(X, \mathcal{O}_X)$를 환 달린 공간이라 하자.
$\mathcal{F}^\bullet$를 $D(\mathcal{O}_X)$의 대상이라 하자. Lemma
\ref{cohomology-lemma-K-flat-resolution}에 따라 K-평탄 분해
$\mathcal{K}^\bullet \to \mathcal{F}^\bullet$를 고르자. Lemma
\ref{cohomology-lemma-derived-tor-exact}에 의해 삼각범주의 완전 함자
$$
K(\mathcal{O}_X)
\longrightarrow
K(\mathcal{O}_X),
\quad
\mathcal{G}^\bullet
\longmapsto
\text{Tot}(\mathcal{G}^\bullet \otimes_{\mathcal{O}_X} \mathcal{K}^\bullet)
$$
를 얻는다. Lemma \ref{cohomology-lemma-K-flat-quasi-isomorphism}에 의해 이 함자는
함자 $D(\mathcal{O}_X) \to D(\mathcal{O}_X)$를 유도한다. 이는
$D(\mathcal{O}_X)$가 $K(\mathcal{O}_X)$를 유사동형들에 대해
국소화한 것이기 때문이다. $\mathcal{F}^\bullet$의 $K$-평탄 분해들의
범주는 공여과되어 있고 Lemma
\ref{cohomology-lemma-derived-tor-quasi-isomorphism-other-side}가 성립하므로,
그 결과 얻는 함자는 동형을 무시하면 $\mathcal{K}^\bullet$의 선택에
의존하지 않는다.

\begin{definition}
\label{cohomology-definition-derived-tor}
$(X, \mathcal{O}_X)$를 환 달린 공간이라 하자.
$\mathcal{F}^\bullet$를 $D(\mathcal{O}_X)$의 대상이라 하자.
{\it 유도 텐서곱}
$$
- \otimes_{\mathcal{O}_X}^{\mathbf{L}} \mathcal{F}^\bullet :
D(\mathcal{O}_X)
\longrightarrow
D(\mathcal{O}_X)
$$
은 위에서 설명한 삼각범주의 완전 함자이다.
\end{definition}

\noindent
명시적인 구성에서 $D(\mathcal{O}_X)$ 안의
$\mathcal{G}^\bullet$와 $\mathcal{F}^\bullet$에 대해 표준 동형
$$
\mathcal{F}^\bullet \otimes_{\mathcal{O}_X}^{\mathbf{L}} \mathcal{G}^\bullet
\cong
\mathcal{G}^\bullet \otimes_{\mathcal{O}_X}^{\mathbf{L}} \mathcal{F}^\bullet
$$
이 존재함이 명백하다. 여기서는 More on Algebra, Section
\ref{more-algebra-section-sign-rules}의 부호 규칙을 사용한다. 따라서
$\mathcal{F}^\bullet \otimes_{\mathcal{O}_X}^{\mathbf{L}} \mathcal{G}^\bullet$
라고 쓸 때에는 보통 어느 변수를 사용해 유도 텐서곱을 정의하는지
구별하지 않을 것이다.

\begin{definition}
\label{cohomology-definition-tor}
$(X, \mathcal{O}_X)$를 환 달린 공간이라 하자.
$\mathcal{F}$, $\mathcal{G}$를 $\mathcal{O}_X$-가군이라 하자.
$\mathcal{F}$와 $\mathcal{G}$의 {\it Tor}들은 위에서 정의한 유도
텐서곱을 사용하여 다음 공식으로 정의한다.
$$
\text{Tor}_p^{\mathcal{O}_X}(\mathcal{F}, \mathcal{G}) =
H^{-p}(\mathcal{F} \otimes_{\mathcal{O}_X}^\mathbf{L} \mathcal{G})
$$
\end{definition}

\noindent
이 정의에 따르면 $\mathcal{O}_X$-가군들의 모든 짧은 완전열
$0 \to \mathcal{F}_1 \to \mathcal{F}_2 \to \mathcal{F}_3 \to 0$과
모든 $\mathcal{O}_X$-가군 $\mathcal{G}$에 대해 코호몰로지 긴 완전열
$$
\xymatrix{
\mathcal{F}_1 \otimes_{\mathcal{O}_X} \mathcal{G} \ar[r] &
\mathcal{F}_2 \otimes_{\mathcal{O}_X} \mathcal{G} \ar[r] &
\mathcal{F}_3 \otimes_{\mathcal{O}_X} \mathcal{G} \ar[r] & 0 \\
\text{Tor}_1^{\mathcal{O}_X}(\mathcal{F}_1, \mathcal{G}) \ar[r] &
\text{Tor}_1^{\mathcal{O}_X}(\mathcal{F}_2, \mathcal{G}) \ar[r] &
\text{Tor}_1^{\mathcal{O}_X}(\mathcal{F}_3, \mathcal{G}) \ar[ull]
}
$$
을 얻는다. 이를 이 상황에 딸린 $\text{Tor}$의 긴 완전열이라 한다.

\begin{lemma}
\label{cohomology-lemma-flat-tor-zero}
\begin{slogan}
Tor는 평탄성에서 벗어나는 정도를 측정한다.
\end{slogan}
$(X, \mathcal{O}_X)$를 환 달린 공간이라 하자.
$\mathcal{F}$를 $\mathcal{O}_X$-가군이라 하자. 다음은 서로 동치이다.
\begin{enumerate}
\item $\mathcal{F}$는 평탄 $\mathcal{O}_X$-가군이다.
\item 모든 $\mathcal{O}_X$-가군 $\mathcal{G}$에 대해
$\text{Tor}_1^{\mathcal{O}_X}(\mathcal{F}, \mathcal{G}) = 0$이다.
\end{enumerate}
\end{lemma}

\begin{proof}
$\mathcal{F}$가 평탄이면 $\mathcal{F} \otimes_{\mathcal{O}_X} -$는 완전
함자이고 위성함자들은 영이다. 역으로 (2)가 성립한다고 가정하자.
$\mathcal{G} \to \mathcal{H}$가 단사이고 그 여핵이
$\mathcal{Q}$이면, $\text{Tor}$의 긴 완전열에 의해
$\mathcal{F} \otimes_{\mathcal{O}_X} \mathcal{G} \to
\mathcal{F} \otimes_{\mathcal{O}_X} \mathcal{H}$의 핵은
$\text{Tor}_1^{\mathcal{O}_X}(\mathcal{F}, \mathcal{Q})$의 몫이다.
이는 가정에 의해 영이다. 따라서 $\mathcal{F}$는 평탄이다.
\end{proof}

\begin{lemma}
\label{cohomology-lemma-factor-through-K-flat}
$(X, \mathcal{O}_X)$를 환 달린 공간이라 하자.
$a : \mathcal{K}^\bullet \to \mathcal{L}^\bullet$를
$\mathcal{O}_X$-가군 복합체들의 사상이라 하자.
$\mathcal{K}^\bullet$가 K-평탄이면 복합체 $\mathcal{N}^\bullet$와
복합체 사상 $b : \mathcal{K}^\bullet \to \mathcal{N}^\bullet$,
$c : \mathcal{N}^\bullet \to \mathcal{L}^\bullet$가 존재하여 다음을
만족한다.
\begin{enumerate}
\item $\mathcal{N}^\bullet$는 K-평탄이다.
\item $c$는 유사동형이다.
\item $a$는 $c \circ b$와 호모토픽이다.
\end{enumerate}
$\mathcal{K}^\bullet$의 항들이 평탄하면 $\mathcal{N}^\bullet$에도
같은 성질이 성립하도록 $\mathcal{N}^\bullet$, $b$, $c$를 고를 수 있다.
\end{lemma}

\begin{proof}
호모토피 범주 $K(\textit{Mod}(\mathcal{O}_X))$가 삼각범주라는 사실을
사용할 것이다. Derived Categories, Proposition
\ref{derived-proposition-homotopy-category-triangulated}을 보라.
특수 삼각형
$\mathcal{K}^\bullet \to \mathcal{L}^\bullet \to
\mathcal{C}^\bullet \to \mathcal{K}^\bullet[1]$을 고르자. Lemma
\ref{cohomology-lemma-K-flat-resolution}에 따라, 항들이 평탄이고
$\mathcal{M}^\bullet$가 K-평탄인 유사동형
$\mathcal{M}^\bullet \to \mathcal{C}^\bullet$를 고르자.
삼각범주의 공리들에 의해 합성
$\mathcal{M}^\bullet \to \mathcal{C}^\bullet \to \mathcal{K}^\bullet[1]$
을 특수 삼각형
$\mathcal{K}^\bullet \to \mathcal{N}^\bullet \to
\mathcal{M}^\bullet \to \mathcal{K}^\bullet[1]$에 넣을 수 있다.
Lemma \ref{cohomology-lemma-K-flat-two-out-of-three}에 의해
$\mathcal{N}^\bullet$가 K-평탄임을 알 수 있다. 다시 삼각범주의 공리들을
사용하여 다음 특수 삼각형의 사상에 들어가는 사상
$\mathcal{N}^\bullet \to \mathcal{L}^\bullet$를 고를 수 있다.
$$
\xymatrix{
\mathcal{K}^\bullet \ar[r] \ar[d] &
\mathcal{N}^\bullet \ar[r] \ar[d] &
\mathcal{M}^\bullet \ar[r] \ar[d] &
\mathcal{K}^\bullet[1] \ar[d] \\
\mathcal{K}^\bullet \ar[r] &
\mathcal{L}^\bullet \ar[r] &
\mathcal{C}^\bullet \ar[r] &
\mathcal{K}^\bullet[1]
}
$$
화살표 셋 가운데 둘이 유사동형이므로 세 번째 화살표
$\mathcal{N}^\bullet \to \mathcal{L}^\bullet$도 이 특수 삼각형들에
딸린 코호몰로지 긴 완전열들에 의해 유사동형이다(원한다면 이 도식의
$D(\mathcal{O}_X)$ 안의 상을 보고 Derived Categories, Lemma
\ref{derived-lemma-third-isomorphism-triangle}를 사용해도 된다).
이로써 (1), (2), (3)의 증명이 끝난다. 마지막 주장을 증명하기 위해
Derived Categories, Lemma
\ref{derived-lemma-improve-distinguished-triangle-homotopy}에 따라
$\mathcal{N}^n \cong \mathcal{M}^n \oplus \mathcal{K}^n$이 되도록
$\mathcal{N}^\bullet$를 고를 수 있다. 따라서
$\mathcal{K}^\bullet$의 항들이 평탄하면 원하는 평탄성을 얻는다.
\end{proof}







\section{유도 당김}
\label{cohomology-section-derived-pullback}

\noindent
$f : (X, \mathcal{O}_X) \to (Y, \mathcal{O}_Y)$를 환 달린 공간의
사상이라 하자. K-평탄 분해를 사용하여 유도 당김 함자
$$
Lf^* : D(\mathcal{O}_Y) \to D(\mathcal{O}_X)
$$
를 정의할 수 있다. 구체적으로, 모든 $\mathcal{O}_Y$-가군 복합체
$\mathcal{G}^\bullet$에 대해 K-평탄 분해
$\mathcal{K}^\bullet \to \mathcal{G}^\bullet$를 고르고
$Lf^*\mathcal{G}^\bullet = f^*\mathcal{K}^\bullet$로 둔다. 이것이
잘 정의됨은 Lemmas \ref{cohomology-lemma-pullback-K-flat},
\ref{cohomology-lemma-K-flat-resolution},
\ref{cohomology-lemma-derived-tor-quasi-isomorphism-other-side}를 사용해 알 수 있다.
그러나 모든 세부를 빠짐없이 정당화하려면 일반 이론을 사용하는 편이
더 편리할 수 있다.

\begin{lemma}
\label{cohomology-lemma-derived-base-change}
위 구성은 선택에 의존하지 않으며 삼각범주의 완전 함자
$Lf^* : D(\mathcal{O}_Y) \to D(\mathcal{O}_X)$를 정의한다.
\end{lemma}

\begin{proof}
Derived Categories, Section \ref{derived-section-derived-functors}에서
전개한 일반 이론을 사용한다.
$\mathcal{D} = K(\mathcal{O}_Y)$, $\mathcal{D}' = D(\mathcal{O}_X)$로
두자. 규칙 $F(\mathcal{G}^\bullet) = f^*\mathcal{G}^\bullet$로 정의되는
삼각범주의 완전 함자를
$F : \mathcal{D} \to \mathcal{D}'$로 쓰자. 또한
$\mathcal{D} = K(\mathcal{O}_Y)$ 안의 유사동형들의 집합을 $S$라 하자.
그러면 Derived Categories, Situation
\ref{derived-situation-derived-functor}의 상황이므로 Derived Categories,
Definition \ref{derived-definition-right-derived-functor-defined}을 적용할
수 있다. $LF$가 모든 곳에서 정의됨을 주장한다. 이는
Derived Categories, Lemma \ref{derived-lemma-find-existence-computes}에서
$\mathcal{P} \subset \Ob(\mathcal{D})$를 $K$-평탄 복합체들의 모임으로
두면 따른다. 실제로 (1)은 Lemma \ref{cohomology-lemma-K-flat-resolution}에서
따른다. (2)를 보이려면 K-평탄 $\mathcal{O}_Y$-가군 복합체 사이의
유사동형
$\mathcal{K}_1^\bullet  \to \mathcal{K}_2^\bullet$에 대해 사상
$f^*\mathcal{K}_1^\bullet  \to f^*\mathcal{K}_2^\bullet$가
유사동형임을 보이면 된다. 이 사상을
$$
f^{-1}\mathcal{K}_1^\bullet \otimes_{f^{-1}\mathcal{O}_Y} \mathcal{O}_X
\longrightarrow
f^{-1}\mathcal{K}_2^\bullet \otimes_{f^{-1}\mathcal{O}_Y} \mathcal{O}_X
$$
로 쓰자. 함자 $f^{-1}$은 완전하므로 사상
$f^{-1}\mathcal{K}_1^\bullet  \to f^{-1}\mathcal{K}_2^\bullet$는
유사동형이다. Lemma \ref{cohomology-lemma-pullback-K-flat}을 사상
$(X, f^{-1}\mathcal{O}_Y) \to (Y, \mathcal{O}_Y)$에 적용하면 복합체
$f^{-1}\mathcal{K}_1^\bullet$와 $f^{-1}\mathcal{K}_2^\bullet$는
K-평탄 $f^{-1}\mathcal{O}_Y$-가군 복합체이다. 따라서 Lemma
\ref{cohomology-lemma-derived-tor-quasi-isomorphism-other-side}에 의해 표시된
사상은 유사동형이다. 이로써 유도 함자
$$
LF :
D(\mathcal{O}_Y) = S^{-1}\mathcal{D}
\longrightarrow
\mathcal{D}' = D(\mathcal{O}_X)
$$
를 얻는다. Derived Categories, Equation
(\ref{derived-equation-everywhere})을 보라. 마지막으로 Derived Categories,
Lemma \ref{derived-lemma-find-existence-computes}에 의해
$\mathcal{K}^\bullet$가 K-평탄일 때
$LF(\mathcal{K}^\bullet) = F(\mathcal{K}^\bullet) = f^*\mathcal{K}^\bullet$
이고, 따라서 $Lf^* = LF$는 실제로 위에서 설명한 방식으로 계산된다.
\end{proof}

\begin{lemma}
\label{cohomology-lemma-derived-pullback-composition}
$f : X \to Y$와 $g : Y \to Z$를 환 달린 공간의 사상이라 하자. 그러면
$D(\mathcal{O}_Z) \to D(\mathcal{O}_X)$인 함자로서
$Lf^* \circ Lg^* = L(g \circ f)^*$이다.
\end{lemma}

\begin{proof}
$E$를 $D(\mathcal{O}_Z)$의 대상이라 하자. 구성에 의해 $Lg^*E$는
$Z$ 위에서 $E$를 나타내는 K-평탄 복합체 $\mathcal{K}^\bullet$를
고르고 $Lg^*E = g^*\mathcal{K}^\bullet$로 두어 계산한다. Lemma
\ref{cohomology-lemma-pullback-K-flat}에 의해 $g^*\mathcal{K}^\bullet$는 $Y$
위에서 K-평탄이다. 따라서 $Lf^*Lg^*E$는
$f^*g^*\mathcal{K}^\bullet = (g \circ f)^*\mathcal{K}^\bullet$로
주어지고, 이는 $L(g \circ f)^*E$도 나타낸다.
\end{proof}

\begin{lemma}
\label{cohomology-lemma-pullback-tensor-product}
$f : (X, \mathcal{O}_X) \to (Y, \mathcal{O}_Y)$를 환 달린 공간의
사상이라 하자. 모든
$\mathcal{F}^\bullet, \mathcal{G}^\bullet \in \Ob(D(\mathcal{O}_Y))$에
대해 표준 이함자적 동형
$$
Lf^*(
\mathcal{F}^\bullet \otimes_{\mathcal{O}_Y}^{\mathbf{L}} \mathcal{G}^\bullet
) =
Lf^*\mathcal{F}^\bullet 
\otimes_{\mathcal{O}_X}^{\mathbf{L}}
Lf^*\mathcal{G}^\bullet 
$$
이 존재한다.
\end{lemma}

\begin{proof}
$\mathcal{F}^\bullet$와 $\mathcal{G}^\bullet$가 K-평탄 복합체라고
가정해도 된다. 이 경우
$\mathcal{F}^\bullet \otimes_{\mathcal{O}_Y}^{\mathbf{L}} \mathcal{G}^\bullet$
는 이중 복합체
$\mathcal{F}^\bullet \otimes_{\mathcal{O}_Y} \mathcal{G}^\bullet$에
딸린 전체 복합체일 뿐이다. Lemma
\ref{cohomology-lemma-tensor-product-K-flat}에 의해
$\text{Tot}(\mathcal{F}^\bullet \otimes_{\mathcal{O}_Y} \mathcal{G}^\bullet)$
도 K-평탄이다. 따라서 보조정리의 동형은 동형
$$
\text{Tot}(f^*\mathcal{F}^\bullet \otimes_{\mathcal{O}_X}
f^*\mathcal{G}^\bullet)
\longrightarrow
f^*\text{Tot}(\mathcal{F}^\bullet \otimes_{\mathcal{O}_Y} \mathcal{G}^\bullet)
$$
에서 오며, 그 각 성분은 Modules, Lemma
\ref{modules-lemma-tensor-product-pullback}의 동형
$f^*\mathcal{F}^p \otimes_{\mathcal{O}_X} f^*\mathcal{G}^q \to
f^*(\mathcal{F}^p \otimes_{\mathcal{O}_Y} \mathcal{G}^q)$이다.
\end{proof}

\begin{lemma}
\label{cohomology-lemma-variant-derived-pullback}
$f : (X, \mathcal{O}_X) \to (Y, \mathcal{O}_Y)$를 환 달린 공간의
사상이라 하자. $D(\mathcal{O}_X)$의 $\mathcal{F}^\bullet$와
$D(\mathcal{O}_Y)$의 $\mathcal{G}^\bullet$에 대해 표준 이함자적 동형
$$
\mathcal{F}^\bullet
\otimes_{\mathcal{O}_X}^{\mathbf{L}}
Lf^*\mathcal{G}^\bullet
=
\mathcal{F}^\bullet 
\otimes_{f^{-1}\mathcal{O}_Y}^{\mathbf{L}}
f^{-1}\mathcal{G}^\bullet 
$$
이 존재한다.
\end{lemma}

\begin{proof}
$\mathcal{F}$를 $\mathcal{O}_X$-가군, $\mathcal{G}$를
$\mathcal{O}_Y$-가군이라 하자. 그러면
$\mathcal{F} \otimes_{\mathcal{O}_X} f^*\mathcal{G} =
\mathcal{F} \otimes_{f^{-1}\mathcal{O}_Y} f^{-1}\mathcal{G}$이다. 이는
$f^*\mathcal{G} =
\mathcal{O}_X \otimes_{f^{-1}\mathcal{O}_Y} f^{-1}\mathcal{G}$이기
때문이다. 이 사실과 정의들에서 보조정리가 따른다.
\end{proof}

\begin{lemma}
\label{cohomology-lemma-tensor-pull-compatibility}
$f : (X, \mathcal{O}_X) \to (Y, \mathcal{O}_Y)$를 환 달린 공간의
사상이라 하자. $\mathcal{K}^\bullet$와 $\mathcal{M}^\bullet$를
$\mathcal{O}_Y$-가군 복합체라 하자. 도식
$$
\xymatrix{
Lf^*(\mathcal{K}^\bullet
\otimes_{\mathcal{O}_Y}^\mathbf{L}
\mathcal{M}^\bullet) \ar[r] \ar[d] &
Lf^*\text{Tot}(\mathcal{K}^\bullet
\otimes_{\mathcal{O}_Y}
\mathcal{M}^\bullet) \ar[d] \\
Lf^*\mathcal{K}^\bullet \otimes_{\mathcal{O}_X}^\mathbf{L}
Lf^*\mathcal{M}^\bullet \ar[d] &
f^*\text{Tot}(\mathcal{K}^\bullet
\otimes_{\mathcal{O}_Y}
\mathcal{M}^\bullet) \ar[d] \\
f^*\mathcal{K}^\bullet \otimes_{\mathcal{O}_X}^\mathbf{L}
f^*\mathcal{M}^\bullet \ar[r] &
\text{Tot}(f^*\mathcal{K}^\bullet \otimes_{\mathcal{O}_X}
f^*\mathcal{M}^\bullet)
}
$$
은 가환한다.
\end{lemma}

\begin{proof}
Lemma \ref{cohomology-lemma-pullback-K-flat}에서처럼 K-평탄 분해의 존재성을
사용한다. 그러한 분해
$\mathcal{P}^\bullet \to \mathcal{K}^\bullet$와
$\mathcal{Q}^\bullet \to \mathcal{M}^\bullet$를 고르면 도식
$$
\xymatrix{
Lf^*\text{Tot}(\mathcal{P}^\bullet
\otimes_{\mathcal{O}_Y}
\mathcal{Q}^\bullet) \ar[r] \ar[d] &
Lf^*\text{Tot}(\mathcal{K}^\bullet
\otimes_{\mathcal{O}_Y}
\mathcal{M}^\bullet) \ar[d] \\
f^*\text{Tot}(\mathcal{P}^\bullet
\otimes_{\mathcal{O}_Y}
\mathcal{Q}^\bullet) \ar[d] \ar[r] &
f^*\text{Tot}(\mathcal{K}^\bullet
\otimes_{\mathcal{O}_Y}
\mathcal{M}^\bullet) \ar[d] \\
\text{Tot}(f^*\mathcal{P}^\bullet \otimes_{\mathcal{O}_X}
f^*\mathcal{Q}^\bullet) \ar[r] &
\text{Tot}(f^*\mathcal{K}^\bullet \otimes_{\mathcal{O}_X}
f^*\mathcal{M}^\bullet)
}
$$
이 가환함을 알 수 있다. 그러나 이제
$\mathcal{P}^\bullet$와 $\mathcal{Q}^\bullet$의 선택과 Lemma
\ref{cohomology-lemma-tensor-product-K-flat}에 의해 이 도식의 왼쪽 변은
이 도식의 왼쪽 변이다.
\end{proof}







\section{비유계 복합체의 코호몰로지}
\label{cohomology-section-unbounded}

\noindent
$(X, \mathcal{O}_X)$를 환 달린 공간이라 하자. 범주
$\textit{Mod}(\mathcal{O}_X)$는 Grothendieck 아벨 범주이다. 즉 모든
여극한을 가지며, 여과 여극한은 완전하고, 생성자
$$
\bigoplus\nolimits_{U \subset X\text{ open}} j_{U!}\mathcal{O}_U,
$$
를 갖는다. Modules, Section \ref{modules-section-kernels}과 Lemmas
\ref{modules-lemma-j-shriek-flat},
\ref{modules-lemma-module-quotient-flat}을 보라. Injectives, Theorem
\ref{injectives-theorem-K-injective-embedding-grothendieck}에 의해
모든 $\mathcal{O}_X$-가군 복합체 $\mathcal{F}^\bullet$에는 K-단사
$\mathcal{O}_X$-가군 복합체로 가는 단사 유사동형
$\mathcal{F}^\bullet \to \mathcal{I}^\bullet$가 존재한다.
그 모든 항은 단사
$\mathcal{O}_X$-가군이며, 이 매장은 복합체 $\mathcal{F}^\bullet$에
대해 함자적으로 고를 수 있다. Derived Categories, Lemma
\ref{derived-lemma-enough-K-injectives-implies}에 의해 다음이 따른다.
\begin{enumerate}
\item 삼각범주 $\mathcal{D}$로 가는 임의의 완전 함자
$F : K(\textit{Mod}(\mathcal{O}_X)) \to \mathcal{D}$는 오른쪽 유도 함자
$RF : D(\mathcal{O}_X) \to \mathcal{D}$를 갖는다.
\item 아벨 범주 $\mathcal{A}$로 가는 임의의 가법 함자
$F : \textit{Mod}(\mathcal{O}_X) \to \mathcal{A}$에 대해, $F$가 유도하는
완전 함자 $F : K(\textit{Mod}(\mathcal{O}_X)) \to K(\mathcal{A})$를
생각하면 오른쪽 유도 함자
$RF : D(\mathcal{O}_X) \to D(\mathcal{A})$를 얻는다.
\end{enumerate}
구성에 의해 $RF(\mathcal{F}^\bullet) = F(\mathcal{I}^\bullet)$이다.
여기서 $\mathcal{F}^\bullet \to \mathcal{I}^\bullet$는 위와 같다.

\medskip\noindent
위 사실의 몇 가지 예를 들면 다음과 같다.
\begin{enumerate}
\item 함자 $\Gamma(X, -) : \textit{Mod}(\mathcal{O}_X) \to
\text{Mod}_{\Gamma(X, \mathcal{O}_X)}$는
$$
R\Gamma(X, -) : D(\mathcal{O}_X) \to D(\Gamma(X, \mathcal{O}_X))
$$
를 준다. 코호몰로지에는 표기 $H^i(X, K) = H^i(R\Gamma(X, K))$를
사용할 것이다.
\item 열린집합 $U \subset X$에 대해 함자
$\Gamma(U, -) : \textit{Mod}(\mathcal{O}_X) \to
\text{Mod}_{\Gamma(U, \mathcal{O}_X)}$를 생각하자. 이는
$$
R\Gamma(U, -) : D(\mathcal{O}_X) \to D(\Gamma(U, \mathcal{O}_X))
$$
를 준다. 코호몰로지에는 표기 $H^i(U, K) = H^i(R\Gamma(U, K))$를
사용할 것이다.
\item 환 달린 공간의 사상
$f : (X, \mathcal{O}_X) \to (Y, \mathcal{O}_Y)$에 대해 함자
$f_* : \textit{Mod}(\mathcal{O}_X) \to \textit{Mod}(\mathcal{O}_Y)$를
생각하자. 이는 비유계 유도범주 위의 전체 직접상
$$
Rf_* : D(\mathcal{O}_X) \longrightarrow D(\mathcal{O}_Y)
$$
을 준다.
\end{enumerate}

\begin{lemma}
\label{cohomology-lemma-adjoint}
$f : (X, \mathcal{O}_X) \to (Y, \mathcal{O}_Y)$를 환 달린 공간의
사상이라 하자. 위에서 정의한 함자 $Rf_*$와 Lemma
\ref{cohomology-lemma-derived-base-change}에서 정의한 함자 $Lf^*$는 서로 수반이다.
즉
$$
\Hom_{D(\mathcal{O}_X)}(Lf^*\mathcal{G}^\bullet, \mathcal{F}^\bullet)
=
\Hom_{D(\mathcal{O}_Y)}(\mathcal{G}^\bullet, Rf_*\mathcal{F}^\bullet)
$$
이며, 이는
$\mathcal{F}^\bullet \in \Ob(D(\mathcal{O}_X))$와
$\mathcal{G}^\bullet \in \Ob(D(\mathcal{O}_Y))$에 대해 이함자적이다.
\end{lemma}

\begin{proof}
이는 $Rf_*$와 $Lf^*$가 존재한다는 사실에서 형식적으로 따른다.
Derived Categories, Lemma
\ref{derived-lemma-derived-adjoint-functors}를 보라.
\end{proof}

\begin{lemma}
\label{cohomology-lemma-derived-pushforward-composition}
$f : X \to Y$와 $g : Y \to Z$를 환 달린 공간의 사상이라 하자. 그러면
$D(\mathcal{O}_X) \to D(\mathcal{O}_Z)$인 함자로서
$Rg_* \circ Rf_* = R(g \circ f)_*$이다.
\end{lemma}

\begin{proof}
Lemma \ref{cohomology-lemma-adjoint}에 의해 $Rg_* \circ Rf_*$는
$Lf^* \circ Lg^*$의 수반이다. Lemma
\ref{cohomology-lemma-derived-pullback-composition}에 의해
$Lf^* \circ Lg^* = L(g \circ f)^*$이고, 따라서 수반함자의 유일성에
의해 $Rg_* \circ Rf_* = R(g \circ f)_*$이다.
\end{proof}

\begin{remark}
\label{cohomology-remark-base-change}
비유계 유도 함자 $Lf^*$와 $Rf_*$의 구성으로 밑변환 사상을 가장
일반적인 상황에서 구성할 수 있다. 구체적으로 환 달린 공간의 가환도식
$$
\xymatrix{
X' \ar[r]_{g'} \ar[d]_{f'} &
X \ar[d]^f \\
S' \ar[r]^g &
S
}
$$
이 주어졌다고 하자. $K$를 $D(\mathcal{O}_X)$의 대상이라 하자. 그러면
$D(\mathcal{O}_{S'})$ 안에 표준 밑변환 사상
$$
Lg^*Rf_*K \longrightarrow R(f')_*L(g')^*K
$$
이 존재한다. 이 사상은 사상
$L(f')^*Lg^*Rf_*K \to L(g')^*K$와 수반인 사상이다.
$L(f')^*Lg^* = L(g')^*Lf^*$이므로 이는 사상
$L(g')^*Lf^*Rf_*K \to L(g')^*K$와 같다. 이 사상으로는 수반 사상
$Lf^*Rf_*K \to K$에 $L(g')^*$를 적용한 것을 취할 수 있다.
\end{remark}

\begin{remark}
\label{cohomology-remark-compose-base-change}
환 달린 공간의 가환도식
$$
\xymatrix{
X' \ar[r]_k \ar[d]_{f'} & X \ar[d]^f \\
Y' \ar[r]^l \ar[d]_{g'} & Y \ar[d]^g \\
Z' \ar[r]^m & Z
}
$$
을 생각하자. 두 정사각형에 대한 Remark \ref{cohomology-remark-base-change}의
밑변환 사상들을 합성하면 바깥 직사각형에 대한 밑변환 사상을 얻는다.
더 정확히 말해, 합성
\begin{align*}
Lm^* \circ R(g \circ f)_*
& =
Lm^* \circ Rg_* \circ Rf_* \\
& \to Rg'_* \circ Ll^* \circ Rf_* \\
& \to Rg'_* \circ Rf'_* \circ Lk^* \\
& = R(g' \circ f')_* \circ Lk^*
\end{align*}
은 그 직사각형에 대한 밑변환 사상이다. 확인은 생략한다.
\end{remark}

\begin{remark}
\label{cohomology-remark-compose-base-change-horizontal}
환 달린 공간의 가환도식
$$
\xymatrix{
X'' \ar[r]_{g'} \ar[d]_{f''} & X' \ar[r]_g \ar[d]_{f'} & X \ar[d]^f \\
Y'' \ar[r]^{h'} & Y' \ar[r]^h & Y
}
$$
을 생각하자. 두 정사각형에 대한 Remark \ref{cohomology-remark-base-change}의
밑변환 사상들을 합성하면 바깥 직사각형에 대한 밑변환 사상을 얻는다.
더 정확히 말해, 합성
\begin{align*}
L(h \circ h')^* \circ Rf_*
& =
L(h')^* \circ Lh_* \circ Rf_* \\
& \to L(h')^* \circ Rf'_* \circ Lg^* \\
& \to Rf''_* \circ L(g')^* \circ Lg^* \\
& = Rf''_* \circ L(g \circ g')^*
\end{align*}
은 그 직사각형에 대한 밑변환 사상이다. 확인은 생략한다.
\end{remark}

\begin{lemma}
\label{cohomology-lemma-adjoints-push-pull-compatibility}
$f : (X, \mathcal{O}_X) \to (Y, \mathcal{O}_Y)$를 환 달린 공간의
사상이라 하고, $\mathcal{K}^\bullet$를
$\mathcal{O}_X$-가군 복합체라 하자. 복합체 위의
$Lf^* \to f^*$, 복합체 위의 $f_* \to Rf_*$, 그리고 수반
$Lf^* \circ Rf_* \to \text{id}$에서 오는 도식
$$
\xymatrix{
Lf^*f_*\mathcal{K}^\bullet \ar[r] \ar[d] &
f^*f_*\mathcal{K}^\bullet \ar[d] \\
Lf^*Rf_*\mathcal{K}^\bullet \ar[r] &
\mathcal{K}^\bullet
}
$$
은 $D(\mathcal{O}_X)$ 안에서 가환한다.
\end{lemma}

\begin{proof}
K-평탄 분해와 K-단사 분해의 존재성을 사용할 것이다. Lemma
\ref{cohomology-lemma-pullback-K-flat}과 위의 논의를 보라. $\mathcal{I}^\bullet$가
$\mathcal{O}_X$-가군 복합체로서 K-단사가 되는 유사동형
$\mathcal{K}^\bullet \to \mathcal{I}^\bullet$를 고르자.
$\mathcal{Q}^\bullet$가 $\mathcal{O}_Y$-가군 복합체로서 K-평탄이
되는 유사동형
$\mathcal{Q}^\bullet \to f_*\mathcal{I}^\bullet$를 고르자. K-평탄
$\mathcal{O}_Y$-가군 복합체 $\mathcal{P}^\bullet$와, 호모토피까지
가환하며 위쪽 수평 화살표가 유사동형인 복합체 사상의 도식
$$
\xymatrix{
\mathcal{P}^\bullet \ar[r] \ar[d] &
f_*\mathcal{K}^\bullet \ar[d] \\
\mathcal{Q}^\bullet \ar[r] & f_*\mathcal{I}^\bullet
}
$$
을 고를 수 있다. 실제로 복합체의 호모토피 범주에서 유사동형들이
곱셈계를 이루므로 먼저 어떤 복합체 $\mathcal{P}^\bullet$에 대해
그러한 도식을 고른 뒤 $\mathcal{P}^\bullet$를 K-평탄 복합체로
바꿀 수 있다. 당김을 취하면 호모토피까지 가환하는 복합체 사상의 도식
$$
\xymatrix{
f^*\mathcal{P}^\bullet \ar[r] \ar[d] &
f^*f_*\mathcal{K}^\bullet \ar[d] \ar[r] &
\mathcal{K}^\bullet \ar[d] \\
f^*\mathcal{Q}^\bullet \ar[r] &
f^*f_*\mathcal{I}^\bullet \ar[r] &
\mathcal{I}^\bullet
}
$$
을 얻는다. 바깥 직사각형이 보조정리의 진술이 참임을 보여 준다.
\end{proof}

\begin{remark}
\label{cohomology-remark-cup-product}
$f : (X, \mathcal{O}_X) \to (Y, \mathcal{O}_Y)$를 환 달린 공간의
사상이라 하자. $Lf^*$와 $Rf_*$의 수반성으로
$D(\mathcal{O}_X)$의 모든 $K, L$에 대해 $D(\mathcal{O}_Y)$ 안에서
상대 컵곱
$$
Rf_*K \otimes_{\mathcal{O}_Y}^\mathbf{L} Rf_*L
\longrightarrow
Rf_*(K \otimes_{\mathcal{O}_X}^\mathbf{L} L)
$$
을 구성할 수 있다. 이 사상은 다음 사상과 수반인 사상이다.
$Lf^*(Rf_*K \otimes_{\mathcal{O}_Y}^\mathbf{L} Rf_*L) \to
K \otimes_{\mathcal{O}_X}^\mathbf{L} L$. 이 사상으로는 동형
$Lf^*(Rf_*K \otimes_{\mathcal{O}_Y}^\mathbf{L} Rf_*L) =
Lf^*Rf_*K \otimes_{\mathcal{O}_X}^\mathbf{L} Lf^*Rf_*L$
(Lemma \ref{cohomology-lemma-pullback-tensor-product})과 여단위
$Lf^* \circ Rf_* \to \text{id}$에서 오는 사상
$Lf^*Rf_*K \otimes_{\mathcal{O}_X}^\mathbf{L} Lf^*Rf_*L
\to K \otimes_{\mathcal{O}_X}^\mathbf{L} L$의 합성을 취할 수 있다.
\end{remark}





\section{여과 복합체의 코호몰로지}
\label{cohomology-section-cohomology-filtered-object}

\noindent
대수다양체의 코호몰로지를 연구할 때 층의 여과 복합체는 흔히 자연스럽게
나타난다. 예를 들어 드람 복합체에는 호지 여과가 딸려 있다. 이 절에서는
매우 일반적인 Injectives, Lemma
\ref{injectives-lemma-K-injective-embedding-filtration}를 사용하여
코호몰로지 위의 스펙트럼열을 구성하고, 이를 앞서 구성한 스펙트럼열들과
연결한다.

\begin{lemma}
\label{cohomology-lemma-spectral-sequence-filtered-object}
$(X, \mathcal{O}_X)$를 환 달린 공간이라 하고, $\mathcal{F}^\bullet$를
$\mathcal{O}_X$-가군들의 여과 복합체라 하자. 이중등급
$\Gamma(X, \mathcal{O}_X)$-가군들의 표준 스펙트럼열
$(E_r, \text{d}_r)_{r \geq 1}$이 존재한다. 여기서 $d_r$의 이중차수는
$(r, -r + 1)$이고
$$
E_1^{p, q} = H^{p + q}(X, \text{gr}^p\mathcal{F}^\bullet)
$$
이다. 모든 $n$에 대해
$$
H^n(X, F^p\mathcal{F}^\bullet) = 0\text{ for }p \gg 0
\quad\text{and}\quad
H^n(X, F^p\mathcal{F}^\bullet) = H^n(X, \mathcal{F}^\bullet)\text{ for }p \ll 0
$$
이면 이 스펙트럼열은 유계이고 $H^*(X, \mathcal{F}^\bullet)$로
수렴한다.
\end{lemma}

\begin{proof}
(복합체가 유한 여과를 갖는 가군들로 이루어진 아래로 유계인 복합체인
경우의 증명은 아래의 주석을 보라.) Injectives, Lemma
\ref{injectives-lemma-K-injective-embedding-filtration}에서와 같은
여과 복합체의 사상
$j : \mathcal{F}^\bullet \to \mathcal{J}^\bullet$을 고르자. 이
스펙트럼열은 여과 복합체
$$
\Gamma(X, \mathcal{J}^\bullet)
\quad\text{with}\quad
F^p\Gamma(X, \mathcal{J}^\bullet) = \Gamma(X, F^p\mathcal{J}^\bullet)
$$
에 딸린 Homology, Section \ref{homology-section-filtered-complex}의
스펙트럼열이다. K-단사 대표에 값을 취하여 코호몰로지를 계산하므로
$E_1$ 쪽은 보조정리에 적힌 것과 같다. 주어진 조건 아래에서의
수렴성과 유계성은 Homology, Lemma
\ref{homology-lemma-ss-converges-trivial}에서 따른다.
\end{proof}

\begin{remark}
\label{cohomology-remark-spectral-sequence-filtered-object}
$(X, \mathcal{O}_X)$를 환 달린 공간이라 하고, $\mathcal{F}^\bullet$를
$\mathcal{O}_X$-가군들의 여과 복합체라 하자.
$\mathcal{F}^\bullet$가 아래로 유계이고 각 $n$에 대해
$\mathcal{F}^n$ 위의 여과가 유한하면 Lemma
\ref{cohomology-lemma-spectral-sequence-filtered-object}의 스펙트럼열을
Injectives, Lemma
\ref{injectives-lemma-K-injective-embedding-filtration}를 사용하지
않고도 구성할 수 있다. 실제로 Derived Categories, Lemma
\ref{derived-lemma-right-resolution-by-filtered-injectives}에 의해
여과 복합체의 여과 유사동형
$i : \mathcal{F}^\bullet \to \mathcal{I}^\bullet$이 존재한다. 여기서
$\mathcal{I}^\bullet$는 아래로 유계이고, 모든 $n$에 대해
$\mathcal{I}^n$ 위의 여과는 유한하며, 각
$\text{gr}^p\mathcal{I}^n$은 단사 $\mathcal{O}_X$-가군이다. 이제
$$
\Gamma(X, \mathcal{I}^\bullet)
\quad\text{with}\quad
F^p\Gamma(X, \mathcal{I}^\bullet) = \Gamma(X, F^p\mathcal{I}^\bullet)
$$
에 딸린 스펙트럼열을 취한다. 단사 대상들로 이루어진 아래로 유계인
복합체에 값을 취하여 코호몰로지를 계산할 수 있으므로 $E_1$ 쪽은
보조정리에 적힌 것과 같다. 주어진 조건 아래에서의 수렴성과 유계성은
Homology, Lemma \ref{homology-lemma-biregular-ss-converges}에서 따른다.
사실 이는 Derived Categories, Lemma
\ref{derived-lemma-ss-filtered-derived}의 스펙트럼열의 특수한 경우이다.
\end{remark}

\begin{example}
\label{cohomology-example-spectral-sequence}
$(X, \mathcal{O}_X)$를 환 달린 공간이라 하고, $\mathcal{F}^\bullet$를
$\mathcal{O}_X$-가군 복합체라 하자.
$F^p\mathcal{F}^\bullet = \tau_{\leq -p}\mathcal{F}^\bullet$로 두고
Lemma \ref{cohomology-lemma-spectral-sequence-filtered-object}를 적용할 수 있다.
($\mathcal{F}^\bullet$가 아래로 유계이면 Remark
\ref{cohomology-remark-spectral-sequence-filtered-object}를 사용할 수 있다.)
그러면 스펙트럼열
$$
E_1^{p, q} = H^{p + q}(X, H^{-p}(\mathcal{F}^\bullet)[p]) =
H^{2p + q}(X, H^{-p}(\mathcal{F}^\bullet))
$$
을 얻는다. $p = -j$, $q = i + 2j$로 번호를 다시 매기면, 모든
$K \in D(\mathcal{O}_X)$에 대해 $d'_r$의 이중차수가
$(r, -r + 1)$이고
$$
(E'_2)^{i, j} = H^i(X, H^j(K))
$$
인 이중등급 가군들의 스펙트럼열
$(E'_r, d'_r)_{r \geq 2}$을 얻는다. 예를 들어 $K$가 아래로
유계이면 이 스펙트럼열은 유계이고 $H^{i + j}(X, K)$로 수렴한다.
아래로 유계인 경우 이 스펙트럼열은 Cartan--Eilenberg 분해를 사용하여
구성한 Derived Categories, Lemma
\ref{derived-lemma-two-ss-complex-functor}의 두 번째 스펙트럼열의
한 예이다.
\end{example}

\begin{example}
\label{cohomology-example-spectral-sequence-bis}
$(X, \mathcal{O}_X)$를 환 달린 공간이라 하고, $\mathcal{F}^\bullet$를
$\mathcal{O}_X$-가군 복합체라 하자.
$F^p\mathcal{F}^\bullet = \sigma_{\geq p}\mathcal{F}^\bullet$로 두고
Lemma \ref{cohomology-lemma-spectral-sequence-filtered-object}를 적용할 수 있다.
그러면 스펙트럼열
$$
E_1^{p, q} = H^{p + q}(X, \mathcal{F}^p[-p]) = H^q(X, \mathcal{F}^p)
$$
을 얻는다. $\mathcal{F}^\bullet$가 아래로 유계이면 다음이 성립한다.
\begin{enumerate}
\item Remark \ref{cohomology-remark-spectral-sequence-filtered-object}를 사용하여
이 스펙트럼열을 구성할 수 있다.
\item 이 스펙트럼열은 유계이고
$H^{i + j}(X, \mathcal{F}^\bullet)$로 수렴한다.
\item 이 스펙트럼열은 Cartan--Eilenberg 분해를 사용하여 구성한
Derived Categories, Lemma \ref{derived-lemma-two-ss-complex-functor}의
첫 번째 스펙트럼열과 같다.
\end{enumerate}
\end{example}

\begin{lemma}
\label{cohomology-lemma-relative-spectral-sequence-filtered-object}
$f : (X, \mathcal{O}_X) \to (Y, \mathcal{O}_Y)$를 환 달린 공간의
사상이라 하고, $\mathcal{F}^\bullet$를
$\mathcal{O}_X$-가군들의 여과 복합체라 하자. 이중등급
$\mathcal{O}_Y$-가군들의 표준 스펙트럼열
$(E_r, \text{d}_r)_{r \geq 1}$이 존재한다. 여기서 $d_r$의 이중차수는
$(r, -r + 1)$이고
$$
E_1^{p, q} = R^{p + q}f_*\text{gr}^p\mathcal{F}^\bullet
$$
이다. 모든 $n$에 대해
$$
R^nf_*F^p\mathcal{F}^\bullet = 0 \text{ for }p \gg 0
\quad\text{and}\quad
R^nf_*F^p\mathcal{F}^\bullet = R^nf_*\mathcal{F}^\bullet \text{ for }p \ll 0
$$
이면 이 스펙트럼열은 유계이고 $Rf_*\mathcal{F}^\bullet$로 수렴한다.
\end{lemma}

\begin{proof}
증명은 Lemma \ref{cohomology-lemma-spectral-sequence-filtered-object}의 증명과
완전히 같다.
\end{proof}





\section{Godement 분해}
\label{cohomology-section-godement}

\noindent
참고문헌은 \cite{Godement}이다.

\medskip\noindent
$(X, \mathcal{O}_X)$를 환 달린 공간이라 하자. $X$와 같은 점들을 갖는
이산 위상공간을 $X_{disc}$로 나타내고, 자명한 연속사상을
$f : X_{disc} \to X$로 나타내자.
$\mathcal{O}_{X_{disc}} = f^{-1}\mathcal{O}_X$로 두자. 그러면
$f : (X_{disc}, \mathcal{O}_{X_{disc}}) \to (X, \mathcal{O}_X)$는
환 달린 공간의 평탄 사상이다. 가군층 위의 수반 함자쌍 $f^*, f_*$에
Simplicial, Section \ref{simplicial-section-standard}의 내용을
{\it 쌍대}하여 적용할 수 있다. 따라서
$\textit{Mod}(\mathcal{O}_X)$에서 자기 자신으로 가는 함자들의 범주에
확대 코심플렉셜 대상
$$
\xymatrix{
\text{id} \ar[r] &
f_*f^* \ar@<1ex>[r] \ar@<-1ex>[r] &
f_*f^*f_*f^* \ar@<0ex>[l] \ar@<-2ex>[r] \ar@<0ex>[r] \ar@<2ex>[r] &
f_*f^*f_*f^*f_*f^* \ar@<1ex>[l] \ar@<-1ex>[l]
}
$$
을 얻는다. Simplicial, Lemma
\ref{simplicial-lemma-standard-simplicial}을 보라. 더욱이 확대
$$
\xymatrix{
f^* \ar[r] &
f^*f_*f^* \ar@<1ex>[r] \ar@<-1ex>[r] &
f^*f_*f^*f_*f^* \ar@<0ex>[l] \ar@<-2ex>[r] \ar@<0ex>[r] \ar@<2ex>[r] &
f^*f_*f^*f_*f^*f_*f^* \ar@<1ex>[l] \ar@<-1ex>[l]
}
$$
는 호모토피 동치이다. Simplicial, Lemma
\ref{simplicial-lemma-standard-simplicial-homotopy}를 보라.

\begin{lemma}
\label{cohomology-lemma-godement-resolution}
$(X, \mathcal{O}_X)$를 환 달린 공간이라 하자. 모든
$\mathcal{O}_X$-가군층 $\mathcal{F}$에는 분해
$$
0 \to
\mathcal{F} \to
f_*f^*\mathcal{F} \to
f_*f^*f_*f^*\mathcal{F} \to
f_*f^*f_*f^*f_*f^*\mathcal{F} \to \ldots
$$
가 존재한다. 이는 $\mathcal{F}$에 대해 함자적이고, 각 항
$f_*f^* \ldots f_*f^*\mathcal{F}$는 플라스크
$\mathcal{O}_X$-가군이며, 모든 $x \in X$에 대해 사상
$$
\mathcal{F}_x[0] \to \Big(
(f_*f^*\mathcal{F})_x \to
(f_*f^*f_*f^*\mathcal{F})_x \to
(f_*f^*f_*f^*f_*f^*\mathcal{F})_x \to \ldots
\Big)
$$
은 $\mathcal{O}_{X, x}$-가군 복합체의 범주에서 호모토피 동치이다.
\end{lemma}

\begin{proof}
복합체
$f_*f^*\mathcal{F} \to  f_*f^*f_*f^*\mathcal{F} \to
f_*f^*f_*f^*f_*f^*\mathcal{F} \to \ldots$는 위에서 설명한,
항들이
$f_*f^*\mathcal{F}, f_*f^*f_*f^*\mathcal{F},
f_*f^*f_*f^*f_*f^*\mathcal{F}, \ldots$인 코심플렉셜 대상에 딸린
복합체이다. Simplicial, Section
\ref{simplicial-section-dold-kan-cosimplicial}을 보라. 확대에서
표시된 사상 $\mathcal{F} \to f_*f^*\mathcal{F}$가 나온다.
$X_{disc}$ 위의 임의의 아벨 층 $\mathcal{H}$에 대해 직접상
$f_*\mathcal{H}$는 플라스크이다. 실제로 $X_{disc}$는 이산공간이고
플라스크 층의 직접상은 플라스크이다. 따라서 이 복합체의 항들은
플라스크 $\mathcal{O}_X$-가군이다.

\medskip\noindent
$x \in X_{disc} = X$가 점이면 모든 $\mathcal{O}_X$-가군
$\mathcal{G}$에 대해 $(f^*\mathcal{G})_x = \mathcal{G}_x$이다.
따라서 $f^*$는 완전 함자이고, $\mathcal{O}_X$-가군 복합체
$\mathcal{G}_1 \to \mathcal{G}_2 \to \mathcal{G}_3$가 완전일
필요충분조건은
$f^*\mathcal{G}_1 \to f^*\mathcal{G}_2 \to f^*\mathcal{G}_3$가
완전인 것이다(Modules, Lemma \ref{modules-lemma-abelian} 참조).
이 절의 도입부에서 언급한 결과에 의해 $f^*$를 적용하면 상수
코심플렉셜 대상 $f^*\mathcal{F}$에서 항들이
$f_*f^*\mathcal{F}, f_*f^*f_*f^*\mathcal{F},
f_*f^*f_*f^*f_*f^*\mathcal{F}, \ldots$인 코심플렉셜 대상으로 가는
호모토피 동치를 얻는다. Simplicial, Lemma
\ref{simplicial-lemma-homotopy-equivalence-s-Q}에 의해
$$
f^*\mathcal{F}[0] \to \Big(
f^*f_*f^*\mathcal{F} \to
f^*f_*f^*f_*f^*\mathcal{F} \to
f^*f_*f^*f_*f^*f_*f^*\mathcal{F} \to \ldots
\Big)
$$
이 호모토피 동치임을 얻는다. 이로부터 보조정리의 나머지 두 명제가
곧바로 따른다.
\end{proof}

\begin{lemma}
\label{cohomology-lemma-godement-resolution-bounded-below}
$(X, \mathcal{O}_X)$를 환 달린 공간이라 하고,
$\mathcal{F}^\bullet$를 아래로 유계인 $\mathcal{O}_X$-가군 복합체라
하자. 아래로 유계인 플라스크 $\mathcal{O}_X$-가군 복합체
$\mathcal{G}^\bullet$로 가는 유사동형
$\mathcal{F}^\bullet \to \mathcal{G}^\bullet$가 존재하고, 모든
$x \in X$에 대해 사상
$\mathcal{F}^\bullet_x \to \mathcal{G}^\bullet_x$은
$\mathcal{O}_{X, x}$-가군 복합체의 범주에서 호모토피 동치이다.
\end{lemma}

\begin{proof}
$\mathcal{O}_X$-가군 복합체의 범주를 $\mathcal{A}$라 하고,
$\mathcal{O}_X$-가군 복합체의 범주를 $\mathcal{B}$라 하자. 그러면
$\mathcal{A}$와 $\mathcal{B}$ 사이의 수반 함자 $f^*$와 $f_*$에 위의
논의를 적용할 수 있다. Lemma \ref{cohomology-lemma-godement-resolution}의 증명과
정확히 같은 논증으로 아벨 범주 $\mathcal{A}$ 안의 분해
$$
0 \to
\mathcal{F}^\bullet \to
f_*f^*\mathcal{F}^\bullet \to
f_*f^*f_*f^*\mathcal{F}^\bullet \to
f_*f^*f_*f^*f_*f^*\mathcal{F}^\bullet \to \ldots
$$
를 얻는다. 여기서 각
$f_*f^*\ldots f_*f^*\mathcal{F}^\bullet$의 각 항은 플라스크
$\mathcal{O}_X$-가군이고, 모든 $x \in X$에 대해 사상
$$
\mathcal{F}^\bullet_x[0] \to \Big(
(f_*f^*\mathcal{F}^\bullet)_x \to
(f_*f^*f_*f^*\mathcal{F}^\bullet)_x \to
(f_*f^*f_*f^*f_*f^*\mathcal{F}^\bullet)_x \to \ldots
\Big)
$$
은 $\mathcal{O}_{X, x}$-가군 복합체들의 복합체 범주에서 호모토피
동치이다. 복합체들의 복합체는 이중 복합체와 같으므로 유도된 사상
$$
\mathcal{F}^\bullet \to
\mathcal{G}^\bullet =
\text{Tot}(
f_*f^*\mathcal{F}^\bullet \to
f_*f^*f_*f^*\mathcal{F}^\bullet \to
f_*f^*f_*f^*f_*f^*\mathcal{F}^\bullet \to \ldots
)
$$
을 생각할 수 있다. 복합체 $\mathcal{F}^\bullet$가 아래로 유계이므로
$\mathcal{G}^\bullet$도 그러하며, 실제로
$\mathcal{G}^\bullet$의 각 항은 복합체
$f_*f^*\ldots f_*f^*\mathcal{F}^\bullet$들의 항들의 유한 직합이므로
플라스크이다. 이제 보조정리의 마지막 주장은 Homology, Lemma
\ref{homology-lemma-homotopy-complex-complexes}에서 따른다. 특히 이는
$\mathcal{F}^\bullet \to \mathcal{G}^\bullet$가 유사동형임을
보이므로 증명이 끝난다.
\end{proof}









\section{컵곱}
\label{cohomology-section-cup-product}

\noindent
$(X, \mathcal{O}_X)$를 환 달린 공간이라 하고, $K, M$을
$D(\mathcal{O}_X)$의 대상이라 하자. $A = \Gamma(X, \mathcal{O}_X)$로
두자. 이 상황에서 (전역) 컵곱은 $D(A)$ 안의 사상
$$
\mu :
R\Gamma(X, K) \otimes_A^\mathbf{L} R\Gamma(X, M)
\longrightarrow
R\Gamma(X, K \otimes_{\mathcal{O}_X}^\mathbf{L} M)
$$
이다. 이를 Remark \ref{cohomology-remark-cup-product}에서처럼 환 달린 공간의 사상
$f : (X, \mathcal{O}_X) \to (pt, A)$에 대한 상대 컵곱으로 정의하되
$D(pt, A) = D(A)$를 사용한다. 특히 이 사상은 짝짓기
$$
\cup :
H^i(X, K) \times H^j(X, M)
\longrightarrow
H^{i + j}(X, K \otimes_{\mathcal{O}_X}^\mathbf{L} M)
$$
를 정의한다. 구체적으로
$\xi \in H^i(X, K) = H^i(R\Gamma(X, K))$와
$\eta \in H^j(X, M) = H^j(R\Gamma(X, M))$가 주어지면 먼저 이들을
“텐서곱”하여
$H^{i + j}(R\Gamma(X, K) \otimes_A^\mathbf{L} R\Gamma(X, M))$의 원소
$\xi \otimes \eta$를 얻을 수 있다. More on Algebra, Section
\ref{more-algebra-section-products-tor}를 보라. 그런 다음 $\mu$를
적용하여 $H^{i + j}(X, K \otimes_{\mathcal{O}_X}^\mathbf{L} M)$의
원하는 원소
$\xi \cup \eta = \mu(\xi \otimes \eta)$를 얻는다.

\medskip\noindent
$\xi$와 $\eta$의 컵곱을 이해하는 또 다른 방법은 다음과 같다.
$\Hom(\mathcal{O}_X, -) = \Gamma(X, -)$이므로
$$
H^i(R\Gamma(X, K)) = \Hom_{D(\mathcal{O}_X)}(\mathcal{O}_X[-i], K)
\quad\text{and}\quad
H^j(R\Gamma(X, M)) = \Hom_{D(\mathcal{O}_X)}(\mathcal{O}_X[-j], M)
$$
로 쓸 수 있다. 따라서 $\xi$와 $\eta$는 각각 사상
$$
\tilde \xi : \mathcal{O}_X[-i] \to K
\quad\text{and}\quad
\tilde \eta : \mathcal{O}_X[-j] \to M
$$
와 “같은” 것이다. 이를 유도 텐서곱의 함자성과 결합하면
$$
\mathcal{O}_X[-i - j] =
\mathcal{O}_X[-i] \otimes_{\mathcal{O}_X}^\mathbf{L} \mathcal{O}_X[-j]
\xrightarrow{\tilde \xi \otimes \tilde \eta}
K \otimes_{\mathcal{O}_X}^\mathbf{L} M
$$
을 얻으며, 위와 같은 이유로 이는
$H^{i + j}(X, K \otimes_{\mathcal{O}_X}^\mathbf{L} M)$의 원소이다.

\begin{lemma}
\label{cohomology-lemma-second-cup-equals-first}
이 구성은 컵곱을 준다.
\end{lemma}

\begin{proof}
위와 같이 $f : (X, \mathcal{O}_X) \to (pt, A)$로 두면
$Rf_*(-) = R\Gamma(X, -)$이고, 사상 $\mu$는 사상
$$
Lf^*(Rf_*K \otimes_A^\mathbf{L} Rf_*M) =
Lf^*Rf_*K \otimes_{\mathcal{O}_X}^\mathbf{L} Lf^*Rf_*M
\xrightarrow{\epsilon_K \otimes \epsilon_M}
K \otimes_{\mathcal{O}_X}^\mathbf{L} M
$$
과 수반이다. 여기서 $\epsilon$은 $Lf^*$와 $Rf_*$ 사이 수반의
여단위이다. $\xi$와 $\eta$를 사상
$\xi : A[-i] \to R\Gamma(X, K)$와
$\eta : A[-j] \to R\Gamma(X, M)$로 생각하면, 텐서곱
$\xi \otimes \eta$는 사상\footnote{여기에는 부호 하나가 숨어 있다.
즉 이 등식은 합성
$$
A[-i - j] \to (A \otimes_A^\mathbf{L} A)[-i - j] \to
A[-i] \otimes_A^\mathbf{L} A[-j]
$$
으로 정의되며, 두 번째 단계에서는 원칙적으로 부호를 사용하는 More on
Algebra, Item (\ref{more-algebra-item-shift-tensor})의 식별을 쓴다.
다만 이 경우 우리 규약에서는 그 부호가 $+1$이다. 설령 $+1$이
아니더라도
$\mathcal{O}_X[-i - j] =
\mathcal{O}_X[-i] \otimes_{\mathcal{O}_X}^\mathbf{L} \mathcal{O}_X[-j]$의
식별에서 같은 부호를 사용했으므로 상관없다.}
$$
A[-i - j] = A[-i] \otimes_A^\mathbf{L} A[-j]
\xrightarrow{\xi \otimes \eta}
R\Gamma(X, K) \otimes_A^\mathbf{L} R\Gamma(X, M)
$$
에 대응한다. 정의에 의해 컵곱 $\xi \cup \eta$는 사상
$A[-i - j] \to R\Gamma(X, K \otimes_{\mathcal{O}_X}^\mathbf{L} M)$이며,
이는
$$
(\epsilon_K \otimes \epsilon_M) \circ Lf^*(\xi \otimes \eta) =
(\epsilon_K \circ Lf^*\xi) \otimes (\epsilon_M \circ Lf^*\eta)
$$
와 수반이다. 그런데
$\epsilon_K \circ Lf^*\xi = \tilde \xi$이고
$\epsilon_M \circ Lf^*\eta = \tilde \eta$임은 쉽게 알 수 있다.
따라서 $\widetilde{\xi \cup \eta} = \tilde \xi \otimes \tilde \eta$이며,
이는 두 구성이 원하는 대로 일치한다는 뜻이다.
\end{proof}

\begin{remark}
\label{cohomology-remark-cup-with-element-map-total-cohomology}
$(X, \mathcal{O}_X)$를 환 달린 공간이라 하고, $K, M$을
$D(\mathcal{O}_X)$의 대상이라 하자. $A = \Gamma(X, \mathcal{O}_X)$로
두자. $\xi \in H^i(X, K)$가 주어지면 이에 딸린 사상
$$
\xi = ``\xi \cup -'' :
R\Gamma(X, M)[-i]
\to
R\Gamma(X, K \otimes_{\mathcal{O}_X}^\mathbf{L} M)
$$
을 얻는다. 실제로 Lemma \ref{cohomology-lemma-second-cup-equals-first}의
증명에서처럼 $\xi$를 사상
$\xi : A[-i] \to R\Gamma(X, K)$로 나타내고, 합성
$$
R\Gamma(X, M)[-i] = A[-i] \otimes_A^\mathbf{L} R\Gamma(X, M)
\xrightarrow{\xi \otimes 1}
R\Gamma(X, K) \otimes_A^\mathbf{L} R\Gamma(X, M)
\to
R\Gamma(X, K \otimes_{\mathcal{O}_X}^\mathbf{L} M)
$$
을 사용한다. 여기서 두 번째 화살표는 위의 전역 컵곱 $\mu$이다.
코호몰로지를 취하면 Lemma \ref{cohomology-lemma-second-cup-equals-first}와 그
증명에서 명백하듯이 $\xi$에 의한 컵곱을 되찾는다.
\end{remark}

\noindent
상대 컵곱의 자연스러운 호환성을 정식화하고 증명하자. 환 달린 공간의
사상 $f : (X, \mathcal{O}_X) \to (Y, \mathcal{O}_Y)$가 주어졌다고
하자. $\mathcal{K}^\bullet$와 $\mathcal{M}^\bullet$를
$\mathcal{O}_X$-가군 복합체라 하자. 소박한 컵곱
$$
\text{Tot}(
f_*\mathcal{K}^\bullet
\otimes_{\mathcal{O}_Y}
f_*\mathcal{M}^\bullet)
\longrightarrow
f_*\text{Tot}(\mathcal{K}^\bullet
\otimes_{\mathcal{O}_X}
\mathcal{M}^\bullet)
$$
이 존재한다. 이것이 상대 컵곱과 관련됨을 주장한다.

\begin{lemma}
\label{cohomology-lemma-cup-compatible-with-naive}
위의 상황에서 다음 도식은 가환한다.
$$
\xymatrix{
f_*\mathcal{K}^\bullet
\otimes_{\mathcal{O}_Y}^\mathbf{L}
f_*\mathcal{M}^\bullet \ar[r] \ar[d]
&
Rf_*\mathcal{K}^\bullet
\otimes_{\mathcal{O}_Y}^\mathbf{L}
Rf_*\mathcal{M}^\bullet \ar[d]^{\text{Remark \ref{cohomology-remark-cup-product}}} \\
\text{Tot}(
f_*\mathcal{K}^\bullet
\otimes_{\mathcal{O}_Y}
f_*\mathcal{M}^\bullet) \ar[d]_{\text{naive cup product}} &
Rf_*(\mathcal{K}^\bullet
\otimes_{\mathcal{O}_X}^\mathbf{L}
\mathcal{M}^\bullet) \ar[d] \\
f_*\text{Tot}(\mathcal{K}^\bullet
\otimes_{\mathcal{O}_X}
\mathcal{M}^\bullet) \ar[r] &
Rf_*\text{Tot}(\mathcal{K}^\bullet
\otimes_{\mathcal{O}_X}
\mathcal{M}^\bullet)
}
$$
\end{lemma}

\begin{proof}
Remark \ref{cohomology-remark-cup-product}의 구성에 의해, 도식을 시계방향으로
돌아 얻는 사상
$$
f_*\mathcal{K}^\bullet
\otimes_{\mathcal{O}_Y}^\mathbf{L}
f_*\mathcal{M}^\bullet 
\longrightarrow
Rf_*\text{Tot}(\mathcal{K}^\bullet
\otimes_{\mathcal{O}_X}
\mathcal{M}^\bullet)
$$
은 사상
\begin{align*}
Lf^*(f_*\mathcal{K}^\bullet
\otimes_{\mathcal{O}_Y}^\mathbf{L}
f_*\mathcal{M}^\bullet)
& =
Lf^*f_*\mathcal{K}^\bullet
\otimes_{\mathcal{O}_Y}^\mathbf{L}
Lf^*f_*\mathcal{M}^\bullet \\
& \to
Lf^*Rf_*\mathcal{K}^\bullet
\otimes_{\mathcal{O}_Y}^\mathbf{L}
Lf^*Rf_*\mathcal{M}^\bullet \\
& \to
\mathcal{K}^\bullet
\otimes_{\mathcal{O}_Y}^\mathbf{L}
\mathcal{M}^\bullet \\
& \to
\text{Tot}(\mathcal{K}^\bullet
\otimes_{\mathcal{O}_X}
\mathcal{M}^\bullet)
\end{align*}
과 수반임을 알 수 있다. Lemma
\ref{cohomology-lemma-adjoints-push-pull-compatibility}에 의해 이는 또한
\begin{align*}
Lf^*(f_*\mathcal{K}^\bullet
\otimes_{\mathcal{O}_Y}^\mathbf{L}
f_*\mathcal{M}^\bullet)
& =
Lf^*f_*\mathcal{K}^\bullet
\otimes_{\mathcal{O}_Y}^\mathbf{L}
Lf^*f_*\mathcal{M}^\bullet \\
& \to
f^*f_*\mathcal{K}^\bullet
\otimes_{\mathcal{O}_Y}^\mathbf{L}
f^*f_*\mathcal{M}^\bullet \\
& \to
\mathcal{K}^\bullet
\otimes_{\mathcal{O}_Y}^\mathbf{L}
\mathcal{M}^\bullet \\
& \to
\text{Tot}(\mathcal{K}^\bullet
\otimes_{\mathcal{O}_X}
\mathcal{M}^\bullet)
\end{align*}
와 같다. 반시계방향으로 돌면 사상
\begin{align*}
Lf^*(f_*\mathcal{K}^\bullet
\otimes_{\mathcal{O}_Y}^\mathbf{L}
f_*\mathcal{M}^\bullet)
& \to
Lf^*\text{Tot}(
f_*\mathcal{K}^\bullet
\otimes_{\mathcal{O}_Y}
f_*\mathcal{M}^\bullet) \\
& \to
Lf^*f_*\text{Tot}(\mathcal{K}^\bullet
\otimes_{\mathcal{O}_X}
\mathcal{M}^\bullet) \\
& \to
Lf^*Rf_*\text{Tot}(\mathcal{K}^\bullet
\otimes_{\mathcal{O}_X}
\mathcal{M}^\bullet) \\
& \to
\text{Tot}(\mathcal{K}^\bullet
\otimes_{\mathcal{O}_X}
\mathcal{M}^\bullet)
\end{align*}
과 수반인 사상을 얻는다. Lemma
\ref{cohomology-lemma-adjoints-push-pull-compatibility}에 의해 이는 또한
\begin{align*}
Lf^*(f_*\mathcal{K}^\bullet
\otimes_{\mathcal{O}_Y}^\mathbf{L}
f_*\mathcal{M}^\bullet)
& \to
Lf^*\text{Tot}(
f_*\mathcal{K}^\bullet
\otimes_{\mathcal{O}_Y}
f_*\mathcal{M}^\bullet) \\
& \to
Lf^*f_*\text{Tot}(\mathcal{K}^\bullet
\otimes_{\mathcal{O}_X}
\mathcal{M}^\bullet) \\
& \to
f^*f_*\text{Tot}(\mathcal{K}^\bullet
\otimes_{\mathcal{O}_X}
\mathcal{M}^\bullet) \\
& \to
\text{Tot}(\mathcal{K}^\bullet
\otimes_{\mathcal{O}_X}
\mathcal{M}^\bullet)
\end{align*}
와 같다. 이제 다음 도식을 살펴보면 증명이 끝난다.
$$
\xymatrix{
Lf^*(f_*\mathcal{K}^\bullet
\otimes_{\mathcal{O}_Y}^\mathbf{L}
f_*\mathcal{M}^\bullet) \ar[d] \ar[rr] & &
Lf^*f_*\mathcal{K}^\bullet \otimes_{\mathcal{O}_X}^\mathbf{L}
Lf^*f_*\mathcal{M}^\bullet \ar[d] \\
Lf^*\text{Tot}(
f_*\mathcal{K}^\bullet
\otimes_{\mathcal{O}_Y}
f_*\mathcal{M}^\bullet) \ar[d]_{naive} \ar[r] &
f^*\text{Tot}(
f_*\mathcal{K}^\bullet
\otimes_{\mathcal{O}_Y}
f_*\mathcal{M}^\bullet) \ar[ldd]^{naive} \ar[dd] &
f^*f_*\mathcal{K}^\bullet \otimes_{\mathcal{O}_X}^\mathbf{L}
f^*f_*\mathcal{M}^\bullet \ar[dd] \ar[ldd] \\
Lf^*f_*\text{Tot}(\mathcal{K}^\bullet
\otimes_{\mathcal{O}_X}
\mathcal{M}^\bullet) \ar[d] \\
f^*f_*\text{Tot}(\mathcal{K}^\bullet \otimes_{\mathcal{O}_X}
\mathcal{M}^\bullet) \ar[rd] &
\text{Tot}(f^*f_*\mathcal{K}^\bullet \otimes_{\mathcal{O}_X}
f^*f_*\mathcal{M}^\bullet) \ar[d] &
\mathcal{K}^\bullet \otimes_{\mathcal{O}_X}^\mathbf{L}
\mathcal{M}^\bullet \ar[ld] \\
& \text{Tot}(\mathcal{K}^\bullet
\otimes_{\mathcal{O}_X}
\mathcal{M}^\bullet)
}
$$
이 도식의 모든 다각형은 가환한다. 위쪽 다각형은 Lemma
\ref{cohomology-lemma-tensor-pull-compatibility}에 의해 가환한다. 두 소박한
컵곱을 포함하는 정사각형은 $Lf^* \to f^*$가 가군 복합체에 대해
함자적이므로 가환한다. 두 사상
$\mathcal{A}^\bullet \otimes^\mathbf{L} \mathcal{B}^\bullet \to
\text{Tot}(\mathcal{A}^\bullet \otimes \mathcal{B}^\bullet)$을 포함하는
정사각형도 마찬가지이다. 마지막으로 남은 정사각형의 가환성은 복합체
수준에서 성립하며, $f^*$와 $f_*$의 수반성에 의한 소박한 컵곱의
정의로 볼 수 있다. 도식의 바깥쪽을 도는 두 사상이 위에서 주어진 두
사상이므로 증명이 끝난다.
\end{proof}

\noindent
$(X, \mathcal{O}_X)$를 환 달린 공간이라 하고,
$\mathcal{K}^\bullet$와 $\mathcal{M}^\bullet$를
$\mathcal{O}_X$-가군 복합체라 하자. 그러면 “소박한” 컵곱
$$
\mu' :
\text{Tot}(
\Gamma(X, \mathcal{K}^\bullet) \otimes_A \Gamma(X, \mathcal{M}^\bullet))
\longrightarrow
\Gamma(X, \text{Tot}(
\mathcal{K}^\bullet \otimes_{\mathcal{O}_X} \mathcal{M}^\bullet))
$$
이 존재한다. Lemma \ref{cohomology-lemma-cup-compatible-with-naive}를 사상
$(X, \mathcal{O}_X) \to (pt, A)$에 적용하면, 이 소박한 컵곱은 이 절의
첫 문단에서 정의한 컵곱 $\mu$와 다음 가환도식으로 관련된다.
$$
\xymatrix{
\Gamma(X, \mathcal{K}^\bullet)
\otimes_A^\mathbf{L}
\Gamma(X, \mathcal{M}^\bullet) \ar[d] \ar[r] &
R\Gamma(X, \mathcal{K}^\bullet)
\otimes_A^\mathbf{L}
R\Gamma(X, \mathcal{M}^\bullet) \ar[d]^-\mu \\
\text{Tot}(\Gamma(X, \mathcal{K}^\bullet)
\otimes_A
\Gamma(X, \mathcal{M}^\bullet)) \ar[d]_-{\mu'} &
R\Gamma(X, \mathcal{K}^\bullet
\otimes_{\mathcal{O}_X}^\mathbf{L}
\mathcal{M}^\bullet) \ar[d] \\
\Gamma(X, \text{Tot}(\mathcal{K}^\bullet
\otimes_{\mathcal{O}_X}
\mathcal{M}^\bullet)) \ar[r] &
R\Gamma(X, \text{Tot}(\mathcal{K}^\bullet
\otimes_{\mathcal{O}_X}
\mathcal{M}^\bullet))
}
$$
이는 $D(A)$ 안의 도식이다. 코호몰로지를 취하면 소박한 컵곱과 실제
컵곱을 연결하는 가환도식
$$
\xymatrix{
H^i(\Gamma(X, \mathcal{K}^\bullet)) \times
H^j(\Gamma(X, \mathcal{M}^\bullet)) \ar[d] \ar[r] &
H^{i + j}(X,
\text{Tot}(\mathcal{K}^\bullet \otimes_{\mathcal{O}_X} \mathcal{M}^\bullet)) \\
H^i(X, \mathcal{K}^\bullet) \times
H^j(X, \mathcal{M}^\bullet) \ar[r]^\cup &
H^{i + j}(X, \mathcal{K}^\bullet \otimes_{\mathcal{O}_X}^\mathbf{L}
\mathcal{M}^\bullet) \ar[u]
}
$$
을 얻는다.

\begin{lemma}
\label{cohomology-lemma-diagrams-commute}
$(X, \mathcal{O}_X)$를 환 달린 공간이라 하자.
$\mathcal{K}^\bullet$와 $\mathcal{M}^\bullet$를 아래로 유계인
$\mathcal{O}_X$-가군 복합체라 하고,
$\mathcal{U} : X = \bigcup_{i \in I} U_i$를 열린 덮개라 하자. 그러면
도식
$$
\xymatrix{
\text{Tot}(\check{\mathcal{C}}^\bullet(\mathcal{U}, \mathcal{K}^\bullet))
\otimes_A^\mathbf{L}
\text{Tot}(\check{\mathcal{C}}^\bullet(\mathcal{U}, \mathcal{M}^\bullet))
\ar[d] \ar[r] &
R\Gamma(X, \mathcal{K}^\bullet)
\otimes_A^\mathbf{L}
R\Gamma(X, \mathcal{M}^\bullet) \ar[d]^\mu \\
\text{Tot}(
\text{Tot}(\check{\mathcal{C}}^\bullet(\mathcal{U}, \mathcal{K}^\bullet))
\otimes_A
\text{Tot}(\check{\mathcal{C}}^\bullet(\mathcal{U}, \mathcal{M}^\bullet)))
\ar[d]^{(\ref{cohomology-equation-needs-signs})} &
R\Gamma(X,
\mathcal{K}^\bullet \otimes_{\mathcal{O}_X}^\mathbf{L} \mathcal{M}^\bullet)
\ar[d] \\
\text{Tot}(
\check{\mathcal{C}}^\bullet({\mathcal U},
\text{Tot}(\mathcal{K}^\bullet \otimes_{\mathcal{O}_X} \mathcal{M}^\bullet)
)) \ar[r] &
R\Gamma(X,
\text{Tot}(\mathcal{K}^\bullet \otimes_{\mathcal{O}_X} \mathcal{M}^\bullet))
}
$$
은 $D(A)$ 안에서 가환한다. 여기서 수평 화살표들은 Lemma
\ref{cohomology-lemma-cech-complex-complex}의 화살표들이다.
\end{lemma}

\begin{proof}
Lemma \ref{cohomology-lemma-godement-resolution-bounded-below}에서와 같은 복합체의
유사동형
$a : \mathcal{K}^\bullet \to \mathcal{K}_1^\bullet$와
$b : \mathcal{M}^\bullet \to \mathcal{M}_1^\bullet$를 고르자. 사상
$a$와 $b$는 줄기에서 호모토피 동치이므로 유도된 사상
$$
\text{Tot}(\mathcal{K}^\bullet \otimes_{\mathcal{O}_X} \mathcal{M}^\bullet)
\to
\text{Tot}(\mathcal{K}_1^\bullet \otimes_{\mathcal{O}_X} \mathcal{M}_1^\bullet)
$$
도 줄기에서 호모토피 동치이며(More on Algebra, Lemma
\ref{more-algebra-lemma-derived-tor-homotopy}), 따라서 유사동형이다.
그러므로 두 도식의 표적
$$
R\Gamma(X,
\text{Tot}(\mathcal{K}^\bullet
\otimes_{\mathcal{O}_X} \mathcal{M}^\bullet)) =
R\Gamma(X,
\text{Tot}(\mathcal{K}_1^\bullet
\otimes_{\mathcal{O}_X} \mathcal{M}_1^\bullet))
$$
은 $D(A)$ 안에서 같다. 따라서 $\mathcal{K}$와 $\mathcal{M}$을
$\mathcal{K}_1$과 $\mathcal{M}_1$으로 바꾼 경우에 도식이 가환함을
증명하면 충분하다. 이는 다음 문단에서 다루는 경우로 환원한다.

\medskip\noindent
$\mathcal{K}^\bullet$와 $\mathcal{M}^\bullet$가 아래로 유계인
플라스크 $\mathcal{O}_X$-가군 복합체라고 가정하고, 컵곱을 체흐
복합체 위의 컵곱 (\ref{cohomology-equation-needs-signs})과 연결하는 도식을
생각하자. 그러면 가환도식
$$
\xymatrix{
\Gamma(X, \mathcal{K}^\bullet)
\otimes_A^\mathbf{L}
\Gamma(X, \mathcal{M}^\bullet) \ar[d] \ar[r] &
\text{Tot}(\check{\mathcal{C}}^\bullet(\mathcal{U}, \mathcal{K}^\bullet))
\otimes_A^\mathbf{L}
\text{Tot}(\check{\mathcal{C}}^\bullet(\mathcal{U}, \mathcal{M}^\bullet))
\ar[d] \\
\text{Tot}(\Gamma(X, \mathcal{K}^\bullet)
\otimes_A
\Gamma(X, \mathcal{M}^\bullet)) \ar[d] \ar[r] &
\text{Tot}(
\text{Tot}(\check{\mathcal{C}}^\bullet(\mathcal{U}, \mathcal{K}^\bullet))
\otimes_A
\text{Tot}(\check{\mathcal{C}}^\bullet(\mathcal{U}, \mathcal{M}^\bullet)))
\ar[d]^{(\ref{cohomology-equation-needs-signs})} \\
\Gamma(X, \text{Tot}(\mathcal{K}^\bullet
\otimes_{\mathcal{O}_X}
\mathcal{M}^\bullet)) \ar[r] &
\text{Tot}(
\check{\mathcal{C}}^\bullet({\mathcal U},
\text{Tot}(\mathcal{K}^\bullet \otimes_{\mathcal{O}_X} \mathcal{M}^\bullet)
))
}
$$
을 생각할 수 있다. 이 도식의 수평 화살표들은 $D(A)$ 안에서
동형이다. 실제로 $\mathcal{K}^\bullet$와 같은 아래로 유계인 플라스크
가군 복합체에 대해 $D(A)$ 안에서
$$
\Gamma(X, \mathcal{K}^\bullet) =
\text{Tot}(\check{\mathcal{C}}^\bullet(\mathcal{U}, \mathcal{K}^\bullet)) =
R\Gamma(X, \mathcal{K}^\bullet)
$$
이다. 이는 Lemma \ref{cohomology-lemma-flasque-acyclic}, Derived Categories,
Lemma \ref{derived-lemma-leray-acyclicity}, 그리고 Lemma
\ref{cohomology-lemma-cech-complex-complex-computes}에서 따른다. 따라서
(\ref{cohomology-equation-needs-signs})를 포함하는 보조정리의 도식의 가환성은
이미 증명한 Lemma \ref{cohomology-lemma-cup-compatible-with-naive}의 가환성에서
따른다. 여기서 $f$는 한 점으로 가는 사상이다(Lemma
\ref{cohomology-lemma-cup-compatible-with-naive} 다음의 논의를 보라).
\end{proof}

\begin{lemma}
\label{cohomology-lemma-cup-product-associative}
$f : (X, \mathcal{O}_X) \to (Y, \mathcal{O}_Y)$를 환 달린 공간의
사상이라 하자. Remark \ref{cohomology-remark-cup-product}의 상대 컵곱은 다음
의미에서 결합적이다. 즉 $D(\mathcal{O}_X)$ 안의 모든 $K, L, M$에
대해 도식
$$
\xymatrix{
Rf_*K \otimes_{\mathcal{O}_Y}^\mathbf{L}
Rf_*L \otimes_{\mathcal{O}_Y}^\mathbf{L}
Rf_*M \ar[r] \ar[d] &
Rf_*(K \otimes_{\mathcal{O}_X}^\mathbf{L} L)
\otimes_{\mathcal{O}_Y}^\mathbf{L} Rf_*M \ar[d] \\
Rf_*K \otimes_{\mathcal{O}_Y}^\mathbf{L}
Rf_*(L \otimes_{\mathcal{O}_X}^\mathbf{L} M) \ar[r] &
Rf_*(K \otimes_{\mathcal{O}_X}^\mathbf{L} 
L \otimes_{\mathcal{O}_X}^\mathbf{L} M)
}
$$
은 $D(\mathcal{O}_Y)$ 안에서 가환한다.
\end{lemma}

\begin{proof}
어느 쪽으로 돌아가도 $D(\mathcal{O}_X)$ 안의 자명한 사상
\begin{align*}
Lf^*(Rf_*K \otimes_{\mathcal{O}_Y}^\mathbf{L}
Rf_*L \otimes_{\mathcal{O}_Y}^\mathbf{L}
Rf_*M) & =
Lf^*(Rf_*K) \otimes_{\mathcal{O}_X}^\mathbf{L}
Lf^*(Rf_*L) \otimes_{\mathcal{O}_X}^\mathbf{L}
Lf^*(Rf_*M) \\
& \to
K \otimes_{\mathcal{O}_X}^\mathbf{L} 
L \otimes_{\mathcal{O}_X}^\mathbf{L} M
\end{align*}
과 수반인 사상을 얻는다.
\end{proof}

\begin{lemma}
\label{cohomology-lemma-cup-product-commutative}
$f : (X, \mathcal{O}_X) \to (Y, \mathcal{O}_Y)$를 환 달린 공간의
사상이라 하자. Remark \ref{cohomology-remark-cup-product}의 상대 컵곱은 다음
의미에서 가환적이다. 즉 $D(\mathcal{O}_X)$ 안의 모든 $K, L$에 대해
도식
$$
\xymatrix{
Rf_*K \otimes_{\mathcal{O}_Y}^\mathbf{L} Rf_*L \ar[r] \ar[d]_\psi &
Rf_*(K \otimes_{\mathcal{O}_X}^\mathbf{L} L) \ar[d]^{Rf_*\psi} \\
Rf_*L \otimes_{\mathcal{O}_Y}^\mathbf{L} Rf_*K \ar[r] &
Rf_*(L \otimes_{\mathcal{O}_X}^\mathbf{L} K)
}
$$
은 $D(\mathcal{O}_Y)$ 안에서 가환한다. 여기서 $\psi$는 유도범주의
가환성 제약이다(Lemma \ref{cohomology-lemma-symmetric-monoidal-derived}).
\end{lemma}

\begin{proof}
생략한다.
\end{proof}

\begin{lemma}
\label{cohomology-lemma-compose-cup-product}
$f : (X, \mathcal{O}_X) \to (Y, \mathcal{O}_Y)$와
$g : (Y, \mathcal{O}_Y) \to (Z, \mathcal{O}_Z)$를 환 달린 공간의
사상이라 하자. Remark \ref{cohomology-remark-cup-product}의 상대 컵곱은 합성과
호환된다. 즉 $D(\mathcal{O}_X)$ 안의 모든 $K, L$에 대해 도식
$$
\xymatrix{
R(g \circ f)_*K \otimes_{\mathcal{O}_Z}^\mathbf{L} R(g \circ f)_*L
\ar@{=}[rr] \ar[d] & &
Rg_*Rf_*K \otimes_{\mathcal{O}_Z}^\mathbf{L} Rg_*Rf_*L \ar[d] \\
R(g \circ f)_*(K \otimes_{\mathcal{O}_X}^\mathbf{L} L) \ar@{=}[r] &
Rg_*Rf_*(K \otimes_{\mathcal{O}_X}^\mathbf{L} L) &
Rg_*(Rf_*K \otimes_{\mathcal{O}_Y}^\mathbf{L}  Rf_*L) \ar[l]
}
$$
은 $D(\mathcal{O}_Z)$ 안에서 가환한다.
\end{lemma}

\begin{proof}
도식의 어느 쪽으로 돌아가도 $D(\mathcal{O}_X)$ 안의 사상
\begin{align*}
& L(g \circ f)^*\left(R(g \circ f)_*K
\otimes_{\mathcal{O}_Z}^\mathbf{L}
R(g \circ f)_*L\right) \\
& =
L(g \circ f)^*R(g \circ f)_*K
\otimes_{\mathcal{O}_X}^\mathbf{L}
L(g \circ f)^*R(g \circ f)_*L) \\
& \to
K \otimes_{\mathcal{O}_X}^\mathbf{L} L
\end{align*}
과 수반인 사상을 얻으므로 성립한다. 이를 확인하려면 여단위들의 합성
$$
L(g \circ f)^*R(g \circ f)_* =
Lf^* Lg^* Rg_* Rf_*  \to
Lf^* Rf_* \to \text{id}
$$
이 $L(g \circ f)^*$와 $R(g \circ f)_*$에 대한 여단위임을 사용한다.
Categories, Lemma \ref{categories-lemma-compose-counits}를 보라.
\end{proof}

\begin{lemma}
\label{cohomology-lemma-base-change-cup-product}
환 달린 공간의 가환 정사각형
$$
\xymatrix{
(X', \mathcal{O}_{X'}) \ar[r]_{g'} \ar[d]_{f'} &
(X, \mathcal{O}_X) \ar[d]^f \\
(Y', \mathcal{O}_{Y'}) \ar[r]^g & (Y, \mathcal{O}_Y) 
}
$$
을 생각하자. $K, L$을 $D(\mathcal{O}_X)$의 대상으로 하자. 상대
컵곱은 다음 의미에서 이 정사각형과 호환된다. 즉 도식
$$
\xymatrix{
Lg^*(Rf_*K \otimes_{\mathcal{O}_Y}^\mathbf{L} Rf_*L)
\ar[r] \ar@{=}[d] &
Lg^*(Rf_*(K \otimes_{\mathcal{O}_X}^\mathbf{L} L)) \ar[d] \\
Lg^*Rf_*K \otimes_{\mathcal{O}_{Y'}}^\mathbf{L} Lg^*Rf_*L \ar[d] &
R(f')_*L(g')^*(K \otimes_{\mathcal{O}_X}^\mathbf{L} L) \ar@{=}[d] \\
R(f')_*(L(g')^*K \otimes_{\mathcal{O}_{Y'}} R(f')_*(L(g')^*L \ar[r] &
R(f')_*(L(g')^*K \otimes_{\mathcal{O}_{X'}}^\mathbf{L} L(g')^*L)
}
$$
은 $D(\mathcal{O}_{Y'})$ 안에서 가환한다. 수평 화살표들은 상대
컵곱(Remark \ref{cohomology-remark-cup-product})으로 주어지고, 수직 화살표들은
밑변환 사상(Remark \ref{cohomology-remark-base-change})과 Lemma
\ref{cohomology-lemma-pullback-tensor-product}로 주어진다.
\end{lemma}

\begin{proof}
생략한다.
\end{proof}












\section{K-단사 복합체의 몇 가지 성질}
\label{cohomology-section-properties-K-injective}

\noindent
$(X, \mathcal{O}_X)$를 환 달린 공간이라 하고, $U \subset X$를 열린 부분집합이라
하자. 이에 대응하는 열린 몰입을
$j : (U, \mathcal{O}_U) \to (X, \mathcal{O}_X)$라 쓰자. 당김 함자
$j^*$는 제한 함자에 불과하므로 완전하다. 따라서 유도 당김 $Lj^*$는
임의의 복합체를 단순히 제한하여 계산된다. 우리는 흔히 이에 대응하는 함자를
$$
D(\mathcal{O}_X) \to D(\mathcal{O}_U),
\quad
E \mapsto j^*E = E|_U
$$
로 간단히 나타낸다. 마찬가지로 영에 의한 연장
$j_! : \textit{Mod}(\mathcal{O}_U) \to \textit{Mod}(\mathcal{O}_X)$
(Sheaves, Section \ref{sheaves-section-open-immersions} 참조)은 완전 함자이다
(Modules, Lemma \ref{modules-lemma-j-shriek-exact}). 그러므로 이는
대상 $F$를 나타내는 임의의 복합체에 $j_!$를 단순히 적용함으로써 함자
$$
j_! : D(\mathcal{O}_U) \to D(\mathcal{O}_X),\quad
F \mapsto j_!F
$$
를 유도한다.

\begin{lemma}
\label{cohomology-lemma-restrict-K-injective-to-open}
$X$를 환 달린 공간이라 하고, $U \subset X$를 열린 부분공간이라 하자.
$\mathcal{O}_X$-가군의 K-단사 복합체를 $U$로 제한하면
$\mathcal{O}_U$-가군의 K-단사 복합체가 된다.
\end{lemma}

\begin{proof}
이는 Derived Categories, Lemma
\ref{derived-lemma-adjoint-preserve-K-injectives}와 제한 함자가 완전한
왼쪽 수반 $j_!$를 갖는다는 사실에서 따른다.
$j_!$의 구성은 Sheaves, Section
\ref{sheaves-section-open-immersions}을, 완전성은 Modules, Lemma
\ref{modules-lemma-j-shriek-exact}를 보라.
\end{proof}

\begin{lemma}
\label{cohomology-lemma-unbounded-cohomology-of-open}
$X$를 환 달린 공간이라 하고, $U \subset X$를 열린 부분공간이라 하자.
$D(\mathcal{O}_X)$의 $K$에 대하여
$H^p(U, K) = H^p(U, K|_U)$이다.
\end{lemma}

\begin{proof}
$K$를 나타내는 $\mathcal{O}_X$-가군의 K-단사 복합체를
$\mathcal{I}^\bullet$라 하자. 코호몰로지의 구성에 의해
$$
H^q(U, K) = H^q(\Gamma(U, \mathcal{I}^\bullet)) =
H^q(\Gamma(U, \mathcal{I}^\bullet|_U))
$$
이다. Lemma \ref{cohomology-lemma-restrict-K-injective-to-open}에 의해
$\mathcal{I}^\bullet|_U$는 $K|_U$를 나타내는 K-단사 복합체이므로
결론이 따른다.
\end{proof}

\begin{lemma}
\label{cohomology-lemma-sheafification-cohomology}
$(X, \mathcal{O}_X)$를 환 달린 공간이라 하고, $K$를
$D(\mathcal{O}_X)$의 대상이라 하자. 준층
$$
U \mapsto H^q(U, K) = H^q(U, K|_U)
$$
의 층화는 $K$의 $q$차 코호몰로지 층 $H^q(K)$이다.
\end{lemma}

\begin{proof}
등식 $H^q(U, K) = H^q(U, K|_U)$는
Lemma \ref{cohomology-lemma-unbounded-cohomology-of-open}에서 따른다.
$K$를 나타내는 K-단사 복합체 $\mathcal{I}^\bullet$를 고르자. 그러면
코호몰로지의 구성에 의해
$$
H^q(U, K) =
\frac{\Ker(\mathcal{I}^q(U) \to \mathcal{I}^{q + 1}(U))}
{\Im(\mathcal{I}^{q - 1}(U) \to \mathcal{I}^q(U))}.
$$
한편
$H^q(K) = \Ker(\mathcal{I}^q \to \mathcal{I}^{q + 1})/
\Im(\mathcal{I}^{q - 1} \to \mathcal{I}^q)$이므로 결과는 명백하다.
\end{proof}

\begin{lemma}
\label{cohomology-lemma-restrict-direct-image-open}
$f : (X, \mathcal{O}_X) \to (Y, \mathcal{O}_Y)$를 환 달린 공간의
사상이라 하자. 열린 부분공간 $V \subset Y$가 주어졌을 때
$U = f^{-1}(V)$라 놓고 유도된 사상을 $g : U \to V$라 쓰자. 그러면
$D(\mathcal{O}_X)$의 $E$에 대하여
$(Rf_*E)|_V = Rg_*(E|_U)$이다.
\end{lemma}

\begin{proof}
$E$를 $\mathcal{O}_X$-가군의 K-단사 복합체
$\mathcal{I}^\bullet$로 나타내자. 그러면
$Rf_*(E) = f_*\mathcal{I}^\bullet$이고, Lemma
\ref{cohomology-lemma-restrict-K-injective-to-open}에 의해
$Rg_*(E|_U) = g_*(\mathcal{I}^\bullet|_U)$이다.
$X$ 위의 임의의 층 $\mathcal{F}$에 대하여
$(f_*\mathcal{F})|_V = g_*(\mathcal{F}|_U)$임은 명백하므로 결과가 따른다.
\end{proof}

\begin{lemma}
\label{cohomology-lemma-Leray-unbounded}
$f : X \to Y$를 환 달린 공간의 사상이라 하자. 그러면 함자
$D(\mathcal{O}_X) \to D(\Gamma(Y, \mathcal{O}_Y))$로서
$R\Gamma(Y, -) \circ Rf_* = R\Gamma(X, -)$이다.
더 일반적으로 열린집합 $V \subset Y$와 $U = f^{-1}(V)$에 대하여
$R\Gamma(U, -) = R\Gamma(V, -) \circ Rf_*$이다.
\end{lemma}

\begin{proof}
$Z$를 한 점으로 이루어지고
$\Gamma(Z, \mathcal{O}_Z) = \Gamma(Y, \mathcal{O}_Y)$인 환 달린 공간이라
하자. 구조층의 대역 단면 위에서 주어진 식별을 유도하는 환 달린 공간의
표준 사상 $Y \to Z$가 있다. 그러면
$D(\mathcal{O}_Z) = D(\Gamma(Y, \mathcal{O}_Y))$이다.
따라서 등식 $R\Gamma(Y, -) \circ Rf_* = R\Gamma(X, -)$는
$X \to Y \to Z$에 Lemma
\ref{cohomology-lemma-derived-pushforward-composition}을 적용하면 따른다.

\medskip\noindent
둘째인 더 일반적인 명제는 Lemma
\ref{cohomology-lemma-restrict-direct-image-open}을 적용한 뒤 첫째 명제에서 따른다.
\end{proof}

\begin{lemma}
\label{cohomology-lemma-unbounded-describe-higher-direct-images}
$f : (X, \mathcal{O}_X) \to (Y, \mathcal{O}_Y)$를 환 달린 공간의
사상이라 하고, $K$를 $D(\mathcal{O}_X)$의 대상이라 하자. 그러면
$H^i(Rf_*K)$는 준층
$$
V \mapsto H^i(f^{-1}(V), K) = H^i(V, Rf_*K)
$$
에 연관된 층이다.
\end{lemma}

\begin{proof}
등식 $H^i(f^{-1}(V), K) = H^i(V, Rf_*K)$는
Lemma \ref{cohomology-lemma-Leray-unbounded}의 둘째 명제에서 코호몰로지를 취하면
따른다. 층화에 관한 명제는 이제 Lemma
\ref{cohomology-lemma-sheafification-cohomology}에서 따른다.
\end{proof}

\begin{lemma}
\label{cohomology-lemma-modules-abelian-unbounded}
$X$를 환 달린 공간이라 하고, $K$를 $D(\mathcal{O}_X)$의 대상이라
하자. $D(\underline{\mathbf{Z}}_X)$에서의 그 상을 $K_{ab}$라 쓰자.
\begin{enumerate}
\item 임의의 열린집합 $U \subset X$에 대하여 표준 사상
$R\Gamma(U, K) \to R\Gamma(U, K_{ab})$가 있으며, 이는
$D(\textit{Ab})$ 안의 동형이다.
\item $f : X \to Y$를 환 달린 공간의 사상이라 하자.
표준 사상 $Rf_*K \to Rf_*(K_{ab})$가 있으며, 이는
$D(\underline{\mathbf{Z}}_Y)$ 안의 동형이다.
\end{enumerate}
\end{lemma}

\begin{proof}
사상은 다음과 같이 구성한다. $K$를 나타내는 K-단사 복합체
$\mathcal{I}^\bullet$를 고르자. $\mathcal{J}^\bullet$가 아벨군의
K-단사 복합체인 유사동형
$\mathcal{I}^\bullet \to \mathcal{J}^\bullet$를 고르자. 그러면 (1)의
사상은 $\Gamma(U, \mathcal{I}^\bullet) \to \Gamma(U, \mathcal{J}^\bullet)$로,
(2)의 사상은
$f_*\mathcal{I}^\bullet \to f_*\mathcal{J}^\bullet$로 주어진다.
이 사상들이 동형임을 보이려면 이들이 코호몰로지 군과 코호몰로지 층에서
동형을 유도함을 보이면 충분하다. Lemma
\ref{cohomology-lemma-unbounded-cohomology-of-open}와
\ref{cohomology-lemma-unbounded-describe-higher-direct-images}에 의해, 사상
$$
H^0(X, K) \longrightarrow H^0(X, K_{ab})
$$
이 동형임을 보이면 충분하다. 다음을 관찰하자.
$$
H^0(X, K) = \Hom_{D(\mathcal{O}_X)}(\mathcal{O}_X, K)
$$
다른 군도 마찬가지이다. $K$를 나타내는 임의의
$\mathcal{O}_X$-가군 복합체 $\mathcal{K}^\bullet$를 고르자.
유도범주를 국소화로 구성한 방식에 의해
$$
\Hom_{D(\mathcal{O}_X)}(\mathcal{O}_X, K) =
\colim_{s : \mathcal{F}^\bullet \to \mathcal{O}_X}
\Hom_{K(\mathcal{O}_X)}(\mathcal{F}^\bullet, \mathcal{K}^\bullet)
$$
이며, 여기서 여극한은 $\mathcal{O}_X$-가군 복합체의 유사동형 $s$ 전체에
걸쳐 취한다. 마찬가지로
$$
\Hom_{D(\underline{\mathbf{Z}}_X)}(\underline{\mathbf{Z}}_X, K) =
\colim_{s : \mathcal{G}^\bullet \to \underline{\mathbf{Z}}_X}
\Hom_{K(\underline{\mathbf{Z}}_X)}(\mathcal{G}^\bullet, \mathcal{K}^\bullet)
$$
이다. 다음으로 $\mathcal{G}^\bullet$가 평탄한
$\underline{\mathbf{Z}}_X$-가군의 위로 유계인 복합체인 유사동형
$s : \mathcal{G}^\bullet \to \underline{\mathbf{Z}}_X$들이 이 계에서
공종임을 관찰하자. (이는 Modules, Lemma
\ref{modules-lemma-module-quotient-flat}와 Derived Categories, Lemma
\ref{derived-lemma-subcategory-left-resolution}에서 따른다.
Section \ref{cohomology-section-flat}의 논의를 보라.)
따라서 원소 $\xi \in H^0(X, K_{ab})$를
$$
(s : \mathcal{G}^\bullet \to \underline{\mathbf{Z}}_X,
a : \mathcal{G}^\bullet \to \mathcal{K}^\bullet)
$$
라는 쌍으로 나타낼 수 있다. 여기서 $\mathcal{G}^\bullet$는 평탄한
$\underline{\mathbf{Z}}_X$-가군의 위로 유계인 복합체이다. 이 쌍을
$$
(\mathcal{G}^\bullet \otimes_{\underline{\mathbf{Z}}_X} \mathcal{O}_X
\to \mathcal{O}_X,
\mathcal{G}^\bullet  \otimes_{\underline{\mathbf{Z}}_X} \mathcal{O}_X
\to \mathcal{K}^\bullet)
$$
로 보내면 $H^0(X, K) \longrightarrow H^0(X, K_{ab})$의 역을 구성할
수 있다. 여기서 유의할 것은 첫째 화살표가 Lemma
\ref{cohomology-lemma-derived-tor-quasi-isomorphism-other-side}와
\ref{cohomology-lemma-bounded-flat-K-flat}에 의해 유사동형이라는 점뿐이다.
이 구성이 실제로 역임을 확인하는 자세한 과정은 생략한다.
\end{proof}

\begin{lemma}
\label{cohomology-lemma-adjoint-lower-shriek-restrict}
$(X, \mathcal{O}_X)$를 환 달린 공간이라 하고, $U \subset X$를 열린
부분집합이라 하자. 이에 대응하는 열린 몰입을
$j : (U, \mathcal{O}_U) \to (X, \mathcal{O}_X)$라 쓰자.
제한 함자 $D(\mathcal{O}_X) \to D(\mathcal{O}_U)$는 영에 의한 연장
$j_! : D(\mathcal{O}_U) \to D(\mathcal{O}_X)$의 오른쪽 수반이다.
\end{lemma}

\begin{proof}
이는 $j_!$와 $j^*$가 수반이고 완전하다는 사실에서 형식적으로 따른다
(따라서 $Lj_! = j_!$와 $Rj^* = j^*$가 존재한다). Derived Categories,
Lemma \ref{derived-lemma-derived-adjoint-functors}를 보라.
\end{proof}

\begin{lemma}
\label{cohomology-lemma-K-injective-flat}
$f : X \to Y$를 환 달린 공간의 평탄 사상이라 하자.
$\mathcal{I}^\bullet$가 $\mathcal{O}_X$-가군의 K-단사 복합체이면,
$f_*\mathcal{I}^\bullet$는 $\mathcal{O}_Y$-가군의 복합체로서
K-단사이다.
\end{lemma}

\begin{proof}
이는 Sheaves, Lemma
\ref{sheaves-lemma-adjoint-pullback-pushforward-modules}에 의해
$$
\Hom_{K(\mathcal{O}_Y)}(\mathcal{F}^\bullet, f_*\mathcal{I}^\bullet)
=
\Hom_{K(\mathcal{O}_X)}(f^*\mathcal{F}^\bullet, \mathcal{I}^\bullet)
$$
이고, $f$가 평탄하다고 가정했으므로 $f^*$가 완전하기 때문이다.
\end{proof}





\section{비유계 Mayer--Vietoris}
\label{cohomology-section-unbounded-mayer-vietoris}

\noindent
비유계 코호몰로지에도 Mayer--Vietoris 열이 있다.

\begin{lemma}
\label{cohomology-lemma-exact-sequence-lower-shriek}
$(X, \mathcal{O}_X)$를 환 달린 공간이라 하자.
$X = U \cup V$가 두 열린 부분공간의 합집합이라 하자.
$D(\mathcal{O}_X)$의 임의의 대상 $E$에 대하여
$D(\mathcal{O}_X)$ 안에 특수 삼각형
$$
j_{U \cap V!}E|_{U \cap V} \to
j_{U!}E|_U \oplus j_{V!}E|_V \to E \to 
j_{U \cap V!}E|_{U \cap V}[1]
$$
이 있다.
\end{lemma}

\begin{proof}
Section \ref{cohomology-section-properties-K-injective}에서 제한 함자와 영에 의한
연장 함자는 임의의 복합체에 해당 함자를 그대로 적용하여 계산됨을 보았다.
$E$를 나타내는 $\mathcal{O}_X$-가군 복합체를
$\mathcal{E}^\bullet$라 하자. 보조정리의 특수 삼각형은
$\mathcal{O}_X$-가군 복합체의 짧은 완전열
$$
0 \to j_{U \cap V!}\mathcal{E}^\bullet|_{U \cap V} \to
j_{U!}\mathcal{E}^\bullet|_U \oplus j_{V!}\mathcal{E}^\bullet|_V
\to \mathcal{E}^\bullet \to 0
$$
에 연관된 특수 삼각형이다(Derived Categories, Section
\ref{derived-section-canonical-delta-functor}, 특히 Lemma
\ref{derived-lemma-derived-canonical-delta-functor}에 의한다).
이 열이 완전함은 Sheaves, Lemma
\ref{sheaves-lemma-j-shriek-modules}를 써서 줄기에서 확인한다
(계산은 생략한다).
\end{proof}

\begin{lemma}
\label{cohomology-lemma-exact-sequence-j-star}
$(X, \mathcal{O}_X)$를 환 달린 공간이라 하자.
$X = U \cup V$가 두 열린 부분공간의 합집합이라 하자.
$D(\mathcal{O}_X)$의 임의의 대상 $E$에 대하여
$D(\mathcal{O}_X)$ 안에 특수 삼각형
$$
E \to 
Rj_{U, *}E|_U \oplus Rj_{V, *}E|_V \to
Rj_{U \cap V, *}E|_{U \cap V} \to
E[1]
$$
이 있다.
\end{lemma}

\begin{proof}
각 항 $\mathcal{I}^n$이 $\textit{Mod}(\mathcal{O}_X)$의 단사 대상인,
$E$를 나타내는 K-단사 복합체 $\mathcal{I}^\bullet$를 고르자.
Injectives, Theorem
\ref{injectives-theorem-K-injective-embedding-grothendieck}을 보라.
$\mathcal{I}^\bullet|U$ 역시 K-단사 복합체임을 이미 보았다
(Lemma \ref{cohomology-lemma-restrict-K-injective-to-open}). 따라서
$Rj_{U, *}E|_U$는 $j_{U, *}\mathcal{I}^\bullet|_U$로 나타내어진다.
$V$와 $U \cap V$에 대해서도 마찬가지이다. 그러므로 보조정리의
특수 삼각형은 복합체의 짧은 완전열
$$
0 \to
\mathcal{I}^\bullet \to
j_{U, *}\mathcal{I}^\bullet|_U \oplus j_{V, *}\mathcal{I}^\bullet|_V \to
j_{U \cap V, *}\mathcal{I}^\bullet|_{U \cap V} \to
0.
$$
에 연관된 특수 삼각형이다(Derived Categories, Section
\ref{derived-section-canonical-delta-functor}, 특히 Lemma
\ref{derived-lemma-derived-canonical-delta-functor}에 의한다).
임의의 열린집합 $W \subset X$와 임의의 $n$에 대하여
$$
0 \to
\mathcal{I}^n(W) \to
\mathcal{I}^n(W \cap U) \oplus \mathcal{I}^n(W \cap V) \to
\mathcal{I}^n(W \cap U \cap V) \to
0
$$
이 완전하기 때문에 이 열은 완전하다(Lemma
\ref{cohomology-lemma-mayer-vietoris}의 증명을 보라).
\end{proof}

\begin{lemma}
\label{cohomology-lemma-mayer-vietoris-hom}
$(X, \mathcal{O}_X)$를 환 달린 공간이라 하고, $X = U \cup V$를
$X$의 두 열린 부분공간의 합집합이라 하자.
$D(\mathcal{O}_X)$의 대상 $E$, $F$에 대하여 Mayer--Vietoris 열
$$
\xymatrix{
& \ldots \ar[r] & \Ext^{-1}(E_{U \cap V}, F_{U \cap V}) \ar[lld] \\
\Hom(E, F) \ar[r] &
\Hom(E_U, F_U) \oplus
\Hom(E_V, F_V) \ar[r] &
\Hom(E_{U \cap V}, F_{U \cap V})
}
$$
이 있다. 여기서 아래첨자는 해당 열린집합으로의 제한을 뜻하고,
$\Hom$과 $\Ext$는 해당 유도범주에서 취한다.
\end{lemma}

\begin{proof}
Lemma \ref{cohomology-lemma-exact-sequence-lower-shriek}의 특수 삼각형을 써서
$\Hom$들의 긴 완전열을 얻고(Derived Categories, Lemma
\ref{derived-lemma-representable-homological}에 의한다), Lemma
\ref{cohomology-lemma-adjoint-lower-shriek-restrict}에 따른 등식
$$
\Hom_{D(\mathcal{O}_X)}(j_{U!}E|_U, F) =
\Hom_{D(\mathcal{O}_U)}(E|_U, F|_U)
$$
을 사용한다.
\end{proof}

\begin{lemma}
\label{cohomology-lemma-unbounded-mayer-vietoris}
$(X, \mathcal{O}_X)$를 환 달린 공간이라 하자.
$X = U \cup V$가 두 열린 부분집합의 합집합이라고 하자.
$D(\mathcal{O}_X)$의 대상 $E$에 대하여 특수 삼각형
$$
R\Gamma(X, E) \to R\Gamma(U, E) \oplus R\Gamma(V, E) \to
R\Gamma(U \cap V, E) \to R\Gamma(X, E)[1]
$$
이 있고, 특히 긴 완전 코호몰로지 열
$$
\ldots \to
H^n(X, E) \to
H^n(U, E) \oplus H^0(V, E) \to
H^n(U \cap V, E) \to
H^{n + 1}(X, E) \to \ldots
$$
이 있다. 이 특수 삼각형과 긴 완전열의 구성은 $E$에 대해 함자적이다.
\end{lemma}

\begin{proof}
$E$를 나타내는 K-단사 복합체 $\mathcal{I}^\bullet$를 고르자.
모든 $n$에 대하여 $\mathcal{I}^n$이
$\textit{Mod}(\mathcal{O}_X)$의 단사 대상이라고 가정해도 된다.
Injectives, Theorem
\ref{injectives-theorem-K-injective-embedding-grothendieck}을 보라.
그러면 $R\Gamma(X, E)$는 $\Gamma(X, \mathcal{I}^\bullet)$로 계산된다.
Lemma \ref{cohomology-lemma-restrict-K-injective-to-open}에 의해 $U$, $V$, 그리고
$U \cap V$에 대해서도 마찬가지이다. 따라서 보조정리의 특수 삼각형은
복합체의 짧은 완전열
$$
0 \to
\mathcal{I}^\bullet(X) \to
\mathcal{I}^\bullet(U) \oplus \mathcal{I}^\bullet(V) \to
\mathcal{I}^\bullet(U \cap V) \to
0.
$$
에 연관된 특수 삼각형이다(Derived Categories, Section
\ref{derived-section-canonical-delta-functor}, 특히 Lemma
\ref{derived-lemma-derived-canonical-delta-functor}에 의한다).
Lemma \ref{cohomology-lemma-mayer-vietoris}의 증명에서 이것이 짧은 완전열임을
보았다. 마지막 명제는 Injectives, Theorem
\ref{injectives-theorem-K-injective-embedding-grothendieck}의 구성의
함자성에서 따른다.
\end{proof}

\begin{lemma}
\label{cohomology-lemma-unbounded-relative-mayer-vietoris}
$f : X \to Y$를 환 달린 공간의 사상이라 하자.
$X = U \cup V$가 두 열린 부분집합의 합집합이라고 하자.
$a = f|_U : U \to Y$, $b = f|_V : V \to Y$, 그리고
$c = f|_{U \cap V} : U \cap V \to Y$라 쓰자.
$D(\mathcal{O}_X)$의 모든 대상 $E$에 대하여 특수 삼각형
$$
Rf_*E \to
Ra_*(E|_U) \oplus Rb_*(E|_V) \to
Rc_*(E|_{U \cap V}) \to
Rf_*E[1]
$$
이 존재한다. 이 삼각형은 $E$에 대해 함자적이다.
\end{lemma}

\begin{proof}
$E$를 나타내는 K-단사 복합체 $\mathcal{I}^\bullet$를 고르자.
모든 $n$에 대하여 $\mathcal{I}^n$이
$\textit{Mod}(\mathcal{O}_X)$의 단사 대상이라고 가정해도 된다.
Injectives, Theorem
\ref{injectives-theorem-K-injective-embedding-grothendieck}을 보라.
그러면 $Rf_*E$는 $f_*\mathcal{I}^\bullet$로 계산된다.
Lemma \ref{cohomology-lemma-restrict-K-injective-to-open}에 의해 $U$, $V$, 그리고
$U \cap V$에 대해서도 마찬가지이다. 따라서 보조정리의 특수 삼각형은
복합체의 짧은 완전열
$$
0 \to
f_*\mathcal{I}^\bullet \to
a_*\mathcal{I}^\bullet|_U \oplus b_*\mathcal{I}^\bullet|_V \to
c_*\mathcal{I}^\bullet|_{U \cap V} \to
0.
$$
에 연관된 특수 삼각형이다(Derived Categories, Section
\ref{derived-section-canonical-delta-functor}, 특히 Lemma
\ref{derived-lemma-derived-canonical-delta-functor}에 의한다).
이는 Lemma \ref{cohomology-lemma-relative-mayer-vietoris}와
$\textit{Mod}(\mathcal{O}_X)$의 단사 대상 $\mathcal{I}$에 대해
$R^1f_*\mathcal{I} = 0$이라는 사실에 의해 복합체의 짧은 완전열이다.
마지막 명제는 Injectives, Theorem
\ref{injectives-theorem-K-injective-embedding-grothendieck}의 구성의
함자성에서 따른다.
\end{proof}

\begin{lemma}
\label{cohomology-lemma-pushforward-restriction}
$(X, \mathcal{O}_X)$를 환 달린 공간이라 하자. 열린 부분공간
$j : U \to X$를 택하자. $T \subset X$를 $U$에 포함되는 닫힌
부분집합이라 하자.
\begin{enumerate}
\item $E$가 코호몰로지 층들이 $T$ 위에 지지되는
$D(\mathcal{O}_X)$의 대상이면, $E \to Rj_*(E|_U)$는 동형이다.
\item $F$가 코호몰로지 층들이 $T$ 위에 지지되는
$D(\mathcal{O}_U)$의 대상이면, $j_!F \to Rj_*F$는 동형이다.
\end{enumerate}
\end{lemma}

\begin{proof}
$V = X \setminus T$ 및 $W = U \cap V$라 놓자. $X = U \cup V$는
$X$의 열린 덮개임에 유의하자. 열린 몰입을
$j_W : W \to V$라 쓰자. 코호몰로지 층들이 $T$ 위에 지지되는
$D(\mathcal{O}_X)$의 대상 $E$를 택하자.
Lemma \ref{cohomology-lemma-restrict-direct-image-open}에 의해
$(Rj_*E|_U)|_V = Rj_{W, *}(E|_W) = 0$이다. 실제로 가정에 의해
$E|_W = 0$이다. 한편 $Rj_*(E|_U)|_U = E|_U$이므로 (1)은 명백하다.
코호몰로지 층들이 $T$ 위에 지지되는 $D(\mathcal{O}_U)$의 대상
$F$를 택하자. Lemma \ref{cohomology-lemma-restrict-direct-image-open}에 의해
$(Rj_*F)|_V = Rj_{W, *}(F|_W) = 0$이다. 실제로 가정에 의해
$F|_W = 0$이다. 또한
$(j_!F)|_V = j_{W!}(F|_W) = 0$이다(첫째 등식은 영에 의한 연장의
정의에서 곧바로 따른다). $(Rj_*F)|_U = F$이고
$(j_!F)|_U = F$이므로 (2)가 성립한다.
\end{proof}

\begin{lemma}
\label{cohomology-lemma-mayer-vietoris-cup}
$(X, \mathcal{O}_X)$를 환 달린 공간이라 하고
$A = \Gamma(X, \mathcal{O}_X)$라 놓자.
$X = U \cup V$가 두 열린 부분집합의 합집합이라고 하자.
$D(\mathcal{O}_X)$의 대상 $K$와 $M$에 대하여 다음 특수 삼각형들의
사상이 있다.
$$
\xymatrix{
R\Gamma(X, K) \otimes_A^\mathbf{L} R\Gamma(X, M) \ar[r] \ar[d] &
R\Gamma(X, K \otimes_{\mathcal{O}_X}^\mathbf{L} M) \ar[d] \\
R\Gamma(X, K) \otimes_A^\mathbf{L}
(R\Gamma(U, M) \oplus R\Gamma(V, M)) \ar[r] \ar[d] &
R\Gamma(U, K \otimes_{\mathcal{O}_X}^\mathbf{L} M)
\oplus R\Gamma(V, K \otimes_{\mathcal{O}_X}^\mathbf{L} M)) \ar[d] \\
R\Gamma(X, K) \otimes_A^\mathbf{L} R\Gamma(U \cap V, M) \ar[r] \ar[d] &
R\Gamma(U \cap V, K \otimes_{\mathcal{O}_X}^\mathbf{L} M) \ar[d] \\
R\Gamma(X, K) \otimes_A^\mathbf{L} R\Gamma(X, M)[1] \ar[r] &
R\Gamma(X, K \otimes_{\mathcal{O}_X}^\mathbf{L} M)[1]
}
$$
여기서
\begin{enumerate}
\item 수평 화살표들은 컵곱으로 주어지고,
\item 오른쪽에는 $K \otimes_{\mathcal{O}_X}^\mathbf{L} M$에 대한
Lemma \ref{cohomology-lemma-unbounded-mayer-vietoris}의 특수 삼각형이 있으며,
\item 왼쪽에는 $M$에 대한 Lemma
\ref{cohomology-lemma-unbounded-mayer-vietoris}의 특수 삼각형에 완전 함자
$R\Gamma(X, K) \otimes_A^\mathbf{L} - $를 적용한 것이 있다.
\end{enumerate}
\end{lemma}

\begin{proof}
$R\Gamma(X, K)$를 나타내는 평탄한 $A$-가군의 K-평탄 복합체
$T^\bullet$를 고르자. More on Algebra, Lemma
\ref{more-algebra-lemma-K-flat-resolution}을 보라.
$T^\bullet \otimes_A \mathcal{O}_X$를 환 달린 공간의 사상
$(X, \mathcal{O}_X) \to (pt, A)$에 의한 $T^\bullet$의 당김이라 쓰자.
$D(\mathcal{O}_X)$ 안에 자연스러운 수반 사상
$\epsilon : T^\bullet \otimes_A \mathcal{O}_X \to K$가 있다.
$T^\bullet \otimes_A \mathcal{O}_X$는 평탄한 항들로 이루어진
$\mathcal{O}_X$-가군의 K-평탄 복합체이다. Lemma
\ref{cohomology-lemma-pullback-K-flat}와 Modules, Lemma
\ref{modules-lemma-pullback-flat}를 보라.
Lemma \ref{cohomology-lemma-factor-through-K-flat}에 의해 $\epsilon$을 나타내는
$\mathcal{O}_X$-가군 복합체의 사상
$$
T^\bullet \otimes_A \mathcal{O}_X \longrightarrow \mathcal{K}^\bullet
$$
을 찾을 수 있으며, 여기서 $\mathcal{K}^\bullet$는 평탄한 항들을 갖는
K-평탄 복합체이다. 구체적으로 $D(\mathcal{O}_X)$의 구성에 의해 먼저
$\epsilon$을 나타내는 $\mathcal{O}_X$-가군 복합체의 사상
$e : T^\bullet \otimes_A \mathcal{O}_X \to \mathcal{L}^\bullet$로
$\epsilon$을 나타낸 뒤, $e$에 이 보조정리를 적용할 수 있다.
$M$을 나타내고 각 항이 단사 $\mathcal{O}_X$-가군인 K-단사 복합체
$\mathcal{I}^\bullet$를 고르자. 마지막으로 각 항이 단사
$\mathcal{O}_X$-가군인 K-단사 복합체 안으로 가는 유사동형
$$
\text{Tot}(\mathcal{K}^\bullet \otimes_\mathcal{O} \mathcal{I}^\bullet)
\longrightarrow
\mathcal{J}^\bullet
$$
을 고르자. 이 화살표의 원천과 표적은
$D(\mathcal{O}_X)$ 안에서 $K \otimes_{\mathcal{O}_X}^\mathbf{L} M$을
나타냄에 유의하자. 이제 임의의 열린집합 $W \subset X$에 대하여
$A$-가군 복합체의 사상
$$
\text{Tot}(T^\bullet \otimes_A \mathcal{I}^\bullet(W))
\to
\text{Tot}(\mathcal{K}^\bullet(W) \otimes_A \mathcal{I}^\bullet(W))
\to
\mathcal{J}^\bullet(W)
$$
을 얻는다. 그 합성은 $D(A)$ 안의 사상
$$
R\Gamma(X, K) \otimes_A^\mathbf{L} R\Gamma(W, M)
\longrightarrow
R\Gamma(W, K \otimes_{\mathcal{O}_X}^\mathbf{L} M)
$$
을 나타낸다. 이 사상들은 제한 사상과 호환됨이 명백하다.
따라서 이제 $A$-가군 복합체의 다음 가환 도식을 생각할 수 있다.
$$
\xymatrix{
0 \ar[d] & 0 \ar[d] \\
\text{Tot}(T^\bullet \otimes_A \mathcal{I}^\bullet(X)) \ar[d] \ar[r] &
\mathcal{J}^\bullet(X) \ar[d] \\
\text{Tot}(T^\bullet \otimes_A
(\mathcal{I}^\bullet(U) \oplus \mathcal{I}^\bullet(V)) \ar[d] \ar[r] &
\mathcal{J}^\bullet(U) \oplus \mathcal{J}^\bullet(V) \ar[d] \\
\text{Tot}(T^\bullet \otimes_A \mathcal{I}^\bullet(U \cap V)) \ar[r] \ar[d] &
\mathcal{J}^\bullet(U \cap V) \ar[d] \\
0 & 0
}
$$
Lemma \ref{cohomology-lemma-mayer-vietoris}의 증명에 의해 열들은
$A$-가군 복합체의 완전열이다(여기서는 $T^\bullet$의 항들이 평탄한
$A$-가군이므로 $\text{Tot}(T^\bullet \otimes_A -)$가
$A$-가군 복합체의 짧은 완전열을 짧은 완전열로 보낸다는 사실도 쓴다).
Lemma \ref{cohomology-lemma-unbounded-mayer-vietoris}의 특수 삼각형들은 이 복합체의
짧은 완전열에 연관된 특수 삼각형이므로, 원하는 결과는
Derived Categories, Section
\ref{derived-section-canonical-delta-functor}에서 논의한
“연관된 특수 삼각형을 취하는 것”의 함자성에서 따른다.
\end{proof}












\section{닫힌 부분집합에 지지를 갖는 코호몰로지, II}
\label{cohomology-section-cohomology-support-bis}

\noindent
Section \ref{cohomology-section-cohomology-support}에서 시작한 논의를 계속한다.

\medskip\noindent
$(X, \mathcal{O}_X)$를 환 달린 공간이라 하고, $Z \subset X$를 닫힌
부분집합이라 하자. 이 상황에서
$\mathcal{F} \mapsto \Gamma_Z(X, \mathcal{F})$로 주어지는 함자
$\textit{Mod}(\mathcal{O}_X) \to \textit{Mod}(\mathcal{O}_X(X))$를
생각할 수 있다. Modules, Definition \ref{modules-definition-support}와
Modules, Lemma \ref{modules-lemma-support-section-closed}를 보라.
Section \ref{cohomology-section-unbounded}의 K-단사 분해를 사용하면 오른쪽 유도함자
$$
R\Gamma_Z(X, - ) : D(\mathcal{O}_X) \to D(\mathcal{O}_X(X))
$$
를 얻는다. $D(\mathcal{O}_X)$의 대상 $K$가 주어졌을 때
$Z$에 지지를 갖는 코호몰로지 가군을
$H^q_Z(X, K) = H^q(R\Gamma_Z(X, K))$라 쓴다. 이것이 나중에
Lemma \ref{cohomology-lemma-sections-support-abelian-unbounded}에서
Section \ref{cohomology-section-cohomology-support}의 구성과 일치함을 볼 것이다.

\medskip\noindent
$\mathcal{O}_X$-가군 $\mathcal{F}$에 대하여 규칙
$$
\mathcal{H}_Z(\mathcal{F})(U) =
\{s \in \mathcal{F}(U) \mid \text{Supp}(s) \subset U \cap Z\} =
\Gamma_{Z \cap U}(U, \mathcal{F}|_U)
$$
으로 정의되는 {\it $Z$에 지지를 갖는 단면들의 부분층}
$\mathcal{H}_Z(\mathcal{F})$를 생각할 수 있다.
Modules, Remark
\ref{modules-remark-sections-support-in-closed-modules}에서 논의했듯이,
$\mathcal{H}_Z(\mathcal{F})$를 $Z$ 위의
$\mathcal{O}_X|_Z$-가군으로 볼 수 있으며 함자
$$
\textit{Mod}(\mathcal{O}_X) \longrightarrow \textit{Mod}(\mathcal{O}_X|_Z),
\quad
\mathcal{F} \longmapsto \mathcal{H}_Z(\mathcal{F})
\text{를 }\mathcal{O}_X|_Z\text{-가군으로 }Z\text{ 위에서 봄}
$$
를 얻는다. 이 함자는 왼쪽 완전이지만 일반적으로 완전하지는 않다.
위와 똑같이 오른쪽 유도함자
$$
R\mathcal{H}_Z : D(\mathcal{O}_X) \longrightarrow D(\mathcal{O}_X|_Z)
$$
를 얻는다. $\mathcal{H}^q_Z(K) = H^q(R\mathcal{H}_Z(K))$로 놓으면
임의의 $\mathcal{O}_X$-가군 층 $\mathcal{F}$에 대하여
$\mathcal{H}^0_Z(\mathcal{F}) = \mathcal{H}_Z(\mathcal{F})$이다.

\begin{lemma}
\label{cohomology-lemma-cohomology-with-support-sheaf-on-support}
$(X, \mathcal{O}_X)$를 환 달린 공간이라 하고, $i : Z \to X$를 닫힌
부분집합의 포함이라 하자.
\begin{enumerate}
\item $R\mathcal{H}_Z : D(\mathcal{O}_X) \to D(\mathcal{O}_X|_Z)$는
$i_* : D(\mathcal{O}_X|_Z) \to D(\mathcal{O}_X)$의 오른쪽 수반이다.
\item $D(\mathcal{O}_X|_Z)$의 $K$에 대하여
$R\mathcal{H}_Z(i_*K) = K$이다.
\item $\mathcal{G}$를 $Z$ 위의 $\mathcal{O}_X|_Z$-가군 층이라 하자.
그러면 $p > 0$에 대하여
$\mathcal{H}^p_Z(i_*\mathcal{G}) = 0$이다.
\end{enumerate}
\end{lemma}

\begin{proof}
함자 $i_*$는 완전하므로 $i_* = Ri_* = Li_*$이다. 따라서 보조정리의
(1)은 Modules, Lemma
\ref{modules-lemma-adjoint-section-with-support}와 Derived Categories,
Lemma \ref{derived-lemma-derived-adjoint-functors}에서 따른다.
$K$가 (2)에서와 같다고 하자. $K$를
$\mathcal{O}_X|_Z$-가군의 K-단사 복합체
$\mathcal{I}^\bullet$로 나타낼 수 있다. Lemma
\ref{cohomology-lemma-K-injective-flat}에 의해 $i_*K$를 나타내는 복합체
$i_*\mathcal{I}^\bullet$는 $\mathcal{O}_X$-가군의 K-단사 복합체이다.
따라서 $R\mathcal{H}_Z(i_*K)$는
$\mathcal{H}_Z(i_*\mathcal{I}^\bullet) = \mathcal{I}^\bullet$로 계산되며,
이것이 (2)를 증명한다. (3)은 (2)의 특수한 경우이다.
\end{proof}

\noindent
$(X, \mathcal{O}_X)$를 환 달린 공간이라 하고, $Z \subset X$를 닫힌
부분집합이라 하자. 지지가 $Z$에 포함되는 $\mathcal{O}_X$-가군들의
범주는 모든 $\mathcal{O}_X$-가군의 범주의 Serre 부분범주이다.
Homology, Definition \ref{homology-definition-serre-subcategory}와
Modules, Lemma \ref{modules-lemma-support-section-closed}를 보라.
코호몰로지 층들이 $Z$ 위에 지지되는 복합체들로 이루어진
$D(\mathcal{O}_X)$의 엄밀히 충만하고 포화된 삼각부분범주를
$D_Z(\mathcal{O}_X)$라 쓴다. Derived Categories, Section
\ref{derived-section-triangulated-sub}을 보라.

\begin{lemma}
\label{cohomology-lemma-complexes-with-support-on-closed}
$(X, \mathcal{O}_X)$를 환 달린 공간이라 하고, $i : Z \to X$를 닫힌
부분집합의 포함이라 하자.
\begin{enumerate}
\item $D(\mathcal{O}_X|_Z)$의 $K$에 대하여
$i_*K$는 $D_Z(\mathcal{O}_X)$에 속한다.
\item 함자 $i_* : D(\mathcal{O}_X|_Z) \to D_Z(\mathcal{O}_X)$는
유사역
$i^{-1}|_{D_Z(\mathcal{O}_X)} = R\mathcal{H}_Z|_{D_Z(\mathcal{O}_X)}$를
갖는 동치이다.
\item 함자
$i_* \circ R\mathcal{H}_Z : D(\mathcal{O}_X) \to D_Z(\mathcal{O}_X)$는
포함 함자 $D_Z(\mathcal{O}_X) \to D(\mathcal{O}_X)$의 오른쪽 수반이다.
\end{enumerate}
\end{lemma}

\begin{proof}
(1)은 정의에서 곧바로 따른다. (3)은 (2)와 Lemma
\ref{cohomology-lemma-cohomology-with-support-sheaf-on-support}의 형식적 귀결이다.
나머지 증명에서는 (2)를 증명한다.

\medskip\noindent
$i$를 환 달린 공간의 사상
$i : (Z, \mathcal{O}_X|_Z) \to (X, \mathcal{O}_X)$로 생각하자.
$i^*$와 $i_*$가 수반 함자쌍임을 상기하자. $i$는 닫힌 몰입이므로
$i_*$는 완전하다. 또한 $i^{-1}\mathcal{O}_X = \mathcal{O}_X|_Z$가
$(Z, \mathcal{O}_X|_Z)$의 구조층이므로 $i^* = i^{-1}$는 완전하고,
$i^*i_* = i^{-1}i_*$는 항등 함자와 동형이다. Modules, Lemma
\ref{modules-lemma-exactness-pushforward-pullback}와
\ref{modules-lemma-i-star-exact}를 보라. 따라서
$i_* : D(\mathcal{O}_X|_Z) \to D_Z(\mathcal{O}_X)$는 완전충실하고
$i^{-1}$는 왼쪽 역을 정한다. 한편 $K$가 $D_Z(\mathcal{O}_X)$의
대상이라고 하고 수반 사상 $K \to i_*i^{-1}K$를 생각하자.
$i_*$와 $i^{-1}$의 완전성을 사용하면 이는 코호몰로지 층에서 수반 사상
$H^n(K) \to i_*i^{-1}H^n(K)$를 유도한다. 이 코호몰로지 층들은
$Z$ 위에 지지되므로 이 수반 사상들은 동형이다. 따라서
$i_* : D(\mathcal{O}_X|_Z) \to D_Z(\mathcal{O}_X)$가 동치임을 얻는다.

\medskip\noindent
증명을 마치려면 $K$가 $D_Z(\mathcal{O}_X)$의 대상일 때
$R\mathcal{H}_Z(K) = i^{-1}K$임을 보이면 충분하다.
방금 $K = i_*i^{-1}K$임을 증명했으므로 이를 사용할 수 있다.
그러면 Lemma \ref{cohomology-lemma-cohomology-with-support-sheaf-on-support}가
원하는 것을 말해 준다.
\end{proof}

\begin{lemma}
\label{cohomology-lemma-sections-with-support-K-injective}
$(X, \mathcal{O}_X)$를 환 달린 공간이라 하고, $i : Z \to X$를 닫힌
부분집합의 포함이라 하자. $\mathcal{I}^\bullet$가
$\mathcal{O}_X$-가군의 K-단사 복합체이면
$\mathcal{H}_Z(\mathcal{I}^\bullet)$는
$\mathcal{O}_X|_Z$-가군의 K-단사 복합체이다.
\end{lemma}

\begin{proof}
$i_* : \textit{Mod}(\mathcal{O}_X|_Z) \to \textit{Mod}(\mathcal{O}_X)$는
완전하고 $\mathcal{H}_Z$의 왼쪽 수반이므로(Modules, Lemma
\ref{modules-lemma-adjoint-section-with-support}), 이는
Derived Categories, Lemma
\ref{derived-lemma-adjoint-preserve-K-injectives}에서 따른다.
\end{proof}

\begin{lemma}
\label{cohomology-lemma-local-to-global-sections-with-support}
$(X, \mathcal{O}_X)$를 환 달린 공간이라 하고, $i : Z \to X$를 닫힌
부분집합의 포함이라 하자. 그러면 함자
$D(\mathcal{O}_X) \to D(\mathcal{O}_X(X))$로서
$R\Gamma(Z, - ) \circ R\mathcal{H}_Z = R\Gamma_Z(X, - )$이다.
\end{lemma}

\begin{proof}
이는 K-단사 분해를 사용한 오른쪽 유도함자의 구성, Lemma
\ref{cohomology-lemma-sections-with-support-K-injective}, 그리고
$\Gamma_Z(X, -) = \Gamma(Z, -) \circ \mathcal{H}_Z$라는 사실에서 따른다.
\end{proof}

\begin{lemma}
\label{cohomology-lemma-triangle-sections-with-support}
$(X, \mathcal{O}_X)$를 환 달린 공간이라 하고, $i : Z \to X$를 닫힌
부분집합의 포함이라 하자. $U = X \setminus Z$라 놓자.
$D(\mathcal{O}_X(X))$ 안에 특수 삼각형
$$
R\Gamma_Z(X, K) \to R\Gamma(X, K) \to R\Gamma(U, K) \to
R\Gamma_Z(X, K)[1]
$$
이 있으며, 이는 $D(\mathcal{O}_X)$의 $K$에 대해 함자적이다.
\end{lemma}

\begin{proof}
각 항이 단사 $\mathcal{O}_X$-가군인, $K$를 나타내는 K-단사 복합체
$\mathcal{I}^\bullet$를 고르자. Section
\ref{cohomology-section-unbounded}를 보라. $\mathcal{I}^\bullet|_U$는
$\mathcal{O}_U$-가군의 K-단사 복합체임을 상기하자. Lemma
\ref{cohomology-lemma-restrict-K-injective-to-open}을 보라. 따라서 특수 삼각형의
각 유도함자는 $\mathcal{I}^\bullet$에 그 바탕 함자를 적용하여 얻어진다.
그러므로 단사 $\mathcal{O}_X$-가군 $\mathcal{I}$에 대하여 짧은 완전열
$$
0 \to \Gamma_Z(X, \mathcal{I}) \to \Gamma(X, \mathcal{I})
\to \Gamma(U, \mathcal{I}) \to 0
$$
이 있음을 증명하면 충분하다. 이는 Lemma
\ref{cohomology-lemma-injective-restriction-surjective}와 정의에서 따른다.
\end{proof}

\begin{lemma}
\label{cohomology-lemma-triangle-sections-with-support-sheaves}
$(X, \mathcal{O}_X)$를 환 달린 공간이라 하고, $i : Z \to X$를 닫힌
부분집합의 포함이라 하자. 여집합의 포함을
$j : U = X \setminus Z \to X$라 쓰자. $D(\mathcal{O}_X)$ 안에
특수 삼각형
$$
i_*R\mathcal{H}_Z(K) \to K \to Rj_*(K|_U) \to
i_*R\mathcal{H}_Z(K)[1]
$$
이 있으며, 이는 $D(\mathcal{O}_X)$의 $K$에 대해 함자적이다.
\end{lemma}

\begin{proof}
각 항이 단사 $\mathcal{O}_X$-가군인, $K$를 나타내는 K-단사 복합체
$\mathcal{I}^\bullet$를 고르자. Section
\ref{cohomology-section-unbounded}를 보라. $\mathcal{I}^\bullet|_U$는
$\mathcal{O}_U$-가군의 K-단사 복합체임을 상기하자. Lemma
\ref{cohomology-lemma-restrict-K-injective-to-open}을 보라. 따라서 특수 삼각형의
각 유도함자는 $\mathcal{I}^\bullet$에 그 바탕 함자를 적용하여 얻어진다.
그러므로 단사 $\mathcal{O}_X$-가군 $\mathcal{I}$에 대하여 짧은 완전열
$$
0 \to i_*\mathcal{H}_Z(\mathcal{I}) \to \mathcal{I}
\to j_*(\mathcal{I}|_U) \to 0
$$
이 있음을 증명하면 충분하다. 이는 Lemma
\ref{cohomology-lemma-injective-restriction-surjective}와 정의에서 따른다.
\end{proof}

\begin{lemma}
\label{cohomology-lemma-sections-support-in-closed-disjoint-open}
$(X, \mathcal{O}_X)$를 환 달린 공간이라 하고, $Z \subset X$를 닫힌
부분집합이라 하자. $j : U \to X$를
$U \cap Z = \emptyset$인 열린 부분집합의 포함이라 하자. 그러면
$D(\mathcal{O}_U)$의 모든 $K$에 대하여
$R\mathcal{H}_Z(Rj_*K) = 0$이다.
\end{lemma}

\begin{proof}
$K$를 나타내는 $\mathcal{O}_U$-가군의 K-단사 복합체
$\mathcal{I}^\bullet$를 고르자. 그러면
$j_*\mathcal{I}^\bullet$는 $Rj_*K$를 나타낸다. Lemma
\ref{cohomology-lemma-K-injective-flat}에 의해 복합체
$j_*\mathcal{I}^\bullet$는 $\mathcal{O}_X$-가군의 K-단사 복합체이다.
따라서 $\mathcal{H}_Z(j_*\mathcal{I}^\bullet)$는
$R\mathcal{H}_Z(Rj_*K)$를 나타낸다. 그러므로 $U$ 위의 임의의 아벨 층
$\mathcal{G}$에 대하여 $\mathcal{H}_Z(j_*\mathcal{G}) = 0$임을
보이면 충분하다. 즉 어떤 열린집합 $W$ 위에서
$W \cap Z$에 지지되는 $j_*\mathcal{G}$의 단면 $s$가 영임을 보여야 한다.
지지 조건은 $s|_{W \setminus W \cap Z} = 0$이라는 뜻이다. 그런데
$j_*\mathcal{G}(W) =
\mathcal{G}(U \cap W) = j_*\mathcal{G}(W \setminus W \cap Z)$이므로
원하는 대로 $s$는 영이다.
\end{proof}

\begin{lemma}
\label{cohomology-lemma-sections-support-abelian-unbounded}
$(X, \mathcal{O}_X)$를 환 달린 공간이라 하고, $Z \subset X$를 닫힌
부분집합이라 하자. $K$를 $D(\mathcal{O}_X)$의 대상이라 하고,
$D(\underline{\mathbf{Z}}_X)$에서의 그 상을 $K_{ab}$라 쓰자.
\begin{enumerate}
\item 표준 사상 $R\Gamma_Z(X, K) \to R\Gamma_Z(X, K_{ab})$가 있으며,
이는 $D(\textit{Ab})$ 안의 동형이다.
\item 표준 사상 $R\mathcal{H}_Z(K) \to R\mathcal{H}_Z(K_{ab})$가 있으며,
이는 $D(\underline{\mathbf{Z}}_Z)$ 안의 동형이다.
\end{enumerate}
\end{lemma}

\begin{proof}
(1)의 증명. 사상은 다음과 같이 구성한다. $K$를 나타내는
$\mathcal{O}_X$-가군의 K-단사 복합체 $\mathcal{I}^\bullet$를 고르자.
$\mathcal{J}^\bullet$가 아벨군의 K-단사 복합체인 유사동형
$\mathcal{I}^\bullet \to \mathcal{J}^\bullet$를 고르자. 그러면 (1)의
사상은
$$
\Gamma_Z(X, \mathcal{I}^\bullet) \to \Gamma_Z(X, \mathcal{J}^\bullet)
$$
로 주어지며, 이는 $\Gamma_Z$가 아벨 층 위의 함자라는 사실로 정해진다.
간단히 확인하면, 이 사상과 Lemma
\ref{cohomology-lemma-modules-abelian-unbounded}의 표준 사상들을 결합한 것이
Lemma \ref{cohomology-lemma-triangle-sections-with-support}의 특수 삼각형들의 사상
$$
\xymatrix{
R\Gamma_Z(X, K) \ar[r] \ar[d] &
R\Gamma(X, K) \ar[r] \ar[d] &
R\Gamma(U, K) \ar[d] \\
R\Gamma_Z(X, K_{ab}) \ar[r] &
R\Gamma(X, K_{ab}) \ar[r] &
R\Gamma(U, K_{ab})
}
$$
에 들어감을 알 수 있다. 인용한 보조정리에 의해 세 화살표 중 둘이
동형이므로, Derived Categories, Lemma
\ref{derived-lemma-third-isomorphism-triangle}에서 결론이 따른다.

\medskip\noindent
(2)의 증명은 생략한다. 힌트: 특수 삼각형에 Lemma
\ref{cohomology-lemma-triangle-sections-with-support-sheaves}를 써서 같은 논증을
사용하라.
\end{proof}

\begin{remark}
\label{cohomology-remark-support-cup-product}
$(X, \mathcal{O}_X)$를 환 달린 공간이라 하고, $i : Z \to X$를 닫힌
부분집합의 포함이라 하자. $D(\mathcal{O}_X)$의 $K$와 $M$이 주어지면
$D(\mathcal{O}_X|_Z)$ 안에 표준 사상
$$
K|_Z \otimes_{\mathcal{O}_X|_Z}^\mathbf{L} R\mathcal{H}_Z(M)
\longrightarrow
R\mathcal{H}_Z(K \otimes_{\mathcal{O}_X}^\mathbf{L} M)
$$
이 있다. 여기서 $K|_Z = i^{-1}K$는 $K$를 $Z$로 제한하여
$D(\mathcal{O}_X|_Z)$의 대상으로 본 것이다. Lemma
\ref{cohomology-lemma-cohomology-with-support-sheaf-on-support}의 $i_*$와
$R\mathcal{H}_Z$의 수반성에 의해 이 사상을 구성하려면 표준 사상
$$
i_*\left(K|_Z \otimes_{\mathcal{O}_X|_Z}^\mathbf{L} R\mathcal{H}_Z(M)\right)
\longrightarrow
K \otimes_{\mathcal{O}_X}^\mathbf{L} M
$$
을 만들면 충분하다. 이를 만들기 위해 $M$을 나타내는
$\mathcal{O}_X$-가군의 K-단사 복합체 $\mathcal{I}^\bullet$와 $K$를
나타내는 $\mathcal{O}_X$-가군의 K-평탄 복합체
$\mathcal{K}^\bullet$를 고르자.
$\mathcal{K}^\bullet|_Z$는 $K|_Z$를 나타내는
$\mathcal{O}_X|_Z$-가군의 K-평탄 복합체이다. Lemma
\ref{cohomology-lemma-pullback-K-flat}을 보라. 따라서
$\mathcal{O}_X$-가군 복합체의 사상
$$
i_*\text{Tot}\left(
\mathcal{K}^\bullet|_Z \otimes_{\mathcal{O}_X|_Z}
\mathcal{H}_Z(\mathcal{I}^\bullet)\right)
\longrightarrow
\text{Tot}(\mathcal{K}^\bullet \otimes_{\mathcal{O}_X} \mathcal{I}^\bullet)
$$
을 만들면 된다. 이를 위해서는 사상
$$
i_*(\mathcal{K}^a|_Z \otimes_{\mathcal{O}_X|_Z}
\mathcal{H}_Z(\mathcal{I}^b))
\longrightarrow
\mathcal{K}^a \otimes_{\mathcal{O}_X} \mathcal{I}^b
$$
들을 만들면 충분하다. 예를 들어 줄기에서 보면, 이 식의 왼쪽 변은
$\mathcal{K}^a \otimes_{\mathcal{O}_X} i_*\mathcal{H}_Z(\mathcal{I}^b)$와
같다. 이제 포함 $\mathcal{H}_Z(\mathcal{I}^b) \to \mathcal{I}^b$를
써서 원하는 사상을 얻는다.
\end{remark}

\begin{remark}
\label{cohomology-remark-support-cup-product-global}
Remark \ref{cohomology-remark-support-cup-product}의 표기 아래 표준 컵곱
\begin{align*}
H^a(X, K) \times H^b_Z(X, M)
& =
H^a(X, K) \times H^b(Z, R\mathcal{H}_Z(M)) \\
& \to
H^a(Z, K|_Z) \times H^b(Z, R\mathcal{H}_Z(M)) \\
& \to
H^{a + b}(Z, K|_Z \otimes_{\mathcal{O}_X|_Z}^\mathbf{L} R\mathcal{H}_Z(M)) \\
& \to
H^{a + b}(Z, R\mathcal{H}_Z(K \otimes_{\mathcal{O}_X}^\mathbf{L} M)) \\
& =
H^{a + b}_Z(X, K \otimes_{\mathcal{O}_X}^\mathbf{L} M)
\end{align*}
을 얻는다. 여기서 등호들은 Lemma
\ref{cohomology-lemma-local-to-global-sections-with-support}로 주어지고, 첫째
화살표는 $Z$로의 제한, 둘째 화살표는 컵곱(Section
\ref{cohomology-section-cup-product}), 셋째 화살표는 Remark
\ref{cohomology-remark-support-cup-product}의 사상이다.
\end{remark}

\begin{lemma}
\label{cohomology-lemma-support-cup-product}
Remark \ref{cohomology-remark-support-cup-product}의 표기 아래 도식
$$
\xymatrix{
H^i(X, K) \times H^j_Z(X, M) \ar[r] \ar[d] &
H^{i + j}_Z(X, K \otimes_{\mathcal{O}_X}^\mathbf{L} M) \ar[d] \\
H^i(X, K) \times H^j(X, M) \ar[r] &
H^{i + j}(X, K \otimes_{\mathcal{O}_X}^\mathbf{L} M)
}
$$
은 가환한다. 여기서 위쪽 수평 화살표는 Remark
\ref{cohomology-remark-support-cup-product-global}의 컵곱이다.
\end{lemma}

\begin{proof}
생략한다.
\end{proof}

\begin{remark}
\label{cohomology-remark-support-functorial}
$f : (X', \mathcal{O}_{X'}) \to (X, \mathcal{O}_X)$를 환 달린 공간의
사상이라 하자. $Z \subset X$를 닫힌 부분집합이라 하고
$Z' = f^{-1}(Z)$라 놓자. 유도된 환 달린 공간의 사상을
$f|_{Z'} : (Z', \mathcal{O}_{X'}|_{Z'}) \to (Z, \mathcal{O}_X|Z)$라
쓰자. 임의의 $D(\mathcal{O}_X)$의 $K$에 대하여
$D(\mathcal{O}_{X'}|_{Z'})$ 안에 표준 사상
$$
L(f|_{Z'})^*R\mathcal{H}_Z(K) \longrightarrow R\mathcal{H}_{Z'}(Lf^*K)
$$
이 있다. 포함 사상들을 $i : Z \to X$와 $i' : Z' \to X'$라 쓰자.
$i'$에 적용한 Lemma \ref{cohomology-lemma-complexes-with-support-on-closed}의
(2)에 의해, 이는 $D_{Z'}(\mathcal{O}_{X'})$ 안의 사상
$$
i'_* L(f|_{Z'})^* R\mathcal{H}_Z(K)
\longrightarrow
i'_*R\mathcal{H}_{Z'}(Lf^*K)
$$
을 주는 것과 같다. Remark \ref{cohomology-remark-base-change}의 함자 사상
$Lf^* \circ i_* \to i'_* \circ L(f|_{Z'})^*$는 이 경우 동형이다
($i_*$와 $i'_*$가 완전하고
$\mathcal{O}_{Z, z} = \mathcal{O}_{X, z}$이며 $Z'$에 대해서도
마찬가지임을 이용하여 줄기에서 확인하면 된다). 따라서 다음 도식의
위쪽 수평 화살표를 구성하면 충분하다.
$$
\xymatrix{
Lf^* i_* R\mathcal{H}_Z(K) \ar[rr] \ar[rd] & &
i'_* R\mathcal{H}_{Z'}(Lf^*K) \ar[ld] \\
& Lf^*K
}
$$
복합체 $Lf^* i_* R\mathcal{H}_Z(K)$는 $Z'$ 위에 지지된다.
남동쪽 화살표는 수반 사상
$i_*R\mathcal{H}_Z(K) \to K$에서 온다(Lemma
\ref{cohomology-lemma-cohomology-with-support-sheaf-on-support}).
수반 사상 $i'_* R\mathcal{H}_{Z'}(Lf^*K) \to Lf^*K$는 Lemma
\ref{cohomology-lemma-complexes-with-support-on-closed}의 (3)에 의해 보편적이므로,
남동쪽 화살표는 남서쪽 화살표를 거쳐 유일하게 인수분해된다. 이로써
원하는 화살표를 얻는다.
\end{remark}

\begin{lemma}
\label{cohomology-lemma-support-functorial}
Remark \ref{cohomology-remark-support-functorial}의 표기와 가정 아래 도식
$$
\xymatrix{
H^p_Z(X, K) \ar[r] \ar[d] & H^p_{Z'}(X', Lf^*K) \ar[d] \\
H^p(X, K) \ar[r] & H^p(X', Lf^*K)
}
$$
은 가환한다. 여기서 위쪽 수평 화살표는 식별
$H^p_Z(X, K) = H^p(Z, R\mathcal{H}_Z(K))$와
$H^p_{Z'}(X', Lf^*K) = H^p(Z', R\mathcal{H}_{Z'}(K'))$,
당김 사상
$H^p(Z, R\mathcal{H}_Z(K)) \to H^p(Z', L(f|_{Z'})^*R\mathcal{H}_Z(K))$,
그리고 Remark \ref{cohomology-remark-support-functorial}에서 구성한 사상으로부터
온다.
\end{lemma}

\begin{proof}
생략한다. 힌트:
$H^p(Z, R\mathcal{H}_Z(K)) = H^p(X, i_*R\mathcal{H}_Z(K))$와
$R\mathcal{H}_{Z'}(Lf^*K)$에 대한 유사한 등식을 사용하면, 당김 사상의
함자성과 Remark \ref{cohomology-remark-support-functorial}의 사상을 정의하는 데
사용한 가환 도식으로부터 결과가 따른다.
\end{proof}






\section{역계와 코호몰로지, I}
\label{cohomology-section-inverse-systems}

\noindent
$A$를 환이라 하고, $I \subset A$를 아이디얼이라 하자.
$A/I^n$-가군 층들의 역계에 관한 몇 가지 결과를 증명한다.

\begin{lemma}
\label{cohomology-lemma-ML-general}
$I$를 환 $A$의 아이디얼이라 하고, $X$를 위상공간이라 하자.
$$
\ldots \to \mathcal{F}_3 \to \mathcal{F}_2 \to \mathcal{F}_1
$$
을 $X$ 위의 $A$-가군 층들의 역계라 하고,
$\mathcal{F}_n = \mathcal{F}_{n + 1}/I^n\mathcal{F}_{n + 1}$이라 하자.
$p \geq 0$라 하자. 등급
$\bigoplus_{n \geq 0} I^n/I^{n + 1}$-가군으로서
$$
\bigoplus\nolimits_{n \geq 0} H^{p + 1}(X, I^n\mathcal{F}_{n + 1})
$$
이 오름 사슬 조건을 만족한다고 가정하자. 그러면 역계
$M_n = H^p(X, \mathcal{F}_n)$은 Mittag--Leffler 조건을
만족한다\footnote{실제로 모든 $n \geq c$에 대하여
$\Im(M_n \to M_{n - c})$가 안정된 상이 되게 하는 $c \geq 0$가 존재한다.}.
\end{lemma}

\begin{proof}
$N_n = H^{p + 1}(X, I^n\mathcal{F}_{n + 1})$로 놓고,
짧은 완전열
$0 \to I^n\mathcal{F}_{n + 1} \to \mathcal{F}_{n + 1} \to \mathcal{F}_n \to 0$
에서 오는 코호몰로지의 경계 사상을
$\delta_n : M_n \to N_n$이라 하자. 그러면
$\bigoplus \Im(\delta_n) \subset \bigoplus N_n$은 등급 부분가군이다.
실제로 $s \in M_n$이고 $f \in I^m$이면 가환 도식
$$
\xymatrix{
0 \ar[r] &
I^n\mathcal{F}_{n + 1} \ar[d]_f \ar[r] &
\mathcal{F}_{n + 1} \ar[d]_f \ar[r] &
\mathcal{F}_n \ar[d]_f \ar[r] & 0 \\
0 \ar[r] &
I^{n + m}\mathcal{F}_{n + m + 1} \ar[r] &
\mathcal{F}_{n + m + 1} \ar[r] &
\mathcal{F}_{n + m} \ar[r] & 0
}
$$
을 얻는다. 가운데 수직 사상은 $\mathcal{F}_{n + 1}$의 국소 단면을
$\mathcal{F}_{n + m + 1}$의 단면으로 올린 뒤 $f$를 곱하여 주어지며,
다른 수직 화살표들도 마찬가지이다. 따라서
$\delta_{n + m}(fs) = f \delta_n(s)$이다.
가정에 의해 $\delta_{n_j}(s_j)$들이
$\bigoplus \Im(\delta_n)$을 등급 가군으로 생성하도록 하는
$s_j \in M_{n_j}$, $j = 1, \ldots, N$을 찾을 수 있다.
$n > c = \max(n_j)$라 하고, $s \in M_n$이라 하자. 그러면
$\delta_n(s) = \sum f_j \delta_{n_j}(s_j)$가 되게 하는
$f_j \in I^{n - n_j}$를 찾을 수 있다. 따라서
$\delta(s - \sum f_j s_j) = 0$이다. 즉 $M_n$ 안에서
$s - \sum f_js_j$로 가는 $s' \in M_{n + 1}$을 찾을 수 있다. 그러므로
$$
\Im(M_{n + 1} \to M_{n - c}) = \Im(M_n \to M_{n - c})
$$
이다. 실제로 원소 $f_js_j$들은 $M_{n - c}$에서 영으로 간다.
이것이 보조정리를 증명한다.
\end{proof}

\begin{lemma}
\label{cohomology-lemma-ML-general-better}
$I$를 환 $A$의 아이디얼이라 하고, $X$를 위상공간이라 하자.
$$
\ldots \to \mathcal{F}_3 \to \mathcal{F}_2 \to \mathcal{F}_1
$$
을 $X$ 위의 $A$-가군의 역계라 하고,
$\mathcal{F}_n = \mathcal{F}_{n + 1}/I^n\mathcal{F}_{n + 1}$이라 하자.
$p \geq 0$라 하자. $n$이 주어졌을 때
$$
N_n =
\bigcap\nolimits_{m \geq n}
\Im\left(
H^{p + 1}(X, I^n\mathcal{F}_{m + 1}) \to H^{p + 1}(X, I^n\mathcal{F}_{n + 1})
\right)
$$
으로 정의하자. 등급 $\bigoplus_{n \geq 0} I^n/I^{n + 1}$-가군으로서
$\bigoplus N_n$이 오름 사슬 조건을 만족하면 역계
$M_n = H^p(X, \mathcal{F}_n)$은 Mittag--Leffler 조건을
만족한다\footnote{실제로 모든 $n \geq c$에 대하여
$\Im(M_n \to M_{n - c})$가 안정된 상이 되게 하는 $c \geq 0$가 존재한다.}.
\end{lemma}

\begin{proof}
증명은 Lemma \ref{cohomology-lemma-ML-general}의 증명과 완전히 같다.
실제로 $\bigoplus N_n$이
$\bigoplus H^{p + 1}(X, I^n\mathcal{F}_{n + 1})$의 등급
$\bigoplus_{n \geq 0} I^n/I^{n + 1}$-부분가군이고, 경계 사상
$\delta_n : M_n \to H^{p + 1}(X, I^n\mathcal{F}_{n + 1})$의 상이
$N_n$에 포함됨을 보이기만 하면 거기서 제시한 논증으로 결과가 따른다.

\medskip\noindent
$\xi \in N_n$이고 $f \in I^k$라고 하자. $m \gg n + k$를 고르고,
$\xi$로 가는
$\xi' \in H^{p + 1}(X, I^n\mathcal{F}_{m + 1})$를 고르자.
Lemma \ref{cohomology-lemma-ML-general}의 증명에서와 같이 구성한 도식
$$
\xymatrix{
0 \ar[r] &
I^n\mathcal{F}_{m + 1} \ar[d]_f \ar[r] &
\mathcal{F}_{m + 1} \ar[d]_f \ar[r] &
\mathcal{F}_n \ar[d]_f \ar[r] & 0 \\
0 \ar[r] &
I^{n + k}\mathcal{F}_{m + 1} \ar[r] &
\mathcal{F}_{m + 1} \ar[r] &
\mathcal{F}_{n + k} \ar[r] & 0
}
$$
을 생각하자. 코호몰로지에서 유도된 사상을 얻고,
$f \xi' \in H^{p + 1}(X, I^{n + k}\mathcal{F}_{m + 1})$가
$f \xi$로 감을 알 수 있다. 이것은 모든 $m \gg n + k$에 대하여
성립하므로 원하는 대로 $f\xi$는 $N_{n + k}$에 속한다.

\medskip\noindent
경계 사상 $\delta_n$의 상이 $N_n$에 포함됨을 보기 위해
$m \geq n$에 대하여 도식
$$
\xymatrix{
0 \ar[r] &
I^n\mathcal{F}_{m + 1} \ar[d] \ar[r] &
\mathcal{F}_{m + 1} \ar[d] \ar[r] &
\mathcal{F}_n \ar[d] \ar[r] & 0 \\
0 \ar[r] &
I^n\mathcal{F}_{n + 1} \ar[r] &
\mathcal{F}_{n + 1} \ar[r] &
\mathcal{F}_n \ar[r] & 0
}
$$
들을 생각한다. 코호몰로지에서 유도된 사상들을 보면 결론이 따른다.
\end{proof}

\begin{lemma}
\label{cohomology-lemma-topology-I-adic-general}
$I$를 환 $A$의 아이디얼이라 하고, $X$를 위상공간이라 하자.
$$
\ldots \to \mathcal{F}_3 \to \mathcal{F}_2 \to \mathcal{F}_1
$$
을 $X$ 위의 $A$-가군 층들의 역계라 하고,
$\mathcal{F}_n = \mathcal{F}_{n + 1}/I^n\mathcal{F}_{n + 1}$이라 하자.
$p \geq 0$라 하자. 등급
$\bigoplus_{n \geq 0} I^n/I^{n + 1}$-가군으로서
$$
\bigoplus\nolimits_{n \geq 0} H^p(X, I^n\mathcal{F}_{n + 1})
$$
이 오름 사슬 조건을 만족한다고 가정하자. 그러면
$M = \lim H^p(X, \mathcal{F}_n)$ 위의 극한 위상은 $I$-진 위상이다.
\end{lemma}

\begin{proof}
$n \geq 1$에 대하여
$F^n = \Ker(M \to H^p(X, \mathcal{F}_n))$으로 놓고 $F^0 = M$으로
놓자. $I F^n \subset F^{n + 1}$임을 관찰하자. 특히
$I^n M \subset F^n$이다. 따라서 $I$-진 위상은 극한 위상보다 더
미세하다. 역을 위해, 주어진 $n$에 대하여
$F^m \subset I^nM$이 되게 하는 $m \geq n$이 존재함을 보일
것이다\footnote{실제로 모든 $n$에 대하여
$F^{n + c} \subset I^nM$이 되게 하는 $c \geq 0$가 존재한다.}.
단사 사상
$$
F^n/F^{n + 1} \longrightarrow H^p(X, \mathcal{F}_{n + 1})
$$
이 있고, 그 상은
$H^p(X, I^n\mathcal{F}_{n + 1}) \to H^p(X, \mathcal{F}_{n + 1})$의
상에 포함된다. $F^n/F^{n + 1}$의 역상을
$$
E_n \subset H^p(X, I^n\mathcal{F}_{n + 1})
$$
이라 쓰자. 그러면 $\bigoplus E_n$은
$\bigoplus H^p(X, I^n\mathcal{F}_{n + 1})$의 등급
$\bigoplus I^n/I^{n + 1}$-부분가군이고,
$\bigoplus E_n \to \bigoplus F^n/F^{n + 1}$은 등급 가군의
준동형이다. 자세한 내용은 생략한다. 가정에 의해 $\bigoplus E_n$은
$\bigoplus I^n/I^{n + 1}$ 위에서 유한 개의 동차 원소들로 생성된다.
$E_n \to F^n/F^{n + 1}$가 전사이므로
$\bigoplus F^n/F^{n + 1}$에 대해서도 같은 것이 성립한다.
따라서 $r$과 $c_1, \ldots, c_r \geq 0$, 그리고 그 상들이
$\bigoplus F^n/F^{n + 1}$을 생성하는 $a_i \in F^{c_i}$를 찾을 수 있다.
$c = \max(c_i)$로 놓자.

\medskip\noindent
$n \geq c$에 대하여 $I F^n = F^{n + 1}$이라고 주장한다.
이 주장은 원하는 대로 $F^{n + c} = I^nF^c \subset I^nM$임을 보인다.
주장을 증명하기 위해 $a \in F^{n + 1}$이라고 하자.
$F^{n + 1}/F^{n + 2}$에서의 $a$의 상은 $a_i$들의 선형결합이다.
따라서 어떤 $f_i \in I^{n + 1 - c_i}$에 대하여
$a  - \sum f_i a_i \in F^{n + 2}$이다.
$n \geq c_i$이므로
$I^{n + 1 - c_i} = I \cdot I^{n - c_i}$이고, 따라서
$g_{i, j} \in I$이고 $h_{i, j}a_i \in F^n$이도록
$f_i = \sum g_{i, j} h_{i, j}$로 쓸 수 있다. 그러므로
$F^{n + 1} = F^{n + 2} + IF^n$이다.
간단한 귀납법으로 모든 $e > 0$에 대하여
$F^{n + 1} = F^{n + e} + IF^n$을 얻는다. 따라서 $IF^n$은
$F^{n + 1}$에서 조밀하다. $I$의 생성자 $k_1, \ldots, k_r$를
고르고 연속 사상
$$
u : (F^n)^{\oplus r} \longrightarrow F^{n + 1},\quad
(x_1, \ldots, x_r) \mapsto \sum k_i x_i
$$
을 극한 위상에서 생각하자. 위의 논의에 의해 모든 $m \geq n$에 대하여
$u$ 아래 $(F^m)^{\oplus r}$의 상은 $F^{m + 1}$에서 조밀하다.
열린 사상 보조정리(More on Algebra, Lemma
\ref{more-algebra-lemma-open-mapping})에 의해 $u$는 열린 사상이다.
따라서 $u$는 전사이다. 그러므로 $n \geq c$에 대하여
$IF^n = F^{n + 1}$이다. 이것으로 증명을 마친다.
\end{proof}

\begin{lemma}
\label{cohomology-lemma-topology-I-adic-general-better}
$I$를 환 $A$의 아이디얼이라 하고, $X$를 위상공간이라 하자.
$$
\ldots \to \mathcal{F}_3 \to \mathcal{F}_2 \to \mathcal{F}_1
$$
을 $X$ 위의 $A$-가군 층들의 역계라 하고,
$\mathcal{F}_n = \mathcal{F}_{n + 1}/I^n\mathcal{F}_{n + 1}$이라 하자.
$p \geq 0$라 하자. $n$이 주어졌을 때
$$
N_n =
\bigcap\nolimits_{m \geq n}
\Im\left(
H^p(X, I^n\mathcal{F}_{m + 1}) \to H^p(X, I^n\mathcal{F}_{n + 1})
\right)
$$
으로 정의하자. 등급 $\bigoplus_{n \geq 0} I^n/I^{n + 1}$-가군으로서
$\bigoplus N_n$이 오름 사슬 조건을 만족하면
$M = \lim H^p(X, \mathcal{F}_n)$ 위의 극한 위상은 $I$-진 위상이다.
\end{lemma}

\begin{proof}
증명은 Lemma \ref{cohomology-lemma-topology-I-adic-general}의 증명과 완전히 같다.
실제로 $\bigoplus N_n$이
$\bigoplus H^{p + 1}(X, I^n\mathcal{F}_{n + 1})$의 등급
$\bigoplus_{n \geq 0} I^n/I^{n + 1}$-부분가군이고,
$F^n/F^{n + 1} \subset H^p(X, \mathcal{F}_{n + 1})$가
$N_n \to H^p(X, \mathcal{F}_{n + 1})$의 상에 포함됨을 보이기만 하면
거기서 제시한 논증으로 결과가 따른다.
Lemma \ref{cohomology-lemma-ML-general-better}의 증명에서 가군 구조에 관한 명제를
보았다.

\medskip\noindent
$t \in F^n$이라 하자. $H^p(X, \mathcal{F}_{n + 1})$에서 $t$의 상으로
가는 원소 $s \in H^p(X, I^n\mathcal{F}_{n + 1})$를 고르자.
$s$가 $N_n$에 속함을 보여야 한다. 이제 $F^n$은 사상
$M \to H^p(X, \mathcal{F}_n)$의 핵이므로, 모든 $m \geq n$에 대하여
$t$를 원소 $t_m \in H^p(X, \mathcal{F}_{m + 1})$으로 보낼 수 있으며
그 원소는 $H^p(X, \mathcal{F}_n)$에서 영으로 간다. 짧은 완전열
$0 \to I^n\mathcal{F}_{m + 1} \to \mathcal{F}_{m + 1} \to \mathcal{F}_n \to 0$
에서 오는 코호몰로지 열
$$
H^{p - 1}(X, \mathcal{F}_n) \to
H^p(X, I^n\mathcal{F}_{m + 1}) \to
H^p(X, \mathcal{F}_{m + 1}) \to
H^p(X, \mathcal{F}_n)
$$
을 생각하자.
$t_m$으로 가는 $s_m \in H^p(X, I^n\mathcal{F}_{m + 1})$을 고를 수 있다.
위의 열과 $m = n$일 때의 열을 비교하면,
$s_m$은
$H^{p - 1}(X, \mathcal{F}_n) \to H^p(X, I^n\mathcal{F}_{n + 1})$의
상에 속하는 원소를 더한 것을 제외하면 $s$로 감을 알 수 있다.
그러나 이 사상은
$H^p(X, I^n\mathcal{F}_{m + 1}) \to H^p(X, I^n\mathcal{F}_{n + 1})$를
거쳐 인수분해되므로 원하는 대로 $s$가 그 상에 속함을 알 수 있다.
\end{proof}





\section{역계와 코호몰로지, II}
\label{cohomology-section-inverse-systems-bis}

\noindent
이 절에서는 아이디얼이 주 아이디얼인 경우에 Section
\ref{cohomology-section-inverse-systems}의 논의를 계속한다.

\begin{lemma}
\label{cohomology-lemma-equivalent-f-good}
$(X, \mathcal{O}_X)$를 환 달린 공간이라 하자.
$f \in \Gamma(X, \mathcal{O}_X)$라 하자.
$$
\ldots \to \mathcal{F}_3 \to \mathcal{F}_2 \to \mathcal{F}_1
$$
을 $\mathcal{O}_X$-가군의 역계라 하자. 다음 조건들을 생각하자.
\begin{enumerate}
\item 모든 $n \geq 1$에 대하여 사상
$f : \mathcal{F}_{n + 1} \to \mathcal{F}_{n + 1}$은
$\mathcal{F}_{n + 1} \to \mathcal{F}_n$을 거쳐 인수분해되어
짧은 완전열
$0 \to \mathcal{F}_n \to \mathcal{F}_{n + 1} \to \mathcal{F}_1 \to 0$을
준다.
\item 모든 $n \geq 1$에 대하여 사상
$f^n : \mathcal{F}_{n + 1} \to \mathcal{F}_{n + 1}$은
$\mathcal{F}_{n + 1} \to \mathcal{F}_1$을 거쳐 인수분해되어
짧은 완전열
$0 \to \mathcal{F}_1 \to \mathcal{F}_{n + 1} \to \mathcal{F}_n \to 0$을
준다.
\item $f$-가분인 $\mathcal{O}_X$-가군 $\mathcal{G}$가 존재하여
$\mathcal{F}_n = \mathcal{G}[f^n]$이다.
\item $f$-꼬임이 없는 $\mathcal{O}_X$-가군 $\mathcal{F}$가 존재하여
$\mathcal{F}_n = \mathcal{F}/f^n\mathcal{F}$이다.
\end{enumerate}
그러면 (4) $\Rightarrow$ (3) $\Leftrightarrow$ (2) $\Leftrightarrow$ (1)이다.
\end{lemma}

\begin{proof}
(1)과 (2)의 동치에 대한 증명은 생략한다. (3)이 (1)을 함의한다는
증명도 생략한다. (1)에서와 같은 $\mathcal{F}_n$이 주어졌다고 하자.
(3)을 증명하기 위해
$\mathcal{G} = \colim \mathcal{F}_n$으로 놓는다. 여기서 사상
$\mathcal{F}_1 \to \mathcal{F}_2 \to \mathcal{F}_3 \to \ldots$은
(1)에서와 같다. 짧은 완전열 (1)에 의해
$\mathcal{F}_{n + 1} \subset \mathcal{G}$의 상은
$\mathcal{F}_n \subset \mathcal{G}$이므로 사상
$f : \mathcal{G} \to \mathcal{G}$는 전사이다.
따라서 $\mathcal{G}$는 $\mathcal{F}_n = \mathcal{G}[f^n]$을 만족하는
$f$-가분 $\mathcal{O}_X$-가군이다.

\medskip\noindent
(4)에서와 같은 $\mathcal{F}$가 주어졌다고 하자. 사상
$\mathcal{F}/f^{n + 1}\mathcal{F} \to \mathcal{F}/f^n\mathcal{F}$는
항상 전사이며, 그 핵은 $f^n$을 곱하여 유도된 사상
$\mathcal{F}/f\mathcal{F} \to \mathcal{F}/f^{n + 1}\mathcal{F}$의 상이다.
(2)를 확인하려면
$f^n : \mathcal{F} \to \mathcal{F}/f^{n + 1}\mathcal{F}$의 핵이
$f\mathcal{F}$임을 보면 충분하다. 이를 위해 열린집합 $U \subset X$
위의 $\mathcal{F}$의 단면 $s, t$가
$f^ns = f^{n + 1}t$를 만족할 때 $s = ft$임을 보이면 충분하다.
$\mathcal{F}$는 $f$-꼬임이 없으므로
$f : \mathcal{F} \to \mathcal{F}$가 단사이고, 따라서 이는 명백하다.
\end{proof}

\begin{lemma}
\label{cohomology-lemma-topology-I-adic-f}
$X$, $f$, $(\mathcal{F}_n)$이 Lemma
\ref{cohomology-lemma-equivalent-f-good}의 조건 (1)을 만족한다고 하자.
$p \geq 0$라 하고 $H^p = \lim H^p(X, \mathcal{F}_n)$으로 놓자.
그러면 모든 $c \geq 1$에 대하여 $f^cH^p$는
$H^p \to H^p(X, \mathcal{F}_c)$의 핵이다. 따라서 $H^p$ 위의
극한 위상은 $f$-진 위상이다.
\end{lemma}

\begin{proof}
$c \geq 1$라 하자. $f^c H^p$가 $H^p(X, \mathcal{F}_c)$에서 영으로
감은 명백하다. $\xi = (\xi_n) \in H^p$가 극한 위상에서 충분히 작으면
$\xi_c = 0$이고, 따라서 $n \geq c$에 대하여 $\xi_n$은
$H^p(X, \mathcal{F}_c)$에서 영으로 간다. 짧은 완전열의 역계
$$
0 \to \mathcal{F}_{n - c} \xrightarrow{f^c} \mathcal{F}_n \to
\mathcal{F}_c \to 0
$$
와 이에 대응하는 긴 완전 코호몰로지 열의 역계
$$
H^{p - 1}(X, \mathcal{F}_c) \to
H^p(X, \mathcal{F}_{n - c}) \to
H^p(X, \mathcal{F}_n) \to
H^p(X, \mathcal{F}_c)
$$
를 생각하자. 항 $H^{p - 1}(X, \mathcal{F}_c)$는 $n$에 의존하지
않으므로 $\xi_n$을 올리는 원소들의 호환 열
$\xi'_n \in H^p(X, \mathcal{F}_{n - c})$을 고를 수 있다.
$\xi' = (\xi'_n)$으로 놓으면 원하는 대로 $\xi = f^c \xi'$임을 알 수 있다.
\end{proof}

\begin{lemma}
\label{cohomology-lemma-limit-finite}
$A$를 주 아이디얼 $(f)$에 관하여 완비인 뇌터 환이라 하자.
$X$를 위상공간이라 하고,
$$
\ldots \to \mathcal{F}_3 \to \mathcal{F}_2 \to \mathcal{F}_1
$$
을 $A$-가군 층들의 역계라 하자. 다음을 가정하자.
\begin{enumerate}
\item $\Gamma(X, \mathcal{F}_1)$은 유한 $A$-가군이다.
\item $X$, $f$, $(\mathcal{F}_n)$은 Lemma
\ref{cohomology-lemma-equivalent-f-good}의 조건 (1)을 만족한다.
\end{enumerate}
그러면
$$
M = \lim \Gamma(X, \mathcal{F}_n)
$$
은 유한 $A$-가군이고, $f$는 $M$ 위의 영인자가 아니며,
$M/fM$은 $\Gamma(X, \mathcal{F}_1)$에서 $M$의 상이다.
\end{lemma}

\begin{proof}
Lemma \ref{cohomology-lemma-topology-I-adic-f}에 의해
$M/fM \subset H^0(X, \mathcal{F}_1)$이다. (1)과 $A$의 뇌터 성질로부터
$M/fM$은 유한 $A$-가군이다. $f^nM$은
$H^0(X, \mathcal{F}_n)$에서 영으로 가므로
$\bigcap f^nM = 0$임을 관찰하자. Algebra, Lemma
\ref{algebra-lemma-finite-over-complete-ring}에 의해 $M$이 $A$ 위에서
유한임을 얻는다. 마지막으로
$s = (s_n) \in M = \lim H^0(X, \mathcal{F}_n)$이고 $fs = 0$이라고 하자.
Lemma \ref{cohomology-lemma-equivalent-f-good}의 조건 (1)에 의해
$s_{n + 1}$은 $\mathcal{F}_{n + 1} \to \mathcal{F}_n$의 핵에 속한다.
따라서 $s_n = 0$이다. $n$은 임의였으므로 $s = 0$이다.
그러므로 $f$는 $M$ 위의 영인자가 아니다.
\end{proof}

\begin{lemma}
\label{cohomology-lemma-ML}
$A$를 환, $f \in A$, $X$를 위상공간이라 하고,
$$
\ldots \to \mathcal{F}_3 \to \mathcal{F}_2 \to \mathcal{F}_1
$$
을 $A$-가군 층들의 역계라 하자. $p \geq 0$라 하자.
다음을 가정하자.
\begin{enumerate}
\item $H^{p + 1}(X, \mathcal{F}_1)$이 유한 길이의 $A$-가군이거나,
$A$가 뇌터 환이고 $H^{p + 1}(X, \mathcal{F}_1)$이 유한 $A$-가군이다.
\item $X$, $f$, $(\mathcal{F}_n)$은 Lemma
\ref{cohomology-lemma-equivalent-f-good}의 조건 (1)을 만족한다.
\end{enumerate}
그러면 역계 $M_n = H^p(X, \mathcal{F}_n)$은 Mittag--Leffler 조건을
만족한다.
\end{lemma}

\begin{proof}
$I = (f)$로 놓자. Lemma \ref{cohomology-lemma-ML-general}의 판정법을 사용할 것이다.
모든 $n \geq 0$에 대하여
$f^n : \mathcal{F}_1 \to I^n\mathcal{F}_{n + 1}$은 동형이다.
따라서
$$
\bigoplus\nolimits_{n \geq 1} H^{p + 1}(X, \mathcal{F}_1) \cdot f^{n + 1}
$$
이 오름 사슬 조건을 만족하는 등급
$S = \bigoplus_{n \geq 0} A/(f) \cdot f^n$-가군임을 보이면 충분하다.
$A$가 뇌터 환이 아니면 $H^{p + 1}(X, \mathcal{F}_1)$은 유한 길이를
가지므로 결과가 성립한다. $A$가 뇌터 환이면 $S$는 뇌터 환이고,
$H^{p + 1}(X, \mathcal{F}_1)$의 가정된 유한성에 의해 이 가군은
$S$ 위에서 유한이므로 결과가 성립한다. 일부 세부 사항은 생략한다.
\end{proof}

\begin{lemma}
\label{cohomology-lemma-ML-better}
$A$를 환, $f \in A$, $X$를 위상공간이라 하고,
$$
\ldots \to \mathcal{F}_3 \to \mathcal{F}_2 \to \mathcal{F}_1
$$
을 $A$-가군 층들의 역계라 하자. $p \geq 0$라 하자.
다음을 가정하자.
\begin{enumerate}
\item 어떤 $m \geq 1$에 대하여
$H^{p + 1}(X, \mathcal{F}_m) \to H^{p + 1}(X, \mathcal{F}_1)$의 상이
유한 길이의 $A$-가군이거나, $A$가 뇌터 환이고
$H^{p + 1}(X, \mathcal{F}_m) \to H^{p + 1}(X, \mathcal{F}_1)$의
상들의 교집합이 유한 $A$-가군이다.
\item $X$, $f$, $(\mathcal{F}_n)$은 Lemma
\ref{cohomology-lemma-equivalent-f-good}의 조건 (1)을 만족한다.
\end{enumerate}
그러면 역계 $M_n = H^p(X, \mathcal{F}_n)$은 Mittag--Leffler 조건을
만족한다.
\end{lemma}

\begin{proof}
$I = (f)$로 놓자. 가군 $N_n$들을 사용하는 Lemma
\ref{cohomology-lemma-ML-general-better}의 판정법을 사용할 것이다.
$m \geq n$에 대하여
$I^n\mathcal{F}_{m + 1} = \mathcal{F}_{m + 1 - n}$이다. 따라서
$$
N_n = \bigcap\nolimits_{m \geq 1} \Im\left(
H^{p + 1}(X, \mathcal{F}_m) \to H^{p + 1}(X, \mathcal{F}_1)
\right)
$$
은 $n$에 의존하지 않고
$\bigoplus N_n = \bigoplus N_1 \cdot f^{n + 1}$이다.
따라서 Lemma \ref{cohomology-lemma-ML}의 증명에서와 정확히 같이 결론을 얻는다.
\end{proof}

\begin{remark}
\label{cohomology-remark-compare-ML}
$(X, \mathcal{O}_X)$를 환 달린 공간이라 하고,
$f \in \Gamma(X, \mathcal{O}_X)$라 하자.
$\mathcal{F}$를 $\mathcal{O}_X$-가군이라 하자.
$\mathcal{F}$가 $f$-꼬임이 없으면 모든 $p \geq 0$에 대하여
역계의 짧은 완전열
$$
0 \to \{H^p(X, \mathcal{F})/f^nH^p(X, \mathcal{F})\}
\to \{H^p(X, \mathcal{F}/f^n\mathcal{F})\}
\to \{H^{p + 1}(X, \mathcal{F})[f^n]\} \to 0
$$
이 있다. 첫째 역계는 Mittag--Leffler 조건(ML)을 만족하므로 여기서
세 가지를 알 수 있다.
\begin{enumerate}
\item 짧은 완전열
$$
0 \to \widehat{H^p(X, \mathcal{F})} \to
\lim H^p(X, \mathcal{F}/f^n\mathcal{F}) \to
T_f(H^{p + 1}(X, \mathcal{F})) \to 0
$$
이 있다. 여기서 $\widehat{\ }$는 통상적인 $f$-진 완비화를 나타내고,
$T_f( - )$는 More on Algebra, Example
\ref{more-algebra-example-spectral-sequence-principal}의
$f$-진 Tate 가군을 나타낸다.
\item 다음 등식이 성립한다.
$R^1\lim H^p(X, \mathcal{F}/f^n\mathcal{F}) =
R^1\lim H^{p + 1}(X, \mathcal{F})[f^n]$.
\item 계 $\{H^{p + 1}(X, \mathcal{F})[f^n]\}$이 ML일 필요충분조건은
$\{H^p(X, \mathcal{F}/f^n\mathcal{F})\}$가 ML인 것이다.
\end{enumerate}
Homology, Lemma \ref{homology-lemma-Mittag-Leffler}와
More on Algebra, Lemmas \ref{more-algebra-lemma-six-term-Rlim} 및
\ref{more-algebra-lemma-Mittag-Leffler}를 보라.
\end{remark}










\section{유도 극한}
\label{cohomology-section-derived-limits}

\noindent
$(X, \mathcal{O}_X)$를 환 달린 공간이라 하자. 삼각범주
$D(\mathcal{O}_X)$는 곱을 가지므로(Injectives, Lemma
\ref{injectives-lemma-derived-products}) $D(\mathcal{O}_X)$는 유도 극한도 가진다.
Derived Categories, Definition
\ref{derived-definition-derived-limit}을 보라.
$(K_n)$이 $D(\mathcal{O}_X)$ 안의 역계이면 그 유도 극한을
$R\lim K_n$이라 쓴다.

\begin{lemma}
\label{cohomology-lemma-RGamma-commutes-with-Rlim}
$(X, \mathcal{O}_X)$를 환 달린 공간이라 하자. 열린집합
$U \subset X$에 대하여 함자 $R\Gamma(U, -)$는 $R\lim$과 가환한다.
더욱이 $D(\mathcal{O}_X)$ 안의 임의의 역계 $(K_n)$과 임의의
$m \in \mathbf{Z}$에 대하여 짧은 완전열
$$
0 \to
R^1\lim H^{m - 1}(U, K_n) \to H^m(U, R\lim K_n) \to
\lim H^m(U, K_n) \to 0
$$
이 있다.
\end{lemma}

\begin{proof}
첫째 명제는 Injectives, Lemma
\ref{injectives-lemma-RF-commutes-with-Rlim}에서 따른다.
그런 다음
$R\lim R\Gamma(U, K_n) = R\Gamma(U, R\lim K_n)$에
More on Algebra, Remark
\ref{more-algebra-remark-compare-derived-limit}을 적용하면 짧은 완전열들을
얻는다.
\end{proof}

\begin{lemma}
\label{cohomology-lemma-Rf-commutes-with-Rlim}
$f : (X, \mathcal{O}_X) \to (Y, \mathcal{O}_Y)$를 환 달린 공간의
사상이라 하자. 그러면 $Rf_*$는 $R\lim$과 가환한다. 즉 $Rf_*$는
유도 극한과 가환한다.
\end{lemma}

\begin{proof}
$(K_n)$을 $D(\mathcal{O}_X)$ 안의 역계라 하자. $D(\mathcal{O}_X)$ 안의
정의 특수 삼각형
$$
R\lim K_n \to \prod K_n \to \prod K_n
$$
을 생각하자. 완전 함자 $Rf_*$를 적용하면 $D(\mathcal{O}_Y)$ 안의
특수 삼각형
$$
Rf_*(R\lim K_n) \to Rf_*\left(\prod K_n\right) \to Rf_*\left(\prod K_n\right)
$$
을 얻는다. 따라서 $Rf_*$가 유도범주에서의 곱과 가환함을 증명하면
충분하다(그 곱은 복합체의 곱으로 단순히 주어지지 않는다. Injectives,
Lemma \ref{injectives-lemma-derived-products}를 보라).
그러나 Lemma \ref{cohomology-lemma-adjoint}에 의해 $Rf_*$는 오른쪽 수반이므로
이는 형식적으로 따른다(Categories, Lemma
\ref{categories-lemma-adjoint-exact} 참조).
주의: $R\lim K_n$은 $D(\mathcal{O}_X)$ 안의 극한이 아니므로
Categories, Lemma \ref{categories-lemma-adjoint-exact}을 직접 적용할
수는 없다.
\end{proof}

\begin{remark}
\label{cohomology-remark-discuss-derived-limit}
$(X, \mathcal{O}_X)$를 환 달린 공간이라 하고, $(K_n)$을
$D(\mathcal{O}_X)$ 안의 역계라 하자. $K = R\lim K_n$으로 놓자.
각 $n$과 $m$에 대하여 $\mathcal{H}^m_n = H^m(K_n)$을 $K_n$의
$m$차 코호몰로지 층이라 하고, 마찬가지로
$\mathcal{H}^m = H^m(K)$으로 놓자. 준층
$$
U \longmapsto \underline{\mathcal{H}}^m_n(U) = H^m(U, K_n)
$$
을 $\underline{\mathcal{H}}^m_n$이라 쓰자. 마찬가지로
$\underline{\mathcal{H}}^m(U) = H^m(U, K)$로 놓는다. Lemma
\ref{cohomology-lemma-sheafification-cohomology}에 의해
$\mathcal{H}^m_n$은 $\underline{\mathcal{H}}^m_n$의 층화이고,
$\mathcal{H}^m$은 $\underline{\mathcal{H}}^m$의 층화이다.
다음 도식이 있다.
$$
\xymatrix{
K \ar@{=}[d] &
\underline{\mathcal{H}}^m \ar[d] \ar[r] & 
\mathcal{H}^m \ar[d] \\
R\lim K_n &
\lim \underline{\mathcal{H}}^m_n \ar[r] & 
\lim \mathcal{H}^m_n
}
$$
일반적으로 $\lim \mathcal{H}^m_n$이
$\lim \underline{\mathcal{H}}^m_n$의 층화인 것은 아닐 수 있다.
$U \subset X$가 열린집합이면 Lemma
\ref{cohomology-lemma-RGamma-commutes-with-Rlim}에 의해 짧은 완전열
\begin{equation}
\label{cohomology-equation-ses-Rlim-over-U}
0 \to
R^1\lim \underline{\mathcal{H}}^{m - 1}_n(U) \to
\underline{\mathcal{H}}^m(U) \to
\lim \underline{\mathcal{H}}^m_n(U) \to 0
\end{equation}
을 얻는다.
\end{remark}

\noindent
다음 보조정리는 스킴 위에서 전이 사상이 전사인 준연접 가군의 역계에
적용된다.

\begin{lemma}
\label{cohomology-lemma-inverse-limit-is-derived-limit}
$(X, \mathcal{O}_X)$를 환 달린 공간이라 하고,
$(\mathcal{F}_n)$을 $\mathcal{O}_X$-가군의 역계라 하자.
$\mathcal{B}$를 $X$의 열린집합들로 이루어진 집합이라 하자.
다음을 가정하자.
\begin{enumerate}
\item $X$의 모든 열린집합은 원소들이 $\mathcal{B}$에 속하는 덮개를 갖는다.
\item $p > 0$ 및 $U \in \mathcal{B}$에 대하여
$H^p(U, \mathcal{F}_n) = 0$이다.
\item $U \in \mathcal{B}$에 대하여 역계 $\mathcal{F}_n(U)$의
$R^1\lim$은 영이다.
\end{enumerate}
그러면 $R\lim \mathcal{F}_n = \lim \mathcal{F}_n$이고,
$p > 0$ 및 $U \in \mathcal{B}$에 대하여
$H^p(U, \lim \mathcal{F}_n) = 0$이다.
\end{lemma}

\begin{proof}
$K_n = \mathcal{F}_n$ 및 $K = R\lim \mathcal{F}_n$으로 놓자.
Remark \ref{cohomology-remark-discuss-derived-limit}의 표기와 가정 (2)에 의해
$U \in \mathcal{B}$에 대하여 $m \not = 0$이면
$\underline{\mathcal{H}}_n^m(U) = 0$이고,
$\underline{\mathcal{H}}_n^0(U) = \mathcal{F}_n(U)$이다.
Equation (\ref{cohomology-equation-ses-Rlim-over-U})과 가정 (3)으로부터
$m \not = 0$이면 $\underline{\mathcal{H}}^m(U) = 0$이고,
$m = 0$이면 $\lim \mathcal{F}_n(U)$와 같음을 알 수 있다.
(1)을 사용하여 층화하면 $m \not = 0$이면
$\mathcal{H}^m = 0$이고, $m = 0$이면 $\lim \mathcal{F}_n$과 같음을
얻는다. 따라서 $K = \lim \mathcal{F}_n$이다.
$H^m(U, K) = \underline{\mathcal{H}}^m(U) = 0$이 $m > 0$에 대하여
성립하므로(위 참조) 둘째 명제도 성립한다.
\end{proof}

\begin{lemma}
\label{cohomology-lemma-cohomology-derived-limit-injective}
$(X, \mathcal{O}_X)$를 환 달린 공간이라 하고, $(K_n)$을
$D(\mathcal{O}_X)$ 안의 역계라 하자. $x \in X$ 및
$m \in \mathbf{Z}$라 하자. 정수 $n(x)$와 $x$의 열린 근방들의 기본계
$\mathfrak{U}_x$가 존재하여 $U \in \mathfrak{U}_x$에 대하여
\begin{enumerate}
\item $R^1\lim H^{m - 1}(U, K_n) = 0$이고,
\item $n \geq n(x)$에 대하여
$H^m(U, K_n) \to H^m(U, K_{n(x)})$가 단사라고 가정하자.
\end{enumerate}
그러면 줄기 위의 사상
$H^m(R\lim K_n)_x \to H^m(K_{n(x)})_x$는 단사이다.
\end{lemma}

\begin{proof}
$H^m(K_{n(x)})_x$에서 영으로 가는
$H^m(R\lim K_n)_x$의 원소를 $\gamma$라 하자.
$H^m(R\lim K_n)$은 준층
$U \mapsto H^m(U, R\lim K_n)$의 층화이므로(Lemma
\ref{cohomology-lemma-sheafification-cohomology}에 의한다),
$\gamma$로 가는 원소 $\tilde \gamma \in H^m(U, R\lim K_n)$와
$U \in \mathfrak{U}_x$를 고를 수 있다. 그러면 $\tilde\gamma$는
$\tilde\gamma_{n(x)} \in H^m(U, K_{n(x)})$로 간다.
$H^m(K_{n(x)})$도 준층 $U \mapsto H^m(U, K_{n(x)})$의 층화이므로
(다시 Lemma \ref{cohomology-lemma-sheafification-cohomology}에 의한다),
$U$를 줄인 뒤 $\tilde\gamma_{n(x)} = 0$이라고 가정해도 된다.
이 $U$에 대하여 Lemma \ref{cohomology-lemma-RGamma-commutes-with-Rlim}의
짧은 완전열
$$
0 \to
R^1\lim H^{m - 1}(U, K_n) \to H^m(U, R\lim K_n) \to
\lim H^m(U, K_n) \to 0
$$
을 생각하자. 가정 (1)에 의해 왼쪽 군은 영이고, 가정 (2)에 의해
오른쪽 군은 $H^m(U, K_{n(x)})$ 안으로 단사로 간다.
따라서 $\tilde\gamma = 0$이고, 원하는 대로 $\gamma = 0$이다.
\end{proof}

\begin{lemma}
\label{cohomology-lemma-is-limit-per-point}
$(X, \mathcal{O}_X)$를 환 달린 공간이라 하고,
$E \in D(\mathcal{O}_X)$라 하자. 모든 $x \in X$에 대하여 함수
$p(x, -) : \mathbf{Z} \to \mathbf{Z}$와 $x$의 열린 근방들의 기본계
$\mathfrak{U}_x$가 존재하여
$$
H^p(U, H^{m - p}(E)) = 0 \text{이고, 이는 }
U \in \mathfrak{U}_x \text{이고 } p > p(x, m)\text{일 때 성립한다}
$$
라고 가정하자. 그러면 Derived Categories, Remark
\ref{derived-remark-map-into-derived-limit-truncations}의 사상
$E \to R\lim \tau_{\geq -n} E$는 $D(\mathcal{O}_X)$ 안의 동형이다.
\end{lemma}

\begin{proof}
$K_n = \tau_{\geq -n}E$ 및 $K = R\lim K_n$으로 놓자.
표준 사상 $E \to K$는 표준 사상
$E \to K_n = \tau_{\geq -n}E$들에서 온다. $E \to K$가 코호몰로지
층들의 동형 $H^m(E) \to H^m(K)$을 유도함을 보여야 한다.
나머지 증명에서는 $m$을 고정한다. $n \geq -m$이면 사상
$E \to \tau_{\geq -n}E = K_n$은 동형
$H^m(E) \to H^m(K_n)$을 유도한다.
증명을 마치려면 모든 $x \in X$에 대하여 사상
$H^m(K)_x \to H^m(K_{n(x)})_x$가 단사가 되게 하는 정수
$n(x) \geq -m$이 존재함을 보이면 충분하다. 실제로 그러면 합성
$$
H^m(E)_x \to H^m(K)_x \to H^m(K_{n(x)})_x
$$
은 전단사이고 둘째 화살표는 단사이므로 첫째 화살표는 전단사이다.
다음과 같이 놓자.
$$
n(x) = 1 + \max\{-m, p(x, m - 1) - m, -1 + p(x, m) - m, -2 + p(x, m + 1) - m\}.
$$
그러면 언제나 $n(x) \geq -m$이다. 주장: $n \geq n(x)$ 및
$U \in \mathfrak{U}_x$에 대하여 사상
$$
H^{m - 1}(U, K_{n + 1}) \to H^{m - 1}(U, K_n)
\quad\text{이고}\quad
H^m(U, K_{n + 1}) \to H^m(U, K_n)
$$
은 동형이다. 이 주장은 Lemma
\ref{cohomology-lemma-cohomology-derived-limit-injective}의 조건 (1)과 (2)가
성립함을 뜻하며, 따라서 원하는 단사성을 함의한다.
다음 특수 삼각형들이 있음을 상기하자(Derived Categories, Remark
\ref{derived-remark-truncation-distinguished-triangle}).
$$
H^{-n - 1}(E)[n + 1] \to
K_{n + 1} \to K_n \to H^{-n - 1}(E)[n + 2]
$$
이에 연관된 긴 완전 코호몰로지 열을 보면,
$n \geq n(x)$ 및 $U \in \mathfrak{U}_x$에 대하여
$$
H^{m + n}(U, H^{-n - 1}(E)),\quad
H^{m + n + 1}(U, H^{-n - 1}(E)),\quad
H^{m + n + 2}(U, H^{-n - 1}(E))
$$
가 영이면 주장이 따른다. 이는 $n(x)$의 선택과 보조정리의 가정에서
따른다.
\end{proof}

\begin{lemma}
\label{cohomology-lemma-is-limit-spaltenstein}
\begin{reference}
\cite[Proposition 3.13]{Spaltenstein}
\end{reference}
$(X, \mathcal{O}_X)$를 환 달린 공간이라 하고,
$E \in D(\mathcal{O}_X)$라 하자. 모든 $x \in X$에 대하여 정수
$d_x \geq 0$와 $x$의 열린 근방들의 기본계 $\mathfrak{U}_x$가
존재하여
$$
H^p(U, H^q(E)) = 0 \text{이고, 이는 }
U \in \mathfrak{U}_x,\ p > d_x, \text{이고 }q < 0\text{일 때 성립한다}
$$
라고 가정하자. 그러면 Derived Categories, Remark
\ref{derived-remark-map-into-derived-limit-truncations}의 사상
$E \to R\lim \tau_{\geq -n} E$는 $D(\mathcal{O}_X)$ 안의 동형이다.
\end{lemma}

\begin{proof}
이는 $p(x, m) = d_x + \max(0, m)$로 둔 Lemma
\ref{cohomology-lemma-is-limit-per-point}에서 따른다.
\end{proof}

\begin{lemma}
\label{cohomology-lemma-is-limit}
$(X, \mathcal{O}_X)$를 환 달린 공간이라 하고,
$E \in D(\mathcal{O}_X)$라 하자. 함수
$p(-) : \mathbf{Z} \to \mathbf{Z}$와 $X$의 열린집합들로 이루어진 집합
$\mathcal{B}$가 존재하여 다음을 만족한다고 가정하자.
\begin{enumerate}
\item $X$의 모든 열린집합은 원소들이 $\mathcal{B}$에 속하는 덮개를 갖는다.
\item $p > p(m)$ 및 $U \in \mathcal{B}$에 대하여
$H^p(U, H^{m - p}(E)) = 0$이다.
\end{enumerate}
그러면 Derived Categories, Remark
\ref{derived-remark-map-into-derived-limit-truncations}의 사상
$E \to R\lim \tau_{\geq -n} E$는 $D(\mathcal{O}_X)$ 안의 동형이다.
\end{lemma}

\begin{proof}
$p(x, m) = p(m)$ 및
$\mathfrak{U}_x = \{U \in \mathcal{B} \mid x \in U\}$로 두고
Lemma \ref{cohomology-lemma-is-limit-per-point}을 적용한다.
\end{proof}

\begin{lemma}
\label{cohomology-lemma-is-limit-dimension}
$(X, \mathcal{O}_X)$를 환 달린 공간이라 하고,
$E \in D(\mathcal{O}_X)$라 하자. 정수 $d \geq 0$와 $X$의 위상의
기저 $\mathcal{B}$가 존재하여
$$
H^p(U, H^q(E)) = 0 \text{이고, 이는 }
U \in \mathcal{B},\ p > d, \text{이고 }q < 0\text{일 때 성립한다}
$$
라고 가정하자. 그러면 Derived Categories, Remark
\ref{derived-remark-map-into-derived-limit-truncations}의 사상
$E \to R\lim \tau_{\geq -n} E$는 $D(\mathcal{O}_X)$ 안의 동형이다.
\end{lemma}

\begin{proof}
$d_x = d$ 및 $\mathfrak{U}_x = \{U \in \mathcal{B} \mid x \in U\}$로
두고 Lemma \ref{cohomology-lemma-is-limit-spaltenstein}을 적용한다.
\end{proof}

\noindent
위의 보조정리들을 사용하면 어떤 상황에서 코호몰로지를 계산할 수 있다.

\begin{lemma}
\label{cohomology-lemma-cohomology-over-U-trivial}
$(X, \mathcal{O}_X)$를 환 달린 공간이라 하고,
$K$를 $D(\mathcal{O}_X)$의 대상이라 하자.
$\mathcal{B}$를 $X$의 열린집합들로 이루어진 집합이라 하자.
다음을 가정하자.
\begin{enumerate}
\item $X$의 모든 열린집합은 원소들이 $\mathcal{B}$에 속하는 덮개를 갖는다.
\item 모든 $p > 0$, $q \in \mathbf{Z}$, $U \in \mathcal{B}$에 대하여
$H^p(U, H^q(K)) = 0$이다.
\end{enumerate}
그러면 $q \in \mathbf{Z}$ 및 $U \in \mathcal{B}$에 대하여
$H^q(U, K) = H^0(U, H^q(K))$이다.
\end{lemma}

\begin{proof}
$d = 0$으로 둔 Lemma \ref{cohomology-lemma-is-limit-dimension}에 의해
$K = R\lim \tau_{\geq -n} K$임을 관찰하자.
$U \in \mathcal{B}$라 하자. Equation
(\ref{cohomology-equation-ses-Rlim-over-U})에 의해 짧은 완전열
$$
0 \to R^1\lim H^{q - 1}(U, \tau_{\geq -n}K) \to
H^q(U, K) \to \lim H^q(U, \tau_{\geq -n}K) \to 0
$$
을 얻는다. Example \ref{cohomology-example-spectral-sequence}의 스펙트럼열을
사용하면 조건 (2)는 모든 $q$에 대하여
$H^q(U, \tau_{\geq -n} K) = H^0(U, H^q(\tau_{\geq -n} K))$임을 함의한다.
$\tau_{\geq -n}K$는 아래로 유계이므로 스펙트럼열은 수렴한다.
$n > -q$이면 $H^q(\tau_{\geq -n}K) = H^q(K)$이다.
따라서 표시한 짧은 완전열의 왼쪽 계와 오른쪽 계는 각각
$H^0(U, H^{q - 1}(K))$와 $H^0(U, H^q(K))$를 값으로 하여 결국
상수이다. 보조정리가 따른다.
\end{proof}

\noindent
유도 극한을 기술할 수 있는 또 다른 경우가 있다.

\begin{lemma}
\label{cohomology-lemma-derived-limit-suitable-system}
$(X, \mathcal{O}_X)$를 환 달린 공간이라 하고, $(K_n)$을
$D(\mathcal{O}_X)$의 대상들의 역계라 하자.
$\mathcal{B}$를 $X$의 열린집합들로 이루어진 집합이라 하자.
다음을 가정하자.
\begin{enumerate}
\item $X$의 모든 열린집합은 원소들이 $\mathcal{B}$에 속하는 덮개를 갖는다.
\item 모든 $U \in \mathcal{B}$와 모든 $q \in \mathbf{Z}$에 대하여
\begin{enumerate}
\item $p > 0$에 대하여 $H^p(U, H^q(K_n)) = 0$이다.
\item 역계 $H^0(U, H^q(K_n))$의 $R^1\lim$은 영이다.
\end{enumerate}
\end{enumerate}
그러면 $q \in \mathbf{Z}$에 대하여
$H^q(R\lim K_n) = \lim H^q(K_n)$이다.
\end{lemma}

\begin{proof}
$K = R\lim K_n$으로 놓고, $U \in \mathcal{B}$라 하자.
Lemma \ref{cohomology-lemma-cohomology-over-U-trivial}과 (2)(a)에 의해
$H^q(U, K_n) = H^0(U, H^q(K_n))$이다. Lemma
\ref{cohomology-lemma-RGamma-commutes-with-Rlim}과 (2)(b)에 의해
$H^q(U, K) = \lim H^0(U, H^q(K_n))$이다. 따라서 $H^q(U, K)$는
$U$ 위의 층 $H^q(K_n)$의 단면들의 역극한이다.
$\lim H^q(K_n)$은 층이므로, 가정 (1)을 사용하면 준층
$U \mapsto H^q(U, K)$의 층화인 $H^q(K)$가
$\lim H^q(K_n)$과 같음을 얻는다. 이것이 보조정리를 증명한다.
\end{proof}








\section{K-단사 분해의 구성}
\label{cohomology-section-K-injective}


\noindent
$(X, \mathcal{O}_X)$를 환 달린 공간이라 하자.
$\mathcal{F}^\bullet$을 $\mathcal{O}_X$-가군의 복합체라 하자.
범주 $\textit{Mod}(\mathcal{O}_X)$는 충분한 단사대상을 가지므로
Derived Categories, Lemma
\ref{derived-lemma-special-inverse-system}을 사용하여 다음 도식을
구성할 수 있다.
$$
\xymatrix{
\ldots \ar[r] &
\tau_{\geq -2}\mathcal{F}^\bullet \ar[r] \ar[d] &
\tau_{\geq -1}\mathcal{F}^\bullet \ar[d] \\
\ldots \ar[r] & \mathcal{I}_2^\bullet \ar[r] & \mathcal{I}_1^\bullet
}
$$
이 도식은 $\mathcal{O}_X$-가군 복합체의 범주 안에 있으며 다음을 만족한다.
\begin{enumerate}
\item 세로 화살표들은 유사동형이다.
\item $\mathcal{I}_n^\bullet$은 단사대상들로 이루어진 아래로 유계인
복합체이다.
\item 화살표 $\mathcal{I}_{n + 1}^\bullet \to \mathcal{I}_n^\bullet$은
항별 분할 전사이다.
\end{enumerate}
$\mathcal{O}_X$-가군의 범주는 극한을 가지며(준층 수준에서 계산된다),
따라서 항별 극한
$\mathcal{I}^\bullet = \lim_n \mathcal{I}_n^\bullet$을 만들 수 있다.
Derived Categories, Lemmas
\ref{derived-lemma-bounded-below-injectives-K-injective} 및
\ref{derived-lemma-limit-K-injectives}
에 의해 이는 K-단사 복합체이다. 일반적으로 표준 사상
\begin{equation}
\label{cohomology-equation-into-candidate-K-injective}
\mathcal{F}^\bullet \to \mathcal{I}^\bullet
\end{equation}
은 유사동형이 아닐 수 있다. 다음 보조정리에서는 이것이 유사동형이 되는
몇 가지 조건을 기술한다.

\begin{lemma}
\label{cohomology-lemma-K-injective}
위에서 기술한 상황에 있다고 하자.
$\mathcal{H}^m = H^m(\mathcal{F}^\bullet)$을 $m$차 코호몰로지 층이라 쓰자.
$\mathcal{B}$를 $X$의 열린 부분집합들로 이루어진 집합이라 하자.
$d \in \mathbf{N}$이라 하자. 다음을 가정하자.
\begin{enumerate}
\item $X$의 모든 열린집합은 원소들이 $\mathcal{B}$에 속하는 덮개를 갖는다.
\item 모든 $U \in \mathcal{B}$에 대하여 $p > d$ 및 $q < 0$이면
$H^p(U, \mathcal{H}^q) = 0$이다\footnote{
$\forall m$, $\exists p(m)$에 대하여 $p > p(m)$이면
$H^p(U. \mathcal{H}^{m - p}) = 0$인 것으로 충분하다. Lemma
\ref{cohomology-lemma-is-limit}을 보라.}.
\end{enumerate}
그러면 (\ref{cohomology-equation-into-candidate-K-injective})은 유사동형이다.
\end{lemma}

\begin{proof}
Derived Categories, Lemma
\ref{derived-lemma-difficulty-K-injectives}에 의해 사상
$\mathcal{F}^\bullet \to R\lim \tau_{\geq -n} \mathcal{F}^\bullet$
이 동형임을 보이면 충분하다. 이는 Lemma
\ref{cohomology-lemma-is-limit-dimension}이다.
\end{proof}

\noindent
다음은 층 복합체들의 계의 역극한이 갖는 코호몰로지 층에 관한 기술적
보조정리이다. 실제 역극한이 아니라 유도 극한을 취해야 한다는 점에서,
어떤 의미로는 이 보조정리는 증명하려고 할 올바른 명제가 아니다.

\begin{lemma}
\label{cohomology-lemma-inverse-limit-complexes}
$(X, \mathcal{O}_X)$를 환 달린 공간이라 하고,
$(\mathcal{F}_n^\bullet)$을 $\mathcal{O}_X$-가군 복합체의 역계라 하자.
$m \in \mathbf{Z}$라 하자. $X$의 열린 부분집합들로 이루어진 집합
$\mathcal{B}$와 정수 $n_0$가 존재하여 다음을 만족한다고 가정하자.
\begin{enumerate}
\item $X$의 모든 열린집합은 원소들이 $\mathcal{B}$에 속하는 덮개를 갖는다.
\item 모든 $U \in \mathcal{B}$에 대하여
\begin{enumerate}
\item 아벨군의 계
$\mathcal{F}_n^{m - 2}(U)$ 및 $\mathcal{F}_n^{m - 1}(U)$
의 $R^1\lim$은 영이다(예를 들어 이 계들은 Mittag--Leffler 조건을
만족한다).
\item 아벨군의 계 $H^{m - 1}(\mathcal{F}_n^\bullet(U))$의
$R^1\lim$은 영이다(예를 들어 이 계는 Mittag--Leffler 조건을 만족한다).
\item 모든 $n \geq n_0$에 대하여
$H^m(\mathcal{F}_n^\bullet(U)) = H^m(\mathcal{F}_{n_0}^\bullet(U))$
이다.
\end{enumerate}
\end{enumerate}
그러면 사상
$H^m(\mathcal{F}^\bullet) \to \lim H^m(\mathcal{F}_n^\bullet) \to
H^m(\mathcal{F}_{n_0}^\bullet)$
들은 층의 동형이다. 여기서
$\mathcal{F}^\bullet = \lim \mathcal{F}_n^\bullet$은 항별 역극한이다.
\end{lemma}

\begin{proof}
$U \in \mathcal{B}$라 하자. $H^m(\mathcal{F}^\bullet(U))$은
$$
\lim_n \mathcal{F}_n^{m - 2}(U) \to
\lim_n \mathcal{F}_n^{m - 1}(U) \to
\lim_n \mathcal{F}_n^m(U) \to
\lim_n \mathcal{F}_n^{m + 1}(U)
$$
의 왼쪽에서 셋째 자리에서의 코호몰로지임을 주목하자.
가정 (2)(a)와 (2)(b)에 의해
More on Algebra, Lemma \ref{more-algebra-lemma-apply-Mittag-Leffler-again}
을 적용하여
$$
H^m(\mathcal{F}^\bullet(U)) = \lim H^m(\mathcal{F}_n^\bullet(U))
$$
을 얻는다. 가정 (2)(c)에 의해 모든 $n \geq n_0$에 대하여
$$
H^m(\mathcal{F}^\bullet(U)) = H^m(\mathcal{F}_n^\bullet(U))
$$
이다. 가정 (1)에 의해 모든 $n \geq n_0$에 대하여
$U \mapsto H^m(\mathcal{F}^\bullet(U))$의 층화와
$U \mapsto H^m(\mathcal{F}_n^\bullet(U))$의 층화가 같음을 얻는다.
따라서 층의 역계 $H^m(\mathcal{F}_n^\bullet)$은 $n \geq n_0$에서
$H^m(\mathcal{F}^\bullet)$을 값으로 하여 상수이다. 이것이 보조정리를
증명한다.
\end{proof}












\section{역계와 코호몰로지, III}
\label{cohomology-section-inverse-systems-tri}

\noindent
이 절에서는 유도 극한을 사용하여 Section
\ref{cohomology-section-inverse-systems-bis}의 논의를 계속한다.

\begin{lemma}
\label{cohomology-lemma-formal-functions-principal}
$(X, \mathcal{O}_X)$를 환 달린 공간이라 하자.
$A \to \Gamma(X, \mathcal{O}_X)$를 환 준동형이라 하고,
$f \in A$라 하자. $E$를 $D(\mathcal{O}_X)$의 대상이라 하자. 다음과 같이 쓰자.
$$
E_n = E \otimes_{\mathcal{O}_X} (\mathcal{O}_X \xrightarrow{f^n} \mathcal{O}_X)
$$
$E^\wedge = R\lim E_n$으로 놓자. $p \in \mathbf{Z}$에 대하여
다음 표준 가환 도식이 있다.
$$
\xymatrix{
& 0 & 0 \\
0 \ar[r] &
\widehat{H^p(X, E)} \ar[r] \ar[u] &
\lim H^p(X, E_n) \ar[r] \ar[u] &
T_f(H^{p + 1}(X, E)) \ar[r] &
0 \\
0 \ar[r] &
H^0(H^p(X, E)^\wedge) \ar[r] \ar[u] &
H^p(X, E^\wedge) \ar[r] \ar[u] &
T_f(H^{p + 1}(X, E)) \ar[r] \ar@{=}[u] &
0 \\
&
R^1\lim H^p(X, E)[f^n] \ar[u] \ar[r]^\cong &
R^1\lim H^{p - 1}(X, E_n) \ar[u] \\
& 0 \ar[u] & 0 \ar[u]
}
$$
그 행과 열은 완전하다. 여기서
$\widehat{H^p(X, E)} = \lim H^p(X, E)/f^n H^p(X, E)$는 통상적인
$f$-진 완비화이고, $H^p(X, E)^\wedge$는 유도 $f$-진 완비화이며,
$T_f(H^{p + 1}(X, E))$는 $f$-진 Tate 가군이다. More on Algebra, Example
\ref{more-algebra-example-spectral-sequence-principal}을 보라.
마지막으로 $H^p(X, E^\wedge) = H^p(R\Gamma(X, E)^\wedge)$이다.
\end{lemma}

\begin{proof}
Lemma \ref{cohomology-lemma-Rf-commutes-with-Rlim}에 의해
$R\Gamma(X, E^\wedge) = R\lim R\Gamma(X, E_n)$임을 관찰하자.
한편
$$
R\Gamma(X, E_n) =
R\Gamma(X, E) \otimes_A^\mathbf{L} (A \xrightarrow{f^n} A)
$$
이다(세부 사항은 생략한다). 따라서 $R\Gamma(X, E^\wedge)$는
$R\Gamma(X, E)^\wedge$라는 유도 $f$-진 완비화이다. 그러므로 도식은
More on Algebra, Lemma
\ref{more-algebra-lemma-compare-completion-derived-completion}에서 따른다.
\end{proof}

\begin{lemma}
\label{cohomology-lemma-stabilizes}
$\mathcal{A}$를 아벨 범주라 하고, $f : M \to M$을
$\mathcal{A}$의 사상이라 하자.
$M[f^n] = \Ker(f^n : M \to M)$이 안정되면 다음 두 역계
$$
(M \xrightarrow{f^n} M)
\quad\text{및}\quad
\Coker(f^n : M \to M)
$$
는 $D(\mathcal{A})$ 안에서 프로-동형이다.
\end{lemma}

\begin{proof}
첫째 역계에서 둘째 역계로 가는 사상이 있음은 명백하다. 다음을 가정하자.
$M[f^c] = M[f^{c + 1}] = M[f^{c + 2}] = \ldots$.
그러면 다음 도식들로 $D(\mathcal{A})$ 안에서 반대 방향의 역계 사상을
정의할 수 있다.
$$
\xymatrix{
M/M[f^c] \ar[r]_-{f^{n + c}} \ar[d]_{f^c} &
M \ar[d]^1 \\
M \ar[r]^{f^n} & M
}
$$
윗 가로 화살표가 단사이므로 윗행의 복합체는
$\Coker(f^{n + c} : M \to M)$와 유사동형이다. 일부 세부 사항은
생략한다.
\end{proof}

\begin{example}
\label{cohomology-example-formal-functions-principal}
$(X, \mathcal{O}_X)$를 환 달린 공간이라 하자.
$A \to \Gamma(X, \mathcal{O}_X)$를 환 준동형이라 하고,
$f \in A$라 하자. $\mathcal{F}$를 $\mathcal{O}_X$-가군이라 하자.
모든 $n \geq c$에 대하여
$\mathcal{F}[f^c] = \mathcal{F}[f^n]$이 되게 하는 $c$가 있다고 가정하자.
$E = \mathcal{F}$로 두고 Lemma
\ref{cohomology-lemma-formal-functions-principal}을 적용한다. Lemma
\ref{cohomology-lemma-stabilizes}에 의해 역계 $(E_n)$은 역계
$(\mathcal{F}/f^n\mathcal{F})$와 프로-동형이다. 따라서
$p \in \mathbf{Z}$에 대하여 다음 가환 도식을 얻는다.
$$
\xymatrix{
& 0 & 0 \\
0 \ar[r] &
\widehat{H^p(X, \mathcal{F})} \ar[r] \ar[u] &
\lim H^p(X, \mathcal{F}/f^n\mathcal{F}) \ar[r] \ar[u] &
T_f(H^{p + 1}(X, \mathcal{F})) \ar[r] &
0 \\
0 \ar[r] &
H^0(H^p(X, \mathcal{F})^\wedge) \ar[r] \ar[u] &
H^p(R\Gamma(X, \mathcal{F})^\wedge) \ar[r] \ar[u] &
T_f(H^{p + 1}(X, \mathcal{F})) \ar[r] \ar@{=}[u] &
0 \\
&
R^1\lim H^p(X, \mathcal{F})[f^n] \ar[u] \ar[r]^\cong &
R^1\lim H^{p - 1}(X, \mathcal{F}/f^n\mathcal{F}) \ar[u] \\
& 0 \ar[u] & 0 \ar[u]
}
$$
그 행과 열은 완전하다. 여기서
$\widehat{H^p(X, \mathcal{F})} =
\lim H^p(X, \mathcal{F})/f^n H^p(X, \mathcal{F})$
는 통상적인 $f$-진 완비화이고, $D(A)$ 안의 $M$에 대하여
$M^\wedge$는 유도 $f$-진 완비화를 나타낸다.
\end{example}
















\section{비유계 복합체의 {\v C}ech 코호몰로지}
\label{cohomology-section-cech-cohomology-of-unbounded-complexes}

\noindent
Section \ref{cohomology-section-cech-cohomology-of-complexes}의 구성은 비유계
복합체에 대해서는 ‘올바른’ 구성이 아니다. 문제는 Stacks project에서
이중 복합체를 전체화할 때 직접합을 사용하므로 이를 직접곱으로
바꾸어야 한다는 데 있다. 이 절에서는 그렇게 하는 대신 덮개가
유한이라고 가정하고 교대 {\v C}ech 복합체를 사용한다.

\medskip\noindent
$(X, \mathcal{O}_X)$를 환 달린 공간이라 하자.
${\mathcal F}^\bullet$를 $\mathcal{O}_X$-가군 전층들의 복합체라 하고,
${\mathcal U} : X = \bigcup_{i \in I} U_i$를 $X$의 {\bf 유한} 열린
덮개라 하자. 교대 {\v C}ech 복합체
$\check{\mathcal{C}}_{alt}^\bullet(\mathcal{U}, \mathcal{F})$
(Section \ref{cohomology-section-alternating-cech})는 전층 $\mathcal{F}$에 대해
함자적이므로 이중 복합체
$\check{\mathcal{C}}^\bullet_{alt}(\mathcal{U}, \mathcal{F}^\bullet)$를
얻는다. 이 절에서는 이에 결부된 전체 복합체를 사용한다. 구성
$\text{Tot}(\check{\mathcal{C}}^\bullet_{alt}({\mathcal U},
{\mathcal F}^\bullet))$
은 ${\mathcal F}^\bullet$에 대해 함자적이다. 또한 다음과 같은
복합체들의 함자적 변환이 있다.
\begin{equation}
\label{cohomology-equation-global-sections-to-alternating-cech}
\Gamma(X, {\mathcal F}^\bullet)
\longrightarrow
\text{Tot}(\check{\mathcal{C}}^\bullet_{alt}({\mathcal U},
{\mathcal F}^\bullet))
\end{equation}
이는 다음 규칙으로 정의된다. 단면
$s\in \Gamma(X, {\mathcal F}^n)$은 원소
$\alpha = \{\alpha_{i_0\ldots i_p}\}$로 보내지며, 여기서
$\alpha_{i_0} = s|_{U_{i_0}}$이고 $\alpha_{i_0\ldots i_p} = 0$은
$p > 0$일 때 성립한다.

\begin{lemma}
\label{cohomology-lemma-alternating-cech-complex-complex}
$(X, \mathcal{O}_X)$를 환 달린 공간이라 하자.
$\mathcal{U} : X = \bigcup_{i \in I} U_i$를 유한 열린 덮개라 하자.
$\mathcal{O}_X$-가군의 복합체 $\mathcal{F}^\bullet$에 대해 표준 사상
$$
\text{Tot}(\check{\mathcal{C}}^\bullet_{alt}(\mathcal{U}, \mathcal{F}^\bullet))
\longrightarrow
R\Gamma(X, \mathcal{F}^\bullet)
$$
이 있으며, 이 사상은 $\mathcal{F}^\bullet$에 대해 함자적이고
(\ref{cohomology-equation-global-sections-to-alternating-cech})와 호환된다.
\end{lemma}

\begin{proof}
각 항이 단사 $\mathcal{O}_X$-가군인 K-단사 복합체
${\mathcal I}^\bullet$를 택하자.
$\mathcal{I}^\bullet$에 대한
(\ref{cohomology-equation-global-sections-to-alternating-cech})의 사상은
$\Gamma(X, {\mathcal I}^\bullet) \to
\text{Tot}(\check{\mathcal{C}}^\bullet_{alt}({\mathcal U},
{\mathcal I}^\bullet))$
이다. 이는 아벨군 복합체들의 유사동형이다. 실제로 Homology, Lemma
\ref{homology-lemma-double-complex-gives-resolution}을 이중 복합체
$\check{\mathcal{C}}^\bullet_{alt}({\mathcal U},
{\mathcal I}^\bullet)$에 적용하고 Lemmas
\ref{cohomology-lemma-injective-trivial-cech} 및 \ref{cohomology-lemma-alternating-usual}을
사용하면 이를 얻는다.
${\mathcal F}^\bullet \to {\mathcal I}^\bullet$가
${\mathcal F}^\bullet$에서 각 항이 단사대상인 K-단사 복합체로 가는
유사동형이라고 하자(Injectives, Theorem
\ref{injectives-theorem-K-injective-embedding-grothendieck}).
$R\Gamma(X, {\mathcal F}^\bullet)$는 복합체
$\Gamma(X, {\mathcal I}^\bullet)$로 나타내어지므로 다음 사상을
사용하여 보조정리의 사상을 얻는다.
$$
\text{Tot}(\check{\mathcal{C}}^\bullet_{alt}({\mathcal U},
{\mathcal F}^\bullet))
\longrightarrow
\text{Tot}(\check{\mathcal{C}}^\bullet_{alt}({\mathcal U},
{\mathcal I}^\bullet)).
$$
함자성과 호환성의 검증은 생략한다.
\end{proof}

\begin{lemma}
\label{cohomology-lemma-alternating-cech-complex-complex-ss}
$(X, \mathcal{O}_X)$를 환 달린 공간이라 하자.
$\mathcal{U} : X = \bigcup_{i \in I} U_i$를 유한 열린 덮개라 하고,
$\mathcal{F}^\bullet$를 $\mathcal{O}_X$-가군의 복합체라 하자.
$\mathcal{B}$를 $X$의 열린 부분집합들의 집합이라 하자. 다음을
가정한다.
\begin{enumerate}
\item $X$의 모든 열린집합에는 그 원소들이 $\mathcal{B}$의 원소인
덮개가 있다.
\item 모든 $i_0, \ldots, i_p \in I$에 대해
$U_{i_0\ldots i_p} \in \mathcal{B}$이다.
\item 모든 $U \in \mathcal{B}$와 $p > 0$에 대해 다음이 성립한다.
\begin{enumerate}
\item $H^p(U, \mathcal{F}^q) = 0$,
\item $H^p(U, \Coker(\mathcal{F}^{q - 1} \to \mathcal{F}^q)) = 0$, 그리고
\item $H^p(U, H^q(\mathcal{F})) = 0$.
\end{enumerate}
\end{enumerate}
그러면 Lemma \ref{cohomology-lemma-alternating-cech-complex-complex}의 사상
$$
\text{Tot}(\check{\mathcal{C}}^\bullet_{alt}(\mathcal{U}, \mathcal{F}^\bullet))
\longrightarrow
R\Gamma(X, \mathcal{F}^\bullet)
$$
은 $D(\textit{Ab})$에서 동형사상이다.
\end{lemma}

\begin{proof}
먼저 $\mathcal{F}^\bullet$가 아래로 유계라고 가정하자. 이 경우 사상
$$
\text{Tot}(\check{\mathcal{C}}^\bullet_{alt}(\mathcal{U}, \mathcal{F}^\bullet))
\longrightarrow
\text{Tot}(\check{\mathcal{C}}^\bullet(\mathcal{U}, \mathcal{F}^\bullet))
$$
은 Lemma \ref{cohomology-lemma-alternating-usual}에 의해 유사동형이다. 즉 이중
복합체들의 사상
$\check{\mathcal{C}}^\bullet_{alt}(\mathcal{U}, \mathcal{F}^\bullet) \to
\check{\mathcal{C}}^\bullet(\mathcal{U}, \mathcal{F}^\bullet)$
은 이 복합체들에 결부된 두 번째 스펙트럼열들의 첫째 쪽 사이의
동형사상을 유도하고
(Homology, Lemma \ref{homology-lemma-ss-double-complex}),
이 스펙트럼열들은 수렴한다
(Homology, Lemma \ref{homology-lemma-first-quadrant-ss}).
따라서 이 경우의 결론은 Lemma
\ref{cohomology-lemma-cech-complex-complex-computes}와 가정 (3)(a)에서 따른다.

\medskip\noindent
일반적으로 가정 (3)(c)에 의해 Lemma \ref{cohomology-lemma-K-injective}에서와
같은 분해
$\mathcal{F}^\bullet \to \mathcal{I}^\bullet = \lim \mathcal{I}_n^\bullet$
을 택할 수 있다. 그러면 보조정리의 사상은
$$
\lim_n 
\text{Tot}(\check{\mathcal{C}}^\bullet_{alt}(\mathcal{U},
\tau_{\geq -n}\mathcal{F}^\bullet))
\longrightarrow
\Gamma(X, \mathcal{I}^\bullet) =
\lim_n \Gamma(X, \mathcal{I}_n^\bullet)
$$
이 된다. 여기서 화살표는 유도범주 안의 사상이지만 오른쪽 등식은
복합체 수준에서 성립한다.
(3)(b)에 의해 $\tau_{\geq -n}\mathcal{F}^\bullet$는 보조정리의 가정을
만족하는 아래로 유계인 복합체이다. 따라서 아래로 유계인 복합체의
경우를 적용하면 각 사상
$$
\text{Tot}(\check{\mathcal{C}}^\bullet_{alt}(\mathcal{U},
\tau_{\geq -n}\mathcal{F}^\bullet))
\longrightarrow
\Gamma(X, \mathcal{I}_n^\bullet)
$$
은 유사동형이다. 왼쪽 복합체들의 주어진 차수에서의 코호몰로지는
결국 상수이다(교대 {\v C}ech 복합체가 유한이기 때문이다).
따라서 오른쪽에서도 마찬가지이다. 그러므로 오른쪽 역극한의
코호몰로지는 Homology, Lemma
\ref{homology-lemma-apply-Mittag-Leffler-again}
(또는 더 강한 More on Algebra, Lemma
\ref{more-algebra-lemma-apply-Mittag-Leffler-again})에 의해 이 상숫값이며,
이로써 보조정리가 증명된다.
\end{proof}






\section{Hom 복합체}
\label{cohomology-section-hom-complexes}

\noindent
$(X, \mathcal{O}_X)$를 환 달린 공간이라 하자.
$\mathcal{L}^\bullet$와 $\mathcal{M}^\bullet$를 두
$\mathcal{O}_X$-가군 복합체라 하자. $\mathcal{O}_X$-가군의 복합체
$\SheafHom^\bullet(\mathcal{L}^\bullet, \mathcal{M}^\bullet)$를
구성하자. 구체적으로 각 $n$에 대해 다음과 같이 둔다.
$$
\SheafHom^n(\mathcal{L}^\bullet, \mathcal{M}^\bullet) =
\prod\nolimits_{n = p + q}
\SheafHom_{\mathcal{O}_X}(\mathcal{L}^{-q}, \mathcal{M}^p)
$$
$\SheafHom^n$은 등급 $\mathcal{O}_X$-가군으로 본
$\mathcal{L}^\bullet$에서 $\mathcal{M}^\bullet$로 가는, 차수가
$n$인 모든 동차 $\mathcal{O}_X$-선형 사상들의
$\mathcal{O}_X$-가군 층으로 생각하면 좋다. 이 용어법에서 미분은
다음 규칙으로 정의한다.
$$
\text{d}(f) =
\text{d}_\mathcal{M} \circ f - (-1)^n f \circ \text{d}_\mathcal{L}
$$
여기서
$f \in \SheafHom^n_{\mathcal{O}_X}(\mathcal{L}^\bullet, \mathcal{M}^\bullet)$
이다. $\text{d}^2 = 0$의 검증은 생략한다. 이 구성은
Differential Graded Algebra, Example
\ref{dga-example-category-complexes}의 특수한 경우이다. 구성에서
곧바로 모든 $n \in \mathbf{Z}$와 모든 열린집합 $U \subset X$에 대해
\begin{equation}
\label{cohomology-equation-cohomology-hom-complex}
H^n(\Gamma(U, \SheafHom^\bullet(\mathcal{L}^\bullet, \mathcal{M}^\bullet))) =
\Hom_{K(\mathcal{O}_U)}(\mathcal{L}^\bullet, \mathcal{M}^\bullet[n])
\end{equation}
을 얻는다.

\begin{lemma}
\label{cohomology-lemma-compose}
$(X, \mathcal{O}_X)$를 환 달린 공간이라 하자.
$\mathcal{O}_X$-가군의 복합체
$\mathcal{K}^\bullet, \mathcal{L}^\bullet, \mathcal{M}^\bullet$가 주어지면
$$
\SheafHom^\bullet(\mathcal{K}^\bullet,
\SheafHom^\bullet(\mathcal{L}^\bullet, \mathcal{M}^\bullet))
=
\SheafHom^\bullet(\text{Tot}(\mathcal{K}^\bullet \otimes_{\mathcal{O}_X}
\mathcal{L}^\bullet), \mathcal{M}^\bullet)
$$
라는 $\mathcal{O}_X$-가군 복합체들의 동형사상이 있으며, 이는
$\mathcal{K}^\bullet, \mathcal{L}^\bullet, \mathcal{M}^\bullet$에
대해 함자적이다.
\end{lemma}

\begin{proof}
생략한다. 힌트: More on Algebra, Lemma
\ref{more-algebra-lemma-compose}와 정확히 같은 방법으로 증명한다.
\end{proof}

\begin{lemma}
\label{cohomology-lemma-composition}
$(X, \mathcal{O}_X)$를 환 달린 공간이라 하자.
$\mathcal{O}_X$-가군의 복합체
$\mathcal{K}^\bullet, \mathcal{L}^\bullet, \mathcal{M}^\bullet$가
주어지면 $\mathcal{O}_X$-가군 복합체들의 표준 사상
$$
\text{Tot}\left(
\SheafHom^\bullet(\mathcal{L}^\bullet, \mathcal{M}^\bullet)
\otimes_{\mathcal{O}_X}
\SheafHom^\bullet(\mathcal{K}^\bullet, \mathcal{L}^\bullet)
\right)
\longrightarrow
\SheafHom^\bullet(\mathcal{K}^\bullet, \mathcal{M}^\bullet)
$$
이 있다.
\end{lemma}

\begin{proof}
생략한다. 힌트: More on Algebra, Lemma
\ref{more-algebra-lemma-composition}과 정확히 같은 방법으로 증명한다.
\end{proof}

\begin{lemma}
\label{cohomology-lemma-diagonal-better}
$(X, \mathcal{O}_X)$를 환 달린 공간이라 하자.
$\mathcal{O}_X$-가군의 복합체
$\mathcal{K}^\bullet, \mathcal{L}^\bullet, \mathcal{M}^\bullet$가
주어지면 표준 사상
$$
\text{Tot}\left(
\mathcal{K}^\bullet \otimes_{\mathcal{O}_X}
\SheafHom^\bullet(\mathcal{M}^\bullet, \mathcal{L}^\bullet)
\right)
\longrightarrow
\SheafHom^\bullet(\mathcal{M}^\bullet,
\text{Tot}(\mathcal{K}^\bullet \otimes_{\mathcal{O}_X} \mathcal{L}^\bullet))
$$
이라는 $\mathcal{O}_X$-가군 복합체들의 사상이 있으며, 이는 세
복합체 모두에 대해 함자적이다.
\end{lemma}

\begin{proof}
생략한다. 힌트: More on Algebra, Lemma
\ref{more-algebra-lemma-diagonal-better}와 정확히 같은 방법으로
증명한다.
\end{proof}

\begin{lemma}
\label{cohomology-lemma-diagonal}
$(X, \mathcal{O}_X)$를 환 달린 공간이라 하자.
$\mathcal{O}_X$-가군의 복합체
$\mathcal{K}^\bullet, \mathcal{L}^\bullet$가 주어지면 표준 사상
$$
\mathcal{K}^\bullet
\longrightarrow
\SheafHom^\bullet(\mathcal{L}^\bullet,
\text{Tot}(\mathcal{K}^\bullet \otimes_{\mathcal{O}_X} \mathcal{L}^\bullet))
$$
이라는 $\mathcal{O}_X$-가군 복합체들의 사상이 있으며, 이는 두
복합체 모두에 대해 함자적이다.
\end{lemma}

\begin{proof}
생략한다. 힌트: More on Algebra, Lemma
\ref{more-algebra-lemma-diagonal}과 정확히 같은 방법으로 증명한다.
\end{proof}

\begin{lemma}
\label{cohomology-lemma-evaluate-and-more}
$(X, \mathcal{O}_X)$를 환 달린 공간이라 하자.
$\mathcal{O}_X$-가군의 복합체
$\mathcal{K}^\bullet, \mathcal{L}^\bullet, \mathcal{M}^\bullet$가
주어지면 표준 사상
$$
\text{Tot}(\SheafHom^\bullet(\mathcal{L}^\bullet,
\mathcal{M}^\bullet) \otimes_{\mathcal{O}_X} \mathcal{K}^\bullet)
\longrightarrow
\SheafHom^\bullet(\SheafHom^\bullet(\mathcal{K}^\bullet,
\mathcal{L}^\bullet), \mathcal{M}^\bullet)
$$
이라는 $\mathcal{O}_X$-가군 복합체들의 사상이 있으며, 이는 세
복합체 모두에 대해 함자적이다.
\end{lemma}

\begin{proof}
생략한다. 힌트: More on Algebra, Lemma
\ref{more-algebra-lemma-evaluate-and-more}와 정확히 같은 방법으로
증명한다.
\end{proof}

\begin{lemma}
\label{cohomology-lemma-RHom-into-K-injective}
$(X, \mathcal{O}_X)$를 환 달린 공간이라 하자. $L$과 $M$을
$D(\mathcal{O}_X)$의 대상이라 하자. $\mathcal{I}^\bullet$를
$M$을 나타내는 $\mathcal{O}_X$-가군의 K-단사 복합체라 하고,
$\mathcal{L}^\bullet$를 $L$을 나타내는 $\mathcal{O}_X$-가군의
복합체라 하자. 그러면 모든 열린집합 $U \subset X$에 대해
$$
H^0(\Gamma(U, \SheafHom^\bullet(\mathcal{L}^\bullet, \mathcal{I}^\bullet))) =
\Hom_{D(\mathcal{O}_U)}(L|_U, M|_U)
$$
이다.
\end{lemma}

\begin{proof}
다음과 같다.
\begin{align*}
H^0(\Gamma(U, \SheafHom^\bullet(\mathcal{L}^\bullet, \mathcal{I}^\bullet)))
& =
\Hom_{K(\mathcal{O}_U)}(\mathcal{L}^\bullet|_U, \mathcal{I}^\bullet|_U) \\
& =
\Hom_{D(\mathcal{O}_U)}(L|_U, M|_U)
\end{align*}
첫째 등식은 (\ref{cohomology-equation-cohomology-hom-complex})이다.
둘째 등식은 Lemma \ref{cohomology-lemma-restrict-K-injective-to-open}에 의해
$\mathcal{I}^\bullet|_U$가 K-단사이므로 성립한다.
\end{proof}

\begin{lemma}
\label{cohomology-lemma-RHom-well-defined}
$(X, \mathcal{O}_X)$를 환 달린 공간이라 하자.
$(\mathcal{I}')^\bullet \to \mathcal{I}^\bullet$를
$\mathcal{O}_X$-가군의 K-단사 복합체들의 유사동형이라 하고,
$(\mathcal{L}')^\bullet \to \mathcal{L}^\bullet$를
$\mathcal{O}_X$-가군 복합체들의 유사동형이라 하자. 그러면
$$
\SheafHom^\bullet(\mathcal{L}^\bullet, (\mathcal{I}')^\bullet)
\longrightarrow
\SheafHom^\bullet((\mathcal{L}')^\bullet, \mathcal{I}^\bullet)
$$
은 유사동형이다.
\end{lemma}

\begin{proof}
$M$을 $\mathcal{I}^\bullet$과 $(\mathcal{I}')^\bullet$로 나타내어지는
$D(\mathcal{O}_X)$의 대상이라 하고, $L$을
$\mathcal{L}^\bullet$과 $(\mathcal{L}')^\bullet$로 나타내어지는
$D(\mathcal{O}_X)$의 대상이라 하자. Lemma
\ref{cohomology-lemma-RHom-into-K-injective}에 의해 두 층
$$
H^0(\SheafHom^\bullet(\mathcal{L}^\bullet, (\mathcal{I}')^\bullet))
\quad\text{and}\quad
H^0(\SheafHom^\bullet((\mathcal{L}')^\bullet, \mathcal{I}^\bullet))
$$
은 모두 전층
$$
U \longmapsto \Hom_{D(\mathcal{O}_U)}(L|_U, M|_U)
$$
에 결부된 층과 같다. 따라서 이 사상은 유사동형이다.
\end{proof}

\begin{lemma}
\label{cohomology-lemma-RHom-from-K-flat-into-K-injective}
$(X, \mathcal{O}_X)$를 환 달린 공간이라 하자.
$\mathcal{I}^\bullet$를 $\mathcal{O}_X$-가군의 K-단사 복합체라 하고,
$\mathcal{L}^\bullet$를 $\mathcal{O}_X$-가군의 K-평탄 복합체라 하자.
그러면 $\SheafHom^\bullet(\mathcal{L}^\bullet, \mathcal{I}^\bullet)$는
$\mathcal{O}_X$-가군의 K-단사 복합체이다.
\end{lemma}

\begin{proof}
실제로 $\mathcal{K}^\bullet$가 $\mathcal{O}_X$-가군의 비순환
복합체이면
\begin{align*}
\Hom_{K(\mathcal{O}_X)}(\mathcal{K}^\bullet,
\SheafHom^\bullet(\mathcal{L}^\bullet, \mathcal{I}^\bullet))
& =
H^0(\Gamma(X,
\SheafHom^\bullet(\mathcal{K}^\bullet,
\SheafHom^\bullet(\mathcal{L}^\bullet, \mathcal{I}^\bullet)))) \\
& =
H^0(\Gamma(X, \SheafHom^\bullet(\text{Tot}(
\mathcal{K}^\bullet \otimes_{\mathcal{O}_X} \mathcal{L}^\bullet),
\mathcal{I}^\bullet))) \\
& =
\Hom_{K(\mathcal{O}_X)}(
\text{Tot}(\mathcal{K}^\bullet \otimes_{\mathcal{O}_X} \mathcal{L}^\bullet),
\mathcal{I}^\bullet) \\
& =
0
\end{align*}
이다. 첫째 등식은 (\ref{cohomology-equation-cohomology-hom-complex})에서 따르고,
둘째 등식은 Lemma \ref{cohomology-lemma-compose}에서 따르며, 셋째 등식은
(\ref{cohomology-equation-cohomology-hom-complex})에서 따른다. 마지막 등식은
$\mathcal{L}^\bullet$가 K-평탄이므로
$\text{Tot}(\mathcal{K}^\bullet \otimes_{\mathcal{O}_X} \mathcal{L}^\bullet)$
가 비순환이고(Definition \ref{cohomology-definition-K-flat}),
$\mathcal{I}^\bullet$가 K-단사이므로 성립한다.
\end{proof}








\section{유도범주에서의 내부 Hom}
\label{cohomology-section-internal-hom}

\noindent
$(X, \mathcal{O}_X)$를 환 달린 공간이라 하고, $L, M$을
$D(\mathcal{O}_X)$의 대상이라 하자. $D(\mathcal{O}_X)$의 임의의 세
번째 대상 $K$에 대해 표준 전단사
\begin{equation}
\label{cohomology-equation-internal-hom}
\Hom_{D(\mathcal{O}_X)}(K, R\SheafHom(L, M))
=
\Hom_{D(\mathcal{O}_X)}(K \otimes_{\mathcal{O}_X}^\mathbf{L} L, M)
\end{equation}
가 존재하도록 $D(\mathcal{O}_X)$의 대상 $R\SheafHom(L, M)$을
구성하고자 한다. Yoneda 보조정리
(Categories, Lemma \ref{categories-lemma-yoneda})에 의해 이 식은
$R\SheafHom(L, M)$을 유일한 동형사상까지 정의한다.

\medskip\noindent
그러한 대상을 구성하기 위해 $M$을 나타내는 K-단사 복합체
$\mathcal{I}^\bullet$와 $L$을 나타내는 임의의
$\mathcal{O}_X$-가군 복합체 $\mathcal{L}^\bullet$를 택한다. 다음과
같이 둔다.
$$
R\SheafHom(L, M) = \SheafHom^\bullet(\mathcal{L}^\bullet, \mathcal{I}^\bullet)
$$
여기서 오른쪽은 Section \ref{cohomology-section-hom-complexes}에서 구성한
$\mathcal{O}_X$-가군 복합체이다. 이는 Lemma
\ref{cohomology-lemma-RHom-well-defined}에 의해 잘 정의된다. 따라서 함자
$$
D(\mathcal{O}_X)^{opp} \times D(\mathcal{O}_X) \longrightarrow D(\mathcal{O}_X),
\quad
(K, L) \longmapsto R\SheafHom(K, L)
$$
를 얻는다. (\ref{cohomology-equation-internal-hom})을 증명하기에 앞서
$R\SheafHom(K, L)$의 코호몰로지 군을 계산하자.

\begin{lemma}
\label{cohomology-lemma-section-RHom-over-U}
$(X, \mathcal{O}_X)$를 환 달린 공간이라 하고, $L, M$을
$D(\mathcal{O}_X)$의 대상이라 하자. 모든 열린집합 $U$에 대해
$$
H^0(U, R\SheafHom(L, M)) =
\Hom_{D(\mathcal{O}_U)}(L|_U, M|_U)
$$
이고, 특히
$H^0(X, R\SheafHom(L, M)) = \Hom_{D(\mathcal{O}_X)}(L, M)$이다.
\end{lemma}

\begin{proof}
$M$을 나타내는 $\mathcal{O}_X$-가군의 K-단사 복합체
$\mathcal{I}^\bullet$와 $L$을 나타내는 K-평탄 복합체
$\mathcal{L}^\bullet$를 택하자. 그러면 Lemma
\ref{cohomology-lemma-RHom-from-K-flat-into-K-injective}에 의해
$\SheafHom^\bullet(\mathcal{L}^\bullet, \mathcal{I}^\bullet)$는
K-단사이다. 따라서 $U$ 위의 단면을 취하는 것만으로 $U$ 위의
코호몰로지를 계산할 수 있고, Lemma
\ref{cohomology-lemma-RHom-into-K-injective}에서 결론이 따른다.
\end{proof}

\begin{lemma}
\label{cohomology-lemma-internal-hom}
$(X, \mathcal{O}_X)$를 환 달린 공간이라 하고, $K, L, M$을
$D(\mathcal{O}_X)$의 대상이라 하자. 위에서 설명한 구성에 대해
$D(\mathcal{O}_X)$ 안의 표준 동형사상
$$
R\SheafHom(K, R\SheafHom(L, M)) =
R\SheafHom(K \otimes_{\mathcal{O}_X}^\mathbf{L} L, M)
$$
이 있다. 이는 $K, L, M$에 대해 함자적이며, $H^0(X, -)$를 취하면
(\ref{cohomology-equation-internal-hom})을 준다.
\end{lemma}

\begin{proof}
$M$을 나타내는 K-단사 복합체 $\mathcal{I}^\bullet$와 $L$을
나타내는 $\mathcal{O}_X$-가군의 K-평탄 복합체
$\mathcal{L}^\bullet$를 택하자. $\mathcal{K}^\bullet$를 $K$를
나타내는 임의의 $\mathcal{O}_X$-가군 복합체라 하자. 그러면
$$
\SheafHom^\bullet(\mathcal{K}^\bullet,
\SheafHom^\bullet(\mathcal{L}^\bullet, \mathcal{I}^\bullet))
=
\SheafHom^\bullet(
\text{Tot}(\mathcal{K}^\bullet \otimes_{\mathcal{O}_X} \mathcal{L}^\bullet),
\mathcal{I}^\bullet)
$$
이다(Lemma \ref{cohomology-lemma-compose}). 왼쪽은
$R\SheafHom(K, R\SheafHom(L, M))$을 나타내고(Lemma
\ref{cohomology-lemma-RHom-from-K-flat-into-K-injective} 사용), 오른쪽은
$R\SheafHom(K \otimes_{\mathcal{O}_X}^\mathbf{L} L, M)$을 나타낸다.
이로써 보조정리의 표시된 식이 증명된다. 대역 단면을 취하고 Lemma
\ref{cohomology-lemma-section-RHom-over-U}를 사용하면
(\ref{cohomology-equation-internal-hom})을 얻는다.
\end{proof}

\begin{lemma}
\label{cohomology-lemma-restriction-RHom-to-U}
$(X, \mathcal{O}_X)$를 환 달린 공간이라 하고, $K, L$을
$D(\mathcal{O}_X)$의 대상이라 하자. $R\SheafHom(K, L)$의 구성은 열린
부분집합으로의 제한과 가환한다. 즉 모든 열린집합 $U$에 대해
$R\SheafHom(K|_U, L|_U) = R\SheafHom(K, L)|_U$이다.
\end{lemma}

\begin{proof}
이는 구성과 Lemma \ref{cohomology-lemma-restrict-K-injective-to-open}에서
명백하다.
\end{proof}

\begin{lemma}
\label{cohomology-lemma-RHom-triangulated}
$(X, \mathcal{O}_X)$를 환 달린 공간이라 하자. 이중함자
$R\SheafHom(- , -)$는 두 변수 각각에서 특수 삼각형을 특수
삼각형으로 보낸다.
\end{lemma}

\begin{proof}
대응
$$
(\mathcal{L}^\bullet, \mathcal{M}^\bullet) \longmapsto
\SheafHom^\bullet(\mathcal{L}^\bullet, \mathcal{M}^\bullet)
$$
은 어느 한 변수에서든 복합체들의 항별 분할 짧은 완전열을 항별
분할 짧은 완전열로 보낸다는 관찰에서 결론이 따른다. 세부 사항은
생략한다.
\end{proof}

\begin{lemma}
\label{cohomology-lemma-internal-hom-composition}
$(X, \mathcal{O}_X)$를 환 달린 공간이라 하자. $D(\mathcal{O}_X)$의
$K, L, M$이 주어지면 표준 사상
$$
R\SheafHom(L, M) \otimes_{\mathcal{O}_X}^\mathbf{L} R\SheafHom(K, L)
\longrightarrow R\SheafHom(K, M)
$$
이 있으며, 이는 $D(\mathcal{O}_X)$ 안에서 $K, L, M$에 대해
함자적이다.
\end{lemma}

\begin{proof}
$M$을 나타내는 K-단사 복합체 $\mathcal{I}^\bullet$와 $L$을 나타내는
K-단사 복합체 $\mathcal{J}^\bullet$, 그리고 $K$를 나타내는 임의의
$\mathcal{O}_X$-가군 복합체 $\mathcal{K}^\bullet$를 택하자. Lemma
\ref{cohomology-lemma-composition}에 의해 복합체들의 사상
$$
\text{Tot}\left(
\SheafHom^\bullet(\mathcal{J}^\bullet, \mathcal{I}^\bullet)
\otimes_{\mathcal{O}_X}
\SheafHom^\bullet(\mathcal{K}^\bullet, \mathcal{J}^\bullet)
\right)
\longrightarrow
\SheafHom^\bullet(\mathcal{K}^\bullet, \mathcal{I}^\bullet)
$$
이 있다. $\mathcal{O}_X$-가군 복합체
$\SheafHom^\bullet(\mathcal{J}^\bullet, \mathcal{I}^\bullet)$,
$\SheafHom^\bullet(\mathcal{K}^\bullet, \mathcal{J}^\bullet)$, 그리고
$\SheafHom^\bullet(\mathcal{K}^\bullet, \mathcal{I}^\bullet)$는 각각
$R\SheafHom(L, M)$, $R\SheafHom(K, L)$, 그리고
$R\SheafHom(K, M)$을 나타낸다. K-평탄 복합체
$\mathcal{H}^\bullet$와 유사동형
$\mathcal{H}^\bullet \to
\SheafHom^\bullet(\mathcal{K}^\bullet, \mathcal{J}^\bullet)$를 택하면,
다음 사상이 있다.
$$
\text{Tot}\left(
\SheafHom^\bullet(\mathcal{J}^\bullet, \mathcal{I}^\bullet)
\otimes_{\mathcal{O}_X} \mathcal{H}^\bullet
\right)
\longrightarrow
\text{Tot}\left(
\SheafHom^\bullet(\mathcal{J}^\bullet, \mathcal{I}^\bullet)
\otimes_{\mathcal{O}_X}
\SheafHom^\bullet(\mathcal{K}^\bullet, \mathcal{J}^\bullet)
\right)
$$
이 사상의 정의역은
$R\SheafHom(L, M) \otimes_{\mathcal{O}_X}^\mathbf{L} R\SheafHom(K, L)$을
나타낸다. 표시된 두 화살표를 합성하면 원하는 사상을 얻는다. 이
구성이 함자적이라는 증명은 생략한다.
\end{proof}

\begin{lemma}
\label{cohomology-lemma-internal-hom-diagonal-better}
$(X, \mathcal{O}_X)$를 환 달린 공간이라 하자. $D(\mathcal{O}_X)$의
$K, L, M$이 주어지면 표준 사상
$$
K \otimes_{\mathcal{O}_X}^\mathbf{L} R\SheafHom(M, L)
\longrightarrow
R\SheafHom(M, K \otimes_{\mathcal{O}_X}^\mathbf{L} L)
$$
이 있으며, 이는 $D(\mathcal{O}_X)$ 안에서 $K, L, M$에 대해
함자적이다.
\end{lemma}

\begin{proof}
$K$를 나타내는 K-평탄 복합체 $\mathcal{K}^\bullet$와 $L$을 나타내는
K-단사 복합체 $\mathcal{I}^\bullet$를 택하고, $M$을 나타내는 임의의
$\mathcal{O}_X$-가군 복합체 $\mathcal{M}^\bullet$를 택하자.
$\mathcal{J}^\bullet$가 K-단사인 유사동형
$\text{Tot}(\mathcal{K}^\bullet \otimes_{\mathcal{O}_X} \mathcal{I}^\bullet)
\to \mathcal{J}^\bullet$
을 택한다. 그러면 사상
$$
\text{Tot}\left(
\mathcal{K}^\bullet \otimes_{\mathcal{O}_X}
\SheafHom^\bullet(\mathcal{M}^\bullet, \mathcal{I}^\bullet)
\right)
\to
\SheafHom^\bullet(\mathcal{M}^\bullet,
\text{Tot}(\mathcal{K}^\bullet \otimes_{\mathcal{O}_X} \mathcal{I}^\bullet))
\to
\SheafHom^\bullet(\mathcal{M}^\bullet, \mathcal{J}^\bullet)
$$
을 사용한다. 여기서 첫째 사상은 Lemma
\ref{cohomology-lemma-diagonal-better}의 사상이다.
\end{proof}

\begin{lemma}
\label{cohomology-lemma-internal-hom-diagonal}
$(X, \mathcal{O}_X)$를 환 달린 공간이라 하자. $D(\mathcal{O}_X)$의
$K, L$이 주어지면 표준 사상
$$
K \longrightarrow R\SheafHom(L, K \otimes_{\mathcal{O}_X}^\mathbf{L} L)
$$
이 있으며, 이는 $D(\mathcal{O}_X)$ 안에서 $K$와 $L$ 모두에 대해
함자적이다.
\end{lemma}

\begin{proof}
$K$를 나타내는 K-평탄 복합체 $\mathcal{K}^\bullet$와 $L$을 나타내는
임의의 $\mathcal{O}_X$-가군 복합체 $\mathcal{L}^\bullet$를 택하자.
K-단사 복합체 $\mathcal{J}^\bullet$와 유사동형
$\text{Tot}(\mathcal{K}^\bullet \otimes_{\mathcal{O}_X} \mathcal{L}^\bullet)
\to \mathcal{J}^\bullet$를 택한다. 그러면
$$
\mathcal{K}^\bullet \to
\SheafHom^\bullet(\mathcal{L}^\bullet,
\text{Tot}(\mathcal{K}^\bullet \otimes_{\mathcal{O}_X} \mathcal{L}^\bullet))
\to
\SheafHom^\bullet(\mathcal{L}^\bullet, \mathcal{J}^\bullet)
$$
을 사용한다. 여기서 첫째 사상은 Lemma \ref{cohomology-lemma-diagonal}에서 온다.
\end{proof}

\begin{lemma}
\label{cohomology-lemma-dual}
$(X, \mathcal{O}_X)$를 환 달린 공간이라 하고, $L$을
$D(\mathcal{O}_X)$의 대상이라 하자.
$L^\vee = R\SheafHom(L, \mathcal{O}_X)$라 두자. $D(\mathcal{O}_X)$의
$M$에 대해 표준 사상
\begin{equation}
\label{cohomology-equation-eval}
M \otimes^\mathbf{L}_{\mathcal{O}_X} L^\vee
\longrightarrow
R\SheafHom(L, M)
\end{equation}
이 있으며, 이는 표준 사상
$$
H^0(X, M \otimes^\mathbf{L}_{\mathcal{O}_X} L^\vee)
\longrightarrow
\Hom_{D(\mathcal{O}_X)}(L, M)
$$
을 유도한다. 이 사상은 $D(\mathcal{O}_X)$의 $M$에 대해 함자적이다.
\end{lemma}

\begin{proof}
(\ref{cohomology-equation-eval})의 사상은 식별
$M = R\SheafHom(\mathcal{O}_X, M)$을 사용하는 Lemma
\ref{cohomology-lemma-internal-hom-composition}의 특수한 경우이다.
\end{proof}

\begin{lemma}
\label{cohomology-lemma-internal-hom-evaluate}
$(X, \mathcal{O}_X)$를 환 달린 공간이라 하고, $K, L, M$을
$D(\mathcal{O}_X)$의 대상이라 하자. 표준 사상
$$
R\SheafHom(L, M) \otimes_{\mathcal{O}_X}^\mathbf{L} K
\longrightarrow
R\SheafHom(R\SheafHom(K, L), M)
$$
이 있으며, 이는 $D(\mathcal{O}_X)$ 안에서 $K, L, M$에 대해
함자적이다.
\end{lemma}

\begin{proof}
$M$을 나타내는 K-단사 복합체 $\mathcal{I}^\bullet$, $L$을 나타내는
K-단사 복합체 $\mathcal{J}^\bullet$, 그리고 $K$를 나타내는 K-평탄
복합체 $\mathcal{K}^\bullet$를 택한다. 이 사상은 다음 사상을
사용하여 정의한다.
$$
\text{Tot}(\SheafHom^\bullet(\mathcal{J}^\bullet,
\mathcal{I}^\bullet) \otimes_{\mathcal{O}_X} \mathcal{K}^\bullet)
\longrightarrow
\SheafHom^\bullet(\SheafHom^\bullet(\mathcal{K}^\bullet,
\mathcal{J}^\bullet), \mathcal{I}^\bullet)
$$
이는 Lemma \ref{cohomology-lemma-evaluate-and-more}의 사상이다. 복합체들을 이렇게
특별히 택했으므로 왼쪽은
$R\SheafHom(L, M) \otimes_{\mathcal{O}_X}^\mathbf{L} K$를 나타내고,
오른쪽은 $R\SheafHom(R\SheafHom(K, L), M)$을 나타낸다. 이것이
$D(\mathcal{O}_X)$의 세 대상 모두에 대해 함자적이라는 증명은
생략한다.
\end{proof}

\begin{remark}
\label{cohomology-remark-tensor-internal-hom}
$(X, \mathcal{O}_X)$를 환 달린 공간이라 하자. $D(\mathcal{O}_X)$의
$K, K', M, M'$에 대해 표준 사상
$$
R\SheafHom(K, K') \otimes_{\mathcal{O}_X}^\mathbf{L}
R\SheafHom(M, M')
\longrightarrow
R\SheafHom(K \otimes_{\mathcal{O}_X}^\mathbf{L} M,
K' \otimes_{\mathcal{O}_X}^\mathbf{L} M')
$$
이 있다. 즉 (\ref{cohomology-equation-internal-hom})에 의해 이 사상을 주는 것은
다음 사상을 주는 것과 같다.
$$
R\SheafHom(K, K') \otimes_{\mathcal{O}_X}^\mathbf{L}
R\SheafHom(M, M') \otimes_{\mathcal{O}_X}^\mathbf{L}
K \otimes_{\mathcal{O}_X}^\mathbf{L} M
\longrightarrow
K' \otimes_{\mathcal{O}_X}^\mathbf{L} M'
$$
이를 위해 먼저 가운데 두 인자를 뒤집고(부호 규칙은 More on Algebra,
Section \ref{more-algebra-section-sign-rules}에서와 같다), 사상
$$
R\SheafHom(K, K') \otimes_{\mathcal{O}_X}^\mathbf{L} K \to K'
\quad\text{and}\quad
R\SheafHom(M, M') \otimes_{\mathcal{O}_X}^\mathbf{L} M \to M'
$$
을 사용하면 된다. 이 사상들은 Lemma
\ref{cohomology-lemma-internal-hom-composition}에서 $K =
R\SheafHom(\mathcal{O}_X, K)$로 생각하여 얻으며, $K'$, $M$, $M'$도
마찬가지이다.
\end{remark}

\begin{remark}
\label{cohomology-remark-projection-formula-for-internal-hom}
$f : X \to Y$를 환 달린 공간의 사상이라 하고, $K, L$을
$D(\mathcal{O}_X)$의 대상이라 하자. 다음과 같은 표준 사상이 있다고
주장한다.
$$
Rf_*R\SheafHom(L, K) \longrightarrow R\SheafHom(Rf_*L, Rf_*K)
$$
즉 (\ref{cohomology-equation-internal-hom})에 의해 이는 사상
$Rf_*R\SheafHom(L, K) \otimes_{\mathcal{O}_Y}^\mathbf{L} Rf_*L \to Rf_*K$
을 주는 것과 같다. 이를 위해 합성
$$
Rf_*R\SheafHom(L, K) \otimes_{\mathcal{O}_Y}^\mathbf{L} Rf_*L \to
Rf_*(R\SheafHom(L, K) \otimes_{\mathcal{O}_X}^\mathbf{L} L) \to
Rf_*K
$$
을 사용할 수 있다. 여기서 첫째 화살표는 상대 컵곱
(Remark \ref{cohomology-remark-cup-product})이고, 둘째 화살표는 Lemma
\ref{cohomology-lemma-internal-hom-composition}에서 오는 표준 사상
$R\SheafHom(L, K) \otimes_{\mathcal{O}_X}^\mathbf{L} L \to K$에
$Rf_*$를 적용한 것이다(한 자리에 $\mathcal{O}_X$를 넣는다).
\end{remark}

\begin{remark}
\label{cohomology-remark-relative-cup-and-composition}
$h : X \to Y$를 환 달린 공간의 사상이라 하고, $K, M$을
$D(\mathcal{O}_Y)$의 대상이라 하자. 도식
$$
\xymatrix{
Rf_*R\SheafHom_{\mathcal{O}_X}(K, M)
\otimes_{\mathcal{O}_Y}^\mathbf{L} Rf_*K
\ar[r] \ar[d] &
Rf_*\left(R\SheafHom_{\mathcal{O}_X}(K, M)
\otimes_{\mathcal{O}_X}^\mathbf{L} K\right)
\ar[d] \\
R\SheafHom_{\mathcal{O}_Y}(Rf_*K, Rf_*M) \otimes_{\mathcal{O}_Y}^\mathbf{L}
Rf_*K \ar[r] &
Rf_*M
}
$$
은 가환한다. 여기서 왼쪽 세로 화살표는 Remark
\ref{cohomology-remark-projection-formula-for-internal-hom}에서 온다. 위쪽 수평
화살표는 Remark \ref{cohomology-remark-cup-product}의 사상이다. 나머지 두
화살표는 Lemma \ref{cohomology-lemma-internal-hom-composition}의 사상에서 한
항목을 $\mathcal{O}_X$ 또는 $\mathcal{O}_Y$로 바꾼 경우이다.
\end{remark}

\begin{remark}
\label{cohomology-remark-prepare-fancy-base-change}
$h : X \to Y$를 환 달린 공간의 사상이라 하고, $K, L$을
$D(\mathcal{O}_Y)$의 대상이라 하자. 표준 사상
$$
Lh^*R\SheafHom(K, L) \longrightarrow R\SheafHom(Lh^*K, Lh^*L)
$$
이 $D(\mathcal{O}_X)$ 안에 존재한다고 주장한다.
(\ref{cohomology-equation-internal-hom})은 Lemma
\ref{cohomology-lemma-internal-hom}에서 증명했으며, 이에 의해 그러한 사상을
주는 것은 사상
$$
Lh^*R\SheafHom(K, L) \otimes^\mathbf{L} Lh^*K \longrightarrow Lh^*L
$$
을 주는 것과 같다. Lemma \ref{cohomology-lemma-pullback-tensor-product}에 의해
이 화살표의 정의역은
$Lh^*(\SheafHom(K, L) \otimes^\mathbf{L} K)$이므로, 표준 사상
$$
R\SheafHom(K, L) \otimes^\mathbf{L} K \longrightarrow L.
$$
을 구성하면 충분하다. 이를 위해
$$
\text{id} :
R\SheafHom(K, L)
\longrightarrow
R\SheafHom(K, L)
$$
에 (\ref{cohomology-equation-internal-hom})으로 대응하는 화살표를 취한다.
\end{remark}

\begin{remark}
\label{cohomology-remark-fancy-base-change}
환 달린 공간들의 가환 도식
$$
\xymatrix{
X' \ar[r]_h \ar[d]_{f'} &
X \ar[d]^f \\
S' \ar[r]^g &
S
}
$$
이 주어졌다고 하자. $K, L$을 $D(\mathcal{O}_X)$의 대상이라 하자.
표준 밑변환 사상
$$
Lg^*Rf_*R\SheafHom(K, L)
\longrightarrow
R(f')_*R\SheafHom(Lh^*K, Lh^*L)
$$
이 $D(\mathcal{O}_{S'})$ 안에 존재한다고 주장한다. 구체적으로 합성
\begin{align*}
L(f')^*Lg^*Rf_*R\SheafHom(K, L)
& =
Lh^*Lf^*Rf_*R\SheafHom(K, L) \\
& \to
Lh^*R\SheafHom(K, L) \\
& \to
R\SheafHom(Lh^*K, Lh^*L)
\end{align*}
에 수반되는 사상을 취한다. 여기서 첫째 화살표는 수반 사상
$Lf^*Rf_* \to \text{id}$를 사용하고, 둘째 화살표는 Remark
\ref{cohomology-remark-prepare-fancy-base-change}에서 구성한 표준 사상이다.
\end{remark}






\section{Ext 층}
\label{cohomology-section-ext}

\noindent
$(X, \mathcal{O}_X)$를 환 달린 공간이라 하고,
$K, L \in D(\mathcal{O}_X)$라 하자. 유도범주에서의 내부 Hom 구성을
사용하고 $D(\mathcal{O}_X)$의 대상 $R\SheafHom(K, L)$의 $n$차
코호몰로지 층을 취하면 잘 정의된 $\mathcal{O}_X$-가군 층
$$
\SheafExt^n(K, L) = H^n(R\SheafHom(K, L))
$$
을 얻는다. 이 대상을 때때로
$\SheafExt^n_{\mathcal{O}_X}(K, L)$로 쓴다. Lemma
\ref{cohomology-lemma-section-RHom-over-U}에 의해 이 $\SheafExt^n$-층은 규칙
$$
U \longmapsto \Ext^n_{D(\mathcal{O}_U)}(K|_U, L|_U)
$$
의 층화임을 알 수 있다. Example \ref{cohomology-example-spectral-sequence}에
의해 항상 스펙트럼열
$$
E_2^{p, q} = H^p(X, \SheafExt^q(K, L))
$$
이 있으며, 이는 좋은 상황에서(예를 들어 $L$이 아래로 유계이고
$K$가 위로 유계이면)
$\Ext^{p + q}_{D(\mathcal{O}_X)}(K, L)$로 수렴한다.





\section{대역 유도 Hom}
\label{cohomology-section-global-RHom}

\noindent
$(X, \mathcal{O}_X)$를 환 달린 공간이라 하고,
$K, L \in D(\mathcal{O}_X)$라 하자. 유도범주에서의 내부 Hom 구성을
사용하면 $D(\Gamma(X, \mathcal{O}_X))$의 잘 정의된 대상
$$
R\Hom_X(K, L) = R\Gamma(X, R\SheafHom(K, L))
$$
을 얻는다. 이 대상을 때때로 $R\Hom_{\mathcal{O}_X}(K, L)$로 쓴다.
Lemma \ref{cohomology-lemma-section-RHom-over-U}에 의해
$$
H^0(R\Hom_X(K, L)) = \Hom_{D(\mathcal{O}_X)}(K, L),
\quad
H^p(R\Hom_X(K, L)) = \Ext_{D(\mathcal{O}_X)}^p(K, L)
$$
이다. $f : Y \to X$가 환 달린 공간의 사상이면, Remark
\ref{cohomology-remark-prepare-fancy-base-change}에서 정의한 사상의 대역 단면을
취하여 $D(\Gamma(X, \mathcal{O}_X))$ 안의 표준 사상
$$
R\Hom_X(K, L) \longrightarrow R\Hom_Y(Lf^*K, Lf^*L)
$$
을 얻는다.






\section{복합체 붙이기}
\label{cohomology-section-glueing-complexes}

\noindent
복합체를 붙일 수 있다! 더 정확히 말하면, 어떤 상황에서는
유도범주의 국소적으로 주어진 대상들을 붙여 대역 대상을 만들 수
있다. 먼저 몇 가지 쉬운 경우를 증명한 뒤, 위상공간과 열린 덮개의
맥락에서 매우 일반적인 \cite[Theorem 3.2.4]{BBD}를 증명한다.

\begin{lemma}
\label{cohomology-lemma-glue}
$(X, \mathcal{O}_X)$를 환 달린 공간이라 하고, $X = U \cup V$를
$X$의 두 열린 부분공간의 합집합이라 하자. 다음이 주어졌다고 하자.
\begin{enumerate}
\item $D(\mathcal{O}_U)$의 대상 $A$,
\item $D(\mathcal{O}_V)$의 대상 $B$, 그리고
\item 동형사상 $c : A|_{U \cap V} \to B|_{U \cap V}$.
\end{enumerate}
그러면 $D(\mathcal{O}_X)$의 대상 $F$와 동형사상
$f : F|_U \to A$, $g : F|_V \to B$가 존재하여
$c = g|_{U \cap V} \circ f^{-1}|_{U \cap V}$를 만족한다.
더 나아가 다음이 주어졌다고 하자.
\begin{enumerate}
\item $D(\mathcal{O}_X)$의 대상 $E$,
\item $D(\mathcal{O}_U)$의 사상 $a : A \to E|_U$,
\item $D(\mathcal{O}_V)$의 사상 $b : B \to E|_V$.
\end{enumerate}
이들이
$$
a|_{U \cap V}  = b|_{U \cap V} \circ c.
$$
를 만족하면, $U$로의 제한이 $a \circ f$이고 $V$로의 제한이
$b \circ g$인 $D(\mathcal{O}_X)$의 사상 $F \to E$가 존재한다.
\end{lemma}

\begin{proof}
대응하는 열린 몰입을 각각 $j_U$, $j_V$, $j_{U \cap V}$로 쓰자.
특수 삼각형
$$
F \to Rj_{U, *}A \oplus Rj_{V, *}B \to Rj_{U \cap V, *}(B|_{U \cap V})
\to F[1]
$$
을 택하되, 사상 $Rj_{V, *}B \to
Rj_{U \cap V, *}(B|_{U \cap V})$는 자명한 사상으로 하고,
$Rj_{U, *}A \to Rj_{U \cap V, *}(B|_{U \cap V})$는
$Rj_{U, *}A \to Rj_{U \cap V, *}(A|_{U \cap V})$와
$Rj_{U \cap V, *}c$의 합성으로 한다. $U$로 제한하면
$$
F|_U \to A \oplus (Rj_{V, *}B)|_U \to (Rj_{U \cap V, *}(B|_{U \cap V}))|_U
\to F|_U[1]
$$
를 얻는다. $j : U \cap V \to U$라 쓰자. 열린집합으로의 제한과
코호몰로지의 호환성에 의해
$(Rj_{V, *}B)|_U$와 $(Rj_{U \cap V, *}(B|_{U \cap V}))|_U$는 모두
$Rj_*(B|_{U \cap V})$와 표준적으로 동형이다. 따라서 마지막 표시
도식의 둘째 화살표에는 단면이 있으므로 사상 $F|_U \to A$가
동형사상임을 얻는다. 마찬가지로 $F|_V \to B$도 동형사상이다.
사상 $F \to E$의 존재는 $\Hom$에 대한 Mayer--Vietoris 열에서
따른다. Lemma \ref{cohomology-lemma-mayer-vietoris-hom}를 보라.
\end{proof}

\begin{lemma}
\label{cohomology-lemma-vanishing-and-glueing}
$f : (X, \mathcal{O}_X) \to (Y, \mathcal{O}_Y)$를 환 달린 공간의
사상이라 하고, $\mathcal{B}$를 $Y$의 위상의 기저라 하자.
\begin{enumerate}
\item $K$가 $D(\mathcal{O}_X)$의 대상이고, $V \in \mathcal{B}$이면
$i < 0$에 대해 $H^i(f^{-1}(V), K) = 0$이라고 가정하자. 그러면
$Rf_*K$의 음의 차수 코호몰로지 층은 사라지고, 모든 열린집합
$V \subset Y$와 $i < 0$에 대해 $H^i(f^{-1}(V), K) = 0$이며,
규칙 $V \mapsto H^0(f^{-1}V, K)$는 $Y$ 위의 층이다.
\item $K, L$이 $D(\mathcal{O}_X)$의 대상이고, $V \in \mathcal{B}$이면
$i < 0$에 대해
$\Ext^i(K|_{f^{-1}V}, L|_{f^{-1}V}) = 0$이라고 가정하자. 그러면 모든
열린집합 $V \subset Y$와 $i < 0$에 대해
$\Ext^i(K|_{f^{-1}V}, L|_{f^{-1}V}) = 0$이고, 규칙
$V \mapsto \Hom(K|_{f^{-1}V}, L|_{f^{-1}V})$는 $Y$ 위의 층이다.
\end{enumerate}
\end{lemma}

\begin{proof}
Lemma \ref{cohomology-lemma-unbounded-describe-higher-direct-images}에 의해
$H^i(Rf_*K)$는 전층
$V \mapsto H^i(f^{-1}(V), K) = H^i(V, Rf_*K)$에 결부된 층이다.
(1)의 가정에 의해 $Rf_*K$의 $< 0$차 코호몰로지 층은 사라진다.
따라서 임의의 열린집합 $V \subset Y$에 대해 코호몰로지 군
$H^i(V, Rf_*K)$는 $i < 0$이면 영이고, $i = 0$이면
$H^0(V, H^0(Rf_*K))$와 같다. 이로써 (1)이 증명된다.

\medskip\noindent
(2)를 증명하려면 Lemma \ref{cohomology-lemma-section-RHom-over-U}를 사용하여
(1)을 복합체 $R\SheafHom(K, L)$에 적용하면 된다.
\end{proof}

\begin{situation}
\label{cohomology-situation-locally-given}
$(X, \mathcal{O}_X)$를 환 달린 공간이라 하자. 다음이 주어져 있다.
\begin{enumerate}
\item $X$의 열린집합들의 모음 $\mathcal{B}$,
\item $U \in \mathcal{B}$에 대해 $D(\mathcal{O}_U)$의 대상 $K_U$,
\item $V \subset U$이고 $V, U \in \mathcal{B}$일 때
$D(\mathcal{O}_V)$의 동형사상
$\rho^U_V : K_U|_V \to K_V$.
\end{enumerate}
이들은 $U, V, W$가 $\mathcal{B}$에 속하고 $W \subset V \subset U$이면
$\rho^U_W = \rho^V_W \circ \rho ^U_V|_W$를 만족한다.
\end{situation}

\noindent
가정을 더 하지 않고서는 이 자료에 관해 아무것도 증명할 수 없다.
흥미로운 한 경우는 $\mathcal{B}$가 $X$의 위상의 기저인 경우이다.
또 다른 경우는 위상공간의 사상 $f : X \to Y$가 있고
$\mathcal{B}$의 원소들이 $Y$의 위상의 한 기저에 속하는 원소들의
역상인 경우이다.

\medskip\noindent
Situation \ref{cohomology-situation-locally-given}에서 {\it 해}란 쌍
$(K, \rho_U)$을 뜻한다. 여기서 $K$는 $D(\mathcal{O}_X)$의 대상이고,
$\rho_U : K|_U \to K_U$, $U \in \mathcal{B}$는 동형사상이며, 모든
$V \subset U$, $U, V \in \mathcal{B}$에 대해
$\rho^U_V \circ \rho_U|_V = \rho_V$를 만족한다. 어떤 경우에는 해가
유일하다.

\begin{lemma}
\label{cohomology-lemma-uniqueness}
Situation \ref{cohomology-situation-locally-given}에서 다음을 가정하자.
\begin{enumerate}
\item $X = \bigcup_{U \in \mathcal{B}} U$이고, $U, V \in \mathcal{B}$이면
$U \cap V = \bigcup_{W \in \mathcal{B}, W \subset U \cap V} W$이다.
\item 임의의 $U \in \mathcal{B}$와 $i < 0$에 대해
$\Ext^i(K_U, K_U) = 0$이다.
\end{enumerate}
해 $(K, \rho_U)$가 존재하면 이는 유일한 동형사상까지 유일하고,
더 나아가 $i < 0$에 대해 $\Ext^i(K, K) = 0$이다.
\end{lemma}

\begin{proof}
$(K, \rho_U)$와 $(K', \rho'_U)$를 두 해라 하자. Topology, Lemma
\ref{topology-lemma-create-map-from-subcollection}에서 구성한 연속사상을
$f : X \to Y$라 하고, $\mathcal{O}_Y = f_*\mathcal{O}_X$로 두자.
그러면 $K, K'$와 $\mathcal{B}$는 Lemma
\ref{cohomology-lemma-vanishing-and-glueing} (2)의 상황에 놓인다. 따라서 $K$의
음의 Ext가 사라지고, 규칙
$$
V \longmapsto \Hom(K|_{f^{-1}V}, K'|_{f^{-1}V})
$$
이 $Y$ 위의 층임을 알 수 있다. $(K, \rho_U)$와
$(K', \rho'_U)$가 모두 해이므로 $U = f^{-1}(f(U))$ 위의 사상
$$
(\rho'_U)^{-1} \circ \rho_U : K|_U \longrightarrow K'|_U
$$
들은 겹침에서 일치한다. 따라서 위 층의 유일한 대역 단면을 얻으며,
이는 주어진 모든 구조와 호환되는 원하는 동형사상 $K \to K'$를
정의한다.
\end{proof}

\begin{remark}
\label{cohomology-remark-uniqueness}
Lemma \ref{cohomology-lemma-uniqueness}의 기호와 가정을 사용하자.
$U, V \in \mathcal{B}$라고 하고, $\mathcal{B}'$를
$U \cap V$에 포함되는 $\mathcal{B}$의 원소들의 집합이라 하자. 그러면
$$
(\{K_{U'}\}_{U' \in \mathcal{B}'},
\{\rho_{V'}^{U'}\}_{V' \subset U'\text{ with }U', V' \in \mathcal{B}'})
$$
은 환 달린 공간 $U \cap V$ 위에서 Lemma
\ref{cohomology-lemma-uniqueness}의 가정을 만족하는 계이다. 또한
$(K_U|_{U \cap V}, \rho^U_{U'})$와
$(K_V|_{U \cap V}, \rho^V_{U'})$는 모두 이 계의 해이다. 보조정리에
의해 유일한 동형사상
$$
\rho_{U, V} : K_U|_{U \cap V} \longrightarrow K_V|_{U \cap V}
$$
을 얻는데, 이는 모든 $U' \subset U \cap V$,
$U' \in \mathcal{B}$에 대해 도식
$$
\xymatrix{
K_U|_{U'} \ar[rr]_{\rho_{U, V}|_{U'}} \ar[rd]_{\rho^U_{U'}} & &
K_V|_{U'} \ar[ld]^{\rho^V_{U'}} \\
& K_{U'}
}
$$
이 가환하도록 하는 것이다. 세 번째 원소 $W \in \mathcal{B}$를
택하자. $\rho_{U, V}$와 같은 성질을 만족하는 동형사상
$\rho_{U, W} : K_U|_{U \cap W} \to K_W|_{U \cap W}$와
$\rho_{V, W} : K_U|_{V \cap W} \to K_W|_{V \cap W}$를 얻는다.
마지막으로
$$
\rho_{U, W}|_{U \cap V \cap W} =
\rho_{V, W}|_{U \cap V \cap W} \circ \rho_{U, V}|_{U \cap V \cap W}
$$
이다. 등식의 양변은 모든 $U'' \subset U \cap V \cap W$,
$U'' \in \mathcal{B}$에 대한 사상 $\rho^U_{U''}$와
$\rho^W_{U''}$에 호환되는 유일한 동형사상이므로, 이는 보조정리의
유일성에서 따른다. 사소한 세부 사항은 생략한다. 모음
$(K_U, \rho_{U, V})$는 $X$의 열린 덮개
$\mathcal{U} : X = \bigcup_{U \in \mathcal{B}} U$에 대한 유도범주의
강하 데이터이다. 이 언어에서 해의 존재를 찾는 것은 ‘강하 데이터의
유효성’을 찾는 것이다.
\end{remark}

\begin{lemma}
\label{cohomology-lemma-solution-in-finite-case}
Situation \ref{cohomology-situation-locally-given}에서 다음을 가정하자.
\begin{enumerate}
\item $U_i \in \mathcal{B}$인
$X = U_1 \cup \ldots \cup U_n$이다.
\item $U, V \in \mathcal{B}$이면
$U \cap V = \bigcup_{W \in \mathcal{B}, W \subset U \cap V} W$이다.
\item 임의의 $U \in \mathcal{B}$와 $i < 0$에 대해
$\Ext^i(K_U, K_U) = 0$이다.
\end{enumerate}
그러면 해가 존재하며 유일한 동형사상까지 유일하다.
\end{lemma}

\begin{proof}
유일성은 Lemma \ref{cohomology-lemma-uniqueness}에서 보았다. $n$에 대한
귀납법으로 보조정리를 증명하자. $n = 1$인 경우는 자명하다.

\medskip\noindent
$n = 2$인 경우. Remark \ref{cohomology-remark-uniqueness}에서 구성한 동형사상
$\rho_{U_1, U_2} : K_{U_1}|_{U_1 \cap U_2} \to K_{U_2}|_{U_1 \cap U_2}$
을 생각하자. Lemma \ref{cohomology-lemma-glue}에 의해 $D(\mathcal{O}_X)$의 대상
$K$와 $\rho_{U_1, U_2}$에 호환되는 동형사상
$\rho_{U_1} : K|_{U_1} \to K_{U_1}$ 및
$\rho_{U_2} : K|_{U_2} \to K_{U_2}$를 얻는다.
$U \in \mathcal{B}$를 택하자. 동형사상 $\rho_U : K|_U \to K_U$를
구성하고 $(K, \rho_U)$가 해라는 검증은 독자에게 맡긴다.
$\mathcal{B}'$를 $U \cap U_1$ 또는 $U \cap U_2$에 포함되는
$\mathcal{B}$의 원소들의 집합이라 하자. 그러면
$(K_U, \rho^U_{U'})$는 환 달린 공간 $U$ 위의 계
$(\{K_{U'}\}_{U' \in \mathcal{B}'},
\{\rho_{V'}^{U'}\}_{V' \subset U'\text{ with }U', V' \in \mathcal{B}'})$
의 해이다. $(K|_U, \tau_{U'})$도 또 하나의 해라고 주장한다. 여기서
$U' \in \mathcal{B}'$에 대한 $\tau_{U'}$는 다음과 같이 택한다.
$U' \subset U_1$이면 합성
$$
K|_{U'} \xrightarrow{\rho_{U_1}|_{U'}}
K_{U_1}|_{U'} \xrightarrow{\rho^{U_1}_{U'}}
K_{U'}
$$
을 취하고, $U' \subset U_2$이면 합성
$$
K|_{U'} \xrightarrow{\rho_{U_2}|_{U'}}
K_{U_2}|_{U'} \xrightarrow{\rho^{U_2}_{U'}}
K_{U'}.
$$
을 취한다. 이것이 해임을 검증하려면 Remark
\ref{cohomology-remark-uniqueness}에 기술된 사상 $\rho_{U_1, U_2}$의 성질과
$\rho_{U_1}$ 및 $\rho_{U_2}$가 $\rho_{U_1, U_2}$와 호환됨을
사용한다. 이제 Lemma \ref{cohomology-lemma-uniqueness}를 적용하면
$U' \in \mathcal{B}'$에 대한 사상 $\tau_{U'}$와
$\rho^U_{U'}$에 호환되는 유일한 동형사상
$K|_{U'} \to K_{U'}$를 얻는다.

\medskip\noindent
$n > 2$인 경우. 열린 부분공간
$X' = U_1 \cup \ldots \cup U_{n - 1}$를 생각하고,
$\mathcal{B}'$를 $X'$에 포함되는 $\mathcal{B}$의 원소들의 집합이라
하자. 그러면 환 달린 공간 $X'$ 위의 계
$(\{K_U\}_{U \in \mathcal{B}'}, \{\rho_V^U\}_{U, V \in \mathcal{B}'})$
를 얻으며, 여기에 귀납 가설을 적용할 수 있다. 해
$(K_{X'}, \rho^{X'}_U)$를 얻는다. 이제 $X$의 열린집합들의 모음
$\mathcal{B}^* = \mathcal{B} \cup \{X'\}$를 생각하면 계
$(\{K_U\}_{U \in \mathcal{B}^*},
\{\rho_V^U\}_{V \subset U\text{ with }U, V \in \mathcal{B}^*})$
를 얻는다. 해 $K_{X'}$에 Lemma \ref{cohomology-lemma-uniqueness}를 적용하면
이 새 계도 조건 (3)을 만족함을 알 수 있다. 이 계에 대해서는
$X = X' \cup U_n$이다. 따라서 위에서 다룬 $n = 2$인 경우로
귀착된다.
\end{proof}

\begin{lemma}
\label{cohomology-lemma-glueing-increasing-union}
$X$를 환 달린 공간이라 하자. $E$를 정렬집합이라 하고
$$
X = \bigcup\nolimits_{\alpha \in E} W_\alpha
$$
를 열린 덮개라 하자. 이때 $W_\alpha \subset W_{\alpha + 1}$이고,
$\alpha$가 후속원이 아니면
$W_\alpha = \bigcup_{\beta < \alpha} W_\beta$라고 하자.
$K_\alpha$를 $D(\mathcal{O}_{W_\alpha})$의 대상으로서
$i < 0$에 대해 $\Ext^i(K_\alpha, K_\alpha) = 0$인 것이라 하자.
모든 $\beta < \alpha$에 대해 $D(\mathcal{O}_{W_\beta})$의
동형사상
$\rho_\beta^\alpha :  K_\alpha|_{W_\beta} \to K_\beta$가 주어져
있고, $\gamma < \beta < \alpha$이면
$\rho_\gamma^\alpha = \rho_\gamma^\beta \circ \rho^\alpha_\beta|_{W_\gamma}$
를 만족한다고 가정하자. 그러면 $D(\mathcal{O}_X)$의 대상 $K$와
모든 $\alpha \in E$에 대한 동형사상
$K|_{W_\alpha} \to K_\alpha$가 존재하며, 이들은 동형사상
$\rho_\beta^\alpha$와 호환된다.
\end{lemma}

\begin{proof}
이 증명에서 $\alpha, \beta, \gamma, \ldots$는 $E$의 원소를 나타낸다.
$W_\alpha$ 위에서 $K_\alpha$를 나타내는 K-단사 복합체
$I_\alpha^\bullet$를 택하자. $\beta < \alpha$이면 포함사상을
$j_{\beta, \alpha} : W_\beta \to W_\alpha$라 쓰자. 초한 재귀를
사용하여 모든 $\beta < \alpha$에 대해 복합체들의 사상
$$
\tau_{\beta, \alpha} :
(j_{\beta, \alpha})_!I_\beta^\bullet
\longrightarrow
I_\alpha^\bullet
$$
을 구성하되, 이는 동형사상
$\rho^\alpha_\beta : K_\alpha|_{W_\beta} \to K_\beta$의 역에
수반되는 사상을 나타내도록 한다. 더 나아가
$\gamma < \beta < \alpha$이면 복합체들의 사상으로서
$$
\tau_{\gamma, \alpha} = \tau_{\beta, \alpha} \circ
(j_{\beta, \alpha})_!\tau_{\gamma, \beta}
$$
가 되도록 구성한다. 실제로 모든 $\gamma < \beta < \alpha$에 대해
올바르게 합성되는 $\tau_{\gamma, \beta}$가 이미 주어졌다고 하자.
$\alpha = \alpha' + 1$이 후속원이면, 동형사상
$\rho^\alpha_{\alpha'} : K_\alpha|_{W_{\alpha'}} \to K_{\alpha'}$의
역에 수반되는 복합체들의 임의의 사상
$$
(j_{\alpha', \alpha})_!I_{\alpha'}^\bullet \to I_\alpha^\bullet
$$
을 택할 수 있다($I_\alpha^\bullet$가 K-단사이므로 가능하다).
임의의 $\beta < \alpha'$에 대해 다음과 같이 둔다.
$$
\tau_{\beta, \alpha} = \tau_{\alpha', \alpha} \circ
(j_{\alpha', \alpha})_!\tau_{\beta, \alpha'}
$$
$\alpha$가 후속원이 아니면 $W_\alpha$ 위의 복합체
$$
C^\bullet = \colim_{\beta < \alpha} (j_{\beta, \alpha})_!I_\beta^\bullet
$$
를 생각한다. 이는 항별 여극한이며, 이 계의 전이 사상은
$\beta' < \beta < \alpha$에 대한 사상
$\tau_{\beta', \beta}$로 주어진다. $C^\bullet$가 $K_\alpha$를
나타낸다고 주장한다. 실제로 $\beta < \alpha$에 대해 여곱 사상
$(j_{\beta, \alpha})_!I_\beta^\bullet \to C^\bullet$을 제한하면 사상
$$
\sigma_\beta : I_\beta^\bullet \longrightarrow C^\bullet|_{W_\beta}
$$
을 얻는데, 이는 유사동형이다. 즉 $x \in W_\beta$이면 줄기를 보아
$$
(C^\bullet)_x =
\colim_{\beta' < \alpha}
\left((j_{\beta', \alpha})_!I_{\beta'}^\bullet\right)_x =
\colim_{\beta \leq \beta' < \alpha} (I_{\beta'}^\bullet)_x
\longleftarrow
(I_\beta^\bullet)_x
$$
를 얻으며, 이는 유사동형이다. 여기서는 줄기를 취하는 것이
여극한과 가환하고, 여과 여극한이 완전하며, 모든
$\beta \leq \beta' < \alpha$에 대해
$(I_\beta^\bullet)_x \to (I_{\beta'}^\bullet)_x$가 유사동형임을
사용했다. 따라서 $(C^\bullet, \sigma_\beta^{-1})$는 계
$(\{K_\beta\}_{\beta < \alpha},
\{\rho^\beta_{\beta'}\}_{\beta' < \beta < \alpha})$의 해이다.
$(K_\alpha, \rho^\alpha_\beta)$도 또 하나의 해이므로, 음의 Ext 군이
사라짐을 사용하는 Lemma \ref{cohomology-lemma-solution-in-finite-case}에 의해
모든 사상과 호환되는 $D(\mathcal{O}_{W_\alpha})$의 유일한 동형사상
$\sigma : K_\alpha \to C^\bullet$을 얻는다. $\sigma$를 나타내는
복합체들의 사상 $\tau : C^\bullet \to I_\alpha^\bullet$를 택하고,
다음과 같이 두어 원하는 사상을 얻는다.
$$
\tau_{\beta, \alpha} = \tau|_{W_\beta} \circ \sigma_\beta
$$
마지막으로 복합체
$$
K^\bullet = \colim_{\alpha \in E} (W_\alpha \to X)_!I_\alpha^\bullet
$$
로 나타내어지는 유도범주의 대상을 $K$로 택한다. 여기서 전이 사상은
모든 $\beta < \alpha$에 대해 세심하게 구성한 사상
$\tau_{\beta, \alpha}$로 주어진다. 위와 정확히 같은 논증으로 모든
$\alpha$에 대해 여곱 사상의 제한이 동형사상
$$
K|_{W_\alpha} \longrightarrow K_\alpha
$$
을 결정하며, 이는 주어진 사상 $\rho^\alpha_\beta$와 호환됨을 알 수
있다.
\end{proof}

\noindent
초한 귀납법을 사용하면 일반적인 경우의 결과를 증명할 수 있다.

\begin{theorem}[BBD 붙이기 보조정리]
\label{cohomology-theorem-glueing-bbd-general}
\begin{reference}
유계성 가정이 없는 \cite[Theorem 3.2.4]{BBD}의 특수한 경우이다.
\end{reference}
Situation \ref{cohomology-situation-locally-given}에서 다음을 가정하자.
\begin{enumerate}
\item $X = \bigcup_{U \in \mathcal{B}} U$이다.
\item $U, V \in \mathcal{B}$이면
$U \cap V = \bigcup_{W \in \mathcal{B}, W \subset U \cap V} W$이다.
\item 임의의 $U \in \mathcal{B}$와 $i < 0$에 대해
$\Ext^i(K_U, K_U) = 0$이다.
\end{enumerate}
그러면 $D(\mathcal{O}_X)$의 대상 $K$와 모든
$U \in \mathcal{B}$에 대한 $D(\mathcal{O}_U)$의 동형사상
$\rho_U : K|_U \to K_U$가 존재하여, $V \subset U$이고
$U, V \in \mathcal{B}$이면
$\rho^U_V \circ \rho_U|_V = \rho_V$를 만족한다. 쌍
$(K, \rho_U)$은 유일한 동형사상까지 유일하다.
\end{theorem}

\begin{proof}
위 글에서 쌍 $(K, \rho_U)$을 해라고 불렀다. 유일성은 Lemma
\ref{cohomology-lemma-uniqueness}에서 따른다. $X$가 $\mathcal{B}$의 원소들로
이루어진 유한 덮개를 가지면(예를 들어 $X$가 준콤팩트이면), 정리는
Lemma \ref{cohomology-lemma-solution-in-finite-case}의 귀결이다. 일반적인
경우에는 초한 귀납법과 Lemma
\ref{cohomology-lemma-glueing-increasing-union}을 사용하여 정확히 같은 방식으로
논증한다.

\medskip\noindent
먼저 초한 재귀를 사용하여 각 순서수 $\alpha$에 대해 열린집합
$W_\alpha \subset X$를 택한다. 구체적으로
$W_0 = \emptyset$로 둔다. $\alpha = \beta + 1$이 후속 순서수이면,
$W_\beta = X$일 때는 $W_\alpha = X$로 두고,
$W_\beta \not = X$일 때는 $W_\beta$에 포함되지 않는
$U_\alpha \in \mathcal{B}$를 택하여
$W_\alpha = W_\beta \cup U_\alpha$로 둔다. $\alpha$가 극한
순서수이면 $W_\alpha = \bigcup_{\beta < \alpha} W_\beta$로 둔다.
충분히 큰 $\alpha$에 대해 $W_\alpha = X$이다. 모든 $\alpha$에 대해
열린집합 $W_\alpha$는 $\mathcal{B}$의 원소들의 합집합이다. 따라서
$\mathcal{B}_\alpha = \{U \in \mathcal{B}, U \subset W_\alpha\}$라
두면
$$
S_\alpha = (\{K_U\}_{U \in \mathcal{B}_\alpha},
\{\rho_V^U\}_{V \subset U\text{ with }U, V \in \mathcal{B}_\alpha})
$$
는 환 달린 공간 $W_\alpha$ 위에서 Lemma
\ref{cohomology-lemma-uniqueness}의 상황에 놓인 계이다.

\medskip\noindent
초한 귀납법으로 모든 $\alpha$에 대해 계 $S_\alpha$가 해를 가짐을
보이자. 충분히 큰 $\alpha$에 대해서는 이 계가 정리에 주어진 계이므로,
이로써 정리가 증명된다.

\medskip\noindent
$\alpha = \beta + 1$이 후속 순서수인 경우.
(이 경우는 위 보조정리의 증명에서 이미 다루었지만, 명확성을 위해
논증을 반복한다.) 이때 어떤 $U_\alpha \in \mathcal{B}$에 대해
$W_\alpha = W_\beta \cup U_\alpha$임을 상기하자. 귀납 가설에 의해
계 $S_\beta$의 해
$(K_{W_\beta}, \{\rho^{W_\beta}_U\}_{U \in \mathcal{B}_\beta})$
가 있다. 이제 $W_\alpha$의 열린집합들의 모음
$\mathcal{B}_\alpha^* = \mathcal{B}_\alpha \cup \{W_\beta\}$를
생각하면 계
$(\{K_U\}_{U \in \mathcal{B}_\alpha^*},
\{\rho_V^U\}_{V \subset U\text{ with }U, V \in \mathcal{B}_\alpha^*})$
를 얻는다. 해 $K_{W_\beta}$에 Lemma
\ref{cohomology-lemma-uniqueness}를 적용하면 이 새 계도 조건 (3)을 만족한다.
이 계에 대해 $W_\alpha = W_\beta \cup U_\alpha$이다. 따라서 Lemma
\ref{cohomology-lemma-solution-in-finite-case}에서 다룬 경우로 귀착된다.

\medskip\noindent
$\alpha$가 극한 순서수인 경우. 이때
$W_\alpha = \bigcup_{\beta < \alpha} W_\beta$임을 상기하자.
$\beta < \alpha$이면
$(K_{W_\beta}, \{\rho^{W_\beta}_U\}_{U \in \mathcal{B}_\beta})$
를 $S_\beta$의 해라 하자. $\gamma < \beta < \alpha$이면 사상
$\rho^{W_\beta}_U$, $U \in \mathcal{B}_\gamma$를 갖춘 제한
$K_{W_\beta}|_{W_\gamma}$는 $S_\gamma$의 해이다. 유일성에 의해
$U \in \mathcal{B}_\gamma$에 대한 사상
$\rho^{W_\beta}_U$와 $\rho^{W_\gamma}_U$에 호환되는 유일한 동형사상
$\rho_{W_\gamma}^{W_\beta} : K_{W_\beta}|_{W_\gamma} \to K_{W_\gamma}$
을 얻는다. 이 사상들은 올바르게 합성된다. 즉
$\delta < \gamma < \beta < \alpha$이면
$\rho_{W_\delta}^{W_\gamma} \circ \rho_{W_\gamma}^{W_\beta}|_{W_\delta}
= \rho^{W_\delta}_{W_\beta}$이다. 따라서 Lemma
\ref{cohomology-lemma-glueing-increasing-union}을 적용할 수 있다. 해
$K_{W_\beta}$에 Lemma \ref{cohomology-lemma-uniqueness}를 적용하면
$K_{W_\beta}$의 음의 Ext가 사라진다는 점에 유의하자. 그 결과
$K_{W_\alpha}$와 동형사상
$$
\rho_{W_\beta}^{W_\alpha} :
K_{W_\alpha}|_{W_\beta}
\longrightarrow
K_{W_\beta}
$$
을 얻으며, 이들은 모든 $\gamma < \beta < \alpha$에 대한 사상
$\rho_{W_\gamma}^{W_\beta}$와 호환된다.

\medskip\noindent
$K_{W_\alpha}$가 해임을 보이려면, 여전히
$U \in \mathcal{B}_\alpha$에 대한 동형사상
$\rho_U^{W_\alpha} : K_{W_\alpha}|_U \to K_U$를 구성하여 일정한
호환성을 만족함을 보여야 한다. 임의의 $\beta < \alpha$와
$V \in \mathcal{B}_\beta$, $V \subset U$에 대해 도식
$$
\xymatrix{
K_{W_\alpha}|_V \ar[r]_{\rho_U^{W_\alpha}|_V}
\ar[d]_{\rho_{W_\beta}^{W_\alpha}|_V}
& K_U|_V \ar[d]^{\rho_U^V} \\
K_{W_\beta} \ar[r]^{\rho_V^{W_\beta}}
& K_V
}
$$
이 가환하도록 하는 유일한 사상으로 $\rho_U^{W_\alpha}$를 택한다.
이 선택이 가능한 까닭은
$$
(\{K_V\}_{V \subset U, V \in \mathcal{B}_\beta\text{ for some }\beta < \alpha},
\{\rho_V^{V'}\}_{V \subset V'\text{ with }V, V' \subset U
\text{ and }V, V' \in \mathcal{B}_\beta\text{ for some }\beta < \alpha})
$$
가 환 달린 공간 $U$ 위에서 Lemma \ref{cohomology-lemma-uniqueness}의 상황에
놓인 계이고, $(K_U, \rho^U_V)$와
$(K_{W_\alpha}|_U,  \rho_V^{W_\beta}\circ \rho_{W_\beta}^{W_\alpha}|_V)$가
모두 이 계의 해이기 때문이다. 이로써 존재성과 유일성을 얻는다.
이 사상들이 원하는 호환성을 만족한다는 증명은 생략한다. 이는 단순히
지표를 추적하는 일이다.
\end{proof}




\section{강한 완전 복합체}
\label{cohomology-section-strictly-perfect}

\noindent
가군의 강한 완전 복합체는 뒤에서 유사코히런트 복합체와 완전
복합체의 개념을 정의하는 데 사용된다. 그 정의는 다음과 같다.

\begin{definition}
\label{cohomology-definition-strictly-perfect}
$(X, \mathcal{O}_X)$를 환 달린 공간이라 하고,
$\mathcal{E}^\bullet$를 $\mathcal{O}_X$-가군의 복합체라 하자. 유한
개를 제외한 모든 $i$에 대해 $\mathcal{E}^i$가 영이고, 모든 $i$에
대해 $\mathcal{E}^i$가 유한 자유 $\mathcal{O}_X$-가군의 직합인자이면
$\mathcal{E}^\bullet$를 {\it 강한 완전}이라고 한다.
\end{definition}

\noindent
주의: $X$가 국소환 달린 공간이라고 가정하지 않으므로, 유한 자유
$\mathcal{O}_X$-가군의 직합인자가 유한 국소자유일 필요는 없다.

\begin{lemma}
\label{cohomology-lemma-cone}
강한 완전 복합체들의 사상의 원뿔은 강한 완전이다.
\end{lemma}

\begin{proof}
정의에서 바로 따른다.
\end{proof}

\begin{lemma}
\label{cohomology-lemma-tensor}
두 강한 완전 복합체의 텐서곱에 결부된 전체 복합체는 강한 완전이다.
\end{lemma}

\begin{proof}
생략한다.
\end{proof}

\begin{lemma}
\label{cohomology-lemma-strictly-perfect-pullback}
$f : (X, \mathcal{O}_X) \to (Y, \mathcal{O}_Y)$를 환 달린 공간의
사상이라 하자. $\mathcal{F}^\bullet$가 $\mathcal{O}_Y$-가군의 강한
완전 복합체이면 $f^*\mathcal{F}^\bullet$는
$\mathcal{O}_X$-가군의 강한 완전 복합체이다.
\end{lemma}

\begin{proof}
유한 자유 가군의 당김은 유한 자유이다. 함자 $f^*$는 가법적이므로
직합인자를 보존한다. 따라서 보조정리가 따른다.
\end{proof}

\begin{lemma}
\label{cohomology-lemma-local-lift-map}
$(X, \mathcal{O}_X)$를 환 달린 공간이라 하자.
$\mathcal{O}_X$-가군들의 실선 도식
$$
\xymatrix{
\mathcal{E} \ar@{..>}[dr] \ar[r] & \mathcal{F} \\
& \mathcal{G} \ar[u]_p
}
$$
이 주어져 있고, $\mathcal{E}$가 유한 자유 $\mathcal{O}_X$-가군의
직합인자이며 $p$가 전사라고 하자. 그러면 도식을 가환하게 하는
점선 화살표가 $X$ 위에서 국소적으로 존재한다.
\end{lemma}

\begin{proof}
어떤 $n$에 대해 $\mathcal{E} = \mathcal{O}_X^{\oplus n}$이라고
가정해도 된다. 이 경우 점선 화살표를 찾는 것은
$\Gamma(X, \mathcal{F})$에 있는 기저 원소들의 상을 들어 올리는
것과 같다. 이는 층의 전사에 대한 특성화에 의해 국소적으로 가능하다
(Sheaves, Section \ref{sheaves-section-exactness-points}).
\end{proof}

\begin{lemma}
\label{cohomology-lemma-local-homotopy}
$(X, \mathcal{O}_X)$를 환 달린 공간이라 하자.
\begin{enumerate}
\item $\alpha : \mathcal{E}^\bullet \to \mathcal{F}^\bullet$가
$\mathcal{O}_X$-가군 복합체들의 사상이고,
$\mathcal{E}^\bullet$가 강한 완전이며 $\mathcal{F}^\bullet$가
비순환이라고 하자. 그러면 $\alpha$는 $X$ 위에서 국소적으로 영과
호모토픽이다.
\item $\alpha : \mathcal{E}^\bullet \to \mathcal{F}^\bullet$가
$\mathcal{O}_X$-가군 복합체들의 사상이고,
$\mathcal{E}^\bullet$가 강한 완전이며, $i < a$에 대해
$\mathcal{E}^i = 0$이고 $i \geq a$에 대해
$H^i(\mathcal{F}^\bullet) = 0$이라고 하자. 그러면 $\alpha$는 $X$
위에서 국소적으로 영과 호모토픽이다.
\end{enumerate}
\end{lemma}

\begin{proof}
첫째 명제는 둘째 명제에서 따르므로 (2)만 증명한다. 복합체
$\mathcal{E}^\bullet$의 길이에 대한 귀납법을 사용한다.
$\mathcal{E}$가 유한 자유 $\mathcal{O}_X$-가군의 직합인자이고 어떤 정수
$n \geq a$에 대해
$\mathcal{E}^\bullet \cong \mathcal{E}[-n]$이면, Lemma
\ref{cohomology-lemma-local-lift-map}과 $H^n(\mathcal{F}^\bullet)$이 사라진다는
가정에 의해
$\mathcal{F}^{n - 1} \to \Ker(\mathcal{F}^n \to \mathcal{F}^{n + 1})$
가 전사라는 사실에서 결론이 따른다. $\mathcal{E}^i$가
$i \in [a, b]$인 경우를 제외하면 영이면, 복합체들의 분할 완전열
$$
0 \to \mathcal{E}^b[-b] \to \mathcal{E}^\bullet \to
\sigma_{\leq b - 1}\mathcal{E}^\bullet \to 0
$$
이 있고, 이는 $K(\mathcal{O}_X)$ 안의 특수 삼각형을 결정한다.
따라서 삼각범주의 공리로 완전열
$$
\Hom_{K(\mathcal{O}_X)}(
\sigma_{\leq b - 1}\mathcal{E}^\bullet, \mathcal{F}^\bullet)
\to
\Hom_{K(\mathcal{O}_X)}(\mathcal{E}^\bullet, \mathcal{F}^\bullet)
\to
\Hom_{K(\mathcal{O}_X)}(\mathcal{E}^b[-b], \mathcal{F}^\bullet)
$$
을 얻는다. 합성 $\mathcal{E}^b[-b] \to \mathcal{F}^\bullet$은
국소적으로 영과 호모토픽이므로, 우리 사상이 위 표시 완전열의
왼쪽 항의 한 원소에서 온다고 가정해도 된다. 이 원소는 귀납 가설에
의해 국소적으로 영이다.
\end{proof}

\begin{lemma}
\label{cohomology-lemma-lift-through-quasi-isomorphism}
$(X, \mathcal{O}_X)$를 환 달린 공간이라 하자.
$\mathcal{O}_X$-가군 복합체들의 실선 도식
$$
\xymatrix{
\mathcal{E}^\bullet \ar@{..>}[dr] \ar[r]_\alpha & \mathcal{F}^\bullet \\
& \mathcal{G}^\bullet \ar[u]_f
}
$$
이 주어졌다고 하자. $\mathcal{E}^\bullet$가 강한 완전이고,
$j < a$에 대해 $\mathcal{E}^j = 0$이며, $j > a$에 대해 $H^j(f)$가
동형사상이고 $j = a$에 대해 전사라고 하자. 그러면 도식을
호모토피까지 가환하게 하는 점선 화살표가 $X$ 위에서 국소적으로
존재한다.
\end{lemma}

\begin{proof}
$f$에 대한 가정에 의해 원뿔 $C(f)^\bullet$의 $\geq a$차
코호몰로지 층은 사라진다. 따라서 Lemma
\ref{cohomology-lemma-local-homotopy}에 의해 열린 덮개 $X = \bigcup U_i$가
존재하여 합성
$\mathcal{E}^\bullet \to \mathcal{F}^\bullet \to C(f)^\bullet$은
$U_i$ 위에서 영과 호모토픽이다. 또한
$$
\mathcal{G}^\bullet \to \mathcal{F}^\bullet \to C(f)^\bullet \to
\mathcal{G}^\bullet[1]
$$
는 $K(\mathcal{O}_{U_i})$ 안의 특수 삼각형으로 제한되므로,
원하는 대로 $\alpha|_{U_i}$를 호모토피까지 사상
$\alpha_i : \mathcal{E}^\bullet|_{U_i} \to
\mathcal{G}^\bullet|_{U_i}$로 들어 올릴 수 있다.
\end{proof}

\begin{lemma}
\label{cohomology-lemma-local-actual}
$(X, \mathcal{O}_X)$를 환 달린 공간이라 하자.
$\mathcal{E}^\bullet$, $\mathcal{F}^\bullet$를
$\mathcal{O}_X$-가군의 복합체라 하고, $\mathcal{E}^\bullet$가 강한
완전이라고 하자.
\begin{enumerate}
\item 임의의 원소
$\alpha \in \Hom_{D(\mathcal{O}_X)}(\mathcal{E}^\bullet, \mathcal{F}^\bullet)$
에 대해 열린 덮개 $X = \bigcup U_i$가 존재하여
$\alpha|_{U_i}$는 복합체들의 사상
$\alpha_i : \mathcal{E}^\bullet|_{U_i} \to
\mathcal{F}^\bullet|_{U_i}$로 주어진다.
\item 복합체들의 사상
$\alpha : \mathcal{E}^\bullet \to \mathcal{F}^\bullet$의 군
$\Hom_{D(\mathcal{O}_X)}(\mathcal{E}^\bullet, \mathcal{F}^\bullet)$
에서의 상이 영이라고 하자. 그러면 열린 덮개 $X = \bigcup U_i$가
존재하여 $\alpha|_{U_i}$는 영과 호모토픽이다.
\end{enumerate}
\end{lemma}

\begin{proof}
(1)의 증명. 유도범주의 구성에 의해 유사동형
$f : \mathcal{F}^\bullet \to \mathcal{G}^\bullet$와 복합체들의 사상
$\beta : \mathcal{E}^\bullet \to \mathcal{G}^\bullet$를
$\alpha = f^{-1}\beta$가 되도록 찾을 수 있다. 따라서 Lemma
\ref{cohomology-lemma-lift-through-quasi-isomorphism}에서 결론이 따른다.
(2)의 증명은 생략한다.
\end{proof}

\begin{lemma}
\label{cohomology-lemma-Rhom-strictly-perfect}
$(X, \mathcal{O}_X)$를 환 달린 공간이라 하자.
$\mathcal{E}^\bullet$, $\mathcal{F}^\bullet$를
$\mathcal{O}_X$-가군의 복합체라 하고, $\mathcal{E}^\bullet$가 강한
완전이라고 하자. 그러면 내부 Hom
$R\SheafHom(\mathcal{E}^\bullet, \mathcal{F}^\bullet)$은 항이
$$
\mathcal{H}^n =
\bigoplus\nolimits_{n = p + q}
\SheafHom_{\mathcal{O}_X}(\mathcal{E}^{-q}, \mathcal{F}^p)
$$
이고 미분이 Section \ref{cohomology-section-hom-complexes}에서 설명된 복합체
$\mathcal{H}^\bullet$로 나타내어진다.
\end{lemma}

\begin{proof}
K-단사 복합체 안으로 가는 유사동형
$\mathcal{F}^\bullet \to \mathcal{I}^\bullet$를 택하자.
$(\mathcal{H}')^\bullet$를 항이
$$
(\mathcal{H}')^n =
\prod\nolimits_{n = p + q}
\SheafHom_{\mathcal{O}_X}(\mathcal{E}^{-q}, \mathcal{I}^p)
$$
인 복합체라 하자. 이는 Section \ref{cohomology-section-internal-hom}의 구성에
의해 $R\SheafHom(\mathcal{E}^\bullet, \mathcal{F}^\bullet)$을
나타낸다. 사상
$$
\mathcal{H}^\bullet \longrightarrow (\mathcal{H}')^\bullet
$$
이 유사동형임을 보이면 충분하다. 열린집합 $U \subset X$가 주어지면
직접 확인하여
$$
H^0(\mathcal{H}^\bullet(U)) =
\Hom_{K(\mathcal{O}_U)}(\mathcal{E}^\bullet|_U, \mathcal{I}^\bullet|_U)
\to
H^0((\mathcal{H}')^\bullet(U)) =
\Hom_{D(\mathcal{O}_U)}(\mathcal{E}^\bullet|_U, \mathcal{I}^\bullet|_U)
$$
를 얻는다. Lemma \ref{cohomology-lemma-local-actual}에 의해
$U \mapsto H^0(\mathcal{H}^\bullet(U))$의 층화는
$U \mapsto H^0((\mathcal{H}')^\bullet(U))$의 층화와 같다. 다른
코호몰로지 층에 대해서도 같은 논증을 할 수 있다. 따라서
$\mathcal{H}^\bullet$는 $(\mathcal{H}')^\bullet$와 유사동형이고,
보조정리가 증명된다.
\end{proof}

\begin{lemma}
\label{cohomology-lemma-Rhom-strictly-perfect-K-flat}
Lemma \ref{cohomology-lemma-Rhom-strictly-perfect}의 상황에서
$\mathcal{F}^\bullet$가 K-평탄이면 $\mathcal{H}^\bullet$도
K-평탄이다.
\end{lemma}

\begin{proof}
강한 완전 복합체 $\mathcal{E}^\bullet$의 유계성 때문에 Hom 복합체
정의의 직접곱이 직접합이 되므로, $\mathcal{H}^\bullet$는 단순히
Hom 복합체
$\SheafHom^\bullet(\mathcal{E}^\bullet, \mathcal{F}^\bullet)$이다.
$\mathcal{K}^\bullet$를 $\mathcal{O}_X$-가군의 비순환 복합체라 하고,
Lemma \ref{cohomology-lemma-diagonal-better}의 사상
$$
\gamma :
\text{Tot}(\mathcal{K}^\bullet \otimes
\SheafHom^\bullet(\mathcal{E}^\bullet, \mathcal{F}^\bullet))
\longrightarrow
\SheafHom^\bullet(\mathcal{E}^\bullet,
\text{Tot}(\mathcal{K}^\bullet \otimes \mathcal{F}^\bullet))
$$
을 생각하자. $\mathcal{F}^\bullet$가 K-평탄이므로 복합체
$\text{Tot}(\mathcal{K}^\bullet \otimes \mathcal{F}^\bullet)$는
비순환이고, 따라서 Lemma \ref{cohomology-lemma-local-actual}
(또는 원한다면 Lemma \ref{cohomology-lemma-Rhom-strictly-perfect})에 의해
$\gamma$의 공역도 비순환이다. 그러므로 보조정리를 증명하려면
$\gamma$가 복합체들의 동형사상임을 보이면 충분하다. 복합체
$\mathcal{E}^\bullet$의 길이에 대한 귀납법을 사용하자. 길이가
$\leq 1$이면 어떤 $n \geq 0$과 $k \in \mathbf{Z}$에 대해
$\mathcal{E}^\bullet$는 $\mathcal{O}_X^{\oplus n}[k]$의 직합인자이고,
이 경우 결과는 직접 확인된다. 길이가 $> 1$이면 적당한
$a \in \mathbf{Z}$에 대한 복합체들의 항별 분할 짧은 완전열
$$
0 \to \sigma_{\geq a + 1} \mathcal{E}^\bullet \to
\mathcal{E}^\bullet \to
\sigma_{\leq a} \mathcal{E}^\bullet \to 0
$$
을 생각하여 더 짧은 경우로 귀착시킨다. Homology, Section
\ref{homology-section-truncations}를 보라. 그러면 $\gamma$는
복합체들의 항별 분할 짧은 완전열들의 사상에 들어맞는다. 귀납
가설에 의해 $\sigma_{\geq a + 1} \mathcal{E}^\bullet$와
$\sigma_{\leq a} \mathcal{E}^\bullet$에 대해서는 $\gamma$가
동형사상이므로, $\mathcal{E}^\bullet$에 대한 결과가 따른다. 일부
세부 사항은 생략한다.
\end{proof}

\begin{lemma}
\label{cohomology-lemma-Rhom-complex-of-direct-summands-finite-free}
$(X, \mathcal{O}_X)$를 환 달린 공간이라 하자.
$\mathcal{E}^\bullet$, $\mathcal{F}^\bullet$를
$\mathcal{O}_X$-가군의 복합체라 하고 다음을 가정하자.
\begin{enumerate}
\item $n \ll 0$이면 $\mathcal{F}^n = 0$이다.
\item $n \gg 0$이면 $\mathcal{E}^n = 0$이다.
\item $\mathcal{E}^n$은 유한 자유 $\mathcal{O}_X$-가군의
직합인자와 동형이다.
\end{enumerate}
그러면 내부 Hom
$R\SheafHom(\mathcal{E}^\bullet, \mathcal{F}^\bullet)$은 항이
$$
\mathcal{H}^n =
\bigoplus\nolimits_{n = p + q}
\SheafHom_{\mathcal{O}_X}(\mathcal{E}^{-q}, \mathcal{F}^p)
$$
이고 미분이 Section \ref{cohomology-section-internal-hom}에서 설명된 복합체
$\mathcal{H}^\bullet$로 나타내어진다.
\end{lemma}

\begin{proof}
$\mathcal{I}^\bullet$가 단사대상들로 이루어진 아래로 유계인 복합체인
유사동형 $\mathcal{F}^\bullet \to \mathcal{I}^\bullet$를 택하자.
$\mathcal{I}^\bullet$는 K-단사임에 유의하자(Derived Categories,
Lemma \ref{derived-lemma-bounded-below-injectives-K-injective}).
따라서 Section \ref{cohomology-section-internal-hom}의 구성에 의해
$R\SheafHom(\mathcal{E}^\bullet, \mathcal{F}^\bullet)$은 항이
$$
(\mathcal{H}')^n =
\prod\nolimits_{n = p + q}
\SheafHom_{\mathcal{O}_X}(\mathcal{E}^{-q}, \mathcal{I}^p) =
\bigoplus\nolimits_{n = p + q}
\SheafHom_{\mathcal{O}_X}(\mathcal{E}^{-q}, \mathcal{I}^p)
$$
인 복합체 $(\mathcal{H}')^\bullet$로 나타내어진다. 영이 아닌 항은
유한 개뿐이므로 등식이 성립한다. $\mathcal{H}^\bullet$는 항이
$\SheafHom_{\mathcal{O}_X}(\mathcal{E}^{-q}, \mathcal{F}^p)$인 이중
복합체에 결부된 전체 복합체이고, $(\mathcal{H}')^\bullet$도
마찬가지이다. 자연스러운 사상
$(\mathcal{H}')^\bullet \to \mathcal{H}^\bullet$은 이중 복합체들의
사상에서 온다. 이 사상이 유사동형임을 보이기 위해 이중 복합체의
스펙트럼열을 사용할 수 있다
(Homology, Lemma \ref{homology-lemma-first-quadrant-ss}).
$$
{}'E_1^{p, q} =
H^p(\SheafHom_{\mathcal{O}_X}(\mathcal{E}^{-q}, \mathcal{F}^\bullet))
$$
이는 $H^{p + q}(\mathcal{H}^\bullet)$로 수렴하며,
$(\mathcal{H}')^\bullet$에 대해서도 마찬가지이다. 보조정리의 증명을
끝내려면 $\mathcal{F}^\bullet \to \mathcal{I}^\bullet$가
$\mathcal{E}$가 유한 자유 $\mathcal{O}_X$-가군의 직합인자일 때마다
코호몰로지 층 위의 동형사상
$$
H^p(\SheafHom_{\mathcal{O}_X}(\mathcal{E}, \mathcal{F}^\bullet))
\longrightarrow
H^p(\SheafHom_{\mathcal{O}_X}(\mathcal{E}, \mathcal{I}^\bullet))
$$
을 유도함을 보이면 충분하다. $\mathcal{E}$가 유한 자유이면 이는
명백하므로 결과가 따른다.
\end{proof}






\section{유사코히런트 가군}
\label{cohomology-section-pseudo-coherent}

\noindent
이 절에서는 유사코히런트 복합체를 논의한다.

\begin{definition}
\label{cohomology-definition-pseudo-coherent}
$(X, \mathcal{O}_X)$를 환 달린 공간이라 하고,
$\mathcal{E}^\bullet$를 $\mathcal{O}_X$-가군의 복합체라 하며,
$m \in \mathbf{Z}$라 하자.
\begin{enumerate}
\item 열린 덮개 $X = \bigcup U_i$가 존재하고, 각 $i$에 대해
복합체들의 사상
$\alpha_i : \mathcal{E}_i^\bullet \to \mathcal{E}^\bullet|_{U_i}$가
있어서, $\mathcal{E}_i^\bullet$가 $U_i$ 위에서 강한 완전이고
$j > m$에 대해 $H^j(\alpha_i)$가 동형사상이며
$H^m(\alpha_i)$가 전사이면 $\mathcal{E}^\bullet$를
{\it $m$-유사코히런트}라고 한다.
\item 모든 $m$에 대해 $m$-유사코히런트이면
$\mathcal{E}^\bullet$를 {\it 유사코히런트}라고 한다.
\item $D(\mathcal{O}_X)$의 대상 $E$가 $m$-유사코히런트
(각각 유사코히런트)인 $\mathcal{O}_X$-가군 복합체로 나타내어질 수
있는 경우, 그리고 그 경우에만 이를 {\it $m$-유사코히런트}
(각각 {\it 유사코히런트})라고 한다.
\end{enumerate}
\end{definition}

\noindent
$X$가 준콤팩트이면 $D(\mathcal{O}_X)$의 $m$-유사코히런트 대상은
$D^-(\mathcal{O}_X)$에 속한다. 그러나 $X$가 준콤팩트가 아니면
그럴 필요가 없다.

\begin{lemma}
\label{cohomology-lemma-pseudo-coherent-independent-representative}
$(X, \mathcal{O}_X)$를 환 달린 공간이라 하고, $E$를
$D(\mathcal{O}_X)$의 대상이라 하자.
\begin{enumerate}
\item 열린 덮개 $X = \bigcup U_i$, $U_i$ 위의 강한 완전 복합체
$\mathcal{E}_i^\bullet$, 그리고 $D(\mathcal{O}_{U_i})$의 사상
$\alpha_i : \mathcal{E}_i^\bullet \to E|_{U_i}$가 존재하여
$j > m$에 대해 $H^j(\alpha_i)$가 동형사상이고
$H^m(\alpha_i)$가 전사이면 $E$는 $m$-유사코히런트이다.
\item $E$가 $m$-유사코히런트이면 $E$를 나타내는 모든 복합체가
$m$-유사코히런트이다.
\end{enumerate}
\end{lemma}

\begin{proof}
$\mathcal{F}^\bullet$를 $E$를 나타내는 임의의 복합체라 하고,
$X = \bigcup U_i$와
$\alpha_i : \mathcal{E}_i^\bullet \to E|_{U_i}$가 (1)에서와 같이
주어졌다고 하자. $\mathcal{F}^\bullet$가 복합체로서
$m$-유사코히런트임을 보이면 (1)과 (2)가 동시에 증명된다. Lemma
\ref{cohomology-lemma-local-actual}에 의해 열린 덮개 $X = \bigcup U_i$를
세분한 뒤 사상 $\alpha_i$를 복합체들의 사상
$\alpha_i : \mathcal{E}_i^\bullet \to \mathcal{F}^\bullet|_{U_i}$로
나타낼 수 있다. 가정에 의해 $j > m$에 대해 $H^j(\alpha_i)$는
동형사상이고 $H^m(\alpha_i)$는 전사이므로,
$\mathcal{F}^\bullet$는 $m$-유사코히런트이다.
\end{proof}

\begin{lemma}
\label{cohomology-lemma-pseudo-coherent-pullback}
$f : (X, \mathcal{O}_X) \to (Y, \mathcal{O}_Y)$를 환 달린 공간의
사상이라 하고, $E$를 $D(\mathcal{O}_Y)$의 대상이라 하자. $E$가
$m$-유사코히런트이면 $Lf^*E$도 $m$-유사코히런트이다.
\end{lemma}

\begin{proof}
$E$를 $\mathcal{O}_Y$-가군 복합체 $\mathcal{E}^\bullet$로 나타내고,
Definition \ref{cohomology-definition-pseudo-coherent}에서와 같은 열린 덮개
$Y = \bigcup V_i$와
$\alpha_i : \mathcal{E}_i^\bullet \to \mathcal{E}^\bullet|_{V_i}$를
택하자. $U_i = f^{-1}(V_i)$라 두자. Lemma
\ref{cohomology-lemma-pseudo-coherent-independent-representative}에 의해
$Lf^*\mathcal{E}^\bullet|_{U_i}$가 $m$-유사코히런트임을 보이면
충분하다. 특수 삼각형
$$
\mathcal{E}_i^\bullet \to
\mathcal{E}^\bullet|_{V_i} \to
C \to
\mathcal{E}_i^\bullet[1]
$$
을 택하자. $\alpha_i$에 대한 가정은 모든 $j \geq m$에 대해
코호몰로지 층 $H^j(C)$가 영이라는 것과 정확히 같다.
$f_i : U_i \to V_i$를 $f$의 제한이라 쓰자. 다음에 유의하자.
$Lf^*\mathcal{E}^\bullet|_{U_i} = 
Lf_i^*(\mathcal{E}|_{V_i})$.
$Lf_i^*$를 적용하면 특수 삼각형
$$
Lf_i^*\mathcal{E}_i^\bullet \to
Lf_i^*\mathcal{E}|_{V_i} \to
Lf_i^*C \to
Lf_i^*\mathcal{E}_i^\bullet[1]
$$
을 얻는다. 왼쪽 유도함자로서의 $Lf_i^*$의 구성에 의해
$j \geq m$에 대해 $H^j(Lf_i^*C) = 0$임을 알 수 있다
(Derived Categories, Lemma
\ref{derived-lemma-negative-vanishing}의 쌍대). 따라서
$j > m$에 대해 $H^j(Lf_i^*\alpha_i)$는 동형사상이고
$H^m(Lf^*\alpha_i)$는 전사이다. 한편
$Lf_i^*\mathcal{E}_i^\bullet = f_i^*\mathcal{E}_i^\bullet$은 Lemma
\ref{cohomology-lemma-strictly-perfect-pullback}에 의해 강한 완전이다. 이로써
결론이 따른다.
\end{proof}

\begin{lemma}
\label{cohomology-lemma-cone-pseudo-coherent}
$(X, \mathcal{O}_X)$를 환 달린 공간이라 하고,
$m \in \mathbf{Z}$라 하자. $(K, L, M, f, g, h)$를
$D(\mathcal{O}_X)$의 특수 삼각형이라 하자.
\begin{enumerate}
\item $K$가 $(m + 1)$-유사코히런트이고 $L$이
$m$-유사코히런트이면 $M$은 $m$-유사코히런트이다.
\item $K$와 $M$이 $m$-유사코히런트이면 $L$도
$m$-유사코히런트이다.
\item $L$이 $(m + 1)$-유사코히런트이고 $M$이
$m$-유사코히런트이면 $K$는 $(m + 1)$-유사코히런트이다.
\end{enumerate}
\end{lemma}

\begin{proof}
(1)의 증명. 열린 덮개 $X = \bigcup U_i$와
$D(\mathcal{O}_{U_i})$의 사상
$\alpha_i : \mathcal{K}_i^\bullet \to K|_{U_i}$를 택하자. 여기서
$\mathcal{K}_i^\bullet$는 강한 완전이고,
$j > m + 1$에 대해 $H^j(\alpha_i)$는 동형사상이며
$j = m + 1$에 대해 전사이다.
$\mathcal{K}_i^\bullet$를
$\sigma_{\geq m + 1}\mathcal{K}_i^\bullet$로 바꿀 수 있으므로,
$j < m + 1$에 대해 $\mathcal{K}_i^j = 0$이라고 가정해도 된다.
열린 덮개를 세분한 뒤 $D(\mathcal{O}_{U_i})$의 사상
$\beta_i : \mathcal{L}_i^\bullet \to L|_{U_i}$를 택하되,
$\mathcal{L}_i^\bullet$는 강한 완전이고, $j > m$에 대해
$H^j(\beta)$는 동형사상이며 $j = m$에 대해 전사라고 할 수 있다.
Lemma \ref{cohomology-lemma-lift-through-quasi-isomorphism}에 의해 덮개를 다시
세분한 뒤 복합체들의 사상
$\gamma_i : \mathcal{K}^\bullet \to \mathcal{L}^\bullet$를 찾아서
도식
$$
\xymatrix{
K|_{U_i} \ar[r] & L|_{U_i} \\
\mathcal{K}_i^\bullet \ar[u]^{\alpha_i} \ar[r]^{\gamma_i} &
\mathcal{L}_i^\bullet \ar[u]_{\beta_i}
}
$$
들이 $D(\mathcal{O}_{U_i})$에서 가환하게 할 수 있다. 여기에는 사상
$\alpha_i$, $\beta_i$, $K|_{U_i} \to L|_{U_i}$를 복합체들의 실제
사상으로 나타내는 일이 필요하다. 일부 세부 사항은 생략한다.
원뿔 $C(\gamma_i)^\bullet$는 강한 완전이다
(Lemma \ref{cohomology-lemma-cone}). 도식의 가환성에 의해 특수 삼각형들의 사상
$$
(\mathcal{K}_i^\bullet, \mathcal{L}_i^\bullet, C(\gamma_i)^\bullet)
\longrightarrow
(K|_{U_i}, L|_{U_i}, M|_{U_i}).
$$
이 존재한다. 유도된 긴 코호몰로지 완전열의 사상과 Homology, Lemmas
\ref{homology-lemma-four-lemma} 및
\ref{homology-lemma-five-lemma}에 의해
$C(\gamma_i)^\bullet \to M|_{U_i}$는 $> m$차 코호몰로지에서
동형사상을, $m$차에서 전사를 유도한다. 따라서 Lemma
\ref{cohomology-lemma-pseudo-coherent-independent-representative}에 의해
$M$은 $m$-유사코히런트이다.

\medskip\noindent
(2)와 (3)은 특수 삼각형을 회전하여 (1)에서 얻는다.
\end{proof}

\begin{lemma}
\label{cohomology-lemma-tensor-pseudo-coherent}
$(X, \mathcal{O}_X)$를 환 달린 공간이라 하고, $K, L$을
$D(\mathcal{O}_X)$의 대상이라 하자.
\begin{enumerate}
\item $K$가 $n$-유사코히런트이고 $i > a$에 대해
$H^i(K) = 0$이며, $L$이 $m$-유사코히런트이고 $j > b$에 대해
$H^j(L) = 0$이면 $K \otimes_{\mathcal{O}_X}^\mathbf{L} L$은
$t = \max(m + a, n + b)$에 대해 $t$-유사코히런트이다.
\item $K$와 $L$이 유사코히런트이면
$K \otimes_{\mathcal{O}_X}^\mathbf{L} L$도 유사코히런트이다.
\end{enumerate}
\end{lemma}

\begin{proof}
(1)의 증명. $X$를 열린 덮개의 원소들로 바꾸어, 강한 완전 복합체
$\mathcal{K}^\bullet$, $\mathcal{L}^\bullet$와 사상
$\alpha : \mathcal{K}^\bullet \to K$,
$\beta : \mathcal{L}^\bullet \to L$이 존재하고, $i > n$에 대해
$H^i(\alpha)$가 동형사상이며 $i = n$에 대해 전사이고,
$i > m$에 대해 $H^i(\beta)$가 동형사상이며 $i = m$에 대해
전사라고 가정할 수 있다. 그러면 사상
$$
\alpha \otimes^\mathbf{L} \beta :
\text{Tot}(\mathcal{K}^\bullet \otimes_{\mathcal{O}_X} \mathcal{L}^\bullet)
\to K \otimes_{\mathcal{O}_X}^\mathbf{L} L
$$
은 $i > t$인 차수 $i$의 코호몰로지 층에서 동형사상을 유도하고,
$i = t$이면 전사를 유도한다. 이는 Tor의 스펙트럼열에서 따른다.
세부 사항은 생략한다.

\medskip\noindent
(2)의 증명. 먼저 $X$를 열린 덮개의 원소들로 바꾸어 $K$와 $L$이
위로 유계인 경우로 귀착시킬 수 있다. 그러면 (1)에서 곧바로
명제가 따른다.
\end{proof}

\begin{lemma}
\label{cohomology-lemma-summands-pseudo-coherent}
$(X, \mathcal{O}_X)$를 환 달린 공간이라 하고,
$m \in \mathbf{Z}$라 하자. $D(\mathcal{O}_X)$에서 $K \oplus L$이
$m$-유사코히런트(각각 유사코히런트)이면 $K$와 $L$도 그러하다.
\end{lemma}

\begin{proof}
$K \oplus L$이 $m$-유사코히런트라고 하자. $X$를 열린 덮개의
원소들로 바꾸어 $K \oplus L \in D^-(\mathcal{O}_X)$, 따라서
$L \in D^-(\mathcal{O}_X)$라고 가정할 수 있다. 특수 삼각형
$$
(K \oplus L, K \oplus L, L \oplus L[1]) =
(K, K, 0) \oplus (L, L, L \oplus L[1])
$$
이 있음에 유의하자. Derived Categories, Lemma
\ref{derived-lemma-direct-sum-triangles}를 보라. Lemma
\ref{cohomology-lemma-cone-pseudo-coherent}에 의해 $L \oplus L[1]$은
$m$-유사코히런트이다. 따라서 $L[1] \oplus L[2]$도
$m$-유사코히런트이고, 귀납법으로
$L[n] \oplus L[n + 1]$은 $m$-유사코히런트이다. $L$이 위로
유계이므로 충분히 큰 $n$에 대해 $L[n]$은 $m$-유사코히런트이다.
따라서 특수 삼각형
$$
(L[n], L[n] \oplus L[n - 1], L[n - 1])
$$
들을 사용하여 거꾸로 내려가면, 원하는 대로
$L[n - 1], L[n - 2], \ldots, L$이 $m$-유사코히런트임을 얻는다.
\end{proof}

\begin{lemma}
\label{cohomology-lemma-complex-pseudo-coherent-modules}
$(X, \mathcal{O}_X)$를 환 달린 공간이라 하고,
$m \in \mathbf{Z}$라 하자. $\mathcal{F}^\bullet$를
$\mathcal{O}_X$-가군의 (국소적으로) 위로 유계인 복합체라 하자.
모든 $i$에 대해 $\mathcal{F}^i$가
$(m - i)$-유사코히런트이면 $\mathcal{F}^\bullet$는
$m$-유사코히런트이다.
\end{lemma}

\begin{proof}
생략한다. 힌트: Lemma \ref{cohomology-lemma-cone-pseudo-coherent}와 절단을
사용하여 More on Algebra, Lemma
\ref{more-algebra-lemma-complex-pseudo-coherent-modules}의 증명에서와
같이 논증한다.
\end{proof}

\begin{lemma}
\label{cohomology-lemma-cohomology-pseudo-coherent}
$(X, \mathcal{O}_X)$를 환 달린 공간이라 하고,
$m \in \mathbf{Z}$라 하자. $E$를 $D(\mathcal{O}_X)$의 대상이라 하자.
$E$가 (국소적으로) 위로 유계이고 모든 $i$에 대해 $H^i(E)$가
$(m - i)$-유사코히런트이면 $E$는 $m$-유사코히런트이다.
\end{lemma}

\begin{proof}
생략한다. 힌트: Lemma \ref{cohomology-lemma-cone-pseudo-coherent}와 절단을
사용하여 More on Algebra, Lemma
\ref{more-algebra-lemma-cohomology-pseudo-coherent}의 증명에서와
같이 논증한다.
\end{proof}

\begin{lemma}
\label{cohomology-lemma-finite-cohomology}
$(X, \mathcal{O}_X)$를 환 달린 공간이라 하고, $K$를
$D(\mathcal{O}_X)$의 대상, $m \in \mathbf{Z}$라 하자.
\begin{enumerate}
\item $K$가 $m$-유사코히런트이고 $i > m$에 대해 $H^i(K) = 0$이면,
$H^m(K)$는 유한형 $\mathcal{O}_X$-가군이다.
\item $K$가 $m$-유사코히런트이고 $i > m + 1$에 대해
$H^i(K) = 0$이면, $H^{m + 1}(K)$는 유한 표시
$\mathcal{O}_X$-가군이다.
\end{enumerate}
\end{lemma}

\begin{proof}
(1)의 증명. $X$ 위에서 국소적으로 작업해도 된다. 따라서 강한 완전
복합체 $\mathcal{E}^\bullet$와 $> m$차 코호몰로지에서 동형사상을,
$m$차에서 전사를 유도하는 사상
$\alpha : \mathcal{E}^\bullet \to K$가 존재한다고 가정해도 된다.
$\mathcal{E}^\bullet$에 대해 결과를 증명하면 충분하다.
$\mathcal{E}^n \not = 0$인 가장 큰 정수를 $n$이라 하자. $n = m$이면
$H^m(\mathcal{E}^\bullet)$는 $\mathcal{E}^n$의 몫이므로 결과가
명백하다. $n > m$이면 $H^n(E^\bullet) = 0$이므로
$\mathcal{E}^{n - 1} \to \mathcal{E}^n$은 전사이다. Lemma
\ref{cohomology-lemma-local-lift-map}에 의해 국소적으로 이 전사의 단면을 찾아
$\mathcal{E}^{n - 1} = \mathcal{E}' \oplus \mathcal{E}^n$로 쓸 수
있다. 따라서 $\mathcal{E}^\bullet$와 같되 $n - 1$차 항이
$\mathcal{E}'$이고 $n$차 항이 $0$인 복합체
$(\mathcal{E}')^\bullet$에 대해 결과를 증명하면 충분하다. $n$에
대한 귀납법으로 결론이 따른다.

\medskip\noindent
(2)의 증명. $X$ 위에서 국소적으로 작업해도 된다. 따라서 강한 완전
복합체 $\mathcal{E}^\bullet$와 $> m$차 코호몰로지에서 동형사상을,
$m$차에서 전사를 유도하는 사상
$\alpha : \mathcal{E}^\bullet \to K$가 존재한다고 가정해도 된다.
(1)의 증명에서와 같이 $i > m + 1$에 대해
$\mathcal{E}^i = 0$인 경우로 귀착할 수 있다. 그러면
$H^{m + 1}(K) \cong H^{m + 1}(\mathcal{E}^\bullet) =
\Coker(\mathcal{E}^m \to \mathcal{E}^{m + 1})$
는 유한 표시이다.
\end{proof}

\begin{lemma}
\label{cohomology-lemma-n-pseudo-module}
$(X, \mathcal{O}_X)$를 환 달린 공간이라 하고, $\mathcal{F}$를
$\mathcal{O}_X$-가군의 층이라 하자.
\begin{enumerate}
\item $D(\mathcal{O}_X)$의 대상으로 본 $\mathcal{F}$가
$0$-유사코히런트인 것은 $\mathcal{F}$가 유한형
$\mathcal{O}_X$-가군일 때, 그리고 그때뿐이다.
\item $D(\mathcal{O}_X)$의 대상으로 본 $\mathcal{F}$가
$(-1)$-유사코히런트인 것은 $\mathcal{F}$가 유한 표시
$\mathcal{O}_X$-가군일 때, 그리고 그때뿐이다.
\end{enumerate}
\end{lemma}

\begin{proof}
한 방향의 함의에는 Lemma \ref{cohomology-lemma-finite-cohomology}를 사용하고,
다른 방향에는 Lemma \ref{cohomology-lemma-cohomology-pseudo-coherent}를
사용한다.
\end{proof}






\section{Tor 차원}
\label{cohomology-section-tor}

\noindent
이 절에서는 평탄 가군에 의한 분해를 더 자세히 살펴본다.

\begin{definition}
\label{cohomology-definition-tor-amplitude}
$(X, \mathcal{O}_X)$를 환 달린 공간이라 하자.
$E$를 $D(\mathcal{O}_X)$의 대상이라 하자.
$a \leq b$인 $a, b \in \mathbf{Z}$를 택하자.
\begin{enumerate}
\item 모든 $\mathcal{O}_X$-가군 $\mathcal{F}$와 모든 $i \not \in [a, b]$에
대하여 $H^i(E \otimes_{\mathcal{O}_X}^\mathbf{L} \mathcal{F}) = 0$이면,
$E$가 {\it $[a, b]$ 안에서 Tor-진폭을 가진다}고 한다.
\item 어떤 $a, b$에 대하여 $E$가 $[a, b]$ 안에서 Tor-진폭을 가지면,
{\it 유한 Tor 차원을 가진다}고 한다.
\item 열린 덮개 $X = \bigcup U_i$가 존재하여 모든 $i$에 대해
$E|_{U_i}$가 유한 Tor 차원을 가지면,
$E$가 {\it 국소적으로 유한 Tor 차원을 가진다}고 한다.
\end{enumerate}
$\mathcal{O}_X$-가군 $\mathcal{F}$에 대하여, $D(\mathcal{O}_X)$의 대상으로
본 $\mathcal{F}[0]$가 $[-d, 0]$ 안에서 Tor-진폭을 가지면
그 가군이 {\it Tor 차원 $\leq d$}를 가진다고 한다.
\end{definition}

\noindent
정의에 나오는 $E$가 유한 Tor 차원을 가지면 $E$는
$D^b(\mathcal{O}_X)$의 대상이다. 실제로 위 정의에서
$\mathcal{F} = \mathcal{O}_X$를 택하면 이를 알 수 있다.

\begin{lemma}
\label{cohomology-lemma-last-one-flat}
$(X, \mathcal{O}_X)$를 환 달린 공간이라 하자.
$\mathcal{E}^\bullet$을 $[a, b]$ 안에서 Tor-진폭을 가지는,
평탄한 $\mathcal{O}_X$-가군들의 위로 유계인 복합체라 하자.
그러면 $\Coker(d_{\mathcal{E}^\bullet}^{a - 1})$은 평탄한
$\mathcal{O}_X$-가군이다.
\end{lemma}

\begin{proof}
$\mathcal{E}^\bullet$은 평탄 가군들의 위로 유계인 복합체이므로,
임의의 $\mathcal{O}_X$-가군 $\mathcal{F}$에 대하여
$\mathcal{E}^\bullet \otimes_{\mathcal{O}_X} \mathcal{F} =
\mathcal{E}^\bullet \otimes_{\mathcal{O}_X}^{\mathbf{L}} \mathcal{F}$이다.
따라서 모든 $\mathcal{O}_X$-가군 $\mathcal{F}$에 대하여 수열
$$
\mathcal{E}^{a - 2} \otimes_{\mathcal{O}_X} \mathcal{F} \to
\mathcal{E}^{a - 1} \otimes_{\mathcal{O}_X} \mathcal{F} \to
\mathcal{E}^a \otimes_{\mathcal{O}_X} \mathcal{F}
$$
은 가운데에서 완전하다. 한편
$\mathcal{E}^{a - 2} \to \mathcal{E}^{a - 1} \to \mathcal{E}^a \to
\Coker(d^{a - 1}) \to 0$
은 평탄 분해이므로, 모든 $\mathcal{O}_X$-가군 $\mathcal{F}$에 대하여
$\text{Tor}_1^{\mathcal{O}_X}(\Coker(d^{a - 1}), \mathcal{F}) = 0$이다.
이는 $\Coker(d^{a - 1})$이 평탄함을 뜻한다.
보조정리 \ref{cohomology-lemma-flat-tor-zero}를 보라.
\end{proof}

\begin{lemma}
\label{cohomology-lemma-tor-amplitude}
$(X, \mathcal{O}_X)$를 환 달린 공간이라 하자.
$E$를 $D(\mathcal{O}_X)$의 대상이라 하자.
$a \leq b$인 $a, b \in \mathbf{Z}$를 택하자. 다음은 동치이다.
\begin{enumerate}
\item $E$는 $[a, b]$ 안에서 Tor-진폭을 가진다.
\item $E$는 평탄한 $\mathcal{O}_X$-가군들의 복합체
$\mathcal{E}^\bullet$로 표시되며,
$i \not \in [a, b]$이면 $\mathcal{E}^i = 0$이다.
\end{enumerate}
\end{lemma}

\begin{proof}
(2)가 성립하면
$E \otimes_{\mathcal{O}_X}^\mathbf{L} \mathcal{F} =
\mathcal{E}^\bullet \otimes_{\mathcal{O}_X} \mathcal{F}$
로 계산할 수 있으므로 (1)이 성립함은 분명하다.

\medskip\noindent
(1)이 성립한다고 하자. 절 \ref{cohomology-section-flat}에 따라 $E$를
평탄한 $\mathcal{O}_X$-가군들의 위로 유계인 복합체
$\mathcal{K}^\bullet$로 표시할 수 있다.
$\mathcal{K}^n \not = 0$인 가장 큰 정수를 $n$이라 하자.
$n > b$이면 $H^n(\mathcal{K}^\bullet) = 0$이므로
$\mathcal{K}^{n - 1} \to \mathcal{K}^n$은 전사이다.
$\mathcal{K}^n$이 평탄하므로
$\Ker(\mathcal{K}^{n - 1} \to \mathcal{K}^n)$도 평탄하다
(Modules, 보조정리 \ref{modules-lemma-flat-ses}).
따라서 $\mathcal{K}^\bullet$을
$\tau_{\leq n - 1}\mathcal{K}^\bullet$로 바꿀 수 있다.
그러므로 $n$에 관한 귀납법에 의해 $K^\bullet$이
$i > b$일 때 $\mathcal{K}^i = 0$인 평탄한
$\mathcal{O}_X$-가군들의 복합체인 경우로 귀착된다.

\medskip\noindent
$\mathcal{E}^\bullet = \tau_{\geq a}\mathcal{K}^\bullet$로 놓자.
$\mathcal{E}^a$가 평탄하다는 점을 제외하면 모두 분명하며,
그 점은 보조정리 \ref{cohomology-lemma-last-one-flat}와 정의에서 곧바로 따른다.
\end{proof}

\begin{lemma}
\label{cohomology-lemma-tor-amplitude-pullback}
$f : (X, \mathcal{O}_X) \to (Y, \mathcal{O}_Y)$를 환 달린 공간의
사상이라 하고, $E$를 $D(\mathcal{O}_Y)$의 대상이라 하자.
$E$가 $[a, b]$ 안에서 Tor-진폭을 가지면 $Lf^*E$도
$[a, b]$ 안에서 Tor-진폭을 가진다.
\end{lemma}

\begin{proof}
$E$가 $[a, b]$ 안에서 Tor-진폭을 가진다고 하자.
보조정리 \ref{cohomology-lemma-tor-amplitude}에 따라 $E$를,
$i \not \in [a, b]$이면 $\mathcal{E}^i = 0$인 평탄한
$\mathcal{O}$-가군들의 복합체 $\mathcal{E}^\bullet$로 표시할 수 있다.
그러면 $Lf^*E$는 $f^*\mathcal{E}^\bullet$로 표시된다.
Modules, 보조정리 \ref{modules-lemma-base-change-flat}에 의해
가군 $f^*\mathcal{E}^i$들은 평탄하다.
따라서 보조정리 \ref{cohomology-lemma-tor-amplitude}에 의해
$Lf^*E$가 $[a, b]$ 안에서 Tor-진폭을 가짐을 얻는다.
\end{proof}

\begin{lemma}
\label{cohomology-lemma-tor-amplitude-stalk}
$(X, \mathcal{O}_X)$를 환 달린 공간이라 하고,
$E$를 $D(\mathcal{O}_X)$의 대상이라 하자.
$a \leq b$인 $a, b \in \mathbf{Z}$를 택하자. 다음은 동치이다.
\begin{enumerate}
\item $E$는 $[a, b]$ 안에서 Tor-진폭을 가진다.
\item 모든 $x \in X$에 대하여 $D(\mathcal{O}_{X, x})$의 대상
$E_x$는 $[a, b]$ 안에서 Tor-진폭을 가진다.
\end{enumerate}
\end{lemma}

\begin{proof}
$x$에서의 줄기를 취하는 것은 환 달린 공간의 사상
$(x, \mathcal{O}_{X, x}) \to (X, \mathcal{O}_X)$에 의한 당김과 같다.
따라서 (1) $\Rightarrow$ (2)는
보조정리 \ref{cohomology-lemma-tor-amplitude-pullback}에서 따른다.
역으로, 줄기를 취하는 것은 텐서곱과 가환한다
(Modules, 보조정리 \ref{modules-lemma-stalk-tensor-product}).
따라서
$$
(E \otimes_{\mathcal{O}_X}^\mathbf{L} \mathcal{F})_x =
E_x \otimes_{\mathcal{O}_{X, x}}^\mathbf{L} \mathcal{F}_x
$$
한편 줄기를 취하는 함자는 완전하므로
$$
H^i(E \otimes_{\mathcal{O}_X}^\mathbf{L} \mathcal{F})_x =
H^i((E \otimes_{\mathcal{O}_X}^\mathbf{L} \mathcal{F})_x) =
H^i(E_x \otimes_{\mathcal{O}_{X, x}}^\mathbf{L} \mathcal{F}_x)
$$
이고, $H^i(E \otimes_{\mathcal{O}_X}^\mathbf{L} \mathcal{F})$의
모든 줄기가 영인지 확인함으로써 이 가군이 영인지 확인할 수 있다
(Modules, 보조정리 \ref{modules-lemma-abelian}). 따라서 (2)는 (1)을 함의한다.
\end{proof}

\begin{lemma}
\label{cohomology-lemma-cone-tor-amplitude}
$(X, \mathcal{O}_X)$를 환 달린 공간이라 하자.
$(K, L, M, f, g, h)$를 $D(\mathcal{O}_X)$의 특수 삼각형이라 하고,
$a, b \in \mathbf{Z}$라 하자.
\begin{enumerate}
\item $K$가 $[a + 1, b + 1]$ 안에서 Tor-진폭을 가지고
$L$이 $[a, b]$ 안에서 Tor-진폭을 가지면,
$M$은 $[a, b]$ 안에서 Tor-진폭을 가진다.
\item $K$와 $M$이 $[a, b]$ 안에서 Tor-진폭을 가지면,
$L$도 $[a, b]$ 안에서 Tor-진폭을 가진다.
\item $L$이 $[a + 1, b + 1]$ 안에서 Tor-진폭을 가지고
$M$이 $[a, b]$ 안에서 Tor-진폭을 가지면,
$K$는 $[a + 1, b + 1]$ 안에서 Tor-진폭을 가진다.
\end{enumerate}
\end{lemma}

\begin{proof}
생략한다. 힌트: 특수 삼각형에 딸린 긴 완전 코호몰로지 수열과
$- \otimes_{\mathcal{O}_X}^{\mathbf{L}} \mathcal{F}$가
특수 삼각형을 보존한다는 사실에서 곧바로 따른다.
(2)를 증명하는 것이 가장 쉽고, 나머지는 평행이동을 통해 그로부터 따른다.
\end{proof}

\begin{lemma}
\label{cohomology-lemma-tensor-tor-amplitude}
$(X, \mathcal{O}_X)$를 환 달린 공간이라 하고,
$K, L$을 $D(\mathcal{O}_X)$의 대상이라 하자.
$K$가 $[a, b]$ 안에서 Tor-진폭을 가지고
$L$이 $[c, d]$ 안에서 Tor-진폭을 가지면,
$K \otimes_{\mathcal{O}_X}^\mathbf{L} L$은
$[a + c, b + d]$ 안에서 Tor-진폭을 가진다.
\end{lemma}

\begin{proof}
생략한다. 힌트: Tor의 스펙트럼열을 사용하라.
\end{proof}

\begin{lemma}
\label{cohomology-lemma-summands-tor-amplitude}
$(X, \mathcal{O}_X)$를 환 달린 공간이라 하고,
$a, b \in \mathbf{Z}$라 하자. $D(\mathcal{O}_X)$의 대상
$K$, $L$에 대하여 $K \oplus L$이 $[a, b]$ 안에서
Tor-진폭을 가지면 $K$와 $L$도 그러하다.
\end{lemma}

\begin{proof}
Tor 함자들이 가법적이라는 사실에서 명백하다.
\end{proof}
\section{완전 복합체}
\label{cohomology-section-perfect}

\noindent
이 절에서는 환 달린 공간 위의 완전 복합체가 가지는 성질을 논한다.

\begin{definition}
\label{cohomology-definition-perfect}
$(X, \mathcal{O}_X)$를 환 달린 공간이라 하고,
$\mathcal{E}^\bullet$을 $\mathcal{O}_X$-가군들의 복합체라 하자.
열린 덮개 $X = \bigcup U_i$가 존재하여 각 $i$마다
복합체의 사상
$\mathcal{E}_i^\bullet \to \mathcal{E}^\bullet|_{U_i}$가 존재하고,
이 사상이 준동형이며 $\mathcal{E}_i^\bullet$이
$\mathcal{O}_{U_i}$-가군들의 강한 완전 복합체이면
$\mathcal{E}^\bullet$을 {\it 완전}이라고 한다.
$D(\mathcal{O}_X)$의 대상 $E$가 완전한
$\mathcal{O}_X$-가군 복합체로 표시될 수 있으면
이 대상을 {\it 완전}이라고 한다.
\end{definition}

\noindent
$X$가 준콤팩트이면 $D(\mathcal{O}_X)$의 완전 대상은
$D^b(\mathcal{O}_X)$에 속한다. 그러나 $X$가 준콤팩트하지 않으면
반드시 그러하지는 않다.

\begin{lemma}
\label{cohomology-lemma-perfect-independent-representative}
$(X, \mathcal{O}_X)$를 환 달린 공간이라 하고,
$E$를 $D(\mathcal{O}_X)$의 대상이라 하자.
\begin{enumerate}
\item 열린 덮개 $X = \bigcup U_i$와 $U_i$ 위의 강한 완전 복합체
$\mathcal{E}_i^\bullet$들이 존재하여
$\mathcal{E}_i^\bullet$이 $D(\mathcal{O}_{U_i})$에서
$E|_{U_i}$를 표시하면, $E$는 완전하다.
\item $E$가 완전하면 $E$를 표시하는 임의의 복합체는 완전하다.
\end{enumerate}
\end{lemma}

\begin{proof}
보조정리 \ref{cohomology-lemma-pseudo-coherent-independent-representative}의
증명과 동일하다.
\end{proof}

\begin{lemma}
\label{cohomology-lemma-perfect-on-locally-ringed}
$(X, \mathcal{O}_X)$를 환 달린 공간이라 하고,
$E$를 $D(\mathcal{O}_X)$의 대상이라 하자.
모든 줄기 $\mathcal{O}_{X, x}$가 국소환이라고 가정하자.
다음은 동치이다.
\begin{enumerate}
\item $E$는 완전하다.
\item 열린 덮개 $X = \bigcup U_i$가 존재하여 $E|_{U_i}$를
유한 국소 자유 $\mathcal{O}_{U_i}$-가군들의 유한 복합체로 표시할 수 있다.
\item 열린 덮개 $X = \bigcup U_i$가 존재하여 $E|_{U_i}$를
유한 자유 $\mathcal{O}_{U_i}$-가군들의 유한 복합체로 표시할 수 있다.
\end{enumerate}
\end{lemma}

\begin{proof}
이는 보조정리 \ref{cohomology-lemma-perfect-independent-representative}와
$X$ 위에서 유한 자유 가군의 모든 직합인자가 유한 국소 자유라는
사실에서 따른다. Modules, 보조정리
\ref{modules-lemma-direct-summand-of-locally-free-is-locally-free}를 보라.
\end{proof}

\begin{lemma}
\label{cohomology-lemma-perfect-precise}
$(X, \mathcal{O}_X)$를 환 달린 공간이라 하고,
$E$를 $D(\mathcal{O}_X)$의 대상이라 하자.
$a \leq b$를 정수라 하자. $E$가 $[a, b]$ 안에서 Tor-진폭을 가지고
$(a - 1)$-유사코히런트이면 $E$는 완전하다.
\end{lemma}

\begin{proof}
$X$를 열린 덮개의 원소들로 바꾸어, 강한 완전 복합체
$\mathcal{E}^\bullet$과 사상
$\alpha : \mathcal{E}^\bullet \to E$가 존재하고
$i \geq a$에 대하여 $H^i(\alpha)$가 동형이라고 가정할 수 있다.
$\mathcal{E}^\bullet$을
$\sigma_{\geq a - 1}\mathcal{E}^\bullet$로 바꾸어 놓는다.
특수 삼각형
$$
\mathcal{E}^\bullet \to E \to C \to \mathcal{E}^\bullet[1]
$$
을 택하자. $E$와 $\mathcal{E}^\bullet$의 코호몰로지 층들의 소멸과
$\alpha$에 관한 가정으로부터
$C \cong \mathcal{K}[2 - a]$를 얻는다. 여기서
$\mathcal{K} = \Ker(\mathcal{E}^{a - 1} \to \mathcal{E}^a)$이다.
$\mathcal{F}$를 $\mathcal{O}_X$-가군이라 하자.
$- \otimes_{\mathcal{O}_X}^\mathbf{L} \mathcal{F}$를 적용하면,
$E$가 $[a, b]$ 안에서 Tor-진폭을 가진다는 가정에 의해
$\mathcal{K} \otimes_{\mathcal{O}_X} \mathcal{F} \to
\mathcal{E}^{a - 1} \otimes_{\mathcal{O}_X} \mathcal{F}$의 상은
$\Ker(\mathcal{E}^{a - 1} \otimes_{\mathcal{O}_X} \mathcal{F}
\to \mathcal{E}^a \otimes_{\mathcal{O}_X} \mathcal{F})$이다.
따라서 $\mathcal{E}' = \Coker(\mathcal{E}^{a - 1} \to \mathcal{E}^a)$라
놓으면 $\text{Tor}_1^{\mathcal{O}_X}(\mathcal{E}', \mathcal{F}) = 0$이다.
그러므로 $\mathcal{E}'$은 평탄하다
(보조정리 \ref{cohomology-lemma-flat-tor-zero}).
따라서 Modules, 보조정리
\ref{modules-lemma-flat-locally-finite-presentation}에 의해
$\mathcal{E}'$은 국소적으로 유한 자유 가군의 직합인자이다.
이에 따라 국소적으로 복합체
$$
\mathcal{E}' \to \mathcal{E}^{a + 1} \to \ldots \to \mathcal{E}^b
$$
는 $E$와 준동형이고 $E$는 완전하다.
\end{proof}

\begin{lemma}
\label{cohomology-lemma-perfect}
$(X, \mathcal{O}_X)$를 환 달린 공간이라 하고,
$E$를 $D(\mathcal{O}_X)$의 대상이라 하자.
다음은 동치이다.
\begin{enumerate}
\item $E$는 완전하다.
\item $E$는 유사코히런트이고 국소적으로 유한 Tor 차원을 가진다.
\end{enumerate}
\end{lemma}

\begin{proof}
(1)을 가정하자. 정의에 의해 이는 열린 덮개
$X = \bigcup U_i$가 존재하여 $E|_{U_i}$가
강한 완전 복합체로 표시됨을 뜻한다. 따라서
보조정리 \ref{cohomology-lemma-pseudo-coherent-independent-representative}에 의해
$E$는 유사코히런트이다(즉, 모든 $m$에 대해
$m$-유사코히런트이다). 또한 유한 자유 가군의 직합인자는 평탄하므로
보조정리 \ref{cohomology-lemma-tor-amplitude}에 의해
$E|_{U_i}$는 유한 Tor 차원을 가진다. 따라서 (2)가 성립한다.

\medskip\noindent
(2)를 가정하자. $X$를 열린 덮개의 원소들로 바꾸어
$E$가 $[a, b]$ 안에서 Tor-진폭을 가지게 하는 정수
$a \leq b$가 존재한다고 가정할 수 있다.
$E$는 모든 $m$에 대해 $m$-유사코히런트이므로
보조정리 \ref{cohomology-lemma-perfect-precise}를 사용하여 완전함을 얻는다.
\end{proof}

\begin{lemma}
\label{cohomology-lemma-perfect-pullback}
$f : (X, \mathcal{O}_X) \to (Y, \mathcal{O}_Y)$를 환 달린 공간의
사상이라 하고, $E$를 $D(\mathcal{O}_Y)$의 대상이라 하자.
$E$가 $D(\mathcal{O}_Y)$에서 완전하면
$Lf^*E$는 $D(\mathcal{O}_X)$에서 완전하다.
\end{lemma}

\begin{proof}
이는 보조정리 \ref{cohomology-lemma-perfect},
\ref{cohomology-lemma-tor-amplitude-pullback}, 그리고
\ref{cohomology-lemma-pseudo-coherent-pullback}에서 따른다.
(다른 증명은 보조정리
\ref{cohomology-lemma-pseudo-coherent-pullback}의 증명을 그대로 따르면 된다.)
\end{proof}

\begin{lemma}
\label{cohomology-lemma-two-out-of-three-perfect}
$(X, \mathcal{O}_X)$를 환 달린 공간이라 하고,
$(K, L, M, f, g, h)$를 $D(\mathcal{O}_X)$의 특수 삼각형이라 하자.
$K, L, M$ 가운데 두 대상이 완전하면 나머지 하나도 완전하다.
\end{lemma}

\begin{proof}
첫 번째 증명: 보조정리 \ref{cohomology-lemma-perfect},
\ref{cohomology-lemma-cone-pseudo-coherent}, 그리고
\ref{cohomology-lemma-cone-tor-amplitude}를 결합한다.
두 번째 증명(개요): $K$와 $L$이 완전하다고 하자.
$X$를 열린 덮개의 원소들로 바꾸어 $K$와 $L$이 각각
강한 완전 복합체 $\mathcal{K}^\bullet$과
$\mathcal{L}^\bullet$로 표시된다고 가정할 수 있다.
다시 $X$를 열린 덮개의 원소들로 바꾸어 사상 $K \to L$이
복합체의 사상
$\alpha : \mathcal{K}^\bullet \to \mathcal{L}^\bullet$로 주어진다고
가정할 수 있다. 보조정리 \ref{cohomology-lemma-local-actual}을 보라.
그러면 $M$은 $\alpha$의 원뿔과 동형이고,
이는 보조정리 \ref{cohomology-lemma-cone}에 의해 강한 완전 복합체이다.
\end{proof}

\begin{lemma}
\label{cohomology-lemma-tensor-perfect}
$(X, \mathcal{O}_X)$를 환 달린 공간이라 하자.
$K, L$이 $D(\mathcal{O}_X)$의 완전 대상이면
$K \otimes_{\mathcal{O}_X}^\mathbf{L} L$도 그러하다.
\end{lemma}

\begin{proof}
보조정리 \ref{cohomology-lemma-perfect}, \ref{cohomology-lemma-tensor-pseudo-coherent}, 그리고
\ref{cohomology-lemma-tensor-tor-amplitude}에서 따른다.
\end{proof}

\begin{lemma}
\label{cohomology-lemma-summands-perfect}
$(X, \mathcal{O}_X)$를 환 달린 공간이라 하자.
$K \oplus L$이 $D(\mathcal{O}_X)$의 완전 대상이면
$K$와 $L$도 그러하다.
\end{lemma}

\begin{proof}
보조정리 \ref{cohomology-lemma-perfect}, \ref{cohomology-lemma-summands-pseudo-coherent}, 그리고
\ref{cohomology-lemma-summands-tor-amplitude}에서 따른다.
\end{proof}

\begin{lemma}
\label{cohomology-lemma-pushforward-perfect}
$(X, \mathcal{O}_X)$를 환 달린 공간이라 하고,
$j : U \to X$를 열린 부분공간이라 하자.
$E$를 $D(\mathcal{O}_U)$의 완전 대상이라 하되,
그 코호몰로지 층들의 지지가 닫힌 부분집합
$T \subset U$에 포함되고 $j(T)$는 $X$에서 닫혀 있다고 하자.
그러면 $Rj_*E$는 $D(\mathcal{O}_X)$의 완전 대상이다.
\end{lemma}

\begin{proof}
완전 복합체라는 성질은 $X$ 위에서 국소적이다.
따라서 $Rj_*E$를 $U$와 $V = X \setminus j(T)$에 제한했을 때
완전함을 확인하면 충분하다. $Rj_*E|_U = E$이고 이는 완전하다.
또한 $E|_{U \setminus T} = 0$이므로 $Rj_*E|_V = 0$이다.
\end{proof}

\begin{lemma}
\label{cohomology-lemma-perfect-max-open-coh-loc-free}
$(X, \mathcal{O}_X)$를 환 달린 공간이라 하고,
$D(\mathcal{O}_X)$의 완전 대상 $E$를 택하자.
모든 줄기 $\mathcal{O}_{X, x}$가 국소환이라고 가정하자.
그러면 집합
$$
U =
\{x \in X \mid
H^i(E)_x\text{ is a finite free }
\mathcal{O}_{X, x}\text{-module for all }i\in \mathbf{Z}\}
$$
은 $X$의 열린집합이며, 모든 $i \in \mathbf{Z}$에 대하여
$H^i(E)|_U$가 유한 국소 자유가 되는 가장 큰 열린집합
$U \subset X$이다.
\end{lemma}

\begin{proof}
$V \subset X$가 모든 $i \in \mathbf{Z}$에 대하여
$H^i(E)|_V$가 유한 국소 자유인 열린집합이면
$V \subset U$임에 유의하자.
$x \in U$라 하자. $x$의 어떤 열린 근방이 $U$에 포함되고,
이 근방 위에서 모든 $i$에 대해 $H^i(E)$가 유한 국소 자유임을
보이면 증명이 끝난다. 증명하는 동안 유한 번
$X$를 $x$의 열린 근방으로 바꾸어도 된다.
따라서 $E$가 강한 완전 복합체
$\mathcal{E}^\bullet$로 표시된다고 가정할 수 있다.
$i \not \in [a, b]$이면 $\mathcal{E}^i = 0$이라고 하자.
$b - a$에 대한 귀납법으로 결과를 증명하겠다. 가군
$H^b(E) = \Coker(d^{b - 1} : \mathcal{E}^{b - 1} \to \mathcal{E}^b)$은
유한 표시이다. $H^b(E)_x$가 유한 자유이므로
Modules, 보조정리 \ref{modules-lemma-finite-presentation-stalk-free}에 의해
$x$의 어떤 열린 근방에서 $H^b(E)$가 유한 자유임을 얻는다.
따라서 $X$를 더 작을 수도 있는 열린 근방으로 바꾸어
직합 분해 $\mathcal{E}^b = \Im(d^{b - 1}) \oplus H^b(E)$를 가지며
$H^b(E)$가 유한 자유라고 가정할 수 있다.
보조정리 \ref{cohomology-lemma-local-lift-map}를 보라.
같은 논증을 다시 적용하면
$\mathcal{E}^{b - 1} = \Ker(d^{b - 1}) \oplus \Im(d^{b - 1})$이라고
가정할 수 있다. 복합체
$\mathcal{E}^a \to \ldots \to \mathcal{E}^{b - 2} \to \Ker(d^{b - 1})$는
완전 대상 $E'$을 표시하는 강한 완전 복합체이며
$i \not = b$에 대하여 $H^i(E) = H^i(E')$이다.
그러므로 귀납 가정으로 결론을 얻는다.
\end{proof}
\section{쌍대}
\label{cohomology-section-duals}

\noindent
이 절에서는 복합체 범주와 유도 범주의 쌍대화 가능한 대상들을
특징짓는다. 특히 $D(\mathcal{O}_X)$의 대상이 쌍대를 가질 필요충분조건은
그 대상이 완전한 것임을 보일 것이다(이는 예
\ref{cohomology-example-dual-derived}와 보조정리
\ref{cohomology-lemma-left-dual-derived}에서 따른다).

\begin{lemma}
\label{cohomology-lemma-symmetric-monoidal-cat-complexes}
$(X, \mathcal{O}_X)$를 환 달린 공간이라 하자.
$\mathcal{O}_X$-가군들의 복합체 범주에 텐서곱을
$\mathcal{F}^\bullet \otimes \mathcal{G}^\bullet =
\text{Tot}(\mathcal{F}^\bullet \otimes_{\mathcal{O}_X} \mathcal{G}^\bullet)$로
정하면 이 범주는 대칭 모노이드 범주이다(부호 규칙은
More on Algebra, 절 \ref{more-algebra-section-sign-rules}를 보라).
\end{lemma}

\begin{proof}
생략한다. 힌트: 단위원 $\mathbf{1}$로는 차수 $0$에서
$\mathcal{O}_X$이고 다른 차수에서 영인 복합체를 택하며,
자명한 동형사상
$\text{Tot}(\mathbf{1} \otimes_{\mathcal{O}_X} \mathcal{G}^\bullet) =
\mathcal{G}^\bullet$과
$\text{Tot}(\mathcal{F}^\bullet \otimes_{\mathcal{O}_X} \mathbf{1}) =
\mathcal{F}^\bullet$을 사용한다.
보조정리를 증명하려면 여러 도식의 가환성을 확인해야 한다.
Categories, 정의 \ref{categories-definition-monoidal-category}와
\ref{categories-definition-symmetric-monoidal-category}를 보라.
각각의 경우 확인은 곧바르다.
\end{proof}

\begin{example}
\label{cohomology-example-dual}
$(X, \mathcal{O}_X)$를 환 달린 공간이라 하자.
$\mathcal{F}^\bullet$을 $\mathcal{O}_X$-가군들의 국소적으로 유계인
복합체라 하되, 각 $\mathcal{F}^n$은 국소적으로 유한 자유
$\mathcal{O}_X$-가군의 직합인자라고 하자. 다시 말해,
열린 덮개 $X = \bigcup U_i$가 존재하여
$\mathcal{F}^\bullet|_{U_i}$가 강한 완전 복합체이다.
절 \ref{cohomology-section-hom-complexes}에서와 같이 복합체
$$
\mathcal{G}^\bullet = \SheafHom^\bullet(\mathcal{F}^\bullet, \mathcal{O}_X)
$$
를 생각하자.
$$
\eta :
\mathcal{O}_X
\to
\text{Tot}(\mathcal{F}^\bullet \otimes_{\mathcal{O}_X} \mathcal{G}^\bullet)
\quad\text{and}\quad
\epsilon :
\text{Tot}(\mathcal{G}^\bullet \otimes_{\mathcal{O}_X} \mathcal{F}^\bullet)
\to
\mathcal{O}_X
$$
를 $\eta = \sum \eta_n$과 $\epsilon = \sum \epsilon_n$으로 정의하자.
여기서 $\eta_n : \mathcal{O}_X \to
\mathcal{F}^n \otimes_{\mathcal{O}_X} \mathcal{G}^{-n}$과
$\epsilon_n : \mathcal{G}^{-n} \otimes_{\mathcal{O}_X} \mathcal{F}^n
\to \mathcal{O}_X$은 Modules, 예
\ref{modules-example-dual}에서와 같다.
그러면 $\mathcal{G}^\bullet, \eta, \epsilon$은
Categories, 정의 \ref{categories-definition-dual}의 의미에서
$\mathcal{F}^\bullet$의 왼쪽 쌍대이다.
등식
$(1 \otimes \epsilon) \circ (\eta \otimes 1) = \text{id}_{\mathcal{F}^\bullet}$과
$(\epsilon \otimes 1) \circ (1 \otimes \eta) =
\text{id}_{\mathcal{G}^\bullet}$의 확인은 생략한다.
More on Algebra, 보조정리
\ref{more-algebra-lemma-left-dual-complex}와 비교하라.
\end{example}

\begin{lemma}
\label{cohomology-lemma-left-dual-complex}
$(X, \mathcal{O}_X)$를 환 달린 공간이라 하고,
$\mathcal{F}^\bullet$을 $\mathcal{O}_X$-가군들의 복합체라 하자.
$\mathcal{F}^\bullet$이 $\mathcal{O}_X$-가군 복합체들의 모노이드
범주에서 왼쪽 쌍대를 가지면
(Categories, 정의 \ref{categories-definition-dual}),
$\mathcal{F}^\bullet$은 각 항이 국소적으로 유한 자유
$\mathcal{O}_X$-가군의 직합인자인 국소적으로 유계인 복합체이고,
그 왼쪽 쌍대는 예 \ref{cohomology-example-dual}에서 구성한 것이다.
\end{lemma}

\begin{proof}
왼쪽 쌍대의 유일성
(Categories, 주석 \ref{categories-remark-left-dual-adjoint})에 의해,
$\mathcal{F}^\bullet$이 명시된 성질을 가짐을 보이면 마지막 명제가 따른다.
$\mathcal{G}^\bullet, \eta, \epsilon$을 왼쪽 쌍대라 하자.
$\eta = \sum \eta_n$과 $\epsilon = \sum \epsilon_n$으로 쓰자.
여기서 $\eta_n : \mathcal{O}_X \to
\mathcal{F}^n \otimes_{\mathcal{O}_X} \mathcal{G}^{-n}$이고
$\epsilon_n : \mathcal{G}^{-n} \otimes_{\mathcal{O}_X} \mathcal{F}^n
\to \mathcal{O}_X$이다.
Categories, 정의 \ref{categories-definition-dual}에 의해
$(1 \otimes \epsilon) \circ (\eta \otimes 1) = \text{id}_{\mathcal{F}^\bullet}$과
$(\epsilon \otimes 1) \circ (1 \otimes \eta) = \text{id}_{\mathcal{G}^\bullet}$이므로,
곧바로
$(1 \otimes \epsilon_n) \circ (\eta_n \otimes 1) = \text{id}_{\mathcal{F}^n}$과
$(\epsilon_n \otimes 1) \circ (1 \otimes \eta_n) =
\text{id}_{\mathcal{G}^{-n}}$을 얻는다.
따라서 Modules, 보조정리
\ref{modules-lemma-left-dual-module}에 의해
$\mathcal{F}^n$은 국소적으로 유한 자유
$\mathcal{O}_X$-가군의 직합인자이다.
합 $\eta = \sum \eta_n$이 국소 유한이므로
$\mathcal{F}^\bullet$은 국소적으로 유계이다.
\end{proof}

\begin{lemma}
\label{cohomology-lemma-internal-hom-evaluate-isom}
$(X, \mathcal{O}_X)$를 환 달린 공간이라 하고,
$K, L, M \in D(\mathcal{O}_X)$라 하자.
$K$가 완전하면 보조정리
\ref{cohomology-lemma-internal-hom-evaluate}의 사상
$$
R\SheafHom(L, M) \otimes_{\mathcal{O}_X}^\mathbf{L} K
\longrightarrow
R\SheafHom(R\SheafHom(K, L), M)
$$
은 동형사상이다.
\end{lemma}

\begin{proof}
이 사상은 전역적으로 정의되고 우변과 좌변의 구성은 국소화와
가환하므로(보조정리 \ref{cohomology-lemma-restriction-RHom-to-U}를 보라),
이를 증명할 때 $X$ 위에서 국소적으로 작업해도 된다.
따라서 $K$가 강한 완전 복합체
$\mathcal{E}^\bullet$로 표시된다고 가정할 수 있다.

\medskip\noindent
$K_1 \to K_2 \to K_3$가 $D(\mathcal{O}_X)$의 특수 삼각형이면
특수 삼각형
$$
R\SheafHom(L, M) \otimes_{\mathcal{O}_X}^\mathbf{L} K_1 \to
R\SheafHom(L, M) \otimes_{\mathcal{O}_X}^\mathbf{L} K_2 \to
R\SheafHom(L, M) \otimes_{\mathcal{O}_X}^\mathbf{L} K_3
$$
과
$$
R\SheafHom(R\SheafHom(K_1, L), M) \to
R\SheafHom(R\SheafHom(K_2, L), M)
R\SheafHom(R\SheafHom(K_3, L), M)
$$
을 얻는다. 절 \ref{cohomology-section-flat}과 보조정리
\ref{cohomology-lemma-RHom-triangulated}를 보라.
보조정리 \ref{cohomology-lemma-internal-hom-evaluate}의 화살표는 $K$에 대해
함자적이므로 이 특수 삼각형들 사이의 사상을 얻는다.
따라서 $K_1$과 $K_3$에 대해 결과가 성립하면,
Derived Categories, 보조정리
\ref{derived-lemma-third-isomorphism-triangle}에 의해
$K_2$에 대해서도 결과가 성립한다.

\medskip\noindent
위의 관찰을 단순 절단의 특수 삼각형
$$
\sigma_{\geq n}\mathcal{E}^\bullet \to \mathcal{E}^\bullet \to
\sigma_{\leq n - 1}\mathcal{E}^\bullet
$$
과 결합하면, $K$가 어떤 차수에 놓인 유한 자유
$\mathcal{O}_X$-가군의 직합인자 하나로 이루어진 경우로 귀착된다.
문제가 직합과 호환된다는 자명한 사실(위에서 말한 것과 유사하다)과
시프트를 사용하면, 다시 어떤 정수 $n$에 대해
$K = \mathcal{O}_X^{\oplus n}$인 경우로 귀착된다.
이 경우는 명백하다.
\end{proof}

\begin{lemma}
\label{cohomology-lemma-dual-perfect-complex}
$(X, \mathcal{O}_X)$를 환 달린 공간이라 하고,
$K$를 $D(\mathcal{O}_X)$의 완전 대상이라 하자.
그러면 $K^\vee = R\SheafHom(K, \mathcal{O}_X)$도 완전 대상이고
$(K^\vee)^\vee \cong K$이다. $D(\mathcal{O}_X)$의 $M$에 대하여
함자적 동형사상
$$
M \otimes^\mathbf{L}_{\mathcal{O}_X} K^\vee = R\SheafHom(K, M)
$$
과
$$
H^0(X, M \otimes^\mathbf{L}_{\mathcal{O}_X} K^\vee) =
\Hom_{D(\mathcal{O}_X)}(K, M)
$$
이 있다.
\end{lemma}

\begin{proof}
보조정리 \ref{cohomology-lemma-internal-hom-evaluate}에 의해 표준 사상
$$
K = R\SheafHom(\mathcal{O}_X, \mathcal{O}_X)
\otimes_{\mathcal{O}_X}^\mathbf{L} K \longrightarrow
R\SheafHom(R\SheafHom(K, \mathcal{O}_X), \mathcal{O}_X) =
(K^\vee)^\vee
$$
이 있으며, 보조정리 \ref{cohomology-lemma-internal-hom-evaluate-isom}에 의해
이는 동형사상이다. 다른 명제들을 확인할 때 내부 Hom의 구성이
열린집합으로의 제한과 가환한다는 사실
(보조정리 \ref{cohomology-lemma-restriction-RHom-to-U})을 더 언급하지 않고 사용하겠다.
$K^\vee$가 완전한지는 $X$ 위에서 국소적으로 확인할 수 있다.
보조정리 \ref{cohomology-lemma-dual}에 의해 마지막 명제를 보이려면 사상
(\ref{cohomology-equation-eval})
$$
M \otimes^\mathbf{L}_{\mathcal{O}_X} K^\vee
\longrightarrow
R\SheafHom(K, M)
$$
이 동형사상임을 확인하면 충분하다. 이 또한 $X$ 위에서 국소적인 문제이다.
따라서 이 두 명제는 $K$가 강한 완전 복합체로 표시되는 경우에
증명하면 충분하다.

\medskip\noindent
$K$가 강한 완전 복합체 $\mathcal{E}^\bullet$로 표시된다고 하자.
그러면 보조정리 \ref{cohomology-lemma-Rhom-strictly-perfect}에 의해
$K^\vee$는 차수 $n$의 항이
$(\mathcal{E}^{-n})^\vee =
\SheafHom_{\mathcal{O}_X}(\mathcal{E}^{-n}, \mathcal{O}_X)$인 복합체로
표시된다. $\mathcal{E}^{-n}$이 유한 자유 $\mathcal{O}_X$-가군의
직합인자이므로 $(\mathcal{E}^{-n})^\vee$도 그러하다.
따라서 $K^\vee$도 강한 완전 복합체로 표시되며,
그 결과 $K^\vee$가 완전함을 알 수 있다.
(\ref{cohomology-equation-eval})이 동형사상임을 보이기 위해
$M$을 복합체 $\mathcal{F}^\bullet$로 표시하자.
보조정리 \ref{cohomology-lemma-Rhom-strictly-perfect}에 의해
$R\SheafHom(K, M)$은 그 항들이
$$
\bigoplus\nolimits_{n = p + q}
\SheafHom_{\mathcal{O}_X}(\mathcal{E}^{-q}, \mathcal{F}^p)
$$
인 복합체로 표시된다. 한편 대상
$M \otimes^\mathbf{L}_{\mathcal{O}_X} K^\vee$는 그 항들이
$$
\bigoplus\nolimits_{n = p + q}
\mathcal{F}^p \otimes_{\mathcal{O}_X} (\mathcal{E}^{-q})^\vee
$$
인 복합체로 표시된다. 따라서 (\ref{cohomology-equation-eval})이
동형사상이라는 명제는 표준 사상
$$
\mathcal{F}
\otimes_{\mathcal{O}_X}
\SheafHom_{\mathcal{O}_X}(\mathcal{E}, \mathcal{O}_X)
\longrightarrow
\SheafHom_{\mathcal{O}_X}(\mathcal{E}, \mathcal{F})
$$
이, $\mathcal{E}$가 유한 자유 $\mathcal{O}_X$-가군의 직합인자이고
$\mathcal{F}$가 임의의 $\mathcal{O}_X$-가군일 때 동형사상이라는
명제로 귀착된다. 이는 $\mathcal{E}$가 유한 자유인 경우의
대응하는 명제에서 곧바로 따른다.
\end{proof}

\begin{lemma}
\label{cohomology-lemma-symmetric-monoidal-derived}
$(X, \mathcal{O}_X)$를 환 달린 공간이라 하자.
유도 텐서곱을 텐서곱으로 하고 통상적인 결합성과 가환성 제약을
갖춘 유도 범주 $D(\mathcal{O}_X)$는 대칭 모노이드 범주이다
(부호 규칙은 More on Algebra, 절
\ref{more-algebra-section-sign-rules}를 보라).
\end{lemma}

\begin{proof}
생략한다. 보조정리
\ref{cohomology-lemma-symmetric-monoidal-cat-complexes}와 비교하라.
\end{proof}

\begin{example}
\label{cohomology-example-dual-derived}
$(X, \mathcal{O}_X)$를 환 달린 공간이라 하고,
$K$를 $D(\mathcal{O}_X)$의 완전 대상이라 하자.
보조정리 \ref{cohomology-lemma-dual-perfect-complex}에서와 같이
$K^\vee = R\SheafHom(K, \mathcal{O}_X)$로 놓자.
그러면 사상
$$
K \otimes_{\mathcal{O}_X}^\mathbf{L} K^\vee \longrightarrow R\SheafHom(K, K)
$$
은 (그 보조정리에 의해) 동형사상이다.
$$
\eta :
\mathcal{O}_X
\longrightarrow
K \otimes_{\mathcal{O}_X}^\mathbf{L} K^\vee
$$
를 위 동형사상 아래에서 $\text{id}_K$에 대응하는 절단으로
$1$을 보내는 사상이라 하자.
$$
\epsilon : 
K^\vee
\otimes_{\mathcal{O}_X}^\mathbf{L} K
\longrightarrow
\mathcal{O}_X
$$
를 평가 사상이라 하자(예를 들어 이를 구성할 때 보조정리
\ref{cohomology-lemma-internal-hom-composition}을 사용할 수 있다).
그러면 $K^\vee, \eta, \epsilon$은 Categories, 정의
\ref{categories-definition-dual}의 의미에서 $K$의 왼쪽 쌍대이다.
등식
$(1 \otimes \epsilon) \circ (\eta \otimes 1) = \text{id}_K$과
$(\epsilon \otimes 1) \circ (1 \otimes \eta) =
\text{id}_{K^\vee}$의 확인은 생략한다.
\end{example}

\begin{lemma}
\label{cohomology-lemma-left-dual-derived}
$(X, \mathcal{O}_X)$를 환 달린 공간이라 하고,
$M$을 $D(\mathcal{O}_X)$의 대상이라 하자.
$M$이 모노이드 범주 $D(\mathcal{O}_X)$에서 왼쪽 쌍대를 가지면
(Categories, 정의 \ref{categories-definition-dual}),
$M$은 완전하고 그 왼쪽 쌍대는 예
\ref{cohomology-example-dual-derived}에서 구성한 것이다.
\end{lemma}

\begin{proof}
$x \in X$라 하자. $M$을 $U$ 위로 제한했을 때 완전 복합체가 되는
$x$의 열린 근방 $U$를 찾으면 충분하다. 따라서 증명하는 동안
유한 번 $X$를 $x$의 열린 근방으로 바꾸어도 된다.
$N, \eta, \epsilon$을 왼쪽 쌍대라 하자.

\medskip\noindent
다음 논증을 여러 번 사용하겠다.
$M$을 표시하는 임의의 $\mathcal{O}_X$-가군 복합체
$\mathcal{M}^\bullet$을 택하자.
$N$을 표시하며 그 항들이 평탄한 $\mathcal{O}_X$-가군인
K-평탄 복합체 $\mathcal{N}^\bullet$을 택하자.
보조정리 \ref{cohomology-lemma-K-flat-resolution}를 보라.
사상
$$
\eta : \mathcal{O}_X \to
\text{Tot}(\mathcal{M}^\bullet \otimes_{\mathcal{O}_X} \mathcal{N}^\bullet)
$$
을 생각하자. $X$를 축소한 뒤, 정수 $N$과
$i = 1, \ldots, N$에 대한 정수 $n_i \in \mathbf{Z}$ 및
$\mathcal{M}^{n_i}$와 $\mathcal{N}^{-n_i}$의 절단 $f_i$, $g_i$를 택하여
$$
\eta(1) = \sum\nolimits_i f_i \otimes g_i
$$
가 되게 할 수 있다. $\mathcal{K}^\bullet \subset \mathcal{M}^\bullet$을
모든 $i = 1, \ldots, N$에 대한 절단 $f_i$들을 포함하는
임의의 $\mathcal{O}_X$-가군 부분복합체라 하자.
가군 $\mathcal{N}^n$들이 평탄하므로
$\text{Tot}(\mathcal{K}^\bullet \otimes_{\mathcal{O}_X} \mathcal{N}^\bullet)
\subset
\text{Tot}(\mathcal{M}^\bullet \otimes_{\mathcal{O}_X} \mathcal{N}^\bullet)$이고,
따라서 $\eta$는
$$
\tilde \eta :
\mathcal{O}_X \to
\text{Tot}(\mathcal{K}^\bullet \otimes_{\mathcal{O}_X} \mathcal{N}^\bullet)
$$
를 통해 인수분해된다.
$\mathcal{K}^\bullet$로 표시되는 $D(\mathcal{O}_X)$의 대상을
$K$라 하면 다음 가환 도식을 얻는다.
$$
\xymatrix{
M \ar[rr]_-{\eta \otimes 1} \ar[rrd]_{\tilde \eta \otimes 1} & &
M \otimes^\mathbf{L} N \otimes^\mathbf{L} M
\ar[r]_-{1 \otimes \epsilon} &
M \\
& &
K \otimes^\mathbf{L} N \otimes^\mathbf{L} M
\ar[u] \ar[r]^-{1 \otimes \epsilon} &
K \ar[u]
}
$$
위 행의 합성이 $M$ 위의 항등사상이므로
$M$은 $D(\mathcal{O}_X)$에서 $K$의 직합인자이다.

\medskip\noindent
위 논증의 첫 번째 응용으로,
모든 $i = 1, \ldots, N$에 대하여 $a < n_i < b$가 되도록
부분복합체
$\mathcal{K}^\bullet = \sigma_{\geq a} \tau_{\leq b}\mathcal{M}^\bullet$을
택할 수 있다. 따라서 $M$은 $D(\mathcal{O}_X)$에서
유계 복합체의 직합인자이며, $M$이 $D^b(\mathcal{O}_X)$에 속한다고
가정할 수 있음을 얻는다. (위 과정에서 $X$를 축소한다는 점을 상기하라.)

\medskip\noindent
$M$이 $D^b(\mathcal{O}_X)$에 속하므로
$\mathcal{M}^\bullet$을 평탄 가군들의 위로 유계인 복합체로
택할 수 있다(Modules, 보조정리
\ref{modules-lemma-module-quotient-flat}과 Derived Categories,
보조정리 \ref{derived-lemma-subcategory-left-resolution}에 의한다).
그러면 위 논증에서 모든 $i = 1, \ldots, N$에 대하여
$a < n_i$가 되도록
$\mathcal{K}^\bullet = \sigma_{\geq a}\mathcal{M}^\bullet$을 택할 수 있다.
따라서 $M$이 $D(\mathcal{O}_X)$에서 평탄 가군들의 유계 복합체의
직합인자라고 가정할 수 있다. 특히 $M$은 유한 Tor-진폭을 가진다.

\medskip\noindent
$M$이 $[a, b]$ 안에서 Tor-진폭을 가진다고 하자.
$M$이 $m$-유사코히런트라고 가정하고, $X$를 축소한 뒤
$M$이 $(m - 1)$-유사코히런트라고 가정할 수 있음을 보이겠다.
그러면 보조정리 \ref{cohomology-lemma-perfect-precise}와
어쨌든 $M$이 $(b + 1)$-유사코히런트라는 사실에 의해 증명이 끝난다.
$X$를 축소하여 강한 완전 복합체 $\mathcal{E}^\bullet$과
$D(\mathcal{O}_X)$의 사상
$\alpha : \mathcal{E}^\bullet \to M$이 존재하고,
$i > m$에 대하여 $H^i(\alpha)$가 동형이며
$i = m$에 대하여 전사라고 가정할 수 있다.
$i < m$이면 $\mathcal{E}^i = 0$이라고 가정해도 된다.
특수 삼각형
$$
\mathcal{E}^\bullet \to M \to L \to \mathcal{E}^\bullet[1]
$$
을 택하자. $i \geq m$이면 $H^i(L) = 0$임에 유의하자.
따라서 $L$을, $i \geq m$이면 $\mathcal{L}^i = 0$인 복합체
$\mathcal{L}^\bullet$로 표시할 수 있다.
사상 $L \to \mathcal{E}^\bullet[1]$은 복합체의 사상
$\mathcal{L}^\bullet \to \mathcal{E}^\bullet[1]$로 주어지며,
이는 차수 $m - 1$을 제외한 모든 차수에서 영이고
그 차수에서는 사상
$\mathcal{L}^{m - 1} \to \mathcal{E}^m$을 준다.
Derived Categories, 보조정리
\ref{derived-lemma-negative-exts}를 보라.
그러면 $M$은 복합체
$$
\mathcal{M}^\bullet :
\ldots \to
\mathcal{L}^{m - 2} \to
\mathcal{L}^{m - 1} \to
\mathcal{E}^m \to
\mathcal{E}^{m + 1} \to \ldots
$$
로 표시된다. 이 복합체에 둘째 문단의 논의를 적용하여
모든 $i = 1, \ldots, N$에 대해
$\mathcal{M}^{n_i}$의 절단 $f_i$를 얻는다.
$n < m$에 대해서는 $\mathcal{K}^n \subset \mathcal{L}^n$을
$n_i = n$인 $f_i$들과 $n_i = n - 1$인 $d(f_i)$들로 생성되는
$\mathcal{O}_X$-부분가군으로 정의하자.
$n \geq m$에 대해서는 $\mathcal{K}^n = \mathcal{E}^n$으로 놓자.
그러면 명백히 다음 특수 삼각형들의 사상이 있다.
$$
\xymatrix{
\mathcal{E}^\bullet \ar[r] &
\mathcal{M}^\bullet \ar[r] &
\mathcal{L}^\bullet \ar[r] &
\mathcal{E}^\bullet[1] \\
\mathcal{E}^\bullet \ar[r] \ar[u] &
\mathcal{K}^\bullet \ar[r] \ar[u] &
\sigma_{\leq m - 1}\mathcal{K}^\bullet \ar[r] \ar[u] &
\mathcal{E}^\bullet[1] \ar[u]
}
$$
여기서 모든 사상은 위에서 지시한 것들이다.
복합체 $\mathcal{K}^\bullet$에 대응하는
$D(\mathcal{O}_X)$의 대상을 $K$라 하자.
증명의 둘째 문단에 나온 논증으로
$D(\mathcal{O}_X)$의 사상 $s : M \to K$를 얻을 수 있으며,
합성 $M \to K \to M$은 $M$ 위의 항등사상이다.
도식
$$
\xymatrix{
\mathcal{E}^\bullet \ar[r] &
\mathcal{K}^\bullet \ar@{=}[r] &
K \\
\mathcal{E}^\bullet \ar[u]^{\text{id}} \ar[r]^i &
\mathcal{M}^\bullet \ar@{=}[r] &
M \ar[u]_s
}
$$
이 가환하는지는 알 수 없지만, 사상 $K \to M$과 합성한 뒤에는
가환함을 알고 있다. 보조정리 \ref{cohomology-lemma-local-actual}에 의해
$X$를 축소하여 $s \circ i$가 복합체의 사상
$\sigma : \mathcal{E}^\bullet \to \mathcal{K}^\bullet$로
주어진다고 가정할 수 있다. 같은 보조정리에 의해,
$\sigma$와 포함사상
$\mathcal{K}^\bullet \subset \mathcal{M}^\bullet$의 합성은 어떤 호모토피
$\{h^i : \mathcal{E}^i \to \mathcal{M}^{i - 1}\}$에 의해
영과 호모토픽이라고 가정할 수 있다.
따라서 $\mathcal{K}^{m - 1}$을
$\mathcal{K}^{m - 1} + \Im(h^m)$으로 바꾼 뒤에도
($\mathcal{K}^{m - 1}$이 유한 개의 전역 절단으로 생성된다는
성질은 여전히 성립한다) $\sigma$ 자체가 영과 호모토픽임을 알 수 있다!
이는 다음 실선 도식이 가환함을 뜻한다.
$$
\xymatrix{
\mathcal{E}^\bullet \ar[r] &
M \ar[r] &
\mathcal{L}^\bullet \ar[r] &
\mathcal{E}^\bullet[1] \\
\mathcal{E}^\bullet \ar[r] \ar[u] &
K \ar[r] \ar[u] &
\sigma_{\leq m - 1}\mathcal{K}^\bullet \ar[r] \ar[u] &
\mathcal{E}^\bullet[1] \ar[u] \\
\mathcal{E}^\bullet \ar[r] \ar[u] &
M \ar[r] \ar[u]^s &
\mathcal{L}^\bullet \ar[r] \ar@{..>}[u] &
\mathcal{E}^\bullet[1] \ar[u]
}
$$
삼각 범주의 공리들에 의해 이 도식에 맞는 점선 화살표를 얻는다.
차수 $m - 1$의 코호몰로지 층을 살펴보면 다음을 얻는다.
$$
\xymatrix{
H^{m - 1}(M) \ar[r] &
H^{m - 1}(\mathcal{L}^\bullet) \ar[r] &
H^m(\mathcal{E}^\bullet) \\
H^{m - 1}(K) \ar[r] \ar[u] &
H^{m - 1}(\sigma_{\leq m - 1}\mathcal{K}^\bullet) \ar[r] \ar[u] &
H^m(\mathcal{E}^\bullet) \ar[u] \\
H^{m - 1}(M) \ar[r] \ar[u] &
H^{m - 1}(\mathcal{L}^\bullet) \ar[r] \ar[u] &
H^m(\mathcal{E}^\bullet) \ar[u]
}
$$
왼쪽 열과 오른쪽 열에서 수직 합성이 모두 항등사상이므로,
가운데의 수직 합성
$H^{m - 1}(\mathcal{L}^\bullet) \to
H^{m - 1}(\sigma_{\leq m - 1}\mathcal{K}^\bullet) \to
H^{m - 1}(\mathcal{L}^\bullet)$은 전사이다!
특히 $H^{m - 1}(\sigma_{\leq m - 1}\mathcal{K}^\bullet) \to
H^{m - 1}(\mathcal{L}^\bullet)$은 전사이다.
위 삼각형 사상에서 유도되는 코호몰로지 층의 긴 완전 수열의
사상을 사용하면, 도식 추적을 통해 이는
$i \geq m$에 대하여 $H^i(K) \to H^i(M)$이 동형이고
$i = m - 1$에 대하여 전사임을 함의한다.
구성에 의해 어떤 $r \geq 0$과 전사
$\mathcal{O}_X^{\oplus r} \to \mathcal{K}^{m - 1}$을 택할 수 있다.
그러면 합성
$$
(\mathcal{O}_X^{\oplus r} \to \mathcal{E}^m \to
\mathcal{E}^{m + 1} \to \ldots ) \longrightarrow
K \longrightarrow M
$$
은 차수 $\geq m$에서 코호몰로지 층에 동형사상을 유도하고
차수 $m - 1$에서 전사를 유도하므로 증명이 끝난다.
\end{proof}
\section{기타 결과}
\label{cohomology-section-misc}

\noindent
다른 곳에 넣기 어려운 몇 가지 결과를 모은다.

\begin{lemma}
\label{cohomology-lemma-colim-and-lim-of-duals}
$(X, \mathcal{O}_X)$를 환 달린 공간이라 하고,
$(K_n)_{n \in \mathbf{N}}$을 $D(\mathcal{O}_X)$의 완전 대상들의
계를 이루는 것이라 하자. $K = \text{hocolim} K_n$을 유도 여극한이라 하자
(Derived Categories, 정의
\ref{derived-definition-derived-colimit}).
그러면 $D(\mathcal{O}_X)$의 임의의 대상 $E$에 대하여
$$
R\SheafHom(K, E) = R\lim E \otimes^\mathbf{L}_{\mathcal{O}_X} K_n^\vee
$$
이다. 여기서 $(K_n^\vee)$은 쌍대 완전 복합체들의 역계이다.
\end{lemma}

\begin{proof}
보조정리 \ref{cohomology-lemma-dual-perfect-complex}에 의해
$R\lim E \otimes^\mathbf{L}_{\mathcal{O}_X} K_n^\vee =
R\lim R\SheafHom(K_n, E)$이며, 이는 특수 삼각형
$$
R\lim R\SheafHom(K_n, E) \to
\prod R\SheafHom(K_n, E) \to
\prod R\SheafHom(K_n, E)
$$
에 놓인다. 마찬가지로 $K$는 특수 삼각형
$\bigoplus K_n \to \bigoplus K_n \to K$에 놓이므로,
$\prod R\SheafHom(K_n, E) = R\SheafHom(\bigoplus K_n, E)$를
보이면 충분하다. 이는 (\ref{cohomology-equation-internal-hom})과
유도 텐서곱이 직합과 가환한다는 사실의 형식적 귀결이다.
\end{proof}

\begin{lemma}
\label{cohomology-lemma-ext-composition-is-cup}
$(X, \mathcal{O}_X)$를 환 달린 공간이라 하자.
$K$와 $E$를 $D(\mathcal{O}_X)$의 대상이라 하되
$E$는 완전하다고 하자. 도식
$$
\xymatrix{
H^0(X, K \otimes_{\mathcal{O}_X}^\mathbf{L} E^\vee) \times H^0(X, E)
\ar[r] \ar[d] &
H^0(X, K \otimes_{\mathcal{O}_X}^\mathbf{L} E^\vee
\otimes_{\mathcal{O}_X}^\mathbf{L} E) \ar[d] \\
\Hom_X(E, K) \times H^0(X, E) \ar[r] &
H^0(X, K)
}
$$
은 가환한다. 여기서 위쪽 수평 화살표는 컵곱이고,
오른쪽 수직 화살표는
$\epsilon : E^\vee \otimes_{\mathcal{O}_X}^\mathbf{L} E \to \mathcal{O}_X$
(예 \ref{cohomology-example-dual-derived})을 사용하며,
왼쪽 수직 화살표는 보조정리
\ref{cohomology-lemma-dual-perfect-complex}를 사용하고,
아래쪽 수평 화살표는 자명한 사상이다.
\end{lemma}

\begin{proof}
$\otimes = \otimes_{\mathcal{O}_X}^\mathbf{L}$와
$\mathcal{O} = \mathcal{O}_X$로 줄여 쓰겠다.
$E$와 $K$를 각각 $R\SheafHom(\mathcal{O}, E)$ 및
$R\SheafHom(\mathcal{O}, K)$와 동일시하고,
$E^\vee$를 $R\SheafHom(E, \mathcal{O})$와 동일시하겠다.

\medskip\noindent
$\xi \in H^0(X, K \otimes E^\vee)$와
$\eta \in H^0(X, E)$를 택하자.
$D(\mathcal{O})$에서 대응하는 사상을
$\tilde \xi : \mathcal{O} \to K \otimes E^\vee$와
$\tilde \eta : \mathcal{O} \to E$로 나타내자.
보조정리 \ref{cohomology-lemma-second-cup-equals-first}에 의해 컵곱
$\xi \cup \eta$는
$\tilde \xi \otimes \tilde \eta : \mathcal{O} \to
K \otimes E^\vee \otimes E$에 대응한다.

\medskip\noindent
보조정리 \ref{cohomology-lemma-dual-perfect-complex}에 의해 $\xi$에 대응하는
사상 $\xi' : E \to K$는 합성
$$
E = \mathcal{O} \otimes E
\xrightarrow{\tilde \xi \otimes 1_E}
K \otimes E^\vee \otimes E
\xrightarrow{1_K \otimes \epsilon}
K
$$
이라고 주장한다. 보조정리
\ref{cohomology-lemma-dual-perfect-complex}의 구성은 평가 사상
(\ref{cohomology-equation-eval})을 사용하는데, 이 사상은 다시
$E$와 $R\SheafHom(\mathcal{O}, E)$의 동일시 및
보조정리 \ref{cohomology-lemma-internal-hom-composition}에서 구성한 합성
$\underline{\circ}$를 사용하여 구성된다.
따라서 $\xi'$은 다음 합성이다.
\begin{align*}
E = \mathcal{O} \otimes
R\SheafHom(\mathcal{O}, E)
& \xrightarrow{\tilde \xi \otimes 1}
R\SheafHom(\mathcal{O}, K) \otimes
R\SheafHom(E, \mathcal{O}) \otimes
R\SheafHom(\mathcal{O}, E) \\
& \xrightarrow{\underline{\circ} \otimes 1}
R\SheafHom(E, K) \otimes R\SheafHom(\mathcal{O}, E) \\
& \xrightarrow{\underline{\circ}}
R\SheafHom(\mathcal{O}, K) = K
\end{align*}
이 사실과 보조정리
\ref{cohomology-lemma-internal-hom-composition}에서 구성한 합성
$\underline{\circ}$가 결합적이라는 사실
(향후 참고문헌을 여기에 삽입), 그리고 예
\ref{cohomology-example-dual-derived}에서 $\epsilon$이 합성
$\underline{\circ} : E^\vee \otimes E \to \mathcal{O}$로
정의된다는 사실로부터 주장이 곧바로 따른다.

\medskip\noindent
앞의 두 문단의 결과를 사용하면, 보조정리의 명제는
$(1_K \otimes \epsilon) \circ (\tilde \xi \otimes \tilde \eta)$가
$(1_K \otimes \epsilon) \circ (\tilde \xi \otimes 1_E)
\circ (1_\mathcal{O} \otimes \tilde \eta)$와 같다는 것이며,
이는 즉시 성립한다.
\end{proof}

\begin{lemma}
\label{cohomology-lemma-pullback-internal-hom}
$h : X \to Y$를 환 달린 공간의 사상이라 하자.
$K, M$을 $D(\mathcal{O}_Y)$의 대상이라 하자.
주석 \ref{cohomology-remark-prepare-fancy-base-change}의 표준 사상
$$
Lh^*R\SheafHom(K, M) \longrightarrow R\SheafHom(Lh^*K, Lh^*M)
$$
은 다음 경우들에 동형사상이다.
\begin{enumerate}
\item $K$가 완전한 경우,
\item $h$가 평탄하고 $K$가 유사코히런트이며
$M$이 (국소적으로) 아래로 유계인 경우,
\item $\mathcal{O}_X$가 $h^{-1}\mathcal{O}_Y$ 위에서 유한 Tor 차원을
가지고, $K$가 유사코히런트이며 $M$이 (국소적으로) 아래로 유계인 경우.
\end{enumerate}
\end{lemma}

\begin{proof}
(1)의 증명. 문제는 $Y$ 위에서 국소적이므로,
절 \ref{cohomology-section-perfect}에 따라 $K$가 강한 완전 복합체
$\mathcal{E}^\bullet$로 표시된다고 가정할 수 있다.
$M$을 표시하는 K-평탄 복합체 $\mathcal{F}^\bullet$을 택하자.
보조정리 \ref{cohomology-lemma-Rhom-strictly-perfect}를 적용하면
$R\SheafHom(K, L)$은 항들이
$\mathcal{H}^n = \bigoplus\nolimits_{n = p + q}
\SheafHom_{\mathcal{O}_X}(\mathcal{E}^{-q}, \mathcal{F}^p)$인 복합체
$\mathcal{H}^\bullet =
\SheafHom^\bullet(\mathcal{E}^\bullet, \mathcal{F}^\bullet)$로 표시된다.
절 \ref{cohomology-section-derived-pullback}의 $Lh^*$ 구성으로부터
$Lh^*K$가 강한 완전 복합체
$h^*\mathcal{E}^\bullet$로 표시됨을 알 수 있다
(보조정리 \ref{cohomology-lemma-strictly-perfect-pullback}).
마찬가지로 대상 $Lh^*M$은 복합체
$h^*\mathcal{F}^\bullet$로 표시된다.
마지막으로 $\mathcal{H}^\bullet$이 보조정리
\ref{cohomology-lemma-Rhom-strictly-perfect-K-flat}에 의해 K-평탄하므로,
대상 $Lh^*R\SheafHom(K, M)$은
$h^*\mathcal{H}^\bullet$로 표시된다.
따라서 증명을 끝내려면
$h^*\mathcal{H}^\bullet =
\SheafHom^\bullet(h^*\mathcal{E}^\bullet, h^*\mathcal{F}^\bullet)$임을
보이면 충분하다. 이를 위해서는 $\mathcal{E}$가 유한 자유
$\mathcal{O}_X$-가군의 직합인자일 때
$h^*\SheafHom(\mathcal{E}, \mathcal{F}) =
\SheafHom(h^*\mathcal{E}, \mathcal{F})$라는 사실에 유의하면 충분하다.

\medskip\noindent
(2)의 증명. $h$가 평탄하므로 임의의
$\mathcal{O}_Y$-가군 복합체에 단순히 $h^*$를 적용하여
$Lh^*$를 계산할 수 있다. 특히 모든 $i \in \mathbf{Z}$에 대하여
$H^i(Lh^*K) = h^*H^i(K)$이다.
$i < a$이면 $H^i(M) = 0$이라고 하자.
$K' \to K$를 모든 $i \geq b$에 대하여
$H^i(K') \to H^i(K)$가 동형사상이 되게 하는
$D(\mathcal{O}_Y)$의 사상이라 하자. 그러면 대응하는 사상
$$
R\SheafHom(K, M) \to R\SheafHom(K', M)
$$
과
$$
R\SheafHom(Lh^*K, Lh^*M) \to R\SheafHom(Lh^*K', Lh^*M)
$$
은 차수 $< a - b$에서 코호몰로지 층에 동형사상을 유도한다
(세부 사항은 생략한다). 따라서 보조정리의 명제에 나오는 사상이
차수 $< a - b$에서 코호몰로지 층에 동형사상을 유도함을 보이려면,
그 차수들에서 $K'$에 대한 결과를 증명하면 충분하다.
또한 (1)의 증명에서처럼 문제는 $Y$ 위에서 국소적이다.
따라서 절 \ref{cohomology-section-pseudo-coherent}에 의해 $K$가
강한 완전 복합체로 표시된다고 가정할 수 있다.
이로써 (1)의 경우로 귀착된다.

\medskip\noindent
(3)의 증명. $Lh^*$가 유계 코호몰로지 차원을 가진다는 사실을
사용하여 필요한 소멸을 얻는다는 점을 제외하면 (2)의 증명과 같다.
세부 사항은 생략한다.
\end{proof}

\begin{lemma}
\label{cohomology-lemma-stalk-internal-hom}
$X$를 환 달린 공간이라 하고,
$K, M$을 $D(\mathcal{O}_X)$의 대상이라 하자.
$x \in X$라 하자. 표준 사상
$$
R\SheafHom(K, M)_x \longrightarrow
R\Hom_{\mathcal{O}_{X, x}}(K_x, M_x)
$$
은 다음 경우들에 동형사상이다.
\begin{enumerate}
\item $K$가 완전한 경우,
\item $K$가 유사코히런트이고 $M$이 (국소적으로) 아래로 유계인 경우.
\end{enumerate}
\end{lemma}

\begin{proof}
$Y = \{x\}$를 $\mathcal{O}_{X, x}$를 구조층으로 가지는 한 점짜리
환 달린 공간이라 하자. 그러면 평탄한 포함사상
$Y \to X$에 보조정리 \ref{cohomology-lemma-pullback-internal-hom}을 적용한다.
\end{proof}
\section{유도 범주의 가역 대상}
\label{cohomology-section-invertible-D-or-R}

\noindent
환 달린 공간의 유도 범주에서 가역인 대상들을 특징짓는다.
구조층의 줄기들이 국소환인 경우와 그렇지 않은 경우를 모두 다룬다.

\begin{lemma}
\label{cohomology-lemma-category-summands-finite-free}
$(X, \mathcal{O}_X)$를 환 달린 공간이라 하자.
$R = \Gamma(X, \mathcal{O}_X)$로 놓자.
유한 자유 $\mathcal{O}_X$-가군의 직합인자인
$\mathcal{O}_X$-가군들의 범주는 유한 사영 $R$-가군들의 범주와 동치이다.
\end{lemma}

\begin{proof}
유한 사영 $R$-가군은 유한 자유 $R$-가군의 직합인자와 같은 것임에
유의하자. 동치는 함자
$\mathcal{E} \mapsto \Gamma(X, \mathcal{E})$로 주어진다.
그 역함자는 Modules, 보조정리
\ref{modules-lemma-construct-quasi-coherent-sheaves}의 구성으로 주어진다.
\end{proof}

\begin{lemma}
\label{cohomology-lemma-invertible-derived}
$(X, \mathcal{O}_X)$를 환 달린 공간이라 하고,
$M$을 $D(\mathcal{O}_X)$의 대상이라 하자. 다음은 동치이다.
\begin{enumerate}
\item $M$은 $D(\mathcal{O}_X)$에서 가역이다.
Categories, 정의 \ref{categories-definition-invertible}를 보라.
\item 국소 유한 직접곱 분해
$$
\mathcal{O}_X = \prod\nolimits_{n \in \mathbf{Z}} \mathcal{O}_n
$$
가 존재하고, 각 $n$에 대해 가역
$\mathcal{O}_n$-가군 $\mathcal{H}^n$이 존재하며
(Modules, 정의 \ref{modules-definition-invertible}),
$D(\mathcal{O}_X)$에서
$M = \bigoplus \mathcal{H}^n[-n]$이다.
\end{enumerate}
(1)과 (2)가 성립하면 $M$은 $D(\mathcal{O}_X)$의 완전 대상이다.
모든 $x \in X$에 대해 $\mathcal{O}_{X, x}$가 국소환이면,
이 조건들은 다음 조건과도 동치이다.
\begin{enumerate}
\item[(3)] 열린 덮개 $X = \bigcup U_i$가 존재하고,
각 $i$에 대해 정수 $n_i$가 존재하여 $M|_{U_i}$가 차수
$n_i$에 놓인 가역 $\mathcal{O}_{U_i}$-가군으로 표시된다.
\end{enumerate}
\end{lemma}

\begin{proof}
(2)를 가정하자. 대상 $R\SheafHom(M, \mathcal{O}_X)$과 합성 사상
$$
R\SheafHom(M, \mathcal{O}_X) \otimes_{\mathcal{O}_X}^\mathbf{L} M \to
\mathcal{O}_X
$$
을 생각하자. 이것이 동형사상임을 보이기 위해 국소적으로 작업해도 된다.
따라서
$\mathcal{O}_X = \prod_{a \leq n \leq b} \mathcal{O}_n$이고
$M = \bigoplus_{a \leq n \leq b} \mathcal{H}^n[-n]$이라고
가정할 수 있다. 그러면
$$
R\SheafHom(\mathcal{H}^m, \mathcal{O}_X)
\otimes_{\mathcal{O}_X}^\mathbf{L} \mathcal{H}^n
$$
이 $n \not = m$이면 영이고 $n = m$이면 $\mathcal{O}_n$과
같음을 보이면 충분하다. $n \not = m$인 경우는
$\mathcal{O}_n$과 $\mathcal{O}_m$이 평탄한
$\mathcal{O}_X$-대수이고
$\mathcal{O}_n \otimes_{\mathcal{O}_X} \mathcal{O}_m = 0$이라는
사실에서 따른다. 가역 $\mathcal{O}_X$-가군의 국소 구조
(Modules, 보조정리 \ref{modules-lemma-invertible})를 사용하고
국소적으로 작업하면 $n = m$인 경우의 동형사상도 곧바로 따른다.
세부 사항은 생략한다. $D(\mathcal{O}_X)$가 대칭 모노이드 범주이므로
$M$이 가역임을 얻는다.

\medskip\noindent
(1)을 가정하자. (2)의 서술은 $\mathcal{O}_n$의 후보가
$\SheafHom_{\mathcal{O}_X}(H^n(M), H^n(M))$임을 보여 준다.
이들이 국소 유한한 환의 층들의 족이고
$\mathcal{O}_X = \prod \mathcal{O}_n$이면, 직접곱 분해에서 오는
$\mathcal{O}_X$의 멱등원들을 사용하여 곧바로 직접합 분해
$M = \bigoplus H^n(M)[-n]$을 얻는다.
따라서 (2)를 증명할 때 $X$ 위에서 국소적으로 작업해도 된다.

\medskip\noindent
$D(\mathcal{O}_X)$의 대상 $N$과 동형사상
$M \otimes_{\mathcal{O}_X}^\mathbf{L} N \cong \mathcal{O}_X$를 택하자.
$x \in X$라 하자. 그러면 $N$은 모노이드 범주
$D(\mathcal{O}_X)$에서 $M$의 왼쪽 쌍대이고,
보조정리 \ref{cohomology-lemma-left-dual-derived}에 의해 $M$은 완전하다.
대칭성에 의해 $N$도 완전하다. $X$를 $x$의 열린 근방으로 바꾸어
$M$과 $N$이 강한 완전 복합체
$\mathcal{E}^\bullet$과 $\mathcal{F}^\bullet$로 표시된다고
가정할 수 있다. 그러면
$M \otimes_{\mathcal{O}_X}^\mathbf{L} N$은
$\text{Tot}(\mathcal{E}^\bullet \otimes_{\mathcal{O}_X} \mathcal{F}^\bullet)$로
표시된다. $X$를 다시 축소하여 서로 역인 동형사상
$\mathcal{O}_X \to M \otimes_{\mathcal{O}_X}^\mathbf{L} N$과
$M \otimes_{\mathcal{O}_X}^\mathbf{L} N \to \mathcal{O}_X$가
복합체의 사상
$$
\alpha : \mathcal{O}_X \to
\text{Tot}(\mathcal{E}^\bullet \otimes_{\mathcal{O}_X} \mathcal{F}^\bullet)
\quad\text{and}\quad
\beta :
\text{Tot}(\mathcal{E}^\bullet \otimes_{\mathcal{O}_X} \mathcal{F}^\bullet)
\to \mathcal{O}_X
$$
으로 주어진다고 가정할 수 있다.
보조정리 \ref{cohomology-lemma-local-actual}을 보라.
그러면 복합체의 사상으로서 $\beta \circ \alpha = 1$이고,
$D(\mathcal{O}_X)$의 사상으로서 $\alpha \circ \beta = 1$이다.
$X$를 축소하여 합성 $\alpha \circ \beta$가 성분
$$
\theta^n :
\text{Tot}^n(\mathcal{E}^\bullet \otimes_{\mathcal{O}_X} \mathcal{F}^\bullet)
\to
\text{Tot}^{n - 1}(
\mathcal{E}^\bullet \otimes_{\mathcal{O}_X} \mathcal{F}^\bullet)
$$
을 가지는 어떤 호모토피 $\theta$에 의해 $1$과 호모토픽이라고
가정할 수 있다. 앞과 같은 보조정리를 사용했다.
$R = \Gamma(X, \mathcal{O}_X)$로 놓자.
보조정리 \ref{cohomology-lemma-category-summands-finite-free}에 의해 다음을 얻는다.
\begin{enumerate}
\item $M^\bullet = \Gamma(X, \mathcal{E}^\bullet)$은
유한 사영 $R$-가군들의 유계 복합체이다.
\item $N^\bullet = \Gamma(X, \mathcal{F}^\bullet)$은
유한 사영 $R$-가군들의 유계 복합체이다.
\item $\alpha$와 $\beta$는 각각 복합체의 사상
$a : R \to \text{Tot}(M^\bullet \otimes_R N^\bullet)$와
$b : \text{Tot}(M^\bullet \otimes_R N^\bullet) \to R$에 대응한다.
\item $\theta^n$은 사상
$h^n : \text{Tot}^n(M^\bullet \otimes_R N^\bullet) \to
\text{Tot}^{n - 1}(M^\bullet \otimes_R N^\bullet)$에 대응한다.
\item $b \circ a = 1$이고 $b \circ a - 1 = dh + hd$이다.
\end{enumerate}
따라서 $M^\bullet$과 $N^\bullet$은 $D(R)$의 서로 역인 대상들을
정의한다. More on Algebra, 보조정리
\ref{more-algebra-lemma-invertible-derived}에 의해
직접곱 분해 $R = \prod_{a \leq n \leq b} R_n$과 가역
$R_n$-가군 $H^n$들이 존재하여
$M^\bullet \cong \bigoplus_{a \leq n \leq b} H^n[-n]$이다.
각 $H^n$이 사영 $R$-가군이므로, $D(R)$에서의 이 동형사상은
복합체의 사상
$$
\bigoplus H^n[-n] \longrightarrow M^\bullet
$$
으로 올릴 수 있다. 이에 대응하여 보조정리
\ref{cohomology-lemma-category-summands-finite-free}를 다시 사용하면
복합체의 사상
$$
\bigoplus H^n \otimes_R \mathcal{O}_X[-n] \to \mathcal{E}^\bullet
$$
을 얻으며, 이는 $D(\mathcal{O}_X)$에서 동형사상이다.
$\mathcal{O}_n = R_n \otimes_R \mathcal{O}_X$으로 놓으면
(2)가 성립함을 얻는다.

\medskip\noindent
$\mathcal{O}_X$의 모든 줄기가 국소환이면 (2)와 (3)의 동치는
곧바로 증명된다. 세부 사항은 생략한다.
\end{proof}
\section{콤팩트 대상}
\label{cohomology-section-compact}

\noindent
이 절에서는 환 달린 공간 위의 가군들의 유도 범주에 있는
콤팩트 대상들을 연구한다. 콤팩트 대상은 Derived Categories,
정의 \ref{derived-definition-compact-object}에서 정의되었음을 상기하자.
적절한 환 달린 공간에서는 완전 대상이 콤팩트하다.

\begin{lemma}
\label{cohomology-lemma-when-jshriek-compact}
$X$를 환 달린 공간이라 하고, $j : U \to X$를 열린집합의
포함사상이라 하자. 다음 조건들을 만족하는 정수 $d$가 존재하면
$\mathcal{O}_X$-가군 $j_!\mathcal{O}_U$는
$D(\mathcal{O}_X)$의 콤팩트 대상이다.
\begin{enumerate}
\item 모든 $p > d$에 대하여 $H^p(U, \mathcal{F}) = 0$이다.
\item 함자 $\mathcal{F} \mapsto H^p(U, \mathcal{F})$들은
직합과 가환한다.
\end{enumerate}
\end{lemma}

\begin{proof}
(1)과 (2)를 가정하자. Sheaves, 보조정리
\ref{sheaves-lemma-j-shriek-modules}에 의해
$\Hom(j_!\mathcal{O}_U, \mathcal{F}) = \mathcal{F}(U)$이므로,
$D(\mathcal{O}_X)$의 $K$에 대하여
$\Hom(j_!\mathcal{O}_U, K) = R\Gamma(U, K)$이다.
따라서 $R\Gamma(U, -)$가 직합과 가환함을 보여야 한다.
첫째 가정은 함자 $F = H^0(U, -)$가 유한 코호몰로지 차원을
가짐을 뜻한다. 또한 둘째 가정에 의해 단사 가군들의 임의의
직합은 $F$에 대해 비순환이다.
$K_i$를 $D(\mathcal{O}_X)$의 대상들의 족이라 하자.
$K_i$를 표시하며 단사 항들을 가지는 K-단사 표시자
$I_i^\bullet$을 택하자. Injectives, 정리
\ref{injectives-theorem-K-injective-embedding-grothendieck}를 보라.
비순환 대상들의 임의의 복합체에 $F$를 적용하여 $RF$를 계산할 수 있고
(Derived Categories, 보조정리
\ref{derived-lemma-unbounded-right-derived}),
$\bigoplus K_i$가 $\bigoplus I_i^\bullet$로 표시되므로
(Injectives, 보조정리 \ref{injectives-lemma-derived-products}),
$R\Gamma(U, \bigoplus K_i)$는
$\bigoplus H^0(U, I_i^\bullet)$로 표시된다.
따라서 원하는 대로 $R\Gamma(U, -)$는 직합과 가환한다.
\end{proof}

\begin{lemma}
\label{cohomology-lemma-perfect-is-compact}
$X$를 환 달린 공간이라 하자. $X$의 바탕 위상공간이 다음 성질들을
가진다고 가정하자.
\begin{enumerate}
\item $X$는 준콤팩트이다.
\item 준콤팩트 열린 부분집합들로 이루어진 기초가 존재한다.
\item 임의의 두 준콤팩트 열린집합의 교집합은 준콤팩트이다.
\end{enumerate}
$K$를 $D(\mathcal{O}_X)$의 완전 대상이라 하자. 그러면
\begin{enumerate}
\item[(a)] $K$는 다음 의미에서 $D^+(\mathcal{O}_X)$의 콤팩트 대상이다.
$M = \bigoplus_{i \in I} M_i$가 아래로 유계이면
$\Hom(K, M) = \bigoplus_{i \in I} \Hom(K, M_i)$이다.
\item[(b)] $X$가 유한 코호몰로지 차원을 가지면, 즉
$i > d$일 때 $H^i(X, \mathcal{F}) = 0$이 되게 하는 $d$가 존재하면,
$K$는 $D(\mathcal{O}_X)$의 콤팩트 대상이다.
\end{enumerate}
\end{lemma}

\begin{proof}
$K^\vee$를 $K$의 쌍대라 하자.
보조정리 \ref{cohomology-lemma-dual-perfect-complex}를 보라.
그러면 $D(\mathcal{O}_X)$의 $M$에 대해 함자적으로
$$
\Hom_{D(\mathcal{O}_X)}(K, M) =
H^0(X, K^\vee \otimes_{\mathcal{O}_X}^\mathbf{L} M)
$$
이다. $K^\vee \otimes_{\mathcal{O}_X}^\mathbf{L} -$가 직합과
가환하므로, $R\Gamma(X, -)$가 관련 직합들과 가환함을 보이면 충분하다.

\medskip\noindent
(b)의 증명. $R\Gamma(X, K) = R\Hom(\mathcal{O}_X, K)$이고
보조정리 \ref{cohomology-lemma-quasi-separated-cohomology-colimit}에 의해
$H^p(X, -)$가 직합과 가환하므로, 이는 보조정리
\ref{cohomology-lemma-when-jshriek-compact}의 특수한 경우이다.

\medskip\noindent
(a)의 증명. $\mathcal{I}_i$, $i \in I$를 단사
$\mathcal{O}_X$-가군들의 족이라 하자. 보조정리
\ref{cohomology-lemma-quasi-separated-cohomology-colimit}에 의해
$$
H^p(X, \bigoplus\nolimits_{i \in I} \mathcal{I}_i) =
\bigoplus\nolimits_{i \in I} H^p(X, \mathcal{I}_i) = 0
$$
이 모든 $p$에 대해 성립한다. 이제 $M = \bigoplus M_i$가
(a)에서와 같으면, $n < a$일 때 $H^n(M_i) = 0$이 되게 하는
$a \in \mathbf{Z}$가 존재한다. 따라서 $M_i$를 표시하며
$n < a$일 때 $\mathcal{I}_i^n = 0$인 단사
$\mathcal{O}_X$-가군들의 복합체 $\mathcal{I}_i^\bullet$을 택할 수 있다.
Derived Categories, 보조정리
\ref{derived-lemma-injective-resolutions-exist}를 보라.
Injectives, 보조정리 \ref{injectives-lemma-derived-products}에 의해
직합 복합체 $\bigoplus \mathcal{I}_i^\bullet$은 $M$을 표시한다.
Leray 비순환성
(Derived Categories, 보조정리 \ref{derived-lemma-leray-acyclicity})에 의해
$$
R\Gamma(X, M) = \Gamma(X, \bigoplus \mathcal{I}_i^\bullet) =
\bigoplus \Gamma(X, \bigoplus \mathcal{I}_i^\bullet) =
\bigoplus R\Gamma(X, M_i)
$$
이며, 이것이 원하는 결론이다.
\end{proof}
\section{사영 공식}
\label{cohomology-section-projection-formula}

\noindent
이 절에서는 사영 공식의 여러 변형을 모은다.
가장 기본적인 판은 보조정리 \ref{cohomology-lemma-projection-formula}이다.
이를 서술하고 증명한 뒤, 완전 복합체가 관련된 더 일반적인 판을 논한다.

\begin{lemma}
\label{cohomology-lemma-injective-tensor-finite-locally-free}
$X$를 환 달린 공간이라 하자.
$\mathcal{I}$를 단사 $\mathcal{O}_X$-가군이라 하고,
$\mathcal{E}$를 $\mathcal{O}_X$-가군이라 하자.
$\mathcal{E}$가 $X$ 위에서 유한 국소 자유라고 가정하자.
Modules, 정의 \ref{modules-definition-locally-free}를 보라.
그러면 $\mathcal{E} \otimes_{\mathcal{O}_X} \mathcal{I}$는
단사 $\mathcal{O}_X$-가군이다.
\end{lemma}

\begin{proof}
보조정리의 가정 아래
$$
\Hom_{\mathcal{O}_X}(\mathcal{F},
\mathcal{E} \otimes_{\mathcal{O}_X} \mathcal{I})
=
\Hom_{\mathcal{O}_X}(
\mathcal{F} \otimes_{\mathcal{O}_X} \mathcal{E}^\vee, \mathcal{I})
$$
이므로 참이다. 여기서
$\mathcal{E}^\vee = \SheafHom_{\mathcal{O}_X}(\mathcal{E}, \mathcal{O}_X)$는
$\mathcal{E}$의 쌍대이며 역시 유한 국소 자유이다.
유한 국소 자유 층과의 텐서곱은 완전 함자이므로
Homology, 보조정리 \ref{homology-lemma-characterize-injectives}에 의해
결론을 얻는다.
\end{proof}

\begin{lemma}
\label{cohomology-lemma-projection-formula}
$f : X \to Y$를 환 달린 공간의 사상이라 하자.
$\mathcal{F}$를 $\mathcal{O}_X$-가군,
$\mathcal{E}$를 $\mathcal{O}_Y$-가군이라 하자.
$\mathcal{E}$가 $Y$ 위에서 유한 국소 자유라고 가정하자.
Modules, 정의 \ref{modules-definition-locally-free}를 보라.
그러면 모든 $q \geq 0$에 대하여 동형사상
$$
\mathcal{E} \otimes_{\mathcal{O}_Y} R^qf_*\mathcal{F}
\longrightarrow
R^qf_*(f^*\mathcal{E} \otimes_{\mathcal{O}_X} \mathcal{F})
$$
이 존재한다. 실제로 $D^{+}(Y)$ 안에서 $\mathcal{F}$에 대해
함자적인 동형사상
$$
\mathcal{E} \otimes_{\mathcal{O}_Y} Rf_*\mathcal{F}
\longrightarrow
Rf_*(f^*\mathcal{E} \otimes_{\mathcal{O}_X} \mathcal{F})
$$
이 존재한다.
\end{lemma}

\begin{proof}
$X$ 위에서 단사 분해
$\mathcal{F} \to \mathcal{I}^\bullet$을 택하자.
$f^*\mathcal{E}$도 유한 국소 자유이므로 분해
$$
f^*\mathcal{E} \otimes_{\mathcal{O}_X} \mathcal{F}
\longrightarrow
f^*\mathcal{E} \otimes_{\mathcal{O}_X} \mathcal{I}^\bullet
$$
를 얻으며, 보조정리
\ref{cohomology-lemma-injective-tensor-finite-locally-free}에 의해 이는 단사 분해이다.
$f_*$를 적용하면
$$
Rf_*(f^*\mathcal{E} \otimes_{\mathcal{O}_X} \mathcal{F})
=
f_*(f^*\mathcal{E} \otimes_{\mathcal{O}_X} \mathcal{I}^\bullet).
$$
따라서 $\mathcal{O}_X$-가군 $\mathcal{F}$에 대해 함자적으로
$f_*(f^*\mathcal{E} \otimes_{\mathcal{O}_X} \mathcal{F}) =
\mathcal{E} \otimes_{\mathcal{O}_Y} f_*(\mathcal{F})$임을 보이면
보조정리가 따른다. $\mathcal{E} = \mathcal{O}_Y^{\oplus n}$이면
이는 명백하고, 일반적인 경우에는 $Y$ 위에서 국소적으로 작업하면 따른다.
세부 사항은 생략한다.
\end{proof}

\noindent
$f : X \to Y$를 환 달린 공간의 사상이라 하고,
$E \in D(\mathcal{O}_X)$와 $K \in D(\mathcal{O}_Y)$라 하자.
아무런 추가 가정 없이도 사상
\begin{equation}
\label{cohomology-equation-projection-formula-map}
Rf_*E \otimes^\mathbf{L}_{\mathcal{O}_Y} K
\longrightarrow
Rf_*(E \otimes^\mathbf{L}_{\mathcal{O}_X} Lf^*K)
\end{equation}
이 있다. 실제로 이는 표준 사상
$$
Lf^*(Rf_*E \otimes^\mathbf{L}_{\mathcal{O}_Y} K) =
Lf^*Rf_*E \otimes^\mathbf{L}_{\mathcal{O}_X} Lf^*K
\longrightarrow
E \otimes^\mathbf{L}_{\mathcal{O}_X} Lf^*K
$$
의 수반 사상이다. 이 표준 사상은
$Lf^*Rf_*E \to E$와 보조정리
\ref{cohomology-lemma-pullback-tensor-product}, \ref{cohomology-lemma-adjoint}에서 온다.
사영 공식의 상당히 일반적인 판은 다음과 같다.

\begin{lemma}
\label{cohomology-lemma-projection-formula-perfect}
$f : X \to Y$를 환 달린 공간의 사상이라 하자.
$E \in D(\mathcal{O}_X)$와 $K \in D(\mathcal{O}_Y)$라 하자.
$K$가 완전하면 $D(\mathcal{O}_Y)$에서
$$
Rf_*E \otimes^\mathbf{L}_{\mathcal{O}_Y} K =
Rf_*(E \otimes^\mathbf{L}_{\mathcal{O}_X} Lf^*K)
$$
이다.
\end{lemma}

\begin{proof}
(\ref{cohomology-equation-projection-formula-map})이 동형사상임을 확인할 때
$Y$ 위에서 국소적으로 작업해도 된다. 즉, 사상을 $V_j$에 제한하면
동형사상이 되는 열린 덮개 $\{V_j \to Y\}$를 찾아야 한다.
완전 대상의 정의에 의해 이는 $K$가
$\mathcal{O}_Y$-가군들의 강한 완전 복합체로 표시된다고
가정해도 됨을 뜻한다. 완전히 일반적으로, 명제가
$K = K_1 \oplus K_2$에 대해 참일 필요충분조건은
$K_1$과 $K_2$ 각각에 대해 참인 것이다.
따라서 $K$가 유한 자유 $\mathcal{O}_Y$-가군들의 유한 복합체라고
가정할 수 있다. 이 경우 단순 절단을 사용하는 간단한 논증으로,
$K$가 유한 자유 $\mathcal{O}_Y$-가군으로 표시되는 경우로 귀착된다.
명제는 $K$ 변수의 유한 직합인자 아래 불변이므로,
$K = \mathcal{O}_Y[n]$인 경우만 증명하면 충분하며 이 경우는 자명하다.
\end{proof}

\noindent
다음은 사영 공식이 완전히 일반적으로 성립하는 경우이다.

\begin{lemma}
\label{cohomology-lemma-projection-formula-closed-immersion}
$f : X \to Y$를 환 달린 공간의 사상이라 하되, $f$가 닫힌
부분집합 위로의 위상동형이라고 하자. 그러면
(\ref{cohomology-equation-projection-formula-map})은 언제나 동형사상이다.
\end{lemma}

\begin{proof}
$f$가 닫힌 부분집합 위로의 위상동형이므로 함자 $f_*$는 완전하다
(Modules, 보조정리 \ref{modules-lemma-i-star-exact}).
따라서 임의의 표시 복합체에 $f_*$를 적용하여 $Rf_*$를 계산한다.
$K$를 표시하는 $\mathcal{O}_Y$-가군들의 K-평탄 복합체
$\mathcal{K}^\bullet$과 $E$를 표시하는 임의의
$\mathcal{O}_X$-가군 복합체 $\mathcal{E}^\bullet$을 택하자.
그러면 $Lf^*K$는 K-평탄 $\mathcal{O}_X$-가군 복합체
$f^*\mathcal{K}^\bullet$로 표시된다
(보조정리 \ref{cohomology-lemma-pullback-K-flat}).
따라서 (\ref{cohomology-equation-projection-formula-map})의 우변은
$$
f_*\text{Tot}(\mathcal{E}^\bullet
\otimes_{\mathcal{O}_X} f^*\mathcal{K}^\bullet)
$$
로 표시된다. 같은 논증으로 좌변은
$$
\text{Tot}(f_*\mathcal{E}^\bullet \otimes_{\mathcal{O}_Y} \mathcal{K}^\bullet)
$$
로 표시된다. $f_*$가 직합과 가환하므로
(Modules, 보조정리 \ref{modules-lemma-i-star-right-adjoint}),
임의의 $\mathcal{O}_X$-가군 $\mathcal{E}$와
$\mathcal{O}_Y$-가군 $\mathcal{K}$에 대해
$$
f_*(\mathcal{E} \otimes_{\mathcal{O}_X} f^*\mathcal{K}) =
f_*\mathcal{E} \otimes_{\mathcal{O}_Y} \mathcal{K}
$$
임을 보이면 충분하다. 줄기에서 확인하겠다.
$y \in Y$라 하자. $y \not \in f(X)$이면 양변의 줄기는 영이다.
$y = f(x)$이면
$$
\mathcal{E}_x \otimes_{\mathcal{O}_{X, x}}
(\mathcal{O}_{X, x} \otimes_{\mathcal{O}_{Y, y}} \mathcal{F}_y) =
\mathcal{E}_x \otimes_{\mathcal{O}_{Y, y}} \mathcal{F}_y
$$
를 보여야 한다
(Sheaves, 보조정리 \ref{sheaves-lemma-stalks-closed-pushforward}와
보조정리 \ref{sheaves-lemma-stalk-pullback-modules}를 사용했다).
이 등식은 성립하므로 보조정리가 증명되었다.
\end{proof}

\begin{remark}
\label{cohomology-remark-compatible-with-diagram}
사상 (\ref{cohomology-equation-projection-formula-map})은 다음 의미에서
주석 \ref{cohomology-remark-base-change}의 밑변환 사상과 호환된다.
즉,
$$
\xymatrix{
X' \ar[r]_{g'} \ar[d]_{f'} &
X \ar[d]^f \\
Y' \ar[r]^g &
Y
}
$$
가 환 달린 공간의 가환 도식이라고 하자.
$E \in D(\mathcal{O}_X)$와 $K \in D(\mathcal{O}_Y)$라 하자.
그러면 도식
$$
\xymatrix{
Lg^*(Rf_*E \otimes^\mathbf{L}_{\mathcal{O}_Y} K) \ar[r]_p \ar[d]_t &
Lg^*Rf_*(E \otimes^\mathbf{L}_{\mathcal{O}_X} Lf^*K) \ar[d]_b \\
Lg^*Rf_*E \otimes^\mathbf{L}_{\mathcal{O}_{Y'}} Lg^*K \ar[d]_b &
Rf'_*L(g')^*(E \otimes^\mathbf{L}_{\mathcal{O}_X} Lf^*K) \ar[d]_t \\
Rf'_*L(g')^*E \otimes^\mathbf{L}_{\mathcal{O}_{Y'}} Lg^*K \ar[rd]_p &
Rf'_*(L(g')^*E \otimes^\mathbf{L}_{\mathcal{O}_{Y'}} L(g')^*Lf^*K) \ar[d]_c \\
&
Rf'_*(L(g')^*E \otimes^\mathbf{L}_{\mathcal{O}_{Y'}} L(f')^*Lg^*K)
}
$$
은 가환한다. 여기서 $t$로 표시된 화살표는 보조정리
\ref{cohomology-lemma-pullback-tensor-product}를 적용하여 얻고,
$b$로 표시된 화살표는 주석 \ref{cohomology-remark-base-change}를 적용하여 얻으며,
$p$로 표시된 화살표는 (\ref{cohomology-equation-projection-formula-map})을
적용하여 얻는다. 또한 $c$는
$L(g')^* \circ Lf^* = L(f')^* \circ Lg^*$에서 온다.
확인은 생략한다.
\end{remark}
\section{Berthelot와 Ogus가 도입한 연산자}
\label{cohomology-section-eta}

\noindent
이 절에서는 More on Algebra, 절
\ref{more-algebra-section-eta}에서 시작한 논의를 이어 간다.
독자는 먼저 그 절을 읽기를 강하게 권한다.

\begin{lemma}
\label{cohomology-lemma-invertible-ideal-sheaf}
$(X, \mathcal{O}_X)$를 환 달린 공간이라 하고,
$\mathcal{I} \subset \mathcal{O}_X$를 아이디얼의 층이라 하자.
다음 두 조건을 생각하자.
\begin{enumerate}
\item 모든 $x \in X$에 대하여 $x$의 열린 근방 $U \subset X$와
$f \in \mathcal{I}(U)$가 존재하여
$\mathcal{I}|_U = \mathcal{O}_U \cdot f$이고
$f : \mathcal{O}_U \to \mathcal{O}_U$가 단사이다.
\item $\mathcal{I}$는 $\mathcal{O}_X$-가군으로서 가역이다.
\end{enumerate}
그러면 (1)은 (2)를 함의한다. 구조층의 모든 줄기
$\mathcal{O}_{X, x}$가 국소환이면 그 역도 참이다.
\end{lemma}

\begin{proof}
생략한다. 힌트: Modules, 보조정리
\ref{modules-lemma-invertible-is-locally-free-rank-1}을 사용하라.
\end{proof}

\begin{situation}
\label{cohomology-situation-eta}
$(X, \mathcal{O}_X)$를 환 달린 공간이라 하자.
$\mathcal{I} \subset \mathcal{O}_X$를 보조정리
\ref{cohomology-lemma-invertible-ideal-sheaf}의 조건 (1)을 만족하는
아이디얼의 층이라 하자\footnote{이 절의 논의는 Modules, 절
\ref{modules-section-invertible}에서 정의된 의미로
$\mathcal{I}$가 가역 $\mathcal{O}_X$-가군이라는 것만 요구하는
경우로 일반화할 수 있다.}.
\end{situation}

\begin{lemma}
\label{cohomology-lemma-I-torsion-free}
상황 \ref{cohomology-situation-eta}에서
$\mathcal{F}$를 $\mathcal{O}_X$-가군이라 하자. 다음은 동치이다.
\begin{enumerate}
\item $\mathcal{I}$에 의해 소멸되는 절단들의 부분층
$\mathcal{F}[\mathcal{I}] \subset \mathcal{F}$은 영이다.
\item 모든 $n \geq 1$에 대하여 부분층
$\mathcal{F}[\mathcal{I}^n]$은 영이다.
\item 곱셈 사상
$\mathcal{I} \otimes_{\mathcal{O}_X} \mathcal{F} \to \mathcal{F}$은
단사이다.
\item 어떤 $f \in \mathcal{I}(U)$에 대하여
$\mathcal{I}|_U = \mathcal{O}_U \cdot f$인 모든 열린집합
$U \subset X$에 대해 사상
$f : \mathcal{F}|_U \to \mathcal{F}|_U$는 단사이다.
\item 모든 $x \in X$와 아이디얼
$\mathcal{I}_x \subset \mathcal{O}_{X, x}$의 생성원 $f$에 대하여,
원소 $f$는 줄기 $\mathcal{F}_x$ 위의 영인자가 아니다.
\end{enumerate}
\end{lemma}

\begin{proof}
생략한다.
\end{proof}

\noindent
상황 \ref{cohomology-situation-eta}에서
$\mathcal{F}$를 $\mathcal{O}_X$-가군이라 하자.
보조정리 \ref{cohomology-lemma-I-torsion-free}의 동치 조건들이 성립하면
$\mathcal{F}$가 {\it $\mathcal{I}$-꼬임이 없다}고 하겠다.
이 경우 임의의 $i \in \mathbf{Z}$에 대하여
$$
\mathcal{I}^i\mathcal{F} =
\mathcal{I}^{\otimes i} \otimes_{\mathcal{O}_X} \mathcal{F}
$$
로 나타내며, 포함관계
$$
\ldots \subset
\mathcal{I}^{i + 1}\mathcal{F} \subset
\mathcal{I}^i\mathcal{F} \subset
\mathcal{I}^{i - 1}\mathcal{F} \subset \ldots
$$
를 가진다. 가군 $\mathcal{I}^i\mathcal{F}$들은
$\mathcal{O}_X$-가군으로서 국소적으로는 $\mathcal{F}$와
동형이지만 전역적으로는 그렇지 않다.

\medskip\noindent
$\mathcal{F}^\bullet$을 $\mathcal{I}$-꼬임 없는
$\mathcal{O}_X$-가군들의 복합체라 하고, 그 미분을
$d^i : \mathcal{F}^i \to \mathcal{F}^{i + 1}$라 하자.
이 경우 $\eta_\mathcal{I}\mathcal{F}^\bullet$을 항들이
\begin{align*}
(\eta_\mathcal{I}\mathcal{F})^i
& =
\Ker\left(
d^i, -1 :
\mathcal{I}^i\mathcal{F}^i \oplus \mathcal{I}^{i + 1}\mathcal{F}^{i + 1}
\to
\mathcal{I}^i\mathcal{F}^{i + 1}
\right) \\
& =
\Ker\left(d^i :
\mathcal{I}^i\mathcal{F}^i
\to
\mathcal{I}^i\mathcal{F}^{i + 1}/
\mathcal{I}^{i + 1}\mathcal{F}^{i + 1}
\right)
\end{align*}
이고 미분이 $d^i$에서 유도되는 복합체로 정의한다.
다시 말해 $(\eta_\mathcal{I}\mathcal{F})^i$의 국소 절단 $s$는
$\mathcal{I}^i\mathcal{F}^i$의 국소 절단 $s$로서,
$\mathcal{I}^i\mathcal{F}^{i + 1}$ 안의 상 $d^i(s)$가 부분층
$\mathcal{I}^{i + 1}\mathcal{F}^{i + 1}$에 속하는 것과 같다.
$\eta_\mathcal{I}\mathcal{F}^\bullet$도
$\mathcal{I}$-꼬임 없는 가군들의 복합체임에 유의하자.

\medskip\noindent
$a^\bullet : \mathcal{F}^\bullet \to \mathcal{G}^\bullet$을
$\mathcal{I}$-꼬임 없는 $\mathcal{O}_X$-가군 복합체들의 사상이라 하자.
그러면 사상들
$\mathcal{I}^i\mathcal{F}^i \to \mathcal{I}^i\mathcal{G}^i$에서
유도되는 복합체의 사상
$$
\eta_\mathcal{I} a^\bullet :
\eta_\mathcal{I}\mathcal{F}^\bullet
\longrightarrow
\eta_\mathcal{I}\mathcal{G}^\bullet
$$
을 얻는다. 독자는 이로써 $\mathcal{I}$-꼬임 없는
$\mathcal{O}_X$-가군 복합체들의 범주 위의 자기함자를 얻는다는
것을 확인할 수 있다.

\medskip\noindent
$a^\bullet, b^\bullet : \mathcal{F}^\bullet \to \mathcal{G}^\bullet$을
$\mathcal{I}$-꼬임 없는 $\mathcal{O}_X$-가군 복합체들의 두 사상이라 하고,
$h = \{h^i : \mathcal{F}^i \to \mathcal{G}^{i - 1}\}$를
$a^\bullet$과 $b^\bullet$ 사이의 호모토피라 하자.
$\eta_\mathcal{I}h$를 $x$를 $h^i(x)$로 보내는 사상들의 족
$(\eta_\mathcal{I}h)^i : (\eta_\mathcal{I}\mathcal{F})^i \to
(\eta_\mathcal{I}\mathcal{G})^{i - 1}$로 정의한다.
이는 $x$가 $\mathcal{I}^i\mathcal{F}^i$의 국소 절단이면
$h^i(x)$가 $\mathcal{I}^i\mathcal{G}^{i - 1}$의 국소 절단이고,
후자는 분명히 $(\eta_\mathcal{I}\mathcal{G})^{i - 1}$에
포함되므로 잘 정의된다. 독자는 $\eta_\mathcal{I}h$가
$\eta_\mathcal{I}a^\bullet$과
$\eta_\mathcal{I}b^\bullet$ 사이의 호모토피임을 확인할 수 있다.
종합하면 가법 범주인 $\mathcal{I}$-꼬임 없는
$\mathcal{O}_X$-가군들의 호모토피 범주
(Derived Categories, 절 \ref{derived-section-homotopy}) 위의 함자
$$
\eta_f :
K(\mathcal{I}\text{-torsion free }\mathcal{O}_X\text{-modules})
\longrightarrow
K(\mathcal{I}\text{-torsion free }\mathcal{O}_X\text{-modules})
$$
를 얻는다. $\eta_\mathcal{I}$가 삼각 범주의 완전 함자라는 의미는 없다.
More on Algebra, 예 \ref{more-algebra-example-eta-not-distinguished}와
비교하라.

\begin{lemma}
\label{cohomology-lemma-eta-stalks}
상황 \ref{cohomology-situation-eta}에서
$\mathcal{F}^\bullet$을 $\mathcal{I}$-꼬임 없는
$\mathcal{O}_X$-가군들의 복합체라 하자.
$x \in X$에 대하여 생성원 $f \in \mathcal{I}_x$를 택하자.
그러면 줄기 $(\eta_\mathcal{I}\mathcal{F}^\bullet)_x$는
More on Algebra, 절 \ref{more-algebra-section-eta}에서 구성한
복합체 $\eta_f\mathcal{F}^\bullet_x$와 표준적으로 동형이다.
\end{lemma}

\begin{proof}
생략한다.
\end{proof}

\begin{lemma}
\label{cohomology-lemma-eta-first-property}
상황 \ref{cohomology-situation-eta}에서
$\mathcal{F}^\bullet$을 $\mathcal{I}$-꼬임 없는
$\mathcal{O}_X$-가군들의 복합체라 하자. 코호몰로지 층의 표준 동형사상
$$
\mathcal{I}^{\otimes i} \otimes_{\mathcal{O}_X}
\left(
H^i(\mathcal{F}^\bullet)/H^i(\mathcal{F}^\bullet)[\mathcal{I}]
\right)
\longrightarrow H^i(\eta_\mathcal{I}\mathcal{F}^\bullet)
$$
이 존재한다.
\end{lemma}

\begin{proof}
다음과 같이 사상
$$
\mathcal{I}^{\otimes i} \otimes_{\mathcal{O}_X} H^i(\mathcal{F}^\bullet)
\longrightarrow
H^i(\eta_\mathcal{I}\mathcal{F}^\bullet)
$$
을 정의한다. $g$를 $\mathcal{I}^{\otimes i}$의 국소 절단이라 하고,
$\overline{s}$를 $H^i(\mathcal{F}^\bullet)$의 국소 절단이라 하자.
그러면 $\overline{s}$는 국소적으로
$\Ker(d^i : \mathcal{F}^i \to \mathcal{F}^{i + 1})$의 국소 절단
$s$의 동치류이다. $g \otimes \overline{s}$를
$(\eta_\mathcal{I}\mathcal{F})^i \subset \mathcal{I}^i\mathcal{F}$의
국소 절단 $gs$로 보낸다. 물론 $gs$는
$\eta_\mathcal{I}\mathcal{F}^\bullet$ 위의 $d^i$의 핵에 속하므로
$H^i(\eta_\mathcal{I}\mathcal{F}^\bullet)$의 국소 절단을 정의한다.
이것이 잘 정의된다는 확인에는 문제가 없다.
이 사상이 보조정리의 동형사상을 통해 인수분해된다고 주장한다.
이는 줄기에서 확인할 수 있고, 따라서 보조정리
\ref{cohomology-lemma-eta-stalks}에 의해 More on Algebra, 보조정리
\ref{more-algebra-lemma-eta-first-property}의 결과로 옮겨진다.
\end{proof}

\begin{lemma}
\label{cohomology-lemma-eta-qis}
상황 \ref{cohomology-situation-eta}에서
$\mathcal{F}^\bullet \to \mathcal{G}^\bullet$을
$\mathcal{I}$-꼬임 없는 $\mathcal{O}_X$-가군 복합체들의 사상이라 하자.
그러면 유도된 사상
$\eta_\mathcal{I}\mathcal{F}^\bullet \to \eta_\mathcal{I}\mathcal{G}^\bullet$도
준동형이다.
\end{lemma}

\begin{proof}
보조정리 \ref{cohomology-lemma-eta-first-property}의 동형사상들이
복합체의 사상과 호환되므로 참이다.
\end{proof}

\begin{lemma}
\label{cohomology-lemma-Leta}
상황 \ref{cohomology-situation-eta}에는 가법 함자\footnote{이 함자는 완전하지 않다.
즉, 특수 삼각형을 특수 삼각형으로 보내지 않는다는 점에 유의하라.}
$L\eta_\mathcal{I} : D(\mathcal{O}_X) \to D(\mathcal{O}_X)$가 존재하며,
$D(\mathcal{O}_X)$의 $M$이 $\mathcal{I}$-꼬임 없는
$\mathcal{O}_X$-가군들의 복합체 $\mathcal{F}^\bullet$로 표시되면
$L\eta_\mathcal{I}M = \eta_\mathcal{I}\mathcal{F}^\bullet$이다.
사상에 대해서도 마찬가지이다.
\end{lemma}

\begin{proof}
$\mathcal{T} \subset \textit{Mod}(\mathcal{O}_X)$를
$\mathcal{I}$-꼬임 없는 $\mathcal{O}_X$-가군들의 충만한 부분범주라 하자.
이에 대응하여 $K(\mathcal{T})$를 $K(\mathcal{O}_X)$의 충만한
삼각 부분범주로 보는 포함
$$
K(\mathcal{T})
\quad\subset\quad
K(\textit{Mod}(\mathcal{O}_X)) = K(\mathcal{O}_X)
$$
이 있다. $S \subset \text{Arrows}(K(\mathcal{T}))$를 준동형들의
모임이라 하자. Derived Categories, 보조정리
\ref{derived-lemma-localization-subcategory}를 적용하여 사상
$$
S^{-1}K(\mathcal{T}) \longrightarrow D(\mathcal{O}_X)
$$
이 삼각 범주의 동치임을 보이겠다. 그 보조정리에 따르면 다음을
증명하면 충분하다. $\mathcal{O}_X$-가군들의 복합체
$\mathcal{G}^\bullet$이 주어지면, $\mathcal{I}$-꼬임 없는
$\mathcal{O}_X$-가군들의 복합체 $\mathcal{F}^\bullet$과
준동형 $\mathcal{F}^\bullet \to \mathcal{G}^\bullet$이 존재한다.
보조정리 \ref{cohomology-lemma-K-flat-resolution}에 의해 복합체
$\mathcal{F}^\bullet$이 K-평탄이고(이 사실은 사용하지 않는다)
평탄한 $\mathcal{O}_X$-가군 $\mathcal{F}^i$들로 이루어지게 하는
준동형 $\mathcal{F}^\bullet \to \mathcal{G}^\bullet$을 찾을 수 있다.
보조정리 \ref{cohomology-lemma-I-torsion-free}의 셋째 특징짓기에 의해
평탄 $\mathcal{O}_X$-가군은 $\mathcal{I}$-꼬임 없는
$\mathcal{O}_X$-가군이므로 증명이 끝난다.

\medskip\noindent
이 예비 작업을 마치고 $L\eta_f$를 정의할 수 있다.
즉, 이 절에서 보조정리 \ref{cohomology-lemma-I-torsion-free} 뒤의 논의에 의해
이미 잘 정의된 함자
$$
K(\mathcal{T}) \xrightarrow{\eta_f} K(\mathcal{T}) \to
K(\mathcal{O}_X) \to D(\mathcal{O}_X)
$$
가 있다. 보조정리 \ref{cohomology-lemma-eta-qis}에 의해 이는 준동형을
준동형으로 보낸다. 따라서 Categories, 보조정리
\ref{categories-lemma-properties-left-localization}에 의해 이 함자는
$S^{-1}K(\mathcal{T}) = D(\mathcal{O}_X)$를 통해 인수분해된다.
\end{proof}
\noindent
상황 \ref{cohomology-situation-eta}에서 Bockstein 연산자들을 구성하자.
먼저 가환 도식
$$
\xymatrix{
0 \ar[r] &
\mathcal{I}^{i + 1} \ar[r] \ar[d] &
\mathcal{I}^i \ar[r] \ar[d] &
\mathcal{I}^i/\mathcal{I}^{i + 1}
\ar[r] \ar@{=}[d] &
0 \\
0 \ar[r] &
\mathcal{I}^{i + 1}/\mathcal{I}^{i + 2} \ar[r] &
\mathcal{I}^i/\mathcal{I}^{i + 2} \ar[r] &
\mathcal{I}^i/\mathcal{I}^{i + 1} \ar[r] &
0
}
$$
이 있음을 관찰하자. 두 행은 $\mathcal{O}_X$-가군들의 짧은 완전
수열이다. $M$을 $D(\mathcal{O}_X)$의 대상이라 하자.
위 도식에 $M$을 텐서하면 특수 삼각형들의 사상
$$
\xymatrix{
M \otimes^\mathbf{L} \mathcal{I}^{i + 1} \ar[r] \ar[d] &
M \otimes^\mathbf{L} \mathcal{I}^i \ar[r] \ar[d] &
M \otimes^\mathbf{L} \mathcal{I}^i/\mathcal{I}^{i + 1} \ar[d]^{\text{id}} \\
M \otimes^\mathbf{L} \mathcal{I}^{i + 1}/\mathcal{I}^{i + 2} \ar[r] &
M \otimes^\mathbf{L} \mathcal{I}^i/\mathcal{I}^{i + 2} \ar[r] &
M \otimes^\mathbf{L} \mathcal{I}^i/\mathcal{I}^{i + 1}
}
$$
을 얻는다. 특히 아래쪽 삼각형에 딸린 코호몰로지 층의 긴 완전
수열은 모든 $i \in \mathbf{Z}$에 대하여 {\it Bockstein 연산자}
$$
\beta = \beta^i :
H^i(M \otimes^\mathbf{L} \mathcal{I}^i/\mathcal{I}^{i + 1})
\longrightarrow
H^{i + 1}(M \otimes^\mathbf{L} \mathcal{I}^{i + 1}/\mathcal{I}^{i + 2})
$$
를 결정한다. 나중에 사용하기 위해, 위 가환 도식에 의해
Bockstein 연산자의 인수분해
\begin{equation}
\label{cohomology-equation-factorization-bockstein}
\vcenter{
\xymatrix{
H^i(M \otimes^\mathbf{L} \mathcal{I}^i/\mathcal{I}^{i + 1})
\ar[r]_\delta \ar[rd]_\beta &
H^{i + 1}(M \otimes^\mathbf{L} \mathcal{I}^{i + 1}) \ar[d] \\
&
H^{i + 1}(M \otimes^\mathbf{L} \mathcal{I}^{i + 1}/\mathcal{I}^{i + 2})
}
}
\end{equation}
가 있음을 기록한다. 여기서 $\delta$는 위 가환 도식의 위쪽
특수 삼각형에서 오는 경계 연산자이다.
복합체
\begin{equation}
\label{cohomology-equation-complex-bocksteins}
H^\bullet(M/\mathcal{I}) =
\left[
\begin{matrix}
\ldots \\
\downarrow \\
H^{i - 1}(M \otimes^\mathbf{L} \mathcal{I}^{i - 1}/\mathcal{I}^i) \\
\downarrow \beta \\
H^i(M \otimes^\mathbf{L} \mathcal{I}^i/\mathcal{I}^{i + 1}) \\
\downarrow \beta \\
H^{i + 1}(M \otimes^\mathbf{L} \mathcal{I}^{i + 1}/\mathcal{I}^{i + 2}) \\
\downarrow \\
\ldots
\end{matrix}
\right]
\end{equation}
를 얻는다. 즉, $\beta \circ \beta = 0$이다.
실제로 이는 줄기에서 확인할 수 있고, 그 경우 More on Algebra,
절 \ref{more-algebra-section-eta}에서 보인 대응하는 대수 결과로부터
따른다. 다른 증명은 다음과 같다. 짧은 완전 수열
$0 \to \mathcal{I}^{i + 1}/\mathcal{I}^{i + 2}
\to \mathcal{I}^i/\mathcal{I}^{i + 2}
\to \mathcal{I}^i/\mathcal{I}^{i + 1} \to 0$은
$D(\mathcal{O}_X)$에서 사상
$b^i : \mathcal{I}^i/\mathcal{I}^{i + 1} \to
(\mathcal{I}^{i + 1}/\mathcal{I}^{i + 2})[1]$을 정의하며,
$M$을 텐서하고 코호몰로지 층을 취하면 위의 사상 $\beta$를 유도한다.
그러면 가군 $\mathcal{I}^i/\mathcal{I}^{i + 3}$이
$\mathcal{I}^i/\mathcal{I}^{i + 1}$에 의한
$\mathcal{I}^{i + 1}/\mathcal{I}^{i + 3}$의 확대라는 사실과
Derived Categories, 보조정리
\ref{derived-lemma-cup-ext-1-zero}의 판정법을 사용하여 합성
$b^{i + 1}[1] \circ b^i : \mathcal{I}^i/\mathcal{I}^{i + 1} \to
(\mathcal{I}^{i + 1}/\mathcal{I}^{i + 2})[1] \to
(\mathcal{I}^{i + 2}/\mathcal{I}^{i + 3})[2]$가
$D(\mathcal{O}_X)$에서 영임을 보인다.

\begin{lemma}
\label{cohomology-lemma-eta-second-property}
상황 \ref{cohomology-situation-eta}에서
$M$을 $D(\mathcal{O}_X)$의 대상이라 하자. 표준 동형사상
$$
L\eta_\mathcal{I}M \otimes^\mathbf{L} \mathcal{O}_X/\mathcal{I}
\longrightarrow
H^\bullet(M/\mathcal{I})
$$
이 $D(\mathcal{O}_X)$에 존재한다. 여기서 우변은 복합체
(\ref{cohomology-equation-complex-bocksteins})이다.
\end{lemma}

\begin{proof}
보조정리 \ref{cohomology-lemma-eta-qis}의 $L\eta_\mathcal{I}$ 구성에 의해,
$M$이 $\mathcal{I}$-꼬임 없는 $\mathcal{O}_X$-가군 복합체
$\mathcal{F}^\bullet$로 표시된다고 가정할 수 있다.
그러면 $L\eta_\mathcal{I}M$은 복합체
$\eta_\mathcal{I}\mathcal{F}^\bullet$로 표시되며,
이 역시 $\mathcal{I}$-꼬임 없는 $\mathcal{O}_X$-가군들의 복합체이다.
따라서 $L\eta_\mathcal{I}M \otimes^\mathbf{L} \mathcal{O}_X/\mathcal{I}$은
복합체
$\eta_\mathcal{I}\mathcal{F}^\bullet \otimes \mathcal{O}_X/\mathcal{I}$로
표시된다. 마찬가지로 복합체 $H^\bullet(M/\mathcal{I})$의 항들은
$H^i(\mathcal{F}^\bullet \otimes \mathcal{I}^i/\mathcal{I}^{i + 1})$이다.

\medskip\noindent
$f$를 $\mathcal{I}$의 국소 생성원이라 하고,
$s$를 $(\eta_\mathcal{I}\mathcal{F})^i$의 국소 절단이라 하자.
그러면 $\mathcal{F}^i$의 어떤 국소 절단 $s'$에 대해
$s = f^is'$로 쓸 수 있고, 마찬가지로
$\mathcal{F}^{i + 1}$의 어떤 국소 절단 $t$에 대해
$d^i(s) = f^{i + 1}t$로 쓸 수 있다.
따라서 $d^i$는 $f^is'$를
$\mathcal{F}^{i + 1} \otimes \mathcal{I}^i/\mathcal{I}^{i + 1}$에서
영으로 보낸다. 그러므로 $s$를
$H^i(\mathcal{F}^\bullet \otimes \mathcal{I}^i/\mathcal{I}^{i + 1})$의
$f^is'$의 동치류로 보낼 수 있다. 이 규칙은
$\mathcal{O}_X$-가군의 사상
$$
(\eta_\mathcal{I}\mathcal{F})^i \otimes \mathcal{O}_X/\mathcal{I}
\longrightarrow
H^i(\mathcal{F}^\bullet \otimes \mathcal{I}^i/\mathcal{I}^{i + 1})
$$
을 정의한다. 계산해 보면 이 사상들은 미분과 호환된다.
본질적으로 $\beta$가 $f^is'$의 동치류를
$f^{i + 1}t$의 동치류로 보내기 때문이다.
따라서 보조정리의 화살표를 표시하는 복합체의 사상을 얻는다.

\medskip\noindent
증명을 끝내기 위해, 앞 문단의 구성이 줄기에서
More on Algebra, 보조정리
\ref{more-algebra-lemma-eta-second-property}에서 구성한 사상들과
일치함을 관찰하면 결론을 얻는다.
\end{proof}

\begin{lemma}
\label{cohomology-lemma-eta-tensor-invertible}
상황 \ref{cohomology-situation-eta}에서
$\mathcal{F}^\bullet$을 $\mathcal{I}$-꼬임 없는
$\mathcal{O}_X$-가군들의 복합체라 하고,
$\mathcal{L}$을 가역 $\mathcal{O}_X$-가군이라 하자.
그러면
$\eta_\mathcal{I}(\mathcal{F}^\bullet \otimes \mathcal{L}) =
(\eta_\mathcal{I}\mathcal{F}^\bullet) \otimes \mathcal{L}$이다.
\end{lemma}

\begin{proof}
구성에서 곧바로 따른다.
\end{proof}

\begin{lemma}
\label{cohomology-lemma-eta-cohomology-locally-free}
상황 \ref{cohomology-situation-eta}에서
$M$을 $D(\mathcal{O}_X)$의 대상이라 하자.
$\mathcal{O}_{X, x}$가 영이 아닌 $x \in X$를 택하자.
$H^i(M)_x$가 $\mathcal{O}_{X, x}$ 위에서 유한 자유이면,
$H^i(L\eta_\mathcal{I}M)_x$도 같은 계수의 유한 자유
$\mathcal{O}_{X, x}$-가군이다.
\end{lemma}

\begin{proof}
$f \in \mathcal{O}_{X, x}$가 줄기 $\mathcal{I}_x$를 생성한다고 하자.
그러면 $f$는 $\mathcal{O}_{X, x}$의 영인자가 아니므로
$H^i(M)_x[f] = 0$이다. 따라서 보조정리
\ref{cohomology-lemma-eta-first-property}에 의해
$H^i(L\eta_\mathcal{I}M)_x$는
$\mathcal{I}^i_x \otimes_{\mathcal{O}_{X, x}} H^i(M)_x$와 동형이며,
이는 원하는 대로 같은 계수의 자유 가군이다.
\end{proof}
